% Options for packages loaded elsewhere
\PassOptionsToPackage{unicode}{hyperref}
\PassOptionsToPackage{hyphens}{url}
\documentclass[
  ukrainian,
]{article}
\usepackage{xcolor}
\usepackage[margin=2.5cm]{geometry}
\usepackage{amsmath,amssymb}
\setcounter{secnumdepth}{-\maxdimen} % remove section numbering
\usepackage{iftex}
\ifPDFTeX
  \usepackage[T1]{fontenc}
  \usepackage[utf8]{inputenc}
  \usepackage{textcomp} % provide euro and other symbols
\else % if luatex or xetex
  \usepackage{unicode-math} % this also loads fontspec
  \defaultfontfeatures{Scale=MatchLowercase}
  \defaultfontfeatures[\rmfamily]{Ligatures=TeX,Scale=1}
\fi
\usepackage{lmodern}
\ifPDFTeX\else
\setmainfont{DejaVu Serif}
\setsansfont{DejaVu Sans}
\setmonofont{DejaVu Sans Mono}
\setmathfont{STIX Two Math}
\fi
\ifPDFTeX\else
  % xetex/luatex font selection
\fi
% Use upquote if available, for straight quotes in verbatim environments
\IfFileExists{upquote.sty}{\usepackage{upquote}}{}
\IfFileExists{microtype.sty}{% use microtype if available
  \usepackage[]{microtype}
  \UseMicrotypeSet[protrusion]{basicmath} % disable protrusion for tt fonts
}{}
\makeatletter
\@ifundefined{KOMAClassName}{% if non-KOMA class
  \IfFileExists{parskip.sty}{%
    \usepackage{parskip}
  }{% else
    \setlength{\parindent}{0pt}
    \setlength{\parskip}{6pt plus 2pt minus 1pt}}
}{% if KOMA class
  \KOMAoptions{parskip=half}}
\makeatother
\usepackage{longtable,booktabs,array}
\newcounter{none} % for unnumbered tables
\usepackage{calc} % for calculating minipage widths
% Correct order of tables after \paragraph or \subparagraph
\usepackage{etoolbox}
\makeatletter
\patchcmd\longtable{\par}{\if@noskipsec\mbox{}\fi\par}{}{}
\makeatother
% Allow footnotes in longtable head/foot
\IfFileExists{footnotehyper.sty}{\usepackage{footnotehyper}}{\usepackage{footnote}}
\makesavenoteenv{longtable}
\usepackage{graphicx}
\graphicspath{{/Users/kalinicenkovitalijmikolajovic/.codex/worktrees/5c15/Script_for_translate_book/Output_Final/}}
\setkeys{Gin}{width=0.82\linewidth,height=0.82\textheight,keepaspectratio}
\makeatletter
\newsavebox\pandoc@box
\newcommand*\pandocbounded[1]{% scales image to fit in text height/width
  \sbox\pandoc@box{#1}%
  \Gscale@div\@tempa{\textheight}{\dimexpr\ht\pandoc@box+\dp\pandoc@box\relax}%
  \Gscale@div\@tempb{\linewidth}{\wd\pandoc@box}%
  \ifdim\@tempb\p@<\@tempa\p@\let\@tempa\@tempb\fi% select the smaller of both
  \ifdim\@tempa\p@<\p@\scalebox{\@tempa}{\usebox\pandoc@box}%
  \else\usebox{\pandoc@box}%
  \fi%
}
% Set default figure placement to htbp
\def\fps@figure{htbp}
\makeatother
\ifLuaTeX
\usepackage[bidi=basic,shorthands=off,provide=*]{babel}
\else
\usepackage[bidi=default,shorthands=off,provide=*]{babel}
\fi
\ifLuaTeX
  \usepackage{selnolig} % disable illegal ligatures
\fi
\setlength{\emergencystretch}{3em} % prevent overfull lines
\providecommand{\tightlist}{%
  \setlength{\itemsep}{0pt}\setlength{\parskip}{0pt}}
\usepackage{bookmark}
\IfFileExists{xurl.sty}{\usepackage{xurl}}{} % add URL line breaks if available
\urlstyle{same}
\hypersetup{
  pdflang={uk},
  hidelinks,
  pdfcreator={LaTeX via pandoc}}

\author{}
\date{}

\begin{document}

ЧАСТКОВО ДИФЕРЕНЦІАЛ РІВНЯННЯ для Вчені та Інженери

Стенлі Дж. Фарлоу, професор математики, Університет Мену

\section{С. Фарлоу}\label{ux441.-ux444ux430ux440ux43bux43eux443}

\section{РІВНЯННЯ З ПОХІДНИМИ ПОХІДНИМИ ДЛЯ НАУКОВЦІВ І
ІНЖЕНЕРІВ}\label{ux440ux456ux432ux43dux44fux43dux43dux44f-ux437-ux43fux43eux445ux456ux434ux43dux438ux43cux438-ux43fux43eux445ux456ux434ux43dux438ux43cux438-ux434ux43bux44f-ux43dux430ux443ux43aux43eux432ux446ux456ux432-ux456-ux456ux43dux436ux435ux43dux435ux440ux456ux432}

Перекладено з англійської \textbf{A.I.} Плісом

Редактор

С. І. Похожаєва

5\textbar{} 66K 22.161.6 φ24 υ ΥΠΚ 517.2

\section{Фарлоу С.}\label{ux444ux430ux440ux43bux43eux443-ux441.}

F 24 Рівняння з частими похідними для науковців та інженерів. Москва,
видавництво Мир, 1985, 384 с.

King's American Mathematician, який є підручником з теорії
диференціальних рівнянь у похідних похідних. Він характеризується
компактністю, ясністю та ясністю подачі, а також неформальним підходом
до подачі матеріалу. Вона містить багато ілюстрацій, графіків і діаграм,
але замість суворих доказів часто даються роздуми, засновані на інтуїції
або аналогіях.

Для інженерів і нематематиків --- біологів, диміків, а також студентів
університету.
\texttt{\textbackslash{}Phi\ =\ \textbackslash{}frac\{1702050000-380\}\{041(01)-85\}\ 26-85,\ \textbackslash{}text\{\ q.\ 1\}}
BBK 22.161.6

Редакційна колегія журналу Literature on Mathematical Sciences

© переклад російською, Мир, 1985

© John Wiley \& Sons, Inc.~1982. Всі права захищені. Авторизований
переклад з англомовного видання, опублікований John Wiley \& Sons, Inc.

\section{З РЕДАКТОРА
ПЕРЕКЛАДУ}\label{ux437-ux440ux435ux434ux430ux43aux442ux43eux440ux430-ux43fux435ux440ux435ux43aux43bux430ux434ux443}

Книга, запропонована для читачів, слід розглядати як підручник для
курсів: «Диференціальні рівняння в похідних похідних» та «Рівняння
математичної фізики» для студентів технічних університетів. Вона також
буде корисною для інженерів різних спеціальностей. У книзі представлено
традиційні методи розв'язання диференціальних рівнянь у похідних
похідних, включаючи деякі чисельні методи. Водночас особлива увага
приділяється методу розділення змінних і методу інтегральних
перетворень, і цю частину книги можна використовувати на практичних
заняттях в університетах

Орієнтація книги на широку «нематематичну» аудиторію визначила стиль
викладання розглянутих питань: інтуїтивний підхід і велика увага до
фізичного значення не лише самих рівнянь, а й меж і початкових умов для
різних задач. Вважається, що читачі знайомі лише з основами
диференціального та інтегрального числення та елементами теорії
звичайних диференціальних рівнянь.

Фахівці можуть не задовольнитися інтуїтивним описом певних понять і
методів автором, але більш суворий виклад неминуче призведе до
ускладнення книги і звузить коло її читачів, що, очевидно, не було
наміром автора.

Більшим запереченням може бути відсутність загальних результатів щодо
задач для диференціальних рівнянь у похідних з змінними коефіцієнтами,
включно з рівняннями вищого порядку. Однак читач має чітко розуміти, що
ця книга є вступним курсом з теорії диференціальних рівнянь у часткових
похідних. Питання лінійної теорії диференціальних рівнянь, які автор не
включив, можна знайти в літературі, наведеній наприкінці книги.

Матеріал книги подається автором з великою методологічною майстерністю.
Однак не все було однаково успішним, і винний читач знайде це,
наприклад, у вивченні принципу Дюамеля, сферичних заповітів та деяких
інших уривків у книзі.

Без сумніву, запропонована книга С. Фарлоу знайде читача не лише серед
студентів і інженерів, а й серед викладачів, які можуть багато чого
навчитися з конкретного способу презентації та розробки спеціальних
математичних курсів.

\section{ПЕРЕДМОВА}\label{ux43fux435ux440ux435ux434ux43cux43eux432ux430}

За останні кілька років спостерігається значне зростання кількості
учнів, які починають вивчати диференціальні рівняння з похідними у своїх
одинадцятих класах. Крім того, багато з них спеціалізуються на сферах,
де інтуїція важливіша за математичну строгість. Тому, пишучи цю книгу, я
прагнув створити текст, який стимулює інтуїтивне мислення читачів, не
втрачаючи при цьому значної математичної строгості. Можливо, на одному
полюсі, подати предмет на високому рівні «епсилон-дельта», але тоді, як
це зазвичай буває, студенти не знатимуть, що робити з отриманою
інформацією. Інша крайність --- повне ігнорування математичних тонкощів,
але тоді ні учні, ні вчителі не знатимуть, як діяти. Я намагався знайти
розумний компроміс між цими крайнощами.

Ця книга виникла з курсу лекцій, які я читав протягом останніх п'яти
років. Розташування матеріалу не є загальноприйнятим --- вони майже
незалежні один від одного

Друг лекцій.

Основна увага приділяється двом найважливішим аналітичним методам:
розділенню змінних та інтегральним перетворенням. Також представлені
деякі нестандартні розділи, такі як метод Монте-Карло, варіаційне
числення, теорія керування, теорія потенціалів, інтегральні рівняння. Це
пов'язано з тим, що більшості студентів доведеться стикатися з цими
питаннями під час навчання, і якщо вони не ознайомляться з ними зараз,
навряд чи коли-небудь систематично їх вивчатимуть.

Цю книгу можна використовувати як навчальний посібник для одно- або
двосеместрового курсу як для початківців, так і для старших студентів.
Читачі повинні знати лише диференціальне та інтегральне числення та
звичайні диференціальні рівняння. Матеріал більшості лекцій подається в
одному-двох уроках. Типовий односеместровий курс може включати лекції
1-13, 15-17, 19-20, 22-23, 25-27, 30-32, 37-39. Усі 47 лекцій легко
читати протягом двох семестрів. Водночас часу буде достатньо для
вирішення проблем.

Автор вдячний Вайлі за пропозицію написати книгу, а також рецензентам
--- професорам Крісу Рорресу та М. Куршиду Алі, які зробили низку
корисних коментарів до тексту книги. Будь-які пропозиції щодо подальшого
покращення цього підручника будуть вдячні як від учнів, так і від
викладачів. Я також вдячний Дороті, Сюзан, Александер і Дезі Фарлоу.

Стенлі Дж. Фарлоу

\section{ВСТУП ДО ТЕОРІЇ ДИФЕРЕНЦІАЛЬНИХ РІВНЯНЬ У ПОХІДНИХ
ПОХІДНИХ}\label{ux432ux441ux442ux443ux43f-ux434ux43e-ux442ux435ux43eux440ux456ux457-ux434ux438ux444ux435ux440ux435ux43dux446ux456ux430ux43bux44cux43dux438ux445-ux440ux456ux432ux43dux44fux43dux44c-ux443-ux43fux43eux445ux456ux434ux43dux438ux445-ux43fux43eux445ux456ux434ux43dux438ux445}

МЕТА ЛЕКЦІЇ: Показати, що таке диференціальні рівняння в похідних
похідних, де і як вони виникають і як їх розв'язують. Коротко розгляньте
методи класифікації рівнянь і надайте список базових понять, які будуть
детально вивчені на наступних лекціях.

Більшість фізичних явищ у таких сферах, як гідродинаміка, електрика та
магнетизм, механіка, оптика та теплопередача, загалом можна описати за
допомогою рівнянь у похідних похідних (PNP). Більшість рівнянь у
математичній фізиці є диференціальними рівняннями з похідними похідними.
Правда, з деякими спрощеними припущеннями ці рівняння зводяться до
звичайних диференціальних рівнянь, але повний опис таких систем неминуче
призводить до використання диференціальних рівнянь у похідних похідних.

\section{Що таке часткові диференціальні
рівняння?}\label{ux449ux43e-ux442ux430ux43aux435-ux447ux430ux441ux442ux43aux43eux432ux456-ux434ux438ux444ux435ux440ux435ux43dux446ux456ux430ux43bux44cux43dux456-ux440ux456ux432ux43dux44fux43dux43dux44f}

Рівняння з частковими похідними --- це рівняння, яке містить
диференціали у похідних похідних. На відміну від звичайних
диференціальних рівнянь (TAC), де невідома функція залежить лише від
однієї змінної, у рівняннях часткових похідних вона залежить від кількох
змінних (наприклад, температура \texttt{u(x,\ t)} залежить від
координати x і часу t).

Тепер наведемо деякі з найважливіших диференціальних рівнянь у похідних
похідних. Для спрощення запису ми використаємо наступні нотації

\begin{verbatim}
u_t \equiv \frac{\partial u}{\partial t}
\end{verbatim}

\texttt{,\ \textquotesingle{}u\_x\ \textbackslash{}equiv\ \textbackslash{}frac\{\textbackslash{}partial\ u\}\{\textbackslash{}partial\ x\}\textquotesingle{}\ ,\ \textquotesingle{}u\_\{xx\}\ \textbackslash{}equiv\ \textbackslash{}frac\{\textbackslash{}partial\^{}2\ u\}\{\textbackslash{}partial\ x\^{}2\}\textquotesingle{}\ ,\ ...}

\section{Деякі диференціальні рівняння в похідних
похідних}\label{ux434ux435ux44fux43aux456-ux434ux438ux444ux435ux440ux435ux43dux446ux456ux430ux43bux44cux43dux456-ux440ux456ux432ux43dux44fux43dux43dux44f-ux432-ux43fux43eux445ux456ux434ux43dux438ux445-ux43fux43eux445ux456ux434ux43dux438ux445}

\begin{verbatim}
u_t = u_{xx}
\end{verbatim}

(одновимірне рівняння теплопровідності),
\texttt{u\_t\ =\ u\_\{xx\}\ +\ u\_\{yy\}} (двовимірне рівняння
теплопровідності),
\texttt{u\_\{rr\}\ +\ \textbackslash{}frac\{1\}\{r\}u\_r\ +\ \textbackslash{}frac\{1\}\{r\^{}2\}u\_\{00\}\ =\ 0}
(рівняння Лапласа за полярними координатами),
\texttt{u\_\{tt\}\ =\ u\_\{xx\}\ +\ u\_\{yy\}\ +\ u\_\{zz\}} (тривимірне
хвильове рівняння),
\texttt{u\_\{tt\}\ =\ u\_\{xx\}\ +\ zu\_t\ +\ \textbackslash{}beta\ u}
(телеграфне рівняння).

\section{Примітки щодо наведених
прикладів}\label{ux43fux440ux438ux43cux456ux442ux43aux438-ux449ux43eux434ux43e-ux43dux430ux432ux435ux434ux435ux43dux438ux445-ux43fux440ux438ux43aux43bux430ux434ux456ux432}

У всіх наведених вище прикладах функція невідома і залежить від більше
ніж однієї змінної. Змінна і (яку ми диференціюємо) називається залежною
змінною. Змінні, за якими відбувається диференціювання, називаються
незалежними змінними. Наприклад, у рівнянні.

\begin{verbatim}
u_t = u_{xx}
\end{verbatim}

Залежна змінна \texttt{u(x,\ t)} є функцією двох незалежних змінних x і
t; та в рівнянні

\begin{verbatim}
u_t = u_{rr} + \frac{1}{r} u_r + \frac{1}{r^2} u_{00}
\end{verbatim}

змінна \texttt{u(r,\ \textbackslash{}theta,\ t)} залежить від
\texttt{r,\ \textbackslash{}theta,\ t} .

\section{Чому необхідно вивчати диференціальні рівняння у похідних
похідних?}\label{ux447ux43eux43cux443-ux43dux435ux43eux431ux445ux456ux434ux43dux43e-ux432ux438ux432ux447ux430ux442ux438-ux434ux438ux444ux435ux440ux435ux43dux446ux456ux430ux43bux44cux43dux456-ux440ux456ux432ux43dux44fux43dux43dux44f-ux443-ux43fux43eux445ux456ux434ux43dux438ux445-ux43fux43eux445ux456ux434ux43dux438ux445}

Більшість фізичних законів природи можна сформулювати через рівняння з
похідними похідними. Прикладами є рівняння Максвелла, закон
теплопередачі Ньютона, рівняння Нав'є-Стокса, рівняння руху Ньютона та
рівняння Шредінгера в кватичній механіці. У всіх цих рівняннях фізичні
явища описуються мовою просторових і часових похідних. Похідні
з'являються в рівняннях, оскільки вони описують найважливіші фізичні
величини (такі як швидкість, прискорення, сила, тертя, потік, струм
тощо). Таким чином, виникають диференціальні рівняння в частинних
похідних, які містять невідому функцію, яку потрібно визначити.

Мета цієї книги --- показати читачеві:

\begin{enumerate}
\def\labelenumi{\arabic{enumi}.}
\item
  Як сформулювати фізичну задачу у формі диференціального рівняння у
  похідних похідних (побудова математичної моделі).
\item
  Як розв'язати диференціальне рівняння в похідних (з урахуванням
  початкових і граничних умов).
\end{enumerate}

Перш ніж перейти до побудови математичних моделей, давайте коротко
розглянемо методи розв'язання рівнянь у похідних похідних.

\section{Як розв'язати диференціальні рівняння в похідних
похідних?}\label{ux44fux43a-ux440ux43eux437ux432ux44fux437ux430ux442ux438-ux434ux438ux444ux435ux440ux435ux43dux446ux456ux430ux43bux44cux43dux456-ux440ux456ux432ux43dux44fux43dux43dux44f-ux432-ux43fux43eux445ux456ux434ux43dux438ux445-ux43fux43eux445ux456ux434ux43dux438ux445}

Гарне питання. Виявляється, існує цілий арсенал методів, придатних для
практичного використання. Найважливішими є ті, у яких диференціальні
рівняння з похідними похідними зводяться до звичайних диференціальних
рівнянь. Перелічимо десять методів розв'язання диференціальних рівнянь у
похідних похідних:

\begin{enumerate}
\def\labelenumi{\arabic{enumi}.}
\tightlist
\item
  Метод розділення змінних. Рівняння з факторами, отриманими з n
  незалежними змінними, зводиться до n звичайних-
\end{enumerate}

диференціальнірівняння 1).

\begin{enumerate}
\def\labelenumi{\arabic{enumi}.}
\setcounter{enumi}{1}
\item
  Метод інтегральних перетворень. Рівняння в похідних з n незалежними
  змінними зводиться до рівняння з похідними у (n-1) незалежних змінних;
  Відповідно, рівняння з частковими похідними з двома незалежними
  змінними можна звести до звичайного диференціального рівняння
  \texttt{\^{}2} ).
\item
  Метод перетворення координат. Початкове диференціальне рівняння в
  похідних похідних зводиться до звичайного рівняння з частковими
  похідними або до іншого, простішого рівняння у похідних похідних
  шляхом відповідного перетворення координат (наприклад, обертання осей
  координат тощо).
\item
  Перетворення залежної змінної. Початкове диференціальне рівняння в
  похідних похідних перетворюється у таке рівняння з похідними похідними
  для іншої невідомої функції, яка
\end{enumerate}

легше розв'язати, ніж початкові 3).

\begin{enumerate}
\def\labelenumi{\arabic{enumi}.}
\setcounter{enumi}{4}
\tightlist
\item
  Чисельні методи. Початкове диференціальне рівняння в похідних похідних
  зводиться до системи різницей, яку розв'язують методом ітерацій на
  комп'ютері. У багатьох випадках це єдиний спосіб розв'язати
  диференціальне рівняння з похідними похідними. Окрім різницьких
  методів для розв'язання рівнянь з частковими про-
\end{enumerate}

\begin{enumerate}
\def\labelenumi{\alph{enumi})}
\setcounter{enumi}{4}
\tightlist
\item
  Ефективніше використовувати загальні перетворення незалежних і
  залежних змінних, включаючи похідні. Ед.
\end{enumerate}

4) Це стосується повного розділення змінних. При частковому розділенні
змінних одне диференціальне рівняння в похідних похідних зводиться до
кількох рівнянь у часткових похідних з меншою кількістю незалежних
змінних. --- Ед. Ед.

2) Це відбувається у випадку одномірного інтегрального перетворення. У
випадку k-вимірного інтегрального перетворення LCP з n незалежними
змінними зводиться до рівняння часткових з n---k незалежними змінними.
--- Ед. Ед.

Існують також інші чисельні методи, зокрема ті, що базуються на
наближенні розв'язків за поліноміальними поверхнями (наближення за
сплайнами).

\begin{enumerate}
\def\labelenumi{\arabic{enumi}.}
\setcounter{enumi}{5}
\tightlist
\item
  Метод теорії збурень. Початкова нелінійна задача зводиться до
  послідовності лінійних задач, апрокенн-
\end{enumerate}

нелінійна задача 1).

\begin{itemize}
\tightlist
\item
  \begin{enumerate}
  \def\labelenumi{\arabic{enumi}.}
  \setcounter{enumi}{6}
  \tightlist
  \item
    Метод функцій Гріна. Початкові та граничні умови замінюються
    системою найпростіших джерел, і задача розв'язується для кожного
    найпростішого джерела. Повне розв'язання початкової задачі
    отримується шляхом підсумовування розв'язків елементарних джерел.
  \end{enumerate}
\item
  \begin{enumerate}
  \def\labelenumi{\arabic{enumi}.}
  \setcounter{enumi}{7}
  \tightlist
  \item
    Метод інтегральних рівнянь. Рівняння з частковими похідними
    зводяться до інтегрального рівняння (рівняння, в якому невідома
    функція знаходиться під знаком інтеграла). Існує багато різних
    методів розв'язання інтегральних рівнянь.
  \end{enumerate}
\item
  \begin{enumerate}
  \def\labelenumi{\arabic{enumi}.}
  \setcounter{enumi}{8}
  \tightlist
  \item
    Варіаційні методи. Замість рівняння з частковими похідними
    диференціальними розв'язують задачу мінімізації. Виявляється, що
    функція, яка дає мінімум певному виразу (наприклад, загальній
    енергії системи), одночасно є розв'язком початкового рівняння
    похідниху похідних 2.
  \end{enumerate}
\item
  \begin{enumerate}
  \def\labelenumi{\arabic{enumi}.}
  \setcounter{enumi}{9}
  \tightlist
  \item
    Метод розкладу за власними функціями. Розв'язок рівняння часткових
    похідних шукається як послідовність власних функцій. Ці власні
    функції знаходяться як розв'язки так званої задачі власних значень,
    що відповідає початковій задачі для рівняння часткових похідних.
  \end{enumerate}
\end{itemize}

\section{Типи рівнянь у похідних
похідних}\label{ux442ux438ux43fux438-ux440ux456ux432ux43dux44fux43dux44c-ux443-ux43fux43eux445ux456ux434ux43dux438ux445-ux43fux43eux445ux456ux434ux43dux438ux445}

Рівняння в похідних похідних можна класифікувати за багатьма
характеристиками. Класифікація рівнянь важлива, оскільки, як виявилося,
кожен клас має свою загальну теорію та методи розв'язання рівнянь.

Тут ми наведемо шість основних методів класифікації

рівняння.

\begin{enumerate}
\def\labelenumi{\arabic{enumi}.}
\tightlist
\item
  Порядок рівняння. Порядок рівняння --- це найвищий порядок часткових
  похідних, включених до рівняння. Наприклад,
\end{enumerate}

Примітка. Ед.

1) Точніше, це метод лінеаризації нелінійних задач. Загальна теорія
збурень містить лінійні та нелінійні теорії збурень, кожна з яких, у
свою чергу, є регулярною та особливою. --- Ед. Ед.

Ці методи застосовні до спеціальних задач, що зумовлено існуванням
відповідних виразів --- функціоналів для цих задач.

Застосовувався раніше

\begin{verbatim}
u_t = u_{xx}
\end{verbatim}

(рівняння другого порядку), \texttt{u\_t\ =\ u\_x} (рівняння першого
порядку), \texttt{u\_1\ =\ uu\_\{xxx\}\ +\ \textbackslash{}sin\ x}
(рівняння трьох порядків).

\begin{enumerate}
\def\labelenumi{\arabic{enumi}.}
\setcounter{enumi}{1}
\tightlist
\item
  Кількість змінних. Кількість змінних --- це кількість незалежних
  змінних. Наприклад,
\end{enumerate}

\begin{verbatim}
u_t = u_{xx}
\end{verbatim}

\texttt{(рівняння\ з\ двома\ змінними\ \textquotesingle{}x\textquotesingle{}\ і\ \textquotesingle{}t\textquotesingle{}\ ),\ \textquotesingle{}u\_t\ =\ u\_\{rr\}\ +\ \textbackslash{}frac\{1\}\{r\}\ u\_r\ +\ \textbackslash{}frac\{1\}\{r\^{}2\}\ u\_\{\textbackslash{}theta\textbackslash{}theta\}\textquotesingle{}\ (рівняння\ з\ трьома\ змінними\ \textquotesingle{}r\textquotesingle{}\ ,}

\begin{enumerate}
\def\labelenumi{\arabic{enumi}.}
\setcounter{enumi}{2}
\tightlist
\item
  Лінійність. Рівняння з частковими похідними можуть бути лінійними та
  нелінійними. У лінійних рівняннях залежна змінна та всі її часткові
  похідні включаються лінійно, зокрема, вони не множаться одна на одну,
  не множаться в квадраті тощо). Точніше, лінійне рівняння другого
  порядку з двома незалежними змінними є рівнянням вигляду
\end{enumerate}

\begin{verbatim}
(1.1) Au_{xx} + Bu_{xy} + Cu_{yy} + Du_x + Eu_y + Fu = G
\end{verbatim}

де A, B, C, D, E, F і G --- це константи або задані функції
пояснювальних змінних x і y.

\begin{verbatim}
u_{tt} = e^{-t}u_{xx} + \sin t
\end{verbatim}

(лінійне рівняння), \texttt{uu\_\{xx\}\ +\ u\_t\ =\ 0} (нелінійне
рівняння), \texttt{u\_\{xx\}\ +\ yu\_\{yy\}\ =\ 0} (лінійне рівняння),
\texttt{xu\_x\ +\ yu\_y\ +\ u\^{}2\ =\ 0} (нелінійне рівняння).

\begin{itemize}
\tightlist
\item
  \begin{enumerate}
  \def\labelenumi{\arabic{enumi}.}
  \setcounter{enumi}{3}
  \tightlist
  \item
    Однорідність. Рівняння (1.1) називають однорідним, якщо права
    частина \texttt{G(x,\ y)} ідентична нулю для всіх x і y. Якщо
    \texttt{G(x,\ y)} не ідентичне нулю, то рівняння називають
    неоднорідним.
  \end{enumerate}
\item
  \begin{enumerate}
  \def\labelenumi{\arabic{enumi}.}
  \setcounter{enumi}{4}
  \tightlist
  \item
    Типи шансів. Якщо коефіцієнти A, B, C, D, E і F рівняння (1.1) є
    сталими, то це рівняння називається рівнянням зі сталими
    коефіцієнтами (інакше рівнянням зі змінними коефіцієнтами).
  \end{enumerate}
\item
  \begin{enumerate}
  \def\labelenumi{\arabic{enumi}.}
  \setcounter{enumi}{5}
  \tightlist
  \item
    Три основні типи лінійних рівнянь. Усі лінійні рівняння з частковими
    похідними другого порядку виду (1.1) належать до одного з трьох
    типів: а) параболічних, б) гіперболічних, в) еліптичних.
  \end{enumerate}
\end{itemize}

Параболічний тип. Рівняння параболічного типу описують процеси
теплопровідності та дифузії і визначаються умовою
\texttt{B\^{}2\ -\ 4AC\ =\ 0} .

Гіперболічний тип. Рівняння гіперболічного типу описів

осцилюючі системи та хвильові рухи визначаються умовою
\texttt{B\^{}2-4AC\textgreater{}0} .

Еліпічний тип. Рівняння еліптичного типу описують усталені процеси і
визначаються умовним \texttt{B\^{}2-4AC\textless{}0} . Приклади

\begin{enumerate}
\def\labelenumi{\alph{enumi})}
\tightlist
\item
\end{enumerate}

\begin{verbatim}
u_i = u_{xx}
\end{verbatim}

, \texttt{B\^{}2\ -\ 4AC\ =\ 0} (параболічна), (гіперболічна),
(гіперболічна), (гіперболічна \texttt{y\ \textgreater{}\ 0} ),
(гіперболічна), (еліптична), (еліптична), параболічна на
\texttt{y\ =\ 0} , гіперболічна на \texttt{y\ \textless{}\ 0} .

(У випадку змінних коефіцієнтів тип рівняння може змінюватися від точки
до точки.)

\section{ЗАУВАЖЕННЯ}\label{ux437ux430ux443ux432ux430ux436ux435ux43dux43dux44f}

\begin{itemize}
\tightlist
\item
  \begin{enumerate}
  \def\labelenumi{\arabic{enumi}.}
  \tightlist
  \item
    Загалом, значення \texttt{B\^{}2-4AC} є функцією незалежних змінних.
    Відповідно, базовий тип рівняння може змінюватися в межах визначення
    рівняння, хоча це не є обов'язковим.
  \end{enumerate}
\item
  \begin{enumerate}
  \def\labelenumi{\arabic{enumi}.}
  \setcounter{enumi}{1}
  \tightlist
  \item
    У рівнянні (1.1) незалежними змінними є x і y. У багатьох задачах
    однією з двох незалежних змінних є час, і рівняння (1.1) можна
    записати через x і t.
  \end{enumerate}
\end{itemize}

\pandocbounded{\includegraphics[keepaspectratio]{images/_page_10_Figure_9.jpeg}}

PHC 1.1. Діаграма класифікації диференціальних рівнянь у похідних
похідних.

\begin{enumerate}
\def\labelenumi{\arabic{enumi}.}
\setcounter{enumi}{2}
\tightlist
\item
  Рис. 1.1. Представлено діаграму класифікації для диференціальних
  рівнянь у похідних похідних.
\end{enumerate}

\section{ЗАВДАННЯ}\label{ux437ux430ux432ux434ux430ux43dux43dux44f}

\begin{itemize}
\tightlist
\item
  \begin{enumerate}
  \def\labelenumi{\arabic{enumi}.}
  \tightlist
  \item
    Класифікуємо наступні рівняння відповідно до всіх характеристик,
    зазначених у днаграмі на рис. 1.1:
  \end{enumerate}
\item
  \begin{enumerate}
  \def\labelenumi{\alph{enumi})}
  \tightlist
  \item
    \texttt{u\_t\ =\ u\_\{xx\}\ +\ 2u\_x\ +\ u} ,
  \end{enumerate}
\end{itemize}

\texttt{6)\ \textquotesingle{}u\_t\ =\ u\_\{xx\}\ +\ e\^{}\{-t\}\textquotesingle{}\ ,\ B)\ \textquotesingle{}u\_\{xx\}\ +\ 3u\_\{xy\}\ +\ u\_\{yy\}\ =\ \textbackslash{}sin\ x\textquotesingle{}\ ,\ \textquotesingle{}u\_\{tt\}\ =\ uu\_\{xxx\}\ +\ e\^{}\{-t\}\textquotesingle{}\ .}

\texttt{-\ 2.\ Скільки\ розв\textquotesingle{}язків\ \textquotesingle{}u\_1\ =\ u\_\{xx\}\textquotesingle{}\ рівняння\ існує?\ Спробуйте\ знайти\ рішення,\ як\ \textquotesingle{}u(x,\ t)\ =\ e\^{}\{a\textbackslash{}lambda\ +\ bt\}\textquotesingle{}\ .}
- 3. Якщо функції \texttt{u\_1(x,\ t)} і \texttt{u\_2(x,\ t)}
задовольняють рівняння (1.1), чи задовольняє сума цих функцій його?
Доведи це. - 4. Ймовірно, найпростіше диференціальне рівняння в похідних
---

\begin{verbatim}
\frac{\partial u(x, y)}{\partial x} = 0.
\end{verbatim}

Чи можете ви розв'язати це рівняння? (Знайдіть усі функції, що
задовольняють це рівняння.)

\begin{enumerate}
\def\labelenumi{\arabic{enumi}.}
\setcounter{enumi}{4}
\tightlist
\item
  Що можна сказати про рівняння
\end{enumerate}

\begin{verbatim}
\frac{\partial^2 u(x, y)}{\partial x \partial y} = 0
\end{verbatim}

?

Чи можете ви знайти всі розв'язки цього рівняння? Скільки їх? Порівняйте
з кількістю розв'язків звичайного диференціального рівняння

\begin{verbatim}
\frac{d^2y}{dx^2} = 0.
\end{verbatim}

\section{ЗАДАЧІ ТИПУ ДИФУЗІЇ (ПАРАБОЛІЧНІ
РІВНЯННЯ)}\label{ux437ux430ux434ux430ux447ux456-ux442ux438ux43fux443-ux434ux438ux444ux443ux437ux456ux457-ux43fux430ux440ux430ux431ux43eux43bux456ux447ux43dux456-ux440ux456ux432ux43dux44fux43dux43dux44f}

МЕТА ЛЕКЦІЇ: Показати, як параболічні рівняння використовуються для
розв'язання задач теплопровідності та дифузії. Наведено кілька прикладів
параболічних рівнянь, а також пояснюється фізичне значення різних членів
\texttt{(u\_1,\ u\_2,\ u\_3,\ u\_4,\ u\_6)} , включених у ці рівняння.

На конкретному прикладі розглядається основна ідея зміщеної задачі для
параболічного рівняння. Одна з головних цілей цієї лекції --- надати
читачеві інтуїтивне уявлення

Розуміння параболічних проблем.

Почнемо з розгляньмо просту фізичну задачу і те, як її можна описати
математичною моделлю з рівнянням у похідних похідних.

Потім ми ускладнимо задачу і покажемо, як нові фізичні умови призводять
до появи нових рівнянь у похідних похідних. У цій лекції диференціальні
рівняння не виводяться і не розв'язуються. Про це ми поговоримо пізніше.

\section{Простий експеримент з
теплопровідністю}\label{ux43fux440ux43eux441ux442ux438ux439-ux435ux43aux441ux43fux435ux440ux438ux43cux435ux43dux442-ux437-ux442ux435ux43fux43bux43eux43fux440ux43eux432ux456ux434ux43dux456ux441ux442ux44e}

Припустимо, що ми провели простий експеримент, який розбили на такі
кроки:

КРОК 1. Візьміть достатньо довгий (скажімо,
\texttt{L=2\textbackslash{},\textbackslash{}mathrm\{m\}} ) стрижень
(наприклад, мідний) діаметром
\texttt{2\textbackslash{},\textbackslash{}mathrm\{cm\}} , бічна поверхня
(але не кінці) якої термічно ізольована. Замість стрижня можна
використовувати мідну трубку, покриту теплоізоляцією зовні та зсередини.

Іншими словами, тепло може проходити лише через кінці стрижня і не може
проходити через його бічну поверхню.

КРОК 2. Тепер помістимо цей стрижень у пристрій із фіксованою
температурою \texttt{T\_\{\textbackslash{}mathfrak\{o\}\}} (градуси
Цельсія) на достатньо довгий час, щоб температура всередині стрижня
стала

Так само, як і температура пристрою. Для простоти припустимо, що
температура пристрою \texttt{T\_0\ =\ 10\^{}\{\textbackslash{}circ\}} C.

КРОК 3. Давайте витягнемо стрижень з нашого пристрою у певний момент
часу, що зручно взяти як нуль t = 0, і приєднаємо до нього дві термопари
з кінців. Завдання цих елементів --- підтримувати фіксовані температури
\texttt{T\_1} і \texttt{T\_2} на кінцях (наприклад,
\texttt{T\_1=0\^{}\{\textbackslash{}circ\}} C і
\texttt{T\_2=50\^{}\{\textbackslash{}circ\}} C). Іншими словами, два
термостати постійно контролюють температуру на кінцях стрижня, і якщо
температура відхиляється від встановлених значень \texttt{T\_1} і
\texttt{T\_2} , активуються потужні нагрівальні або охолоджувальні
елементи (пристрої) для відповідного регулювання температури. Наш
експеримент ілюструє RPS, 2 1,

\pandocbounded{\includegraphics[keepaspectratio]{images/_page_13_Figure_4.jpeg}}

Рис. 2.1. Схема експерименту: a --- це термопара, яка підтримує
встановлену температуру \texttt{T\_1} на лівому кінці стрижня;
\texttt{\textbackslash{}delta} --- це термопара, яка підтримує
встановлену температуру \texttt{T\_2} на правому кінці стрижня.

КРОК 4. Тепер ми будемо слідкувати за температурним профілем у стрижні
за допомогою пристрою візуалізації. Чому ми хочемо провести такий
експеримент? Це інше питання, про яке ми поговоримо пізніше.

Це кінець опису нашого експерименту.

Головна мета цієї лекції --- показати, як цю фізичну проблему та її
модифікації можна змоделювати мовою параболічних рівнянь.

\section{Математична модель
теплопровідності}\label{ux43cux430ux442ux435ux43cux430ux442ux438ux447ux43dux430-ux43cux43eux434ux435ux43bux44c-ux442ux435ux43fux43bux43eux43fux440ux43eux432ux456ux434ux43dux43eux441ux442ux456}

Опис нашої фізичної проблеми вимагає трьох типів відносин.

\begin{itemize}
\tightlist
\item
  \begin{enumerate}
  \def\labelenumi{\arabic{enumi}.}
  \tightlist
  \item
    Рівняння з частковими похідними, які описують фізичне явище
    теплопровідності.
  \end{enumerate}
\item
  \begin{enumerate}
  \def\labelenumi{\arabic{enumi}.}
  \setcounter{enumi}{1}
  \tightlist
  \item
    Граничні умови, що описують процеси теплопередачі на межах.
  \end{enumerate}
\item
  \begin{enumerate}
  \def\labelenumi{\arabic{enumi}.}
  \setcounter{enumi}{2}
  \tightlist
  \item
    Початкові твердження, які описують стан системи на початку процесу.
  \end{enumerate}
\end{itemize}

\section{Рівняння
теплопровідності}\label{ux440ux456ux432ux43dux44fux43dux43dux44f-ux442ux435ux43fux43bux43eux43fux440ux43eux432ux456ux434ux43dux43eux441ux442ux456}

Базове одновимірне рівняння теплопровідності записується у вилці

(2.1)

\begin{verbatim}
u_t = \alpha^2 u_{xx}, \quad 0 < x < L, \quad 0 < t < \infty.
\end{verbatim}

Це рівняння пов'язує значення \texttt{u\_t} --- швидкість зміни
температури з часом (виміряна в градусах/с)

та величину \texttt{u\_\{xx\}} --- це увігнутий профілю температури
\texttt{u(x,\ t)} (який є мірою різниці між температурою в заданій точці
та температурою в сусідніх точках).

\pandocbounded{\includegraphics[keepaspectratio]{images/_page_14_Figure_14.jpeg}}

Рис. 2.2. Стрілки показують зміну температури відповідно до рівняння
\texttt{u\_t\ =\ \textbackslash{}alpha\^{}2\ u\_\{xx\}} .

\texttt{Це\ рівняння\ буде\ виведено\ з\ базового\ рівняння\ збереження\ тепла\ в\ наступних\ лекціях,\ але\ наразі\ ми\ розглянемо\ його\ заданим.\ Це\ рівняння\ каже,\ що\ температура\ \textquotesingle{}u\textbackslash{}left(x,\textbackslash{}\ t\textbackslash{}right)\textquotesingle{}\ (у\ певній\ точці\ t\ у\ певній\ точці\ x\ стрижня)\ зростає\ \textquotesingle{}(u\_t\textgreater{}0)\textquotesingle{}\ або\ зменшується\ \textquotesingle{}(u\_t\textless{}0)\textquotesingle{}\ залежно\ від\ того,\ чи\ є\ друга\ похідна\ \textquotesingle{}u\_\{xx\}\textquotesingle{}\ позитивною\ чи\ негативною.\ Фіг.\ Рисунок\ 2.2\ ілюструє\ зміну\ температури\ в\ різних\ точках\ на\ елементі.}

Давайте подивимось, як значення \texttt{u\_\{xx\}} можна інтерпретувати
мовою теплопровідності. Припустимо, ми наблизимо значення
\texttt{u\_\{xx\}} фінальних постів

\begin{verbatim}
u_{xx}(x, t) \cong \frac{1}{\Delta x^2} [u(x + \Delta x, t) - 2u(x, t) + u(x - \Delta x, t)].
\end{verbatim}

Це співвідношення можна переписати у вигляді

\begin{verbatim}
u_{\lambda x}(x, t) \simeq \frac{2}{\Delta x^2} \left[ \frac{u(x + \Delta x, t) + u(x - \Delta x, t)}{2} - u(x, t) \right].
\end{verbatim}

Тепер ми можемо дати наступну інтерпретацію значення \texttt{u\_\{xx\}}
.

\begin{enumerate}
\def\labelenumi{\arabic{enumi}.}
\item
  Якщо температура \texttt{u(x,\ t)} менша за середнє значення
  температури у двох сусідніх точках, то
  \texttt{u\_\{xx\}\ \textgreater{}\ 0} (тут загальний тепловий потік по
  осі x додатний).
\item
  Якщо температура \texttt{u(x,\ t)} дорівнює середньому значенню
  температур у двох сусідніх точках, то \texttt{u\_\{xx\}\ =\ 0} (тут
  загальний тепловий потік по осі x дорівнює нулю).
\end{enumerate}

\texttt{3.\ Якщо\ температура\ \textquotesingle{}u\textbackslash{}left(x,\textbackslash{}\ t\textbackslash{}right)\textquotesingle{}\ вища\ за\ середню\ температуру\ у\ двох\ сусідніх\ точках,\ то\ \textquotesingle{}u\_\{xx\}\textless{}0\textquotesingle{}\ (тут\ —\ повний\ потік}

тепла вздовж осі x є від'ємною).

\texttt{Усе\ це\ ілюструється\ на\ рис.\ 2.2.\ Іншими\ словами,\ якщо\ температура\ в\ точці\ x\ більша\ за\ середню\ температуру\ у\ двох\ сусідніх\ точках\ \textquotesingle{}x\ -\ \textbackslash{}Delta\ x\textquotesingle{}\ і\ \textquotesingle{}x\ +\ \textbackslash{}Delta\ x\textquotesingle{}\ ,\ то\ температура\ в\ точці\ x\ зменшиться.\ Відповідно,\ точна\ швидкість\ зниження\ температури,\ а\ саме\ величина\ \textquotesingle{}u\_t\textquotesingle{}\ ,\ пропорційна\ цій\ різниці.\ Коефіцієнт\ пропорційності\ \textquotesingle{}\textbackslash{}alpha\^{}2\textquotesingle{}\ визначається\ властивостями\ матеріалу.\ Ми\ розглянемо\ цю\ константу\ детальніше\ у\ майбутніх\ лекціях.}

\section{Граничні умови
(ГІ)}\label{ux433ux440ux430ux43dux438ux447ux43dux456-ux443ux43cux43eux432ux438-ux433ux456}

Усі фізичні проблеми характеризуються наявністю певних меж. Отже, щоб
рухатися вперед, ми повинні включити межі в математичну модель, щоб
адекватно описати фізичну проблему. У нашому експерименті граничні умови
(ГІ) дуже прості. Оскільки температура на кінцях x=0 і x=L завжди
залишається сталою і дорівнює \texttt{T\_1} та \texttt{T\_2} , можливо
записати

(2.2)

\begin{verbatim}
\left\{ \begin{array}{l} u\left(0,\ t\right) = T_{t}, \\ u\left(L,\ t\right) = T_{2}, \end{array} \right. \quad 0 < t < \infty,
\end{verbatim}

\section{Початкові умови
(NU)}\label{ux43fux43eux447ux430ux442ux43aux43eux432ux456-ux443ux43cux43eux432ux438-nu}

Усі фізичні процеси мають починатися у певний момент часу (зазвичай
приймається як нуль t=0), тому ми повинні встановити фізичні умови саме
в цю точку. Відтоді, як ми почали контролювати температуру з моменту,
коли паличка мала постійну температуру
\texttt{T\_\{\textbackslash{}mathrm\{\ \}o\}} , ми

(2.3)

\begin{verbatim}
u(x, 0) = T_0, \quad 0 \le x \le L
\end{verbatim}

Тепер ми створили математичну модель експерименту. Система співвідношень
(2.1), (2.2) та (2.3) називається задачею зміщення для рівняння
теплопровідності і зазвичай записується як

(2.4)

\begin{verbatim}
\begin{aligned} & (\text{YUII}) \quad u_t = \alpha^2 u_{xx}, & 0 < x < L, \ 0 < t < \infty, \\ & \left( \Gamma \text{Y} \right) \quad \begin{cases} u \ (0, \ t) = T_1, \\ u \ (L, \ t) = T_2, \end{cases} & 0 < t < \infty, \\ & (\text{HY}) \quad u \ (x, \ 0) = T_0, \quad 0 \leqslant x \leqslant L. \end{aligned}
\end{verbatim}

Цікаво, що це зовсім не очевидно: чи існує одна функція
\texttt{u(x,\ t)}, яка задовольняє умови (2.4), і чи описує ця функція
температуру в стрижні? Отже, наше безпосереднє завдання --- знайти цю
функцію \texttt{u(x,\ t)} з (2.4).

Перш ніж завершити цю лекцію, давайте розглянемо деякі варіації основної
проблеми. Почнемо з невеликих змін рівняння теплопровідності.

\section{Деякі рівняння дифузійного
типу}\label{ux434ux435ux44fux43aux456-ux440ux456ux432ux43dux44fux43dux43dux44f-ux434ux438ux444ux443ux437ux456ux439ux43dux43eux433ux43e-ux442ux438ux43fux443}

Теплопередача через бічну поверхню пропорційна різниці температур

Рівняння

\begin{verbatim}
u_t = \alpha^2 u_{xx} - \beta (u - u_0)
\end{verbatim}

Описує теплопровідність елемента, враховуючи не лише дифузію теплових
\texttt{\textbackslash{}alpha\^{}2\ u\_\{xx\}} вздовж елемента, а й
теплопередачу через бічну поверхню елемента. Відтік тепла (
\texttt{\textbackslash{}beta\ \textgreater{}\ 0} ) або його приплив (
\texttt{\textbackslash{}beta\ \textless{}\ 0} ) пропорційний різниці між
температурою u стрижня та температурою навколишнього середовища
\texttt{u\_0} ( \texttt{\textbackslash{}beta} --- це сталий коефіцієнт
пропорційності). Якщо коефіцієнт \texttt{\textbackslash{}beta} великий
порівняно з \texttt{\textbackslash{}alpha\^{}2} , то теплопотік уздовж
стрижня буде малим порівняно з

Протікання через бічну поверхню, і тому тепло буде проходити через неї
відповідно до закону
\texttt{u\_t\ =\ -\textbackslash{}beta\ (u\ -\ u\_a)} .

У хімії, де величина відіграє роль концентрації, рівняння

\begin{verbatim}
u_1 = \alpha^2 u_{xx} - \beta (u - u_0)
\end{verbatim}

означає, що швидкість зміни \texttt{u\_i} кількості речовини залежить як
від дифузії \texttt{\textbackslash{}alpha\^{}2\ u\_\{xx\}} (у напрямку
осі x), так і від появи ( \texttt{\textbackslash{}beta\ \textless{}\ 0}
) або розпаду ( \texttt{\textbackslash{}beta\ \textgreater{}\ 0} )
речовини в хімічній реакції, і пропорційна різниці між двома
концентраціями u і \texttt{u\_\{\textbackslash{}theta\}} .

\section{Внутрішнє джерело
тепла}\label{ux432ux43dux443ux442ux440ux456ux448ux43dux454-ux434ux436ux435ux440ux435ux43bux43e-ux442ux435ux43fux43bux430}

Неоднородне рівняння

\begin{verbatim}
u_t = \alpha^2 u_{xx} + f(x, t)
\end{verbatim}

відповідає випадку, коли всередині елемента є джерело тепла (розташоване
вздовж усього елемента і діє в будь-який момент t). Наприклад, може бути
так, що дріт з електричним струмом проходить через стрижень і опір є
постійним джерелом тепла \texttt{f(x,t)} = K.

\section{Рівняння конвективної
дифузії}\label{ux440ux456ux432ux43dux44fux43dux43dux44f-ux43aux43eux43dux432ux435ux43aux442ux438ux432ux43dux43eux457-ux434ux438ux444ux443ux437ux456ux457}

\texttt{Припустимо,\ що\ домішка\ поширюється\ вздовж\ потоку,\ що\ рухається\ зі\ швидкістю\ v.\ Очевидно,\ що\ концентрація\ \textquotesingle{}u\textbackslash{}left(x,\textbackslash{},t\textbackslash{}right)\textquotesingle{}\ домішок\ змінюється\ залежно\ від\ x\ (позитивний\ напрямок\ осі\ x\ обирається\ вздовж\ потоку)\ і\ часу\ t.\ Швидкість\ зміни\ концентрації\ \textquotesingle{}u\_t\textquotesingle{}\ визначається\ рівнянням\ конвективної\ дифузії}

\begin{verbatim}
u_t = \alpha^2 u_{xx} - v u_x.
\end{verbatim}

Термін \texttt{\textbackslash{}alpha\^{}2\ u\_\{xx\}} описує внесок
дифузії, а \texttt{vu\_x} --- конвективний компонент. Що важливіше,
дифузія або конвекція залежить від співвідношення
\texttt{\textbackslash{}alpha\^{}2} до v. Димові частинки піднімаються
конвективно через рух гарячого повітря і одночасно розсіюють через
вихорові рухи повітря. Окрім цих модифікацій рівняння теплопровідності,
можливо змінювати граничні умови на кінцях стрижня відповідно до
фізичної ситуації. Деякі з цих змін ми обговоримо у третій лекції.

\section{ЗАУВАЖЕННЯ}\label{ux437ux430ux443ux432ux430ux436ux435ux43dux43dux44f-1}

Рівняння теплопровідності
\texttt{u\_t\ =\ \textbackslash{}alpha\^{}2(x)\ u\_\{xx\}} з змінним
коефіцієнтом \texttt{\textbackslash{}alpha(x)} відповідає задачі, в якій
дифузія залежить від x (тобто матеріал є гетерогенним). Наприклад, мідь
і сталь---

\pandocbounded{\includegraphics[keepaspectratio]{images/_page_18_Picture_4.jpeg}}

PHC. 2.3.

Пластини контактують, як показано на рис. 2.3. Припустимо, що
температура \texttt{u(0,t)=0\^{}\{\textbackslash{}circ\}} C
підтримується на лівому боці міді, а температура
\texttt{u(L,t)=20\^{}\{\textbackslash{}circ\}} C --- з правого боку
сталі, тоді рівняння, що описує теплопровідність у цій системі, виглядає
так:

\begin{verbatim}
u_t = \alpha^2(x) \, u_{_{\! A \, X}}, \quad 0 < x < L,
\end{verbatim}

де

\begin{verbatim}
\alpha(x) = \begin{cases} \alpha_1 & \text{(коэффициент диффузии в меди)}, \quad 0 < x < L/2, \\ \alpha_2 & \text{(коэффициент диффузии в стали)}, \quad L/2 < x < L. \end{cases}
\end{verbatim}

\# ЗАВДАННЯ

\begin{enumerate}
\def\labelenumi{\arabic{enumi}.}
\tightlist
\item
  Якщо початкова температура в стрижні
\end{enumerate}

\begin{verbatim}
u(x, 0) = \sin \pi x, \quad 0 \leqslant x \leqslant 1,
\end{verbatim}

а граничні умови такі:

\begin{verbatim}
u(0, t) = 0,
\end{verbatim}

\texttt{u(1,\ t)\ =\ 0,} Тоді як поводитиметься температура в стрижні
\texttt{U(x,\ t)} при T \textgreater{} 0? 2. Припустимо, що в стрижні є
постійне внутрішнє джерело тепла, так що рівняння теплопровідності в
стрижні записується як

\begin{verbatim}
u_t = \alpha^2 u_{xx} + 1, \quad 0 < x < 1.
\end{verbatim}

Припустимо, що ми зафіксували температуру на межі з значеннями
\texttt{u(0,\ t)} = 0 і \texttt{u(1,\ t)} = 1. Якою буде стаціонарна
температура у стрижні? Іншими словами, який c-незалежний розподіл
\texttt{U(x,\ t)} перешкоджує \texttt{u(x,\ t)}?

\begin{enumerate}
\def\labelenumi{\arabic{enumi}.}
\setcounter{enumi}{2}
\tightlist
\item
  Припустимо, що металевий стрижень втрачає тепло через бічну поверхню
  згідно з рівнянням
\end{enumerate}

\begin{verbatim}
u_1 = \alpha^2 u_{xx} - \beta u
\end{verbatim}

, \texttt{0\ \textless{}\ x\ \textless{}\ 1} .

І припустимо, що ми підтримуємо кінці стрижня при температурах
\texttt{U(0,\ T)} = 1, \texttt{U(1,\ T)} = 1. Знайдіть стаціонарний
розподіл температури по стрижню. Побудуйте графік цього розподілу. У
яких моментах відбувається тепловий потік?

\texttt{4.\ Припустимо,\ що\ термоізольований\ металевий\ стрижень\ довжини\ L\ =\ 1\ має\ початкову\ температуру\ \textquotesingle{}\textbackslash{}sin(3\textbackslash{}pi\ x)\textquotesingle{}\ ,\ а\ фіксовану\ температуру\ —\ 0\ і\ \textquotesingle{}10\textbackslash{},\^{}\{\textbackslash{}circ\}\textquotesingle{}\ C\ на\ лівому\ та\ правому\ кінцях.}

\section{Лекція 3}\label{ux43bux435ux43aux446ux456ux44f-3}

\section{ГРАНИЧНІ УМОВИ В ЗАДАЧАХ ДИФУЗІЙНОГО
ТИПУ}\label{ux433ux440ux430ux43dux438ux447ux43dux456-ux443ux43cux43eux432ux438-ux432-ux437ux430ux434ux430ux447ux430ux445-ux434ux438ux444ux443ux437ux456ux439ux43dux43eux433ux43e-ux442ux438ux43fux443}

МЕТА ЛЕКЦІЇ: Показати, як виникають різні граничні умови у задачах
теплопровідності та дифузії, а також ввести важливе поняття теплового
потоку.

Розглядаються три найважливіші типи граничних умов:

\begin{itemize}
\tightlist
\item
  \begin{enumerate}
  \def\labelenumi{\arabic{enumi}.}
  \tightlist
  \item
    u = \texttt{g(t)} (температура задається на межі),
    \texttt{-\ \textquotesingle{}2.\ \textbackslash{}\ \textbackslash{}frac\{\textbackslash{}partial\ u\}\{\textbackslash{}partial\ n\}\ +\ \textbackslash{}lambda\ u\ =\ g\textbackslash{}left(t\textbackslash{}right)\textquotesingle{}\ (встановлення\ температури\ навколишнього\ середовища);}
  \end{enumerate}
\end{itemize}

n --- вектор зовнішньої нормалі до межі),
\texttt{3.\ \textbackslash{}frac\{\textbackslash{}partial\ u\}\{\textbackslash{}partial\ n\}\ =\ g(t)}
(визначений тепловий потік через межу).

У задачах теплопровідності зазвичай використовують три основні типи
граничних умов. Ця лекція наводить приклади того, як різні фізичні умови
експерименту призводять до різних типів граничних умов.

\section{Граничні умови першого виду (температура задається на
межі)}\label{ux433ux440ux430ux43dux438ux447ux43dux456-ux443ux43cux43eux432ux438-ux43fux435ux440ux448ux43eux433ux43e-ux432ux438ux434ux443-ux442ux435ux43cux43fux435ux440ux430ux442ux443ux440ux430-ux437ux430ux434ux430ux454ux442ux44cux441ux44f-ux43dux430-ux43cux435ux436ux456}

Розглянемо тепловий потік в одномірному стрижні, показаний на рис. 3 1,
і припустимо, що температура кінців змінюється відповідно до законів
\texttt{g\_1(t)} і \texttt{g\_2(t)} .

\pandocbounded{\includegraphics[keepaspectratio]{images/_page_20_Picture_4.jpeg}}

Рис. 31. Температура встановлюється на межі.

Як згадувалося в попередній лекції, кожен пристрій, що підтримує
встановлену температуру на кінцях, складається з термостата та
нагрівального елемента для підтримки відповідних теплових умов. Проблеми
з граничними умовами першого виду дуже поширені. У деяких випадках
складність полягає у пошуку контролерів граничних умов, тобто температур
\texttt{g\_1(t)} та \texttt{g\_2(t)} меж, які призводять до певної зміни
температури всередині стрижня. Наприклад, у металургії часто необхідно
обирати граничні умови контролю, щоб температура металу всередині печі
змінювалася з часом, а температурний градієнт був невеликим.

\pandocbounded{\includegraphics[keepaspectratio]{images/_page_20_Picture_7.jpeg}}

Рис. 3.2. Коливання температури на межі.

Граничні умови першого виду також виникають у багатовимірних задачах
теплопровідності. Як приклад можна навести цікаву задачу визначення
температурного поля всередині кругового диска з радіусом \emph{R} (див.
рис. 3.2), коли гранична температура задається у полярних координатах
формулою
\texttt{u(R,\ \textbackslash{}theta,\ t)\ =\ \textbackslash{}cos\ t\ \textbackslash{}sin\ \textbackslash{}theta}
.

Звісно, щоб розв'язати задачу, потрібно знати початкову температуру, але
в цьому випадку початкові умови перестають впливати на розв'язок через
дуже короткий час, і отримана температура всередині диска визначається
лише значеннями температури на межі.

\section{Граничні умови другого виду (задана температура навколишнього
середовища)}\label{ux433ux440ux430ux43dux438ux447ux43dux456-ux443ux43cux43eux432ux438-ux434ux440ux443ux433ux43eux433ux43e-ux432ux438ux434ux443-ux437ux430ux434ux430ux43dux430-ux442ux435ux43cux43fux435ux440ux430ux442ux443ux440ux430-ux43dux430ux432ux43aux43eux43bux438ux448ux43dux44cux43eux433ux43e-ux441ux435ux440ux435ux434ux43eux432ux438ux449ux430}

Припустимо, що ми знову розглядаємо термічно ізольований мідний
стержень, але тепер відмовляємося від заданого температурного режиму на
кінцях стрижня. Натомість давайте наблизимо кінці стрижня до контакту з
двома середовищами. Нехай температура одного з них \texttt{g\_1(t)} , а
іншого \texttt{g\_2(t)} . Іншими словами, припустимо, що лівий кінець
стрижня поміщений у контейнер із рідиною, температура якої змінюється
відповідно до закону \texttt{g\_1(t)} . А правий кінець поміщається в
іншу рідину з температурою \texttt{g\_2(t)} (рис. 3, 3)

\pandocbounded{\includegraphics[keepaspectratio]{images/_page_21_Picture_6.jpeg}}

Рис. 3.3. Конвективна теплопередача через межі, a --- це рідина при
температурі \texttt{g\_1(t)} , \texttt{\textbackslash{}delta} --- рідина
при температурі \texttt{g\_2(t)} .

Задаючи граничні умови такого типу, ми не можемо вважати, що граничні
температури стрижня збігаються з температурами рідини. Але ми знаємо
закон Ньютона: якщо температура одного з кінців стрижня менша за
температуру відповідної рідини, тепло потече в стрижень зі швидкістю,
пропорційною різницям температур. Іншими словами, для одномірного
стрижня з межами x=0 і x=1 закон теплопередачі Ньютона формулюється так:

(3.1)

\begin{verbatim}
\{
\end{verbatim}

\texttt{вихідний\ тепловий\ потік\ на\ \textquotesingle{}x=0\textquotesingle{}\ дорівнює\ \textquotesingle{}h\textbackslash{}left{[}u\textbackslash{}left(0,\textbackslash{},t\textbackslash{}right)-g\_\{1\}\textbackslash{}left(t\textbackslash{}right)\textbackslash{}right{]}\textquotesingle{}\ ,\ вихідний\ тепловий\ потік\ при\ \textquotesingle{}x=L\textquotesingle{}\ дорівнює\ \textquotesingle{}h\textbackslash{}left{[}u\textbackslash{}left(L,\textbackslash{},t\textbackslash{}right)-g\_\{2\}\textbackslash{}left(t\textbackslash{}right)\textbackslash{}right{]}\textquotesingle{}\ ,}

де h --- коефіцієнт теплопередачі, який показує, скільки калорій
проходить через межу за одну секунду при різниці температур у один
градус. Вихідний потік дорівнює кількості кало-

Пройшов крізь палочний спис за одну секунду. Зверніть увагу, що вихідний
тепловий потік буде позитивним на кінці стрижня, де температура стрижня
вища за температуру навколишнього середовища. Рівняння (3.1) разом із
відомим законом теплопровідності Фур'є тепер можна використовувати для
отримання Р и с

\pandocbounded{\includegraphics[keepaspectratio]{images/_page_22_Figure_3.jpeg}}

. 3.4. Ілюстрація закону Фур'є:
\texttt{\textbackslash{}frac\{\textbackslash{}partial\ u\}\{\textbackslash{}partial\ n\}\ \textgreater{}\ 0}
означає, що тепло надходить у цю область;
\texttt{\textbackslash{}frac\{\textbackslash{}partial\ u\}\{\textbackslash{}partial\ n\}\ \textless{}\ 0}
означає, що тепло виходить із цієї зони; n --- напрямок зовнішньої
нормалі;
\texttt{\textbackslash{}frac\{\textbackslash{}partial\ u\}\{\textbackslash{}partial\ n\}\ -}
Зміна температури в напрямку N;
\texttt{\textbackslash{}frac\{\textbackslash{}partial\ u\}\{\textbackslash{}partial\ n\}\ =\ -\textbackslash{}frac\{\textbackslash{}partial\ u\}\{\textbackslash{}partial\ (-n)\}}
є похідною у напрямку внутрішньої нормальної --- це похідна, взята
\texttt{\textbackslash{}mathbf\{e\}} протилежному знаку протилежного
знака пормалі.

граничні умови. Закон Фур'є дає друге представлення (перше міститься у
(3.1)) для вихідного теплового потоку. Прирівнюючи ці два вирази до
потоку, ми отримуємо граничні умови, які шукаємо. Сформулюємо закон
Фур'є (встановлений експериментально):

(3.2) Тепловий потік, що проходить через межу площі, пропорційний
нормальній похідній температури у напрямку внутрішньої норми.

Цей закон стверджує, що якщо температура швидко підвищується у напрямку
зовнішньої температури до межі D (рисунок 3.4), то потік буде текти з
навколишнього середовища до області D.

У нашій одновимірній задачі закон Фур'є має вигляд:

(3.3)

\begin{verbatim}
\begin{cases} \text{вытекающий поток тепла при } x=0 \text{ равен } k \frac{\partial u}{\partial x}, \\ \text{вытекающий поток тепла при } x=' \text{ равен } -k \frac{\partial u}{\partial x}, \end{cases}
\end{verbatim}

де k --- теплопровідність матеріалу, яка є мірою того, наскільки добре
матеріал проводить тепло. Погано провідні матеріали мають коефіцієнт
близький до нуля в одиницях GHS, тоді як мідь та алюміній мають значення
значно вище нуля.

\pandocbounded{\includegraphics[keepaspectratio]{images/_page_23_Figure_5.jpeg}}

Рис. 3.5. Ще одна ілюстрація закону Фур'є.

Закон Фур'є (3.3) фактично описує теплопровідність всюди всередині
стрижня, а не лише на його межі; Наприклад, (див. рис. 3.5).

(3.4) Потік, що проходить через точку

x

\texttt{зліва\ направо\ \textquotesingle{}=-k\textbackslash{}frac\{\textbackslash{}partial\ n\}\{\textbackslash{}partial\ x\}\textquotesingle{}\ Закон\ Фур\textquotesingle{}є\ (3.4)\ стверджує:\ якщо\ \textquotesingle{}u\_x\ \textless{}\ 0\textquotesingle{}\ ,\ то\ в\ точці\ \textquotesingle{}x\_0\textquotesingle{}\ тепло\ тече\ зліва\ направо,\ якщо\ \textquotesingle{}u\_x\ \textgreater{}\ 0\textquotesingle{}\ ,\ то\ в\ точці\ \textquotesingle{}x\_0\textquotesingle{}\ тепло\ тече\ справа\ наліво.\ (Тепло\ завжди\ переходить\ від\ вищих\ температур\ до\ нижчих.)}

Нарешті, якщо використовувати вирази (3.1) і (3.3) для теплового потоку,
то отримуємо бажані граничні умови для задачі, описаної на рис. 3.3, у
чисто математичній формі

\begin{verbatim}
(\Gamma Y) \quad \begin{cases} \frac{\partial u\left(0,\,t\right)}{\partial x} = \frac{h}{k}\left[u\left(0,\,t\right) - g_{1}\left(t\right)\right], \\ \frac{\partial u\left(L,\,t\right)}{\partial x} = -\frac{h}{k}\left[u\left(L,\,t\right) - g_{2}\left(t\right)\right]. \end{cases} \quad 0 < t < \infty,
\end{verbatim}

Дуже часто константа h/k позначається як \texttt{\textbackslash{}lambda}
, а граничні умови теплового потоку на межі записуються у вигляді

(3.5)

\begin{verbatim}
u_{x}(0, t) = \lambda [u(0, t) - g_{1}(t)], u_{x}(L, t) = -\lambda [u(L, t) - g_{2}(t)].
\end{verbatim}

Для випадку вищих розмірностей отримують подібні граничні умови.
Наприклад, якщо межа кругового диска осяяна рухомою рідиною з
температурою \texttt{g(\textbackslash{}theta,t)} , то граничні умови
записуються як

\begin{verbatim}
\frac{\partial u}{\partial r}(R, \theta, t) = -\frac{h}{k} [u(R, \theta, t) - g(\theta, t)].
\end{verbatim}

Тут
\texttt{\textbackslash{}frac\{\textbackslash{}partial\ u\}\{\textbackslash{}partial\ r\}(R,\ \textbackslash{}theta,\ t)}
позначає зовнішню нормальну похідну (у позитивному напрямку r-осі)
значення u, обчислену в точці \texttt{(R,\ \textbackslash{}theta)} межі.
Ми назвемо такі граничні умови лінійними (оскільки вони лінійні в U і
\texttt{u\_r} ), але неоднорідними, оскільки функція
\texttt{g(\textbackslash{}theta,\ t)} знаходиться праворуч.

\section{Граничні умови третього виду (враховуючи потік; зокрема,
враховується випадок термічно ізольованих
меж)}\label{ux433ux440ux430ux43dux438ux447ux43dux456-ux443ux43cux43eux432ux438-ux442ux440ux435ux442ux44cux43eux433ux43e-ux432ux438ux434ux443-ux432ux440ux430ux445ux43eux432ux443ux44eux447ux438-ux43fux43eux442ux456ux43a-ux437ux43eux43aux440ux435ux43cux430-ux432ux440ux430ux445ux43eux432ux443ux454ux442ux44cux441ux44f-ux432ux438ux43fux430ux434ux43eux43a-ux442ux435ux440ux43cux456ux447ux43dux43e-ux456ux437ux43eux43bux44cux43eux432ux430ux43dux438ux445-ux43cux435ux436}

Жоден потік не проходить через термоізольовані межі, тому нормальна
похідна (зовнішня чи внутрішня) має зрівнятися з нулем на межі, оскільки
потік пропорційний нормальній похідній. У випадку одномірного елемента з
ізольованими кінцями x=0 і x=L, граничні умови мають вигляд

\begin{verbatim}
u_{\lambda}(0, t) = 0,
\end{verbatim}

\texttt{\textquotesingle{}u\_\{\textbackslash{}lambda\}(L,\ t)\ =\ 0.\textquotesingle{}\ \textquotesingle{}0\ \textless{}\ t\ \textless{}\ \textbackslash{}infty,\textquotesingle{}\ У\ двовимірній\ області\ ізоляція\ межі\ означає,\ що\ нормальна\ похідна\ температури\ на\ межі\ перетворюється\ на\ нуль.\ Припустимо,\ наприклад,\ що\ круглий\ диск\ термічно\ ізольований\ уздовж\ межі.\ Тоді\ граничну\ умову\ слід\ записати\ так:\ \textquotesingle{}u\_r(R,\ \textbackslash{}theta,\ t)\ =\ 0\textquotesingle{}\ для\ всіх\ \textquotesingle{}0\ \textbackslash{}le\ \textbackslash{}theta\ \textless{}\ 2\textbackslash{}pi\textquotesingle{}\ і\ всіх\ \textquotesingle{}0\ \textless{}\ t\ \textless{}\ \textbackslash{}infty\textquotesingle{}\ .}

З іншого боку, якщо вказати кількість тепла, що проходить через межу
нашого диска, то граничні умови набувають вигляду

\begin{verbatim}
u_r(R, \theta, t) = f(\theta, t),
\end{verbatim}

де \texttt{f(0,\ t)} --- це кількість тепла, що переходить у круглий
диск із зовнішнього джерела.

Тепер проілюструємо різні типи граничних умов.

\section{Типові граничні умови для задачі одномірної
теплопровідності}\label{ux442ux438ux43fux43eux432ux456-ux433ux440ux430ux43dux438ux447ux43dux456-ux443ux43cux43eux432ux438-ux434ux43bux44f-ux437ux430ux434ux430ux447ux456-ux43eux434ux43dux43eux43cux456ux440ux43dux43eux457-ux442ux435ux43fux43bux43eux43fux440ux43eux432ux456ux434ux43dux43eux441ux442ux456}

Припустимо, що у нас є мідний стрижень довжиною 200 см, бікова поверхня
якого термічно ізольована, а початкова температура становить 0 °C.
Припустимо також, що верхній кінець стрижня x=0 термічно ізольований, а
нижній кінець x=200 змивається рухомою водою, яка має постійну
температуру \texttt{g\_2(t)=20} °C (див. рис. 3.6).

\pandocbounded{\includegraphics[keepaspectratio]{images/_page_25_Picture_6.jpeg}}

Рис. 3.6. Змішане завдання.

Математична модель цієї задачі описується наступними чотирма
відношеннями:

\begin{verbatim}
(3.6) \qquad \begin{array}{llll} (\mathrm{Y}\mathrm{H}\Pi) & u_t = \alpha^2 u_{xx}, & 0 < x < 200, & 0 < t < \infty, \\ u_x (0, \, t) = 0, & \\ u_x (200, \, t) = -\frac{h}{k} \left[ u \, (200, \, t) - 20 \right], & 0 < t < \infty, \\ (\mathrm{HY}) & u \, (x, \, 0) = 0 \, ^{\mathrm{o}}\mathrm{C}, & 0 \leqslant x \leqslant 200, & \end{array}
\end{verbatim}

де \texttt{\textbackslash{}alpha\^{}2\ =\ 1,16} см 2/с --- теплова
дифузія міді;

k = 0,93 кал/см·с·°C --- теплопровідність міді;

h --- коефіцієнт теплопередачі. Визначення цього співвідношення ---
велика проблема. Він вимірюється швидкістю виходу тепла від другого
кінця стрижня до води. Це залежить від швидкості потоку води, стану
поверхні, яку слід змивати тощо.

Читач може спробувати провести експеримент, щоб визначити

ділення значення цього коефіцієнта.

\section{ЗАУВАЖЕННЯ}\label{ux437ux430ux443ux432ux430ux436ux435ux43dux43dux44f-2}

\begin{enumerate}
\def\labelenumi{\arabic{enumi}.}
\tightlist
\item
  Типова задача внутрішньої теплопровідності для квадрата показана на
  рис. 3.7.
\end{enumerate}

\texttt{У\ цій\ задачі\ початкова\ температура\ \textquotesingle{}u\textbackslash{}left(x,\textbackslash{},y,\textbackslash{},t\textbackslash{}right)\textquotesingle{}\ при\ t\ =\ 0\ задається\ всередині\ квадрата,\ і\ рівняння\ з\ похідними\ з\ частковими\ має\ вигляд\ Р\ и\ с}

\pandocbounded{\includegraphics[keepaspectratio]{images/_page_26_Figure_9.jpeg}}

. 3.7. Типові граничні умови для задачі дифузії в квадраті: a ---
кількість тепла, \texttt{f\_1(t)} проходить через межу; b ---
термоізольована межа; a --- це навколишня температура Равії
\texttt{10\^{}\{\textbackslash{}circ\}} C; z --- температура
підтримується рівною \texttt{10\^{}\{\textbackslash{}circ\}} C.

а граничні умови будуть виконані на
\texttt{0\ \textless{}\ t\ \textless{}\ \textbackslash{}infty} і
визначатимуть температуру в будь-який момент. Якою б не була ця
температура, вона повинна відповідати граничним умовам, зазначеним на
рис. 3.7.

\begin{enumerate}
\def\labelenumi{\arabic{enumi}.}
\setcounter{enumi}{1}
\tightlist
\item
  Зверніть увагу, що гранична умова
\end{enumerate}

\begin{verbatim}
u_r(R, \theta, t) = -\frac{h}{k} [u(R, \theta, t) - g(\theta, t)]
\end{verbatim}

На колі не означає, що гранична температура дорівнює
\texttt{g(\textbackslash{}theta,t)} , але якщо коефіцієнт теплопередачі
h великий, то ці умови фактично еквівалентні вимогі
\texttt{g(\textbackslash{}theta,t)} температури на межі.

\section{ЗАВДАННЯ}\label{ux437ux430ux432ux434ux430ux43dux43dux44f-1}

\begin{enumerate}
\def\labelenumi{\arabic{enumi}.}
\item
  Намалюйте контури графіків для розв'язання змішаної задачі (3.6) для
  різних моментів часу. Чи відповідають ваші графи визначеним граничним
  умовам? Яка стаціонарна температура стрижня? Очевидно, ваша відповідь
  базується на інтуїції?
\item
  Як би ви інтерпретували наступне змішане завдання?
\end{enumerate}

\begin{verbatim}
\begin{aligned}
(\text{УЧП})\quad & u_t = \alpha^2 u_{xx}, && 0 < x < 1,\ 0 < t < \infty, \\
(\text{ГУ})\quad & u(0,t) = 0,\quad u_x(1,t) = 1, && 0 < t < \infty, \\
(\text{НУ})\quad & u(x,0) = \sin(\pi x), && 0 \leqslant x \leqslant 1.
\end{aligned}
\end{verbatim}

Чи можете ви графічно показати рішення цієї проблеми в різні моменти
часу? Чи буде це рішення, як правило, стаціонарним? Це очевидно?

\begin{enumerate}
\def\labelenumi{\arabic{enumi}.}
\setcounter{enumi}{2}
\tightlist
\item
  Яку фізичну інтерпретацію ви можете дати проблемі
\end{enumerate}

\begin{verbatim}
\begin{array}{llll} (\text{УЧП}) & u_t = \alpha^2 u_{xx}, & 0 < x < 1, & 0 < t < \infty, \\ (\text{ГУ}) & \begin{cases} u_x \, (0, \, t) = 0, & 0 < t < \infty, \\ u_x \, (1, \, t) = 0, & 0 < t < \infty, \end{cases} \\ (\text{НУ}) & u \, (x, \, 0) = \sin (\pi x), & 0 \leqslant x \leqslant 1? \\ \end{array}
\end{verbatim}

Чи можете ви графічно показати рішення цієї проблеми в різні моменти
часу?

Що можна сказати про стабільну температуру?

\begin{enumerate}
\def\labelenumi{\arabic{enumi}.}
\setcounter{enumi}{3}
\tightlist
\item
  Припустимо, що бічна поверхня металевого стрижня не ізольована і має
  початкову температуру 20°C, але температура одного кінця стає 50°C.
  Потім стрижень промивають рідиною з температурою 30°C. Як виглядає
  змішана задача, що відповідає цьому випадку?
\end{enumerate}

\section{ВИВЕДЕННЯ РІВНЯННЯ
ТЕПЛОПРОВІДНОСТІ}\label{ux432ux438ux432ux435ux434ux435ux43dux43dux44f-ux440ux456ux432ux43dux44fux43dux43dux44f-ux442ux435ux43fux43bux43eux43fux440ux43eux432ux456ux434ux43dux43eux441ux442ux456}

МЕТА ЛЕКЦІЇ: Показати, як можна отримати одномірне рівняння
теплопровідності

\begin{verbatim}
u_t = \alpha^2 u_{xx} + f(x, t),
\end{verbatim}

спираючись лише на фундаментальний закон збереження кількості тепла.
Розглядається залежність швидкості теплопередачі від основних
термофізичних параметрів: теплопровідності, теплоємності та густини
матеріалу. Розглядаються різні модифікації базового рівняння
теплопровідності.

У кожній галузі знань існує набір базових принципів, правильність яких
очевидна. Усі інші твердження мають бути виведені з цих базових
принципів. Звісно, те, що очевидно одній людині, може здатися сумнівним
для іншої. Історія будь-якої науки полягає у постійному пошуку більш
фундаментальних принципів, що лежать в основі цієї науки. Основні
принципи мають бути настільки універсальними, щоб ніхто не сумнівався.

? --- Припущення A --- Припущення B --- Припущення C --- Припущення D
--- ? Рис. 4.1. Аксіоматичний метод.

Наприклад, одна людина може вважати, що всі факти певної науки можна
вивести з одного базового припущення, яке ми позначимо як B. На основі
припущення B вона або вона може довести теорему C, що призведе до
теореми D, а яка, у свою чергу, призведе до доведення багатьох інших
теорем (див. рисунок 4.1). Ось як виглядає прогрес у науці --- дослідник
може отримувати все більше нових результатів. Найвидатніші фізики,
хіміки та біологи або саме на цьому шляху.

З іншого боку, замість доведення нових теорем можна спробувати, якщо
можливо, знайти новий принцип (назвемо його A), який є більш
фундаментальним за B. Припущення A природно вважається більш
фундаментальним, ніж B.

Якщо B можна вивести з A. Отже, якщо можна знайти новий принцип A, тоді
межа наших знань зміщується до сфери більш фундаментальних ідей про
певну науку. У теорії теплопровідності основним принципом є закон
збереження енергії (теплова енергія). Усі інші твердження випливають із
цього базового принципу (рис. 4.2).

Закон збереження енергії

\begin{verbatim}
u_t = \alpha^2 u_{xx} + f(x,t)
\end{verbatim}

Інші властивості теплопровідності Рис. 4.2. Енергозбереження: наріжний
камінь теплопровідності.

Звісно, ми могли б забути цю лекцію і використати рівняння
теплопровідності як фундаментальне поняття (дехто може так подумати,
якщо це очевидно), але тоді ми обдуримо більш серйозних студентів, адже
закон збереження вважається найфундаментальнішим принципом у науці.
Науковці, моделюючи свої задачі, часто записують відношення, що
виражають закони збереження, а потім переписують їх у вигляді рівнянь у
похідних похідних.

Повернемося до головної мети нашої лекції --- вивести рівняння
теплопровідності з рівняння збереження енергії.

\section{Виведення рівняння
теплопровідності}\label{ux432ux438ux432ux435ux434ux435ux43dux43dux44f-ux440ux456ux432ux43dux44fux43dux43dux44f-ux442ux435ux43fux43bux43eux43fux440ux43eux432ux456ux434ux43dux43eux441ux442ux456-1}

Розглянемо однорідний стрижень довжини L, щодо якого зробимо такі
припущення:

\begin{enumerate}
\def\labelenumi{\arabic{enumi}.}
\tightlist
\item
  Стрижень виготовлений з одного однорідного провідного матеріалу.
\end{enumerate}

\pandocbounded{\includegraphics[keepaspectratio]{images/_page_29_Figure_10.jpeg}}

Рис. 4.3. Тонкий теплопровідний стержень.

\begin{enumerate}
\def\labelenumi{\arabic{enumi}.}
\setcounter{enumi}{1}
\tightlist
\item
  Бікова поверхня стрижня термічно ізольована (тепло може поширюватися
  лише вздовж осі x).
\end{enumerate}

\texttt{3.\ Стрижень\ тонка,\ тобто\ температура\ всіх\ точок\ у\ кожному\ поперечному\ перерізі\ стрижня\ є\ сталою.\ Якщо\ розглядати\ частину\ вудки\ на\ сегменті\ \textquotesingle{}{[}x,\ x+\textbackslash{}Delta\ x{]}\textquotesingle{}\ і\ використовувати}

закон збереження кількості тепла, тоді можна написати:

\texttt{(4.1)\ Загальна\ зміна\ кількості\ тепла\ на\ сегменті\ \textquotesingle{}{[}x,\ x\ +\ \textbackslash{}Delta\ x{]}\ =\textquotesingle{}\ дорівнює\ загальній\ кількості\ тепла,\ що\ пройшло\ через\ межі\ +\ Нульова\ кількість\ тепла,\ утвореного\ всередині\ сегмента\ \textquotesingle{}{[}x,\ x\ +\ \textbackslash{}Delta\ x{]}\textquotesingle{}\ .}

\texttt{Оскільки\ загальна\ кількість\ тепла\ (у\ калоріях)\ всередині\ сегмента\ \textquotesingle{}\textbackslash{}\{x,\ x\ +\ \textbackslash{}Delta\ x\textbackslash{}\}\textquotesingle{}\ в\ будь-який\ момент\ часу\ t\ обчислюється\ формулою:}

\texttt{Загальна\ кількість\ тепла\ всередині\ сегмента\ \textquotesingle{}{[}x,\ x+\textbackslash{}Delta\ x{]}\ =\ \textbackslash{}int\textbackslash{}limits\_\{-\textbackslash{}infty\}\^{}\{x+\textbackslash{}Delta\ x\}\ coAu\textquotesingle{}\ (s,\ t)\ ds,}

де c --- питома теплоємність матеріалу (показує здатність матеріалу
накопичувати тепло),

Щільність матеріалу,

A --- це площа поперечного перерізу стрижня.

\textbf{va}con збереження енергії (4.1) можна надати у такій
математичній формі:

(4.2)

\begin{verbatim}
\frac{d}{dt} \int_{x}^{x+\Delta x} c\rho Au(s, t) ds = c\rho A \int_{x}^{x+\Delta x} u_{t}(s, t) ds = eA \left[ u_{x}(x+\Delta x, t) - u_{x}(x, t) \right] + A \int_{x+\Delta x}^{x+\Delta x} f(s, t) ds.
\end{verbatim}

\textbf{де} k --- теплопровідність матеріалу (вказує на здатність
матеріалу проводити тепло), \texttt{f(x,\ t)} --- об'ємна потужність
зовнішнього джерела тепла (калорії/см·с).

Задача полягає в тому, щоб записати рівняння (4.2) у формі, яка не
містить інтегралів. Щоб розв'язати цю задачу, нагадаємо читачеві теорему
про середнє значення з ходу інтегральної нечисельності.

\section{Теорема про
значення}\label{ux442ux435ux43eux440ux435ux43cux430-ux43fux440ux43e-ux437ux43dux430ux447ux435ux43dux43dux44f}

Якщо функція \texttt{f(x)} є неперервною на сегменті {[}a, b{]}, то
існує принаймні одна точка
\texttt{\textbackslash{}xi\ \textbackslash{}in\ {[}a,\ b{]}} така, що

\begin{verbatim}
\int_{a}^{b} f(x) dx = f(\xi) (b-a).
\end{verbatim}

Застосовуючи цей результат до рівняння (4.2), ми отримуємо наступний
результат

\begin{verbatim}
coAu_{t}(\xi, t) \Delta x = kA\{u_{x}(x + \Delta x, t) - u_{x}(x, t)\} + Af(\xi, t) \Delta x,
x < \xi < x + \Delta x,
\end{verbatim}

або

\begin{verbatim}
u_t(\xi, t) = \frac{k}{c\rho} \left\{ \frac{u_x(x + \Delta x, t) - u_x(x, t)}{\Delta x} \right\} + \frac{1}{c\rho} f(\xi, t).
\end{verbatim}

\begin{verbatim}
Спрямовуючи \Delta x до 0, отримуємо потрібне рівняння
\end{verbatim}

(4.3)

\begin{verbatim}
u_t(x, t) = \alpha^2 u_{xx}(x, t) + F(x, t),
\end{verbatim}

де

\begin{verbatim}
\alpha^{2}=\frac{k}{c\rho}
\end{verbatim}

\begin{verbatim}
— коефіцієнт теплової дифузії, F\left(x,\ t\right)=\frac{1}{c\rho}f\left(x,\ t\right) — густина джерел тепла.
\end{verbatim}

Ми могли б зупинитися на цьому. Однак ми також розглянемо випадок, коли
бічна поверхня стрижня не термічно ізольована. Припустимо, що кількість
теплового потоку через бічну поверхню стрижня в цьому випадку
пропорційна різниці між температурою стрижня \texttt{u(x,t)} і
температурою навколишнього середовища, яка залишається сталою і дорівнює
нулю. У цьому випадку закон збереження кількості тепла веде до рівняння

\begin{verbatim}
(4.4) u_t = \alpha^2 u_{xx} - \beta u + F(x, t),
\end{verbatim}

\begin{verbatim}
де \beta — коефіцієнт пропорційності потоку через бічну поверхню.
\end{verbatim}

\section{ЗАУВАЖЕННЯ}\label{ux437ux430ux443ux432ux430ux436ux435ux43dux43dux44f-3}

\begin{enumerate}
\def\labelenumi{\arabic{enumi}.}
\tightlist
\item
  Як зазначалося вище, константа k називається теплопровідністю речовини
  і чисельно дорівнює кількості тепла (у калоріях), що проходить за
  секунду через пластину товщиною 1 см і площею поперечного перерізу 1
  см³ при різниці температур на протилежних гранях пластини рівною 1°C.
  Числові значення теплопровідності для різних матеріалів можна знайти в
  довідниках з фізики або хімії. Типові значення теплопровідності
  варіюються від 1 (для міді) до значень, близьких до нуля для хороших
  теплоізоляторів. Якщо стрижень зроблений із однорідного матеріалу, то
  значення k не залежить від x. Для деяких матеріалів значення k
  залежить від температури, і тому рівняння теплопровідності в цьому
  випадку стає нелінійним:
\end{enumerate}

\begin{verbatim}
u_{t} = \frac{1}{c\rho} \frac{\partial}{\partial x} \{ k (u) | u_{x} \}.
\end{verbatim}

Однак найчастіше теплопровідність k дуже слабо впливає на температуру u,
і цю нелінійність можна ігнорувати.

\begin{enumerate}
\def\labelenumi{\arabic{enumi}.}
\setcounter{enumi}{1}
\item
  Значення c зазвичай називають питомою теплоємністю речовини. Він
  характеризує здатність речовини зберігати теплову енергію. Наприклад,
  запечена картопля має високу теплоємність, оскільки може зберігати
  велику кількість тепла на одиницю маси (саме тому вона зберігає високу
  температуру так довго). Чисельна теплоємність дорівнює кількості
  калорій тепла, які потрібно передати 1 граму речовини, щоб підвищити
  її температуру на 1°C. Для більшості розглянутих задач значення c
  можна вважати сталим, не записуючи ні x, ні i. Значення питомої
  теплоємності для конкретних матеріалів можна знайти у фізичних
  довідниках.
\item
  Одиниці вимірювання основних величин, включених у рівняння
\end{enumerate}

теплопровідність (у системі GHS) наступна:

i-температура (градуси Цельсія), \texttt{u\_t} --- швидкість зміни
температури (°C/s), \texttt{u\_x} --- нахил температурної кривої
(°C/см), \texttt{u\_\{xx\}} --- опуклість температурної кривої (°C/см²),

\begin{verbatim}
c = щільність тепла (калор/г \cdot °C)
\end{verbatim}

\begin{verbatim}
k — теплопровідність (кал/см \cdot при \cdot °C), \rho — щільність (г/<sup>см3</sup>), \alpha^2 — це теплова дифузія (см<sup>²</sup>/с).
\end{verbatim}

\begin{verbatim}
4. Слід зазначити, що теплова дифузія матеріалу \alpha^2 = \frac{k}{\epsilon \rho} прямо пропорційна теплопровідності матеріалу і обернено пропорційна густині \rho та питому теплоємності \epsilon ; Ми сподіваємося, що цей зв'язок відповідає інтуїтивним ідеям читача.
\end{verbatim}

\section{ЗАВДАННЯ}\label{ux437ux430ux432ux434ux430ux43dux43dux44f-2}

\begin{verbatim}
1. Підставити розмірність кожної з u,\ u_t,\ \dots величин у рівняння
\end{verbatim}

\begin{verbatim}
u_t = \alpha^2 u_{xx} - \beta u
\end{verbatim}

\begin{verbatim}
і переконайтеся, що всі терміни мають однаковий розмір {}^{\circ}\mathrm{C}/c .
\end{verbatim}

\begin{enumerate}
\def\labelenumi{\arabic{enumi}.}
\setcounter{enumi}{1}
\tightlist
\item
  Підставити розмірність кожної з величин у рівняння
\end{enumerate}

\begin{verbatim}
u_t = \alpha^2 u_{xx} - v u_x,
\end{verbatim}

\begin{verbatim}
де \upsilon вимірюється в одиницях швидкості, і переконайтеся, що всі терми мають однакову розмірність.
\end{verbatim}

3 Отримаємо рівняння теплопровідності

\begin{verbatim}
u_{t} = \frac{1}{c\rho} \frac{\partial}{\partial x} [h(x) u_{x}] + f(x, t)
\end{verbatim}

\begin{verbatim}
Для випадку, коли теплопровідність k\left( x\right) залежить від координати X.
\end{verbatim}

\begin{enumerate}
\def\labelenumi{\arabic{enumi}.}
\setcounter{enumi}{3}
\tightlist
\item
  Припустимо, що \texttt{u(x,\ t)} --- це значення концентрації речовини
  в потоку, що рухається зі швидкістю v. Припустимо, що концентрація
  змінюється внаслідок процесів дифузії та конвекції. Виведемо рівняння
\end{enumerate}

\begin{verbatim}
u_1 = \alpha^2 u_{xx} - v u_x,
\end{verbatim}

\begin{verbatim}
Виходячи з того, що в будь-який момент часу в певному [x, x + \Delta x] сегменті матерії не з'являється і не зникає.
\end{verbatim}

ПРИМІТКА. Закон збереження матерії виглядає так:

\begin{verbatim}
Зміни маси речовини на сегменті [x, x + \Delta x] = = Зміна внаслідок дифузії через межі + Зміна внаслідок переносу через межі.
\end{verbatim}

Лекція 5

\section{ЗМІННЕ
РОЗДІЛЕННЯ}\label{ux437ux43cux456ux43dux43dux435-ux440ux43eux437ux434ux456ux43bux435ux43dux43dux44f}

МЕТА ЛЕКЦІЇ: Познайомити читача з потужним методом розділення змінних і
показати, як цей метод можна використати для розв'язання добре відомої
проблеми дифузії. Оскільки цей метод створює труднощі для деяких учнів
(головним чином через громіздкі алгебранічні перетворення), ми
супроводжуємо презентацію інтуїтивними міркуваннями.

Основна ідея методу полягає в тому, щоб розкласти початкову умову на
найпростіші компоненти, знайти відгук системи на кожен найпростіший
компонент і підсумувати всі відповіді. Таким чином можна знайти
відповідь на довільну початкову умову.

Існуючий метод поділу змінних певною мірою зменшує можливість такої
інтерпретації методу, але все ж не виключає його.

Метод розділення змінних є одним із найвідоміших методів розв'язання
зсувних задач і використовується, якщо:

\begin{itemize}
\tightlist
\item
  \begin{enumerate}
  \def\labelenumi{\arabic{enumi}.}
  \tightlist
  \item
    Рівняння є лінійним і однорідним (не обов'язково з постійними
    коефіцієнтами).
  \end{enumerate}
\item
  \begin{enumerate}
  \def\labelenumi{\arabic{enumi}.}
  \setcounter{enumi}{1}
  \tightlist
  \item
    Граничні умови визначаються як
  \end{enumerate}
\end{itemize}

\begin{verbatim}
\alpha u_x(0, t) + \beta u(0, t) = 0,
\end{verbatim}

\begin{verbatim}
\gamma u_x(1, t) + \delta u(1, t) = 0, де \alpha , \beta , \gamma і \delta — константи (граничні умови, задані в такій формі, називаються лінійними однорідними граничними умовами).
\end{verbatim}

Метод був створений за часів Фур'є (зазвичай відомий як метод Фур'є) і,
здається, є найпопулярнішим (де це застосовно).

\pandocbounded{\includegraphics[keepaspectratio]{images/_page_34_Figure_8.jpeg}}

Рис. 5.1. Схематичне представлення задачі дифузії.

Замість того, щоб вивчати метод у загальному випадку, давайте спочатку
розглянемо конкретну задачу (загальний випадок обговоримо пізніше).
Розглянемо змішану задачу типу дифузії: щоб знайти розв'язок

(PPP)

\begin{verbatim}
u_t = \alpha^2 u_{xx}
, 0 < x < 1 , 0 < t < \infty ,
\end{verbatim}

що задовольняє граничні умови

(GI)

\begin{verbatim}
\begin{cases} u(0, t) = 0, & 0 < t < \infty, \\ u(1, t) = 0, & 0 < t < \infty, \end{cases}
\end{verbatim}

та початкова умова

(ГУ)

\begin{verbatim}
u(x, 0) = \varphi(x), 0 \le x \le 1.
\end{verbatim}

\begin{verbatim}
Перш ніж почати розділяти змінні, давайте дамо фізичну інтерпретацію нашої проблеми. Отже, існує стрижень скінченної довжини, кінці якого підтримуються при сталій температурі, рівній нулю (насправді кінці можна підтримувати при значно вищій температурі, значення якої вважається початковою точкою). Додаткові дані про завдання подаються як початкова умова. Наша мета — знайти розподіл температури u\left(x,t\right) в наступний час.
\end{verbatim}

\section{Загальні принципи методу розділення
змінних}\label{ux437ux430ux433ux430ux43bux44cux43dux456-ux43fux440ux438ux43dux446ux438ux43fux438-ux43cux435ux442ux43eux434ux443-ux440ux43eux437ux434ux456ux43bux435ux43dux43dux44f-ux437ux43cux456ux43dux43dux438ux445}

Для найпростішого диференціального рівняння в похідних похідних
розділення змінних є пошуком розв'язків у вигляді

\begin{verbatim}
u(x, t) = X(x)T(t),
\end{verbatim}

де \texttt{X(x)} --- це функція, що залежить лише від змінної x, а
\texttt{T(t)} --- функція, що залежить лише від t. Це розв'язок у
певному сенсі є найпростішим, оскільки температура \texttt{u(x,\ t)},
представлена у такій формі, зберігатиме «форму» профілю у різний час
часу t (див. рисунок 5.2).

\pandocbounded{\includegraphics[keepaspectratio]{images/_page_35_Figure_6.jpeg}}

РИСУНОК 5.2 Графік функції X(x)

\begin{verbatim}
T(t) у різних точках часу t.

~~
Загальна ідея полягає в тому, щоб знайти нескінченну кількість таких часткових диференціальних розв'язків (які задовольняють граничні умови). Ці найпростіші функції u_n\left(x,\ t\right) = X_n\left(x\right)T_n\left(t\right) (які називаються фундаментальними розв'язками) є, так би мовити, елементарними будівельними блоками, з яких будується розв'язок нашої проблеми. Розв'язок нашої задачі u\left(x,\ t\right) у вигляді такої лінійної комбінації фундаментальних розв'язків X_n\left(x\right)T_n\left(t\right) , що отримана сума
~~
\end{verbatim}

\texttt{\textbackslash{}sum\_\{n=1\}\^{}\{\textbackslash{}infty\}\ A\_n\ X\_n(x)\ T\_n(t)}

\begin{verbatim}
задовольняє початкові умови. Оскільки ця сума задовольняє рівняння та граничні умови, вона є розв'язком початкової задачі. Ми маємо робити всі ці розрахунки достатньо детально.

# Розділення змінних

КРОК 1. (Пошук елементарних розв'язків рівняння з частковими похідними.)

Ми хочемо знайти функцію u(x, t), яка є розв'язком

Наступного завдання:
\end{verbatim}

\texttt{\textbackslash{}begin\{aligned\}}
\texttt{(\textbackslash{}text\{УЧП\})\textbackslash{}quad\ \&\ u\_t\ =\ \textbackslash{}alpha\^{}2\ u\_\{xx\},\ \&\&\ 0\ \textless{}\ x\ \textless{}\ 1,\textbackslash{}\ 0\ \textless{}\ t\ \textless{}\ \textbackslash{}infty,\ \textbackslash{}\textbackslash{}}
\texttt{(\textbackslash{}text\{ГУ\})\textbackslash{}quad\ \&\ u(0,t)\ =\ 0,\textbackslash{}quad\ u(1,t)\ =\ 0,\ \&\&\ 0\ \textless{}\ t\ \textless{}\ \textbackslash{}infty,\ \textbackslash{}\textbackslash{}}
\texttt{(\textbackslash{}text\{НУ\})\textbackslash{}quad\ \&\ u(x,0)\ =\ \textbackslash{}varphi(x),\ \&\&\ 0\ \textless{}\ x\ \textbackslash{}leqslant\ 1.}
\texttt{\textbackslash{}end\{aligned\}}

\begin{verbatim}
Пошукаємо розв'язки, які представлені у вигляді u(x, t) = X(x) T(t). Для цього підставляємо вираз X(x) T(t) у рівняння. Внаслідок підстановки отримуємо
\end{verbatim}

X(x) T'(t) = q\^{}2 X'\,'(x) T(t).

\begin{verbatim}
~~
**Тепер** виконаємо операцію, властиву цьому методу: *поділимо* обидві сторони останнього рівняння на \alpha^2 X(x) T(t) , у результаті чого отримуємо
~~
\end{verbatim}

\texttt{\textbackslash{}frac\{T\textquotesingle{}(t)\}\{\textbackslash{}alpha\^{}2\ T(t)\}\ =\ \textbackslash{}frac\{X\textquotesingle{}\textquotesingle{}(x)\}\{X(x)\}.}

\begin{verbatim}
Цей вираз вважається таким, що розділив змінні, тобто ліва частина рівняння залежить лише від t, а права — лише від x. Оскільки x і t не є розділенням комунікації один від одного, кожна частина рівняння має бути сталою. Позначимо цю константу k, тоді
\end{verbatim}

\texttt{\textbackslash{}frac\{T\textquotesingle{}\}\{\textbackslash{}alpha\^{}2\ T\}\ =\ \textbackslash{}frac\{X\textquotesingle{}\textquotesingle{}\}\{X\}\ =\ k}

\begin{verbatim}
або
\end{verbatim}

\texttt{T\textquotesingle{}-k\textbackslash{}alpha\^{}2T=0,} X'\,'-kX=0.

\begin{verbatim}
Тепер можливо розв'язати кожне з цих звичайних диференціальних рівнянь. Добуток відповідних розв'язків задовольнятиме початкове рівняння з похідними похідними. (Зверніть увагу, що ми спростили оригінальне диференціальне рівняння другого порядку до двох звичайних диференціальних рівнянь.)

Тепер звернемо увагу на таку важливу обставину: константа відстань k має бути від'ємною ^1 ) (іншими словами, функції T(t) мають прагнути до нуля в точці t \to \infty ). З огляду на це, давайте розглянемо позначення k = -\lambda^2 , де \lambda не дорівнює нулю (у цьому випадку вираз -\lambda^2 завжди буде від'ємним). Враховуючи нову нотацію для константи відстань двох звичайних диференціальних рівнянь,

<sup>1)</sup> Інакше рівняння X'' - kX = 0 з граничними умовами X(0) = X(1) = 0 має лише гривну розв'язку X(x) = 0. Ед.

записуються у формі
\end{verbatim}

\texttt{T\textquotesingle{}\ +\ \textbackslash{}lambda\^{}2\ \textbackslash{}alpha\^{}2\ T\ =\ 0,}
\texttt{X\textquotesingle{}\textquotesingle{}\ +\ \textbackslash{}lambda\^{}2\ X\ =\ 0.}

\begin{verbatim}
Розв'яжемо ці рівняння. Це рівняння стандартного типу. Їхні загальні рішення фіксуються у формі
\end{verbatim}

\texttt{T(t)\ =\ Ae\^{}\{-\textbackslash{}gamma\ *\_\{O\}\ *\_\{I\}\}}

\begin{verbatim}
(A — довільна константа),
\end{verbatim}

\texttt{X(x)\ =\ A\ \textbackslash{}sin(\textbackslash{}lambda\ x)\ +\ B\ \textbackslash{}cos(\textbackslash{}lambda\ x)}

\begin{verbatim}
(A, B — довільні константи).

Відповідно, функції виду
\end{verbatim}

\texttt{u(x,\ t)\ =\ e\^{}\{-\textbackslash{}lambda\^{}2\ \textbackslash{}alpha\^{}2\ t\}\ {[}A\ \textbackslash{}sin(\textbackslash{}lambda\ x)\ +\ B\ \textbackslash{}cos(\textbackslash{}lambda\ x){]}}

\begin{verbatim}
~~
(де A, B і \lambda — довільні константи) задовольняють NPP u_t = \alpha^2 u_{xx} (пропонується перевірити це у задачі 1, розміщеній наприкінці цієї лекції). Отже, маємо нескінченну множину функцій, які задовольняють початкове диференціальне рівняння в похідних похідних.
~~

КРОК 2. (Пошук розв'язків, що задовольняють граничним умовам.)

Зараз ситуація така: ми маємо нескінченну кількість розв'язків початкового рівняння, але не всі вони задовольняють граничні або початкові умови. Наступний крок — вибрати таку підмножину розв'язків вигляду

(5.1)
\end{verbatim}

\texttt{e\^{}\{-\textbackslash{}lambda\^{}2\ \textbackslash{}alpha\^{}\{2\textbackslash{}ell\}\}\ {[}A\ \textbackslash{}sin(\textbackslash{}lambda\ x)\ +\ B\ \textbackslash{}cos(\textbackslash{}lambda\ x){]},}

\begin{verbatim}
які задовольняють граничні умови
\end{verbatim}

u(0, t) = 0,

\begin{verbatim}
u(1, t) = 0. Для цього підставляємо розв'язки (5.1) у ці граничні умови. У результаті ми отримуємо
\end{verbatim}

\texttt{u(0,\ t)\ =\ Be\^{}\{-\textbackslash{}lambda\^{}2\ \textbackslash{}alpha\^{}\{2t\}\}\ =\ 0\ \textbackslash{}Rightarrow\ B\ =\ 0,}
\texttt{u(1,\ t)\ =\ Ae\^{}\{-\textbackslash{}lambda\^{}2\ \textbackslash{}alpha\^{}\{2t\}\}\ \textbackslash{}sin\ \textbackslash{}lambda\ =\ 0\ \textbackslash{}Rightarrow\ \textbackslash{}sin\ \textbackslash{}lambda\ =\ 0.}

\begin{verbatim}
Друга гранична умова накладає обмеження на можливі значення константи \lambda розділення: вона має бути коренем рівняння \sin\lambda=0 . Іншими словами, щоб виконати умову u(1, t)=0, необхідно вимагати виконання
\end{verbatim}

\texttt{\textbackslash{}lambda\ =\ \textbackslash{}pm\ \textbackslash{}pi}

\begin{verbatim}
, \pm 2\pi , ... або
\end{verbatim}

\texttt{\textbackslash{}lambda\_n\ =\ \textbackslash{}pm\ n\textbackslash{}pi}

\begin{verbatim}
, n = 1, 2, \ldots Зверніть увагу, що друга гранична умова може бути виконана, якщо поставити A = 0, але в цьому випадку розв'язок (5.1) буде ідентичним нулю.

![](images/_page_38_Figure_3.jpeg)

Рис. 5.3. Фундаментальні рішення n_n(x, t) = A_n e^{-(\pi n \alpha)^{n/t}} \sin{(\pi n x)} .

Тепер, коли ми зробили другий крок, маємо нескінченний набір ознак

(5.2)
\end{verbatim}

\texttt{u\_n(x,\ t)\ =\ A\_n\ e\^{}\{-(n\textbackslash{}pi\textbackslash{}alpha)\^{}2\ t\}\ \textbackslash{}sin(n\textbackslash{}pi\ x),\ \textbackslash{}quad\ n\ =\ 1,\ 2\ \textbackslash{}dots}

\begin{verbatim}
кожна з яких задовольняє рівняння з частковими похідними та граничними умовами). Отримані функції — це будівельні блоки, з яких ми будемо створювати рішення нашої проблеми. Це рішення буде певною сумою

<sup>2)</sup> Зверніть увагу, що функції u_n і u_{-n} відрізняються лише знаком.

цих найпростіших функцій. Конкретний тип суми залежить від початкового стану. Фіг. Рисунок 5.3 показує графіки кількох таких фундаментальних функцій

КРОК 3. (Пошук розв'язку, який задовольняє рівняння, межу та початкові умови.)

Останній (і, ймовірно, найцікавіший з математичної точки зору) крок — знайти таку суму фундаментальних розв'язків
\end{verbatim}

\texttt{u(x,\ t)\ =\ \textbackslash{}sum\_\{n=1\}\^{}\{\textbackslash{}infty\}\ A\_n\ e\^{}\{-(n\textbackslash{}pi\textbackslash{}alpha)\^{}n\ t\}\ \textbackslash{}sin(n\textbackslash{}pi\ x),}

\begin{verbatim}
тобто при виборі таких коефіцієнтів A_n , що функція задовольнятиме початкову умову
\end{verbatim}

\texttt{u\textbackslash{}left(\ x,\textbackslash{}\ 0\textbackslash{}right)\ =\textbackslash{}varphi\textbackslash{}left(\ x\textbackslash{}right)\ .}

\begin{verbatim}
Підставляючи суму в початковому стані, отримуємо

(5.3)
\end{verbatim}

\texttt{q(x)\ =\ \textbackslash{}sum\_\{n=1\}\^{}\{\textbackslash{}infty\}\ A\_n\ \textbackslash{}sin(n\textbackslash{}pi\ x).}

\begin{verbatim}
Це рівняння підводить нас до цікавого питання: чи можна розкласти початкову температуру \phi(x) на ряд елементарних функцій вигляду
\end{verbatim}

\texttt{A\_1\ \textbackslash{}sin(\textbackslash{}pi\ x)\ +\ A\_2\ \textbackslash{}sin(2\textbackslash{}pi\ x)\ +\ A\_3\ \textbackslash{}sin(3\textbackslash{}pi\ x)\ +\ \textbackslash{}dots}

\begin{verbatim}
?

Позитивну відповідь на це питання дав французький математик Жозеф Фур'є. Виявилося, що для достатньо хороших функцій така декомпозиція можлива ^1 ). Тоді виникає нове питання: знайти коефіцієнти розкладу A_n ?

Насправді це дуже просто. Для цього використаємо властивість системи функцій
\end{verbatim}

\texttt{\textbackslash{}\{\textbackslash{}sin(n\textbackslash{}pi\ x);\ n=1,\ 2,\ \textbackslash{}ldots\textbackslash{}\},}

\begin{verbatim}
відомий як ортогональність.

У нашому випадку (див. завдання 2) ці функції задовольняють умови
\end{verbatim}

\texttt{\textbackslash{}int\_\{0\}\^{}\{1\}\ \textbackslash{}sin(m\textbackslash{}pi\ x)\ \textbackslash{}sin(n\textbackslash{}pi\ x)\ dx\ =\ \textbackslash{}begin\{cases\}\ 0,\ \&\ m\ \textbackslash{}neq\ n,\ \textbackslash{}\textbackslash{}\ 1/2,\ \&\ m\ =\ n.\ \textbackslash{}end\{cases\}}

\begin{verbatim}
Умови розкладу функцій у рядах Фур'є наведені в лекції 11. Для глибшого ознайомлення з цим питанням можна звернутися до численних освітніх літератур (див., наприклад: Г. П. Толстов, Серія Фур'є. — Москва: Наука, 1980). — Ед. Ед.

Ця властивість чітко показана на рис. 5.4, яка показує графіки деяких функцій із цієї системи.

Отже, ми хочемо знайти коефіцієнти у розкладі
\end{verbatim}

\texttt{\textbackslash{}varphi(x)\ =\ A\_1\ \textbackslash{}sin(\textbackslash{}pi\ x)\ +\ A\_2\ \textbackslash{}sin(2\textbackslash{}pi\ x)\ +\ A\_3\ \textbackslash{}sin(3\textbackslash{}pi\ x)\ +\ \textbackslash{}dots}

\begin{verbatim}
![](images/_page_40_Figure_5.jpeg)

РИСУНОК 5.4. Послідовність ортогональних функцій

Помножте обидві сторони цього співвідношення на \sin{(m\pi x)} (m — довільне ціле число) і інтегруйте від нуля до одиниці. У результаті ми отримуємо
\end{verbatim}

\texttt{\textbackslash{}int\_\{0\}\^{}\{1\}\ \textbackslash{}varphi(x)\ \textbackslash{}sin(m\textbackslash{}pi\ x)\ \textbackslash{},\ dx\ =\ A\_\{m\}\ \textbackslash{}int\_\{0\}\^{}\{1\}\ \textbackslash{}sin\^{}\{2\}(m\textbackslash{}pi\ x)\ \textbackslash{},\ dx\ =\ \textbackslash{}frac\{A\_\{m\}\}\{2\}}

\begin{verbatim}
(усі інші члени зрівнялися з нулем через ортогональність). Розв'язуючи рівняння відносно A_m , отримуємо
\end{verbatim}

\texttt{A\_m\ =\ 2\ \textbackslash{}int\_0\^{}1\ \textbackslash{}varphi(x)\ \textbackslash{}sin(m\textbackslash{}pi\ x)\ dx.}

\begin{verbatim}
Отже, отримано, що розв'язок записується у вигляді

(5.4)
\end{verbatim}

\texttt{u(x,\ t)\ =\ \textbackslash{}sum\_\{n=1\}\^{}\{\textbackslash{}infty\}\ A\_n\ e\^{}\{-(n\textbackslash{}pi\textbackslash{}alpha)\^{}2\ t\}\ \textbackslash{}sin(\textbackslash{}pi\ n\ x),}

\begin{verbatim}
де коефіцієнти A_n визначаються формулами
\end{verbatim}

\texttt{A\_n\ =\ 2\textbackslash{}int\_0\^{}1\ \textbackslash{}varphi(x)\ \textbackslash{}sin(\textbackslash{}pi\ nx)\ dx}

\begin{verbatim}
Ви можете переконатися, що отримане рішення задовольняє всі умови початкової задачі. На цьому завершується третій крок.

Багато студентів розчаровані, коли бачать, наскільки складним виглядає рішення, і майже не намагаються дивитися на нього знову (виглядає дуже погано). Однак це рішення не здаватиметься таким складним, якщо ви витратите час на його аналіз. Насправді, складніша форма розв'язку робить його більш інформативним. Нижче наведено кілька нотаток, які допоможуть вам інтерпретувати це рішення.

# ЗАУВАЖЕННЯ

Зверніть увагу, що єдина різниця між синусоїдальним розкладом функції φ(x) і розв'язком (5.4) — це наявність часового фактора e<sup>-(nπα)<sup>3</sup>1</sup> у кожному члені ряду. Отже, якщо початкова умова має просту декомпозицію, наприклад, вигляду
\end{verbatim}

\texttt{\textbackslash{}varphi(x)\ =\ \textbackslash{}sin(\textbackslash{}pi\ x)\ +\ \textbackslash{}frac\{1\}\{2\}\textbackslash{}sin(3\textbackslash{}pi\ x),}

\begin{verbatim}
тоді рішення можна негайно записати
\end{verbatim}

\texttt{u(x,\ t)\ =\ e\^{}\{-(\textbackslash{}pi\textbackslash{}alpha)\^{}2\ t\}\ \textbackslash{}sin(\textbackslash{}pi\ x)\ +\ \textbackslash{}frac\{1\}\{2\}\ e\^{}\{-(3\textbackslash{}pi\textbackslash{}alpha)\^{}2\ t\}\ \textbackslash{}sin(3\textbackslash{}pi\ x).}

\begin{verbatim}
У цьому випадку очевидно, що якщо ми розкладемо \phi(x) на ряд Фур'є вздовж синусів, то отримуємо
\end{verbatim}

A\_1 = 1,

\begin{verbatim}
A_2 = 0, A_3 = 1/2, A_4 = A_5 = \dots = 0. - 2. Розв'язок (5.4) можна інтерпретувати так: ми представляємо початкову температуру \varphi(x) як суму найпростіших функцій виду A_n \sin(\pi nx) ; кожна така функція породжує відповідь A_n e^{-(\pi n\alpha)^2 t} \sin(\pi nx) . Сумуючи всі такі відповіді, отримуємо розв'язок, який відповідає початковій умові u(x, 0) = \varphi(x) .
- 3. Кожен член у декомпозиції
\end{verbatim}

\texttt{u(x,\ t)\ =\ A\_1\ e\^{}\{-(\textbackslash{}pi\textbackslash{}alpha)\^{}2\ t\}\ \textbackslash{}sin(\textbackslash{}pi\ x)\ +\ A\_2\ e\^{}\{-(2\textbackslash{}pi\textbackslash{}alpha)\^{}2\ t\}\ \textbackslash{}sin(2\textbackslash{}pi\ x)\ +\ \textbackslash{}cdots}

\begin{verbatim}
є функцією від x і t. Зверніть увагу, що внесок членів з великими числами в t>0 дуже малий через множник e^{-(n\pi\alpha)^2 t} . Відповідно, після закінчення досить великого Р и с

![](images/_page_42_Figure_2.jpeg)

. 5.5. У задачах дифузії вищі порядки терміни швидко затухають; \phi(x) — початкова температура. Слід зазначити, що по-перше, у розчині зникають високі частоти — по-перше, зникли найбільші коливання міді.

Повний розв'язок приблизно збігається з першим членом, який є півхвилі синусоїди, що з часом згасає (рис. 1). 5.5.)

# ЗАВДАННЯ

1. Показати, що функції
\end{verbatim}

\texttt{u(x,\ t)\ =\ e\^{}\{-\textbackslash{}lambda\^{}2\ \textbackslash{}alpha\textquotesingle{}\ t\}\ \textbackslash{}left{[}\ A\ \textbackslash{}sin(\textbackslash{}lambda\ x)\ +\ B\ \textbackslash{}cos(\textbackslash{}lambda\ x)\ \textbackslash{}right{]}}

\begin{verbatim}
задовольняють рівняння u_l = \alpha^2 u_{xx} для довільних значень A, B і \lambda .

2. Покажи це
\end{verbatim}

\texttt{\textbackslash{}int\_0\^{}1\ \textbackslash{}sin(m\textbackslash{}pi\ x)\ \textbackslash{}sin(n\textbackslash{}pi\ x)\ dx\ =\ \textbackslash{}begin\{cases\}\ 0,\ \&\ m\ \textbackslash{}neq\ n,\ \textbackslash{}\textbackslash{}\ 1/2,\ \&\ m\ =\ n.\ \textbackslash{}end\{cases\}}

\begin{verbatim}
ПРИМІТКА. Використання ідентичності
\end{verbatim}

\texttt{\textbackslash{}sin(mx)\textbackslash{}sin(nx)\ =\ \textbackslash{}frac\{1\}\{2\}{[}\textbackslash{}cos(m-n)x\ -\ \textbackslash{}cos(m+n)x{]}.}

\begin{verbatim}
3. Знайти розклад у рядах Фур'є синусами функції \varphi(x) = 1 на сегменті [0, 1]. Побудуйте перші три терміни розкладу.

Використовуючи результати завдання 3, знайдіть розв'язок наступної змішаної задачі:

(PPP)
\end{verbatim}

u\_t = u\_\{xx\}

\begin{verbatim}
0 < x < 1 , 
(GI)
\end{verbatim}

\texttt{\textbackslash{}begin\{cases\}\ u(0,\ t)\ =\ 0,\ \textbackslash{}\textbackslash{}\ u(1,\ t)\ =\ 0,\ \textbackslash{}end\{cases\}}

\begin{verbatim}
0 < t < \infty , 
(НУ) u(x, 0) = 1, \quad 0 \le x \le 1 .

(Зверніть увагу, що ця проблема фізично не має значення, оскільки припускається, що температура на кінцях стрижня миттєво знижується з одного до нуля \phi(x) .

5. Знайти розв'язок проблеми 4, якщо початкова умова встановлена у вигляді
\end{verbatim}

\texttt{u(x,\ 0)\ =\ \textbackslash{}sin(2\textbackslash{}pi\ x)\ +\ \textbackslash{}frac\{1\}\{3\}\textbackslash{}sin(4\textbackslash{}pi\ x)\ +\ \textbackslash{}frac\{1\}\{5\}\textbackslash{}sin(6\textbackslash{}pi\ x).}

\begin{verbatim}
6. Розв'яжіть задачу 4, якщо початкова умова u\left( x,\ 0\right) =x-x^{2},\quad 0< x<1. Лекція 6

# ПЕРЕТВОРЕННЯ НЕОДНОРІДНИХ ГРАНИЧНИХ УМОВ У ОДНОРІДНІ ГРАНИЧНІ УМОВИ

МЕТА ЛЕКЦІЇ: Показати, як виникає змішана задача
\end{verbatim}

\texttt{\textbackslash{}begin\{array\}\{lll\}\ (\textbackslash{}text\{УЧП\})\ \&\ u\_t\ -\ a\^{}2\ u\_\{xx\}\ =\ f\textbackslash{}left(x,\ \textbackslash{}\ t\textbackslash{}right)\ \&\ \textbackslash{}frac\{\textbackslash{}beta\_t\}\{97\}\ \textbackslash{}\textbackslash{}\ (\textbackslash{}text\{ГУ\})\ \&\ \textbackslash{}begin\{cases\}\ \textbackslash{}alpha\_1\ u\_x\textbackslash{}left(0,\ \textbackslash{}\ t\textbackslash{}right)\ +\ \textbackslash{}beta\_1\ u\textbackslash{}left(0,\ \textbackslash{}\ t\textbackslash{}right)\ =\ g\_1\textbackslash{}left(t\textbackslash{}right),\ \textbackslash{}\textbackslash{}\ a\_2\ u\_x\textbackslash{}left(L,\ \textbackslash{}\ t\textbackslash{}right)\ +\ \textbackslash{}beta\_2\ u\textbackslash{}left(L,\ \textbackslash{}\ t\textbackslash{}right)\ =\ g\_2\textbackslash{}left(t\textbackslash{}right),\ \textbackslash{}end\{cases\}\ \textbackslash{}\textbackslash{}\ (\textbackslash{}text\{HУ\})\ \&\ u\textbackslash{}left(x,\ \textbackslash{}\ 0\textbackslash{}right)\ =\ \textbackslash{}varphi\textbackslash{}left(x\textbackslash{}right)\ \textbackslash{}end\{array\}}

\begin{verbatim}
можна звести до задачі вигляду
\end{verbatim}

\texttt{U\_\{t\}\ -\ \textbackslash{}alpha\^{}\{2\}U\_\{xx\}\ =\ F(x,\ t),}
\texttt{\textbackslash{}alpha\_\{1\}U\_\{x\}(0,\ t)\ +\ \textbackslash{}beta\_\{1\}U(0,\ t)\ =\ 0,}
\texttt{\textbackslash{}alpha\_\{2\}U\_\{x\}(L,\ t)\ +\ \textbackslash{}beta\_\{2\}U(L,\ t)\ =\ 0,}
\texttt{U(x,\ 0)\ =\ \textbackslash{}varphi(x)}

\begin{verbatim}
з нульовими граничними умовами. Цей новий виклик можна подолати

- 1. Шляхом ділення змінних (якщо нова задача виявиться однорідною, тобто F(x, t) = 0).

~~
- 2. Методом інтегральних перетворень і розкладу власними функціями, якщо F(x, t) \neq 0 .
~~

Метод розділення змінних, про який йшлося в попередній лекції, є досить потужним, а розв'язки, отримані за його допомогою, подаються у зручній формі. Однак читач повинен розуміти, що цей метод не застосовний до всіх задач. Щоб метод розділення змінних був застосовним, граничні умови мають бути лінійними, однорідними, тобто представленими як

(6.1)
\end{verbatim}

\texttt{\textbackslash{}begin\{aligned\}\ \textbackslash{}alpha\_1\ u\_x(0,\ t)\ +\ \textbackslash{}beta\_1\ u(0,\ t)\ \&=\ 0,\ \textbackslash{}\textbackslash{}\ \textbackslash{}alpha\_2\ u\_x(L,\ t)\ +\ \textbackslash{}beta\_2\ u(L,\ t)\ \&=\ 0.\ \textbackslash{}end\{aligned\}}

\begin{verbatim}
Мета цієї лекції — показати, як виникає задача з неоднорідними граничними умовами форми

(6.2)
\end{verbatim}

\texttt{\textbackslash{}begin\{aligned\}\ \&\ (\textbackslash{}mathsf\{У\}\textbackslash{}mathsf\{Ч\}\textbackslash{}Pi)\ \&\ u\_t\ =\ \textbackslash{}alpha\^{}2\ u\_\{xx\},\ \textbackslash{}\textbackslash{}\ \&\ (\textbackslash{}mathsf\{\textbackslash{}Gamma\}\textbackslash{}mathsf\{Y\})\ \&\ \textbackslash{}begin\{cases\}\ \textbackslash{}alpha\_1\ u\_x\ \textbackslash{},\ (0,\ \textbackslash{},\ t)\ +\ \textbackslash{}beta\_1\ u\ \textbackslash{},\ (0,\ \textbackslash{},\ t)\ =\ g\_1\ \textbackslash{},\ (t)\ \textbackslash{}\textbackslash{}\ \textbackslash{}alpha\_2\ u\_x\ \textbackslash{},\ (L,\ \textbackslash{},\ t)\ +\ \textbackslash{}beta\_2\ u\ \textbackslash{},\ (L,\ \textbackslash{},\ t)\ =\ g\_2\ \textbackslash{},\ (t)\ \textbackslash{}end\{cases\}\ \textbackslash{}text\{\ (неоднородные\ \}\ \textbackslash{}mathsf\{\textbackslash{}Gamma\}\textbackslash{}mathsf\{Y\}),\ \textbackslash{}\textbackslash{}\ \&\ (\textbackslash{}mathsf\{H\}\textbackslash{}mathsf\{Y\})\ \&\ u\ \textbackslash{},\ (x,\ \textbackslash{},\ 0)\ =\ \textbackslash{}varphi\ \textbackslash{},\ (x)\ \textbackslash{}end\{aligned\}}

\begin{verbatim}
можна розв'язати, зведивши її до задачі з однорідними граничними умовами.

Нове завдання можна розв'язати іншими методами

(наприклад, шляхом розкладу за власними функціями).

Почнемо з надзвичайно простого прикладу перетворення гетерогенних граничних умов на однорідні.

# Перетворення неоднорідних граничних умов у однорідні граничні умови

Розглянемо проблему поширення тепла в термоізольованому стрижні, кінці якого підтримуються при сталих температурах k_1 і k_2, тобто
\end{verbatim}

\texttt{\textbackslash{}begin\{aligned\}}
\texttt{(\textbackslash{}text\{УЧП\})\textbackslash{}quad\ \&\ u\_t\ =\ \textbackslash{}alpha\^{}2\ u\_\{xx\},\ \&\&\ 0\ \textless{}\ x\ \textless{}\ L,\textbackslash{}\ 0\ \textless{}\ t\ \textless{}\ \textbackslash{}infty,\ \textbackslash{}\textbackslash{}}
\texttt{(6.3)\textbackslash{}\ (\textbackslash{}text\{ГУ\})\textbackslash{}quad\ \&\ u(0,t)\ =\ k\_1,\textbackslash{}quad\ u(L,t)\ =\ k\_2,\ \&\&\ 0\ \textless{}\ t\ \textless{}\ \textbackslash{}infty,\ \textbackslash{}\textbackslash{}}
\texttt{(\textbackslash{}text\{НУ\})\textbackslash{}quad\ \&\ u(x,0)\ =\ \textbackslash{}varphi(x),\ \&\&\ 0\ \textbackslash{}leqslant\ x\ \textbackslash{}leqslant\ L.}
\texttt{\textbackslash{}end\{aligned\}}

\begin{verbatim}
Складність цієї задачі полягає в тому, що, оскільки граничні умови в ній є гетерогенними, ми не можемо розв'язати її шляхом розділення змінних. Однак очевидно, що в t \to \infty розв'язок нашої задачі схиляється до стаціонарного розв'язку, який лінійно (вздовж x) змінюється від температури k_1 до температури k_2 (рис. 6.1). Іншими словами, розумно припустити, що температуру в нашій задачі можна представити як суму
~~

# з двох термів
\end{verbatim}

\texttt{u\textbackslash{}left(x,\textbackslash{},t\textbackslash{}right)=}

\begin{verbatim}
~~
стаціонарна + тимчасова = \uparrow \uparrow Частина розв'язку, що залежить від HY і має тенденцію до нуля з помноженням на =\left[k_1+\frac{x}{L}\left(k_2-k_1\right)\right]+U\left(x,\,t\right). У цьому випадку наше завдання — знайти температуру переходу U\left( x,\,t\right) . Заміна
~~
\end{verbatim}

\texttt{u(x,\ t)\ =\ \textbackslash{}left{[}k\_1\ +\ \textbackslash{}frac\{x\}\{L\}(k\_2\ -\ k\_1)\textbackslash{}right{]}\ +\ U(x,\ t)}

\begin{verbatim}
~~
У початковій задачі (6.3) ми отримуємо нову задачу відносно невідомої функції U\left(x,\ t\right) . Ви можете вирішити проблему Р и с
~~

![](images/_page_45_Figure_7.jpeg)

~~
. 6.1. Розв'язання задачі (6.3) у різний час часу a- початкова температура u\left(x,\ 0\right)=\phi\left(x\right); b - стаціонарна температура u\left(x,\ \infty\right)=k_{1}+\frac{x}{L}\left(k_{2}-k_{1}\right); b - перехідна температура.
~~

~~
Відносно нова невідома функція U\left(x,\,t\right) і додає її до стаціонарного розв'язку, що дає бажану функцію u\left(x,\,t\right) . Виконавши ці прості перетворення з (6.3),
~~

Ми отримаємо

(6.4)
\end{verbatim}

\texttt{U\_\{t\}\ =\ \textbackslash{}alpha\^{}\{2\}U\_\{xx\}\ \textbackslash{}quad\ 0\ \textless{}\ x\ \textless{}\ L,}
\texttt{(\textbackslash{}Gamma\ Y)\ \textbackslash{}quad\ \textbackslash{}begin\{cases\}\ U(0,\ t)\ =\ 0,\ \textbackslash{}\textbackslash{}\ U(L,\ t)\ =\ 0,\ \textbackslash{}end\{cases\}\ \textbackslash{}quad\ 0\ \textless{}\ t\ \textless{}\ \textbackslash{}infty,}
\texttt{(HY)\ \textbackslash{}quad\ U(x,\ 0)\ =\ \textbackslash{}varphi(x)\ -\ \textbackslash{}left{[}k\_\{1\}\ +\ \textbackslash{}frac\{x\}\{L\}(k\_\{2\}\ -\ k\_\{1\})\textbackslash{}right{]}\ =\ \textbackslash{}overline\{\textbackslash{}varphi\}(x),}

\begin{verbatim}
~~
де \overline{\phi}(x) — нова, але відома початкова умова.
~~

Ця задача стосується не лише однорідного рівняння, а, на щастя, однорідних граничних умов, що дозволяє розв'язувати її методом розділення змінних. Насправді, читач, ймовірно, пам'ятає рішення:

(6.5)
\end{verbatim}

\texttt{U(x,\ t)\ =\ \textbackslash{}sum\_\{n=1\}\^{}\{\textbackslash{}infty\}\ a\_n\ e\^{}\{-(n\textbackslash{}pi\textbackslash{}alpha)\^{}4\ t\}\ \textbackslash{}sin(n\textbackslash{}pi\ x/L),}

\begin{verbatim}
рей
\end{verbatim}

\texttt{a\_n\ =\ \textbackslash{}frac\{2\}\{L\}\ \textbackslash{}int\_\{-\textbackslash{}infty\}\^{}\{L\}\ \textbackslash{}overline\{\textbackslash{}phi\}(\textbackslash{}xi)\ \textbackslash{}sin(\textbackslash{}pi\ n\textbackslash{}xi/L)\ d\textbackslash{}xi.}

\begin{verbatim}
Це все для сердечників із фіксованою температурою на кінцях. А як щодо більш реалістичних GU з правою стороною, залежною від часу? Основні ідеї тут ті ж, що й у попередній задачі, але вони набагато складніші.

# Перетворення граничних умов, залежних від часу, у нульові

Розглянемо типову задачу

(PPP)
\end{verbatim}

\texttt{u\_t\ =\ \textbackslash{}alpha\^{}2\ u\_\{xx\}}

\begin{verbatim}
~~
, 0 < x < L , 0 < t < \infty , 0 < t < \infty , 0 < t < \infty , 0 < t < \infty , 0 < t < \infty , 0 < t < \infty , 0 < t < \infty , 0 < t < \infty , 0 < t < \infty , 0 < t < \infty , 0 < t < \infty , 0 < t < \infty , 0 < t < \infty , 0 < t < \infty , 0 < t < \infty , 0 < t < \infty , 0 < t < \infty , 0 < t < \infty , 0 < t < \infty , 0 < t < \infty , 0 < t < \infty , 0 < t < \infty , 0 < t < \infty , 0 < t < \infty , 0 < t < \infty , 0 < t < \infty , , 0 < t < \infty , 0 < t < \infty , 0 < t < \infty , 0 < t < \infty , 0 < t < \infty , 0 < t < \infty , 0 < t < \infty , 0 < t < \infty , 0 < t < \infty , 0 < t < \infty , 0 < t < \infty , 0 < t < \infty , 0 < t < \infty , 0 < t < \infty , 0 < t < \infty , 0 < t < \infty , 0 < t < \infty , 0 < t < \infty , 0 < t < \infty , 0 < t < \infty , 0 < t < \infty , 0 < t < \infty , 0 < t < \infty , 0 < t < \infty , 0 < t < \infty , 0 < t < \infty , 0 < t < \infty , 0 < t < \infty , 0 < t < \infty , 0 < t < \infty , 0 < t < \infty , 0 < t < \infty , 0 < t < \infty , 0 < t < \infty , 0 < t < \infty , 0 < t < \infty , 0 < t < \infty , 0 < t < \infty , 0 < t < \infty , 0 < t < \infty , 0 < t < \infty , 0 < t < \infty , 0 < t < \infty , 0 < t < \infty , 0 < t < \infty , 0 < t < \infty , 0 < t < \infty , 0 < t < \infty , 0 < t < \infty , 0 < t < \infty , 0 < t < \infty , 0 < t < \infty , 0 < t < \infty , 0 < t < \infty , 0 < t < \infty , 0 < t < \infty , 0 < t < \infty , 0 < t < \infty , 0 < t < \infty , 0 < t < \infty , 0 < t < \infty , 0 < t < \infty , 0 < t < \infty , 0 < t < \infty , 0 < t < \infty , 0 < t < \infty , 0 < t < \infty , 0 < t < \infty , 0 < t < \infty , 0 < t < \infty , 0 < t < \infty , 0 < t < \infty , 0 < t < \infty , 0 < t < \infty , 0 < t < \infty , 0 < t < \infty , 0 < t < \infty , 0 < t < \infty , 0 < t < \infty , 0 < t < \infty , 0 < t < \infty , 0 < t < \infty , 0 < t < \infty , 0 < t < \infty , 0 < t < \infty , 0 < t < \infty , 0 < t < \infty , 0 < t < \infty , 0 < t < \infty , 0 < t < \infty , 0 < t < \infty , 0 < t < \infty , 0 < t < \infty , 0 < t < \infty , 0 < t < \infty , 0 < t < \infty , 0 < t < \infty , 0 < t < \infty , 0 < t < \infty , 0 < t < \infty , 0 < t < \infty , 0 < t < \infty , 0 < t < \infty , 0 < t < \infty , 0 < t < \infty , 0 < t < \infty , 0 < t < \infty , 0 < t < \infty , 0 < t < \infty , 0 < t < \infty , , 0 < t < \infty , 0 < t < \infty , 0 < t < \infty , 0 < t < \infty , 0 < t < \infty , 0 < t < \infty , 0 < t < \infty , 0 < t < \infty , 0 < t < \infty , 0 < t < \infty , 0 < t < \infty , 0 < t < \infty , 0 < t < \infty , 0 < t < \infty , 0 < t < \infty , 0 < t < \infty , 0 < t < \infty , 0 < t < \infty , 0 < Щоб перетворити ці граничні умови на нульові умови, ми (після деяких спроб і помилок) зупинилися на такій формі розв'язку:
~~

(6.7)
\end{verbatim}

u(x, t) = A(t){[}1-x/L{]} + B(t){[}x/L{]} + U(x, t),

\begin{verbatim}
~~
де функції A\left(t\right) і B\left(t\right) обираються так, що «квазі-стаціонарна» частина розв'язку (6.7)
~~

(6.8)
\end{verbatim}

S(x, t) = A(t){[}1 - x/L{]} + B(t){[}x/L{]}

\begin{verbatim}
~~
задовольняло граничні умови початкової задачі. У цьому випадку функція U\left(x,\ t\right) задовольнятиме однорідні граничні умови.
~~

Підставляючи функцію S(x, t) у граничні умови
\end{verbatim}

S(0, t) = g\_1(t),

\begin{verbatim}
S_x(L, t) + hS(L, t) = g_2(t) призводить до двох рівнянь, з яких можна визначити A(t) і B(t). У результаті ми отримуємо

(6.9)
\end{verbatim}

\texttt{A(t)\ =\ g\_1(t),\ B(t)\ =\ \textbackslash{}frac\{g\_1(t)\ +\ Lg\_2(t)\}\{1\ +\ Lh\}.}

\begin{verbatim}
Отже.
\end{verbatim}

\texttt{u(x,\ t)\ =\ g\_1(t)\ {[}1\ -\ x/L{]}\ +\ \textbackslash{}frac\{g\_1(t)\ +\ Lg\_2(t)\}\{1\ +\ Lh\}\ {[}x/L{]}\ +\ U(x,\ t).}

\begin{verbatim}
Якщо ми підставимо ці вирази для you(x, t) у початкову задачу (це просять читача), отримаємо нову задачу для невідомої функції U(x, t):

~~
\begin{array}{llllllllllllllllllllllllllllllllllll
~~

Тепер ми стикаємося з новою проблемою однорідних граничних умов (на жаль, рівняння стало гетерогенним). Ми не можемо розв'язати цю задачу методом розділення змінних, але якщо читач погоджується чекати, то за кілька лекцій він ознайомиться з розв'язком цієї задачі методом цілісних перетворень і розкладання за власними функціями.

# ЗАУВАЖЕННЯ .

1. Як зазначалося раніше, мета цієї лекції — навчити читача, як перетворювати задачі з неунірольними граничними умовами на задачі з нульовими граничними умовами. Після цього може статися, що нове диференціальне рівняння в похідних похідних стане однорідним. Якщо це станеться, ми можемо вважати себе щасливцями, оскільки в цьому випадку проблему можна розв'язати методом розділення змінних (див. перший приклад з цієї лекції). Якщо перетворене рівняння виявиться гетерогенним, тоді нову задачу потрібно розв'язати іншим методом.

2. Лінійні неоднорідні граничні умови найзагальної форми
\end{verbatim}

\texttt{\textbackslash{}alpha\_1\ u\_x(0,\ t)\ +\ \textbackslash{}beta\_1\ u(0,\ t)\ =\ g\_1(t),}

\begin{verbatim}
~~
\alpha_2 u_x(L, t) + \beta_2 u(L, t) = g_2(t) Ви також можете перетворити умови на нульові граничні умови, використовуючи метод, описаний у другому прикладі. Звісно, нове рівняння, ймовірно, буде гетерогенним.
~~

- 3. Деякі методи розв'язку не містять вимог щодо унікальності граничних умов і, отже, не потребують попередніх перетворень. Пізніше, коли ми ознайомимося з перетворенням Лапласа, побачимо, що нульові граничні умови є необов'язковими (у випадку нульових граничних умов розв'язок дещо простіший, ніж в інших випадках).

- 4. Для граничних умов типу
\end{verbatim}

u(0, t) = g\_1(t),

\begin{verbatim}
u(L, t) = g_2(t) Метод, описаний у другому прикладі, призводить до наступної форми розв'язання задач:
\end{verbatim}

\texttt{u(x,\ t)\ =\ \textbackslash{}left\textbackslash{}\{\ g\_1(t)\ +\ \textbackslash{}frac\{x\}\{L\}\ \textbackslash{}left{[}\ g\_2(t)\ -\ g\_1(t)\ \textbackslash{}right{]}\ \textbackslash{}right\textbackslash{}\}\ +\ U(x,\ t).}

\begin{verbatim}
# ЗАВДАННЯ

1. Знайти розв'язок наступної змішаної задачі:
\end{verbatim}

\texttt{\textbackslash{}begin\{array\}\{lll\}\ (\textbackslash{}text\{УЧП\})\ \&\ u\_t\ =\ \textbackslash{}alpha\^{}2\ u\_\{xx\}\ +\ x\ \textbackslash{}cos\ t,\ \&\ 0\ \textless{}\ x\ \textless{}\ 1,\ \textbackslash{}\textbackslash{}\ (\textbackslash{}Gamma\ \textbackslash{}text\{У\})\ \&\ \textbackslash{}begin\{cases\}\ u\ \textbackslash{}\ (0,\ \textbackslash{}\ t)\ =\ 1,\ \textbackslash{}\textbackslash{}\ u\_x\ \textbackslash{}\ (1,\ \textbackslash{}\ t)\ +\ hu\ \textbackslash{}\ (1,\ \textbackslash{}\ t)\ =\ 1,\ \textbackslash{}end\{cases\}\ \&\ 0\ \textless{}\ t\ \textless{}\ \textbackslash{}infty,\ \textbackslash{}\textbackslash{}\ (\textbackslash{}text\{HY\})\ \&\ u\ \textbackslash{}\ (x,\ \textbackslash{}\ 0)\ =\ \textbackslash{}sin\ (\textbackslash{}pi\ x),\ \&\ 0\ \textbackslash{}leqslant\ x\ \textbackslash{}leqslant\ 1,\ \textbackslash{}end\{array\}}

\begin{verbatim}
попередньо перетворюючи його на задачу з нульовими граничними умовами. Чи відповідає рішення, яке ви отримуєте, вашим задумам щодо цієї проблеми?

Завдання

(PPP)
\end{verbatim}

\texttt{u\_t\ =\ u\_\{xx\},\ \textbackslash{}quad\ 0\ \textless{}\ x\ \textless{}\ 1,}

\begin{verbatim}
(GI)
\end{verbatim}

\texttt{\textbackslash{}begin\{cases\}\ u(0,\ t)\ =\ 0,\ \&\ 0\ \textless{}\ t\ \textless{}\ \textbackslash{}infty,\ \textbackslash{}\textbackslash{}\ u(1,\ t)\ =\ 1,\ \&\ 0\ \textless{}\ t\ \textless{}\ \textbackslash{}infty,\ \textbackslash{}end\{cases\}}

\begin{verbatim}
~~
(НУ) u(x, 0) = x^2, \quad 0 \le x \le 1, Перетворіть у задачу з нульовими граничними умовами і розв'яжіть її. Як виглядатиме рішення у різних
~~

Миті часу? Чи відповідає це вашим інтуїтивним ідеям? Що таке стаціонарне рішення? Як виглядає перехідне рішення?

3. Завдання
\end{verbatim}

\texttt{\textbackslash{}begin\{aligned\}}
\texttt{(\textbackslash{}text\{УЧП\})\textbackslash{}quad\ \&\ u\_t\ =\ u\_\{xx\},\ \&\&\ 0\ \textless{}\ x\ \textless{}\ 1,\textbackslash{}\ 0\ \textless{}\ t\ \textless{}\ \textbackslash{}infty,\ \textbackslash{}\textbackslash{}}
\texttt{(\textbackslash{}text\{ГУ\})\textbackslash{}quad\ \&\ u\_x(0,t)\ =\ 0,\textbackslash{}quad\ u\_x(1,t)\ +\ h\ u(1,t)\ =\ 1,\ \&\&\ 0\ \textless{}\ t\ \textless{}\ \textbackslash{}infty,\ \textbackslash{}\textbackslash{}}
\texttt{(\textbackslash{}text\{НУ\})\textbackslash{}quad\ \&\ u(x,0)\ =\ \textbackslash{}sin(\textbackslash{}pi\ x),\ \&\&\ 0\ \textbackslash{}leqslant\ x\ \textbackslash{}leqslant\ 1.}
\texttt{\textbackslash{}end\{aligned\}}

\begin{verbatim}
Перетворимо на задачу з нульовими граничними умовами. Чи буде нове рівняння однорідним?

Лекція 7

# РОЗВ'ЯЗАННЯ СКЛАДНІШИХ ЗАДАЧ ШЛЯХОМ РОЗДІЛЕННЯ ЗМІННИХ

МЕТА ЛЕКЦІЇ: Показати, як більш складні проблеми теплопровідності можна розв'язати методом розділення змінних. У цій лекції ми головним чином розглядаємо такі задачі, розв'язання яких дозволить читачеві краще ознайомитися з методом ділення змінних. Сподіваємося, що читач зможе екстраполювати викладені тут ідеї для власних завдань.

У тій же лекції читач ознайомиться з задачею за її власними функціями, або, як її ще називають, проблемою Штурма-Ліувіля. Обговорюються деякі властивості проблеми Штурма-Ліувіля.

Як згадувалося вище, мета цієї лекції — розв'язати змішану задачу шляхом розділення змінних. Можливо, саме таким методом читач вирішить свою проблему. Водночас ми сподіваємося, що він зможе перенести метод, розглянутий у цій лекції, для власних завдань.

Почнемо з одномірної задачі теплопровідності, в якій одна з граничних умов містить похідну.

# Задача теплопровідності з похідною в граничній умові

Розглянемо пристрій, показаний на рис. 7.1. Припустимо, що верхній кінець стрижня підтримується при нульовій температурі, тобто u(0, t) = 0, а нижній кінець занурений у воду, температура якої також стала і дорівнює нулю (у цьому випадку — нуль Р и с

![](images/_page_50_Picture_4.jpeg)

. 7.1. Схематичне представлення зсувної задачі.

![](images/_page_50_Figure_6.jpeg)

Рис. 7.2. Природа кривих, які задовольняють
\end{verbatim}

\texttt{(\textbackslash{}Gamma\ Y)\ =}
\texttt{\textbackslash{}begin\{cases\}\ u(0,\ t)\ =\ 0,\ \textbackslash{}\textbackslash{}\ u\_x(1,\ t)\ =\ -hu(1,\ t).\ \textbackslash{}end\{cases\}}

\begin{verbatim}
відповідає певній еталонній температурі). Теплообмін між стрижнем і водою відбувається згідно із законом Ньютона, який записується у x=1 як наступна гранична умова:
\end{verbatim}

u\_x(1, t) = -hu(1, t).

\begin{verbatim}
Тепер припустимо, що початкова температура стрижня дорівнює u(x, 0) = x, але одразу при t > 0 включаємо граничні умови. Щоб визначити температуру в стрижні, необхідно розв'язати проблему зміщення

~~
\begin{array}{llll} (\mathrm{УЧ}\Pi) & u_t = \alpha^2 u_{xx}, & 0 < x < 1, & 0 < t < \infty, \\ (7.1) & (\Gamma\mathrm{У}) & \begin{cases} u_t(0, t) = 0 \\ u_x(1, t) + hu_t(1, t) = 0 \end{cases} & \text{(однородные } \Gamma\mathrm{У}), \\ (\mathrm{H}\mathrm{У}) & u_t(x, 0) = x, & 0 \leq x \leq 1. \end{array}
~~

 Щоб використати метод розділення змінних, розбиємо процес розв'язання на кілька етапів.

**КРОК** 1. (Перетворення диференціального рівняння в похідних у два звичайних диференціальних рівняння)

Підставляючи u(x, t) = X(x)
\end{verbatim}

T(t)\texttt{у\ рівняння\ (7.1),\ отримуємо}XT' = \alpha\^{}2
X'\,'T\texttt{Поділивши\ інші\ частини\ цього\ рівняння\ на}\alpha\^{}2
XT отримуємо

\begin{verbatim}
\frac{T'}{\alpha^2 T} = \frac{X''}{X} .
\end{verbatim}

Оскільки ліва частина останнього рівняння залежить лише від часу, а
права --- лише від x (і оскільки x і t незалежні), обидві сторони
рівняння мають бути константами. Прирівнюючи обидві частини до однієї
сталої \texttt{\textbackslash{}mu} , отримуємо

(7.2)

\begin{verbatim}
T' - \mu \alpha^{2} T = 0, X'' - \mu X = 0.
\end{verbatim}

Процес розщеплення змінних завершено.

КРОК 2. (Визначення константи відокремлення)

\texttt{По-перше,\ константа\ і\ не\ повинна\ бути\ додатною,\ інакше\ \textquotesingle{}\^{}\{1\}\textquotesingle{}\ )\ функція\ \textquotesingle{}T\textbackslash{}left(\ t\textbackslash{}right)\textquotesingle{}\ буде\ зростати\ експоненційно,\ що\ призведе\ до\ нескінченного\ підвищення\ температури\ u=XT.\ Очевидно,\ що\ таке\ рішення\ не\ має\ фізичного\ значення.}

По-друге, припустимо, що \texttt{\textbackslash{}mu\ =\ 0} . Тоді

\begin{verbatim}
X'' = 0
\end{verbatim}

і, отже.

\begin{verbatim}
X(x) = A + \hat{B}x.
\end{verbatim}

Однак, оскільки граничні умови вимагають, що

\begin{verbatim}
u(0, t) = X(0) T(t) = 0,
\end{verbatim}

\texttt{u(1,\ t)\ =\ (hX(1)\ +\ X\textquotesingle{}(1))\ T(t)\ =\ 0,} Ми
одразу дійшли висновку, що

\begin{verbatim}
X(0) = 0
\end{verbatim}

\texttt{\textquotesingle{}\textbackslash{}Rightarrow\ A\ =\ 0\textquotesingle{}\ \textquotesingle{}X\textquotesingle{}(1)\ +\ hX(1)\ =\ 0\textquotesingle{}\ \textquotesingle{}\textbackslash{}Rightarrow\ B\ =\ 0\textquotesingle{}\ .}

Це означає, що \texttt{u(x,\ t)\ \textbackslash{}equiv\ 0} . Іншими
словами, лише нульове розв'язання відповідає нульовій константі
розділення. Ми відкинемо його, оскільки нас цікавлять ненульові
розв'язки.

Якщо, нарешті, \texttt{\textbackslash{}mu\ \textless{}\ 0} , можна
поставити \texttt{\textbackslash{}mu\ =\ -\textbackslash{}lambda\^{}2} і
переписати рівняння (7.2) у вигляді

\begin{verbatim}
T' + \lambda^2 \alpha^2 T = 0,
X' + \lambda^2 X = 0.
\end{verbatim}

1) Розв'язок \texttt{X(x)} відповідної задачі з граничним значенням буде
лише тривіальним, \texttt{X(x)} = 0. --- Ед. Ед.

Загальні розв'язки цих рівнянь легко знайти:

\begin{verbatim}
T(t) = Ae^{-(\lambda \alpha)^2 t},
X(x) = B\sin(\lambda x) + C\cos(\lambda x).
\end{verbatim}

Отже, будь-яка функція провідника

(7.3)

\begin{verbatim}
u(x, t) = e^{-(\lambda \alpha)^2 t} [A \sin(\lambda x) + B \cos(\lambda x)]
\end{verbatim}

при будь-яких значеннях \texttt{\textbackslash{}lambda} A і B
задовольняє рівняння (читач може перевірити це самостійно простою
заміною).

Тепер нам потрібно з'ясувати, які з цих функцій задовольняють граничним
умовам

(7.4)

\begin{verbatim}
u(0, t) = 0, u_x(1, t) + hu(1, t) = 0.
\end{verbatim}

Підстановка розв'язків виду (7.3) у граничні умови (7.4) дозволяє знайти
відношення, яким мають задовольняти значення
\texttt{\textbackslash{}lambda} , \texttt{\textbackslash{}Lambda} та B:

\begin{verbatim}
Be^{-(\lambda\alpha)^2 t} = 0 \Rightarrow B = 0,
\end{verbatim}

\texttt{\textquotesingle{}A\textbackslash{}lambda\ e\^{}\{-(\textbackslash{}lambda\textbackslash{}alpha)\^{}2\ t\}\ \textbackslash{}cos\ \textbackslash{}lambda\ +\ hAe\^{}\{-(\textbackslash{}lambda\textbackslash{}alpha)\^{}2\ t\}\ \textbackslash{}sin\ \textbackslash{}lambda\ =\ 0.\textquotesingle{}\ Після\ простих\ алгебранхічних\ перетворень\ останнє\ рівняння\ зводиться\ до\ вигляду}

\begin{verbatim}
tg \lambda = -\lambda/\hbar
\end{verbatim}

.

Іншими словами, щоб визначити \texttt{\textbackslash{}lambda} , потрібно
знайти точки перетину графів функцій tg \texttt{\textbackslash{}lambda}
і --- \texttt{\textbackslash{}lambda/h} (рисунок 7.3).

Числові значення значень \texttt{\textbackslash{}lambda\_1} ,
\texttt{\textbackslash{}lambda\_2} , \ldots{} можна знайти за допомогою
комп'ютера, якщо задано значення h Р и с

\pandocbounded{\includegraphics[keepaspectratio]{images/_page_52_Figure_15.jpeg}}

. 7.3. Перетин графів функцій \texttt{tg(\textbackslash{}lambda)} і
\texttt{-\textbackslash{}lambda/h} .

викликати власні значення задачі граничного значення

(7.5)

\begin{verbatim}
X'' + \lambda^{2}X = 0,
\begin{cases} X(0) = 0 \\ X'(1) + hX(1) = 0. \end{cases}
\end{verbatim}

\begin{verbatim}
Іншими словами, з цими значеннями параметра \lambda існують нетривіальні розв'язки задачі (7.5). Власні значення задачі (7.5) визначаються як корені рівняння \log \lambda = -\lambda/h . Їхні числові значення знаходилися на комп'ютері для випадку h = 1. Таблиця 7.1 показує результати розрахунків перших кількох власних значень.
\end{verbatim}

\begin{verbatim}
| n | \lambda_B | rt . | \lambda_n |
\end{verbatim}

\textbar-------------\textbar----------------------\textbar--------\textbar----------------\textbar{}

123 \textbar{} 2,024,917,98 \textbar{} 45 \textbar{} 11,0814,20
\textbar{}

\begin{verbatim}
Розв'язки задачі (7.5), що відповідають власним значенням \lambda_n , називаються власними функціями X_n(x) . Для розглянутого завдання власні функції виглядають так:
\end{verbatim}

\begin{verbatim}
X_n(x) = \sin(\lambda_n x)
\end{verbatim}

і зображені на рис. 7.4.

\pandocbounded{\includegraphics[keepaspectratio]{images/_page_53_Figure_10.jpeg}}

Рис. 7.4. Власні функції \texttt{X\_n}(x) завдання (7.5) при h=1; Усі
власні функції задовольняють граничні умови.

КРОК 3, (Пошук фундаментальних рішень)

Тепер у нас є нескінченний набір функцій (фундаментальних розв'язків) у
нашому розпорядженні

\begin{verbatim}
u_n(x, t) = X_n(x) T_n(t) = e^{-(\lambda_n x)^2 t} \sin(\lambda_n x),
\end{verbatim}

кожна з яких задовольняє рівняння та граничні умови. Залишився ще один
крок --- загалом

\begin{verbatim}
u(x, t) = \sum_{n=1}^{\infty} a_n X_n(x) T_n(t) =
= \sum_{n=1}^{\infty} a_n e^{-(\lambda_n c t^2 t)} \sin(\lambda_n x)
\end{verbatim}

Вибрати коефіцієнти \texttt{a\_n} так, щоб функція \texttt{u(x,\ t)}
задовольняла початкову умову \texttt{u(x,\ 0)} = x (у цьому випадку
функція \texttt{u(x,\ t)} задовольнятиме рівняння та граничні умови,
оскільки, по-перше, кожен член суми їх задовольняє, а по-друге, і
рівняння, і граничні умови є лінійними та однорідними).

Отже, ймовірність \texttt{a\_n} слід визначати з цієї умови

(7 6)

\begin{verbatim}
u(x, 0) = x = \sum_{n=1}^{\infty} a_n \sin(\lambda_n x).
\end{verbatim}

Ми маємо зробити останній крок.

КРОК 4. (Розклад початкової умови на ряд за власними функціями)

Щоб знайти коефіцієнти \texttt{a\_n} у розкладі (7,6) на власні функції,
помножте обидві частини цього співвідношення на
\texttt{\textbackslash{}sin(\textbackslash{}lambda\_m\ x)} і інтегруйте
від 0 до 1. У результаті ми отримуємо

\begin{verbatim}
\int_{0}^{2} \xi \sin(\lambda_{m} \xi) d\xi = \sum_{n=1}^{\infty} a_{n} \int_{0}^{1} \sin(\lambda_{n} \xi) \sin(\lambda_{m} \xi) d\xi =
= a_{m} \int_{0}^{1} \sin^{2}(\lambda_{m} \xi) d\xi =
= a_{m} \left( \frac{\lambda_{m} - \sin \lambda_{m} \cos \lambda_{m}}{2\lambda_{m}} \right).
\end{verbatim}

Розв'язавши останнє відношення відносно \texttt{a\_m} і замінивши m на
n, отримуємо бажаний результат:

(7.7)

\begin{verbatim}
a_n = \frac{2\lambda_n}{(\lambda_n - \sin \lambda_n \cos \lambda_n)} \int_0^{\infty} \xi \sin (\lambda_n \xi) d\xi.
\end{verbatim}

Таблиця 7.2.

\section{\texorpdfstring{Шанси \texttt{a\_n} у формулі
(7.8)}{Шанси a\_n у формулі (7.8)}}\label{ux448ux430ux43dux441ux438-a_n-ux443-ux444ux43eux440ux43cux443ux43bux456-7.8}

{\def\LTcaptype{none} % do not increment counter
\begin{longtable}[]{@{}llll@{}}
\toprule\noalign{}
n & an & n & \texttt{a\_n} \\
\midrule\noalign{}
\endhead
\bottomrule\noalign{}
\endlastfoot
123 & 0,240,22-0,03 & 45 & -0.11-0.09  \\
\end{longtable}
}

Отже, розв'язок задачі (7.1) записується у вигляді

(7.8)

\begin{verbatim}
u(x, t) = \sum_{n=1}^{\infty} a_n e^{-i\lambda_n \alpha_n^{2} t} \sin(\lambda_n x),
\end{verbatim}

де коефіцієнти \texttt{a\_n} задаються формулами (7.7). Перші п'ять
коефіцієнтів \texttt{a\_n} нашої задачі були знайдені чисельно. Їхні
значення наведені в Таблиці 1. 7.2.

\pandocbounded{\includegraphics[keepaspectratio]{images/_page_55_Figure_8.jpeg}}

Рис. 7.5. Графічне представлення функції (7.8); межа умовно
задовольняється у всі моменти часу.

Відповідно, перші три члени розв'язання змішаної задачі

(7.9)

\begin{verbatim}
\begin{aligned} u_t &= u_{xx}, & 0 < x < 1, & 0 < t < \infty, \\ u_t &= 0, & 0 < t < \infty, \\ u_x &= 0, & 0 < t < \infty, \\ u_x &= 0, & 0 < t < \infty, \\ u_x &= 0, & 0 < t < \infty, \\ u_x &= 0, & 0 < t < \infty, \\ u_x &= 0, & 0 < t < \infty, \\ u_x &= 0, & 0 < t < \infty, \\ u_x &= 0, & 0 < t < \infty, \\ u_x &= 0, & 0 < t < \infty, \\ u_x &= 0, & 0 < t < \infty, \\ u_x &= 0, & 0 < t < \infty, \\ u_x &= 0, & 0 < t < \infty, \\ u_x &= 0, & 0 < t < \infty, \\ u_x &= 0, & 0 < t < \infty, \\ u_x &= 0, & 0 < t < \infty, \\ u_x &= 0, & 0 < t < \infty, \\ u_x &= 0, & 0 < t < \infty, \\ u_x &= 0, & 0 < t < \infty, \\ u_x &= 0, & 0 < t < \infty, \\ u_x &= 0, & 0 < t < \infty, \\ u_x &= 0, & 0 < t < \infty, \\ u_x &= 0, & 0 < t < \infty, \\ u_x &= 0, & 0 < t < \infty, \\ u_x &= 0, & 0 < t < \infty, \\ u_x &= 0, & 0 < t < \infty, \\ u_x &= 0, & 0 < t < \infty, \\ u_x &= 0, & 0 < t < \infty, \\ u_x &= 0, & 0 < t < \infty, \\ u_x &= 0, & 0 < t < \infty, \\ u_x &= 0, & 0 < t < \infty, \\ u_x &= 0, & 0 < t < \infty, \\ u_x &= 0, & 0 < t < \infty, \\ u_x &= 0, & 0 < t < \infty, \\ u_x &= 0, & 0 < t < \infty, \\ u_x &= 0, & 0 < t < \infty, \\ u_x &= 0, & 0 < t < \infty, \\ u_x &= 0, & 0 < t < \infty, \\ u_x &= 0, & 0 < t < \infty, \\ u_x &= 0, & 0 < t < \infty, \\ u_x &= 0, & 0 < t < \infty, \\ u_x &= 0, & 0 < t < \infty, \\ u_x &= 0, & 0 < t < \infty, \\ u_x &= 0, & 0 < t < \infty, \\ u_x &= 0, & 0 < t < \infty, \\ u_x &= 0, & 0 < t < \infty, \\ u_x &= 0, & 0 < t < \infty, \\ u_x &= 0, & 0 < t < \infty, \\ u_x &= 0, & 0 < t < \infty, \\ u_x &= 0, & 0 < t < \infty, \\ u_x &= 0, & 0 < t < \infty, \\ u_x &= 0, & 0 < t < \infty, \\ u_x &= 0, & 0 < t < \infty, \\ u_x &= 0, & 0 < t < \infty, \\ u_x &= 0, & 0 < t < \infty, \\ u_x &= 0, & 0 < t < \infty, \\ u_x &= 0, & 0 < t < \infty, \\ u_x &= 0, & 0 < t < \infty, \\ u_x &= 0, & 0 < t < \infty, \\ u_x &= 0, & 0 < t < \infty, \\ u_x &= 0, & 0 < t < \infty, \\ u_x &= 0, & 0 < t < \infty, \\ u_x &= 0, & 0 < t < \infty, \\ u_x &= 0, & 0 < t < \infty, \\ u_x &= 0, & 0 < t < \infty, \\ u_x &= 0, & 0 < t < \infty, \\ u_x &= 0, & 0 < t < \infty, \\ u_x &= 0, & 0 < t < \infty, \\ u_x &= 0, & 0 < t < \infty, \\ u_x &= 0, & 0 < t < \infty, \\ u_x &= 0, & 0 < t < \infty, \\ u_x &= 0, & 0 < t < \infty, \\ u_x &= 0, & 0 < t < \infty, \\ u_x &= 0, & 0 < t < \infty, \\ u_x &= 0, & 0 < t < \infty, \\ u_x &= 0, & 0 < t < \infty, \\ u_x &= 0, & 0 < t < \infty, \\ u_x &= 0, & 0 < t < \infty, \\ u_x &= 0, & 0 < t < \infty, \\ u_x &= 0, & 0 < t < \infty, \\ u_x &= 0, & 0 < t \end{aligned}
\end{verbatim}

мають вигляд

\begin{verbatim}
u(x, t) = = 0.24e^{-4t} \sin(2.02x) + 0.22e^{-24t} \sin(4.9x) - 0.03e^{-48.3t} \sin(7.98x) + \dots
\end{verbatim}

Фіг. На рисунку 7-5 показані графіки цього розв'язку для кількох точок
часу. Рекомендуємо читачеві відповісти на такі питання: Чи відповідає
тип рішення, яке ви отримуєте, з вашими інтуїтивними ідеями? Чи
виконуються граничні умови?

\section{ЗАУВАЖЕННЯ}\label{ux437ux430ux443ux432ux430ux436ux435ux43dux43dux44f-4}

Задача власних значень (7.5) є особливим випадком більш загальної
задачі, яку зазвичай називають проблемою Штурма-Ліувіля:

(7.10)

\begin{verbatim}
\begin{aligned} & (\text{ОДУ}) \quad [p(x)y']' - q(x)y + \lambda r(x)y = 0, \quad 0 < x < 1, \\ & (\alpha_1 y(0) + \beta_1 y'(0) = 0, \\ & (\alpha_2 y(1) + \beta_1 y'(1) = 0. \end{aligned}
\end{verbatim}

Якщо розв'язати задачу з однорідними граничними умовами методом
розділення змінних, то функції \texttt{X(x)} завжди знаходяться як
розв'язки деякої задачі Штурма--Ліувіля. Очевидно, що задача власних
значень (7.5) є особливим випадком задачі (7.10).

Перелічимо основні властивості задачі (7.10), які були встановлені
Штурмом і Ліувілем за певних природних обмежень.

функції \texttt{p(x)}, \texttt{q(x)} та \texttt{r(x)}.

\begin{enumerate}
\def\labelenumi{\arabic{enumi}.}
\tightlist
\item
  Існує нескінченна послідовність власних значень, що задовольняють
  нерівності
\end{enumerate}

\begin{verbatim}
\lambda_1 < \lambda_2 < \lambda_3 < \ldots < \lambda_n < \ldots \to \infty
\end{verbatim}

.

\begin{enumerate}
\def\labelenumi{\arabic{enumi}.}
\setcounter{enumi}{1}
\item
  Кожне власне значення \texttt{\textbackslash{}lambda\_n} відповідає
  одній (точній до константного множника) ненульовій власній функції
  \texttt{y\_n(x)} .
\item
  Якщо \texttt{y\_n(x)} і \texttt{y\_m(x)} дві різні власні функції, що
  відповідають різним власним значенням
  \texttt{\textbackslash{}lambda\_n} і
  \texttt{\textbackslash{}lambda\_m} , то вони є ортогональними з вагою
  \texttt{r(x)} на сегменті {[}0, 1{]}, тобто вони задовольняють це
  відношення
\end{enumerate}

\begin{verbatim}
\int_{0}^{1} r(x) y_{n}(x) y_{m}(x) dx = 0.
\end{verbatim}

Більше деталей про проблему Штурма-Ліувіля можна знайти в літературі,
наведеній наприкінці книги.

1) Див., наприклад, Тіхоног А. Н., Самарський А. А. Рівняння
математичної фізики. --- Москва: Наука, 1972. --- ред. Ед.

\section{ЗАВДАННЯ}\label{ux437ux430ux432ux434ux430ux43dux43dux44f-3}

\begin{enumerate}
\def\labelenumi{\arabic{enumi}.}
\tightlist
\item
  Розв'язати наступну задачу для рівняння теплопровідності методом
  розкладу
\end{enumerate}

(PPP)

\begin{verbatim}
u_1 = u_{xx}, \quad 0 < x < 1, \quad 0 < t < \infty,
\end{verbatim}

(GI)

\begin{verbatim}
\begin{cases} u(0, t) = 0, \\ u_x(1, t) = 0, \\ u(x, 0) = x, \quad 0 \le x \le 1. \end{cases}
\end{verbatim}

Чи відповідає рішення, яке ви знаходите, з інтонативними
представленнями? Що таке стаціонарне рішення?

\begin{enumerate}
\def\labelenumi{\arabic{enumi}.}
\setcounter{enumi}{1}
\tightlist
\item
  Знайдіть власні цінності та функції проблеми Ассолта-Ліувіля
\end{enumerate}

(ODY)

\begin{verbatim}
X' + \lambda X = 0, \quad 0 < x < 1,
\end{verbatim}

(FY)

\begin{verbatim}
\begin{cases} X(0) = 0, \\ X'(1) = 0. \end{cases}
\end{verbatim}

Задайте специфічні вирази для функцій \texttt{p(x)}, \texttt{q(x)} та
\texttt{r(x)} із загальної задачі Штурма--Ліувіля для цього випадку.

\begin{enumerate}
\def\labelenumi{\arabic{enumi}.}
\setcounter{enumi}{2}
\tightlist
\item
  Розв'язати задачу теплопровідності для планки з термоізольованими
  межами
\end{enumerate}

(PPP)

\begin{verbatim}
u_t = u_{xx}, \quad 0 < x < 1, \quad 0 < t < \infty,
\end{verbatim}

(GI)

\begin{verbatim}
\begin{cases} u_x (0, t) = 0, \\ u_x (1, t) = 0, \end{cases} \quad 0 < t < \infty,
\end{verbatim}

(ГУ)
\texttt{u(x,\ 0)\ =\ x,\ \textbackslash{}quad\ 0\ \textbackslash{}le\ x\ \textbackslash{}le\ 1.}
Чи відповідає рішення, яке ви отримуєте, з вашою інтерпретацією цієї
проблеми? Що таке стаціонарне рішення? Яке його фізичне значення?

\begin{enumerate}
\def\labelenumi{\arabic{enumi}.}
\setcounter{enumi}{3}
\tightlist
\item
  Знайдіть власні цінності та функції завдань
\end{enumerate}

(ОДА)

\begin{verbatim}
X'' + \lambda X = 0, \quad 0 < x < 1,
\end{verbatim}

(GI)

\begin{verbatim}
\begin{cases} X'(0) = 0, \\ X'(1) = 0. \end{cases}
\end{verbatim}

\#\# ПЕРЕТВОРЕННЯ СКЛАДНИХ РІВНЯНЬ У ПРОСТУ ФОРМУ

МЕТА ЛЕКЦІЇ Показати, як PPP для функції \texttt{u(x,\ t)} можна
перетворити на простіший PPP для нової функції \texttt{w(x,\ t)}. Вибір
правильної трансформації зазвичай базується на ненавмисному намірі. У
нашій лекції

\begin{verbatim}
u_t = \alpha^2 u_{xx} - \beta u,
u_t = \alpha^2 u_{xx} - v u_x, \quad (v = \text{const})
\end{verbatim}

переходимо до найпростішого рівняння теплопровідності

\begin{verbatim}
w_1 = \alpha^2 w_{xx}
\end{verbatim}

внаслідок відповідних перетворень

\begin{verbatim}
u(x, t) = e^{-\beta t} w(x, t)
u(x, t) = e^{\epsilon t (1 - \delta t - 2) 2\alpha^2} w(x, t).
\end{verbatim}

Розв'язавши просте рівняння теплопровідності за функцією
\texttt{\textbackslash{}omega(x,t)} , ми можемо знайти розв'язки
початкових рівнянь, перетворивши функцію
\texttt{\textbackslash{}omega(x,t)} відповідно до цих формул (звісно, ці
перетворення слід застосовувати і до початкових і граничних умов).

Після прочитання двох останніх лекцій читач може створити враження, що
метод розділення змінних застосовується \emph{лише} до рівняння

\begin{verbatim}
u_1 = \alpha^2 u_{xx}
\end{verbatim}

Насправді рівняння теплопотоку є найпростішим, але не єдиним
параболічним рівнянням, яке можна розв'язати методом розділення змінних.
Як уже згадувалося, якщо рівняння є лінійним і однорідним, його можна
розв'язати, поділивши змінні. Наприклад, двовимірна задача
теплопровідності в колі описується рівнянням

\begin{verbatim}
u_t = \alpha^2 \left[ u_{rr} + \frac{1}{r} u_r + \frac{1}{r^2} u_{\theta\theta} \right],
\end{verbatim}

І хоча це рівняння з \emph{змінними коефіцієнтами}, його можна звести до
трьох звичайних диференціальних рівнянь.

У цій лекції ми покажемо читачеві, що іноді, замість того, щоб
безпосередньо розв'язувати диференціальне рівняння в похідних, краще
перетворити його на простіше рівняння, яке можна розв'язати методом
розділення змінних або іншим методом. Наступний приклад ілюструє цей
підхід.

\section{Перетворення проблеми з перенесенням бічної поверхні у проблему
з термічно ізольованою бічною
поверхнею}\label{ux43fux435ux440ux435ux442ux432ux43eux440ux435ux43dux43dux44f-ux43fux440ux43eux431ux43bux435ux43cux438-ux437-ux43fux435ux440ux435ux43dux435ux441ux435ux43dux43dux44fux43c-ux431ux456ux447ux43dux43eux457-ux43fux43eux432ux435ux440ux445ux43dux456-ux443-ux43fux440ux43eux431ux43bux435ux43cux443-ux437-ux442ux435ux440ux43cux456ux447ux43dux43e-ux456ux437ux43eux43bux44cux43eux432ux430ux43dux43eux44e-ux431ux456ux447ux43dux43eux44e-ux43fux43eux432ux435ux440ux445ux43dux435ux44e}

Розглянемо наступне завдання:

(8.1)

\begin{verbatim}
u_{t} = \alpha^{2}u_{xx} - \beta u, \quad 0 < x < 1, \quad 0 < t < \infty,
\begin{cases} u(0, t) = 0, \\ u(1, t) = 0, \\ u(x, 0) = \varphi(x), \quad 0 \le x \le 1, \end{cases}
\end{verbatim}

(8.1)

\begin{verbatim}
(FY) \quad
\begin{cases} u(0, t) = 0, \\ u(x, 0) = \varphi(x), \quad 0 \le x \le 1, \end{cases}
\end{verbatim}

де термін \texttt{-\textbackslash{}beta\ u} описує потік тепла через
бічну поверхню стрижня (рис. 8.1).

\pandocbounded{\includegraphics[keepaspectratio]{images/_page_59_Figure_7.jpeg}}

\texttt{Рис.\ 81.\ Задача\ теплопровідності,\ описана\ рівнянням\ \textquotesingle{}u\_t\ =\ \textbackslash{}alpha\^{}t\ u\_\{\textbackslash{}lambda\textbackslash{}lambda\}\ -\ \textbackslash{}beta\ u\textquotesingle{}\ :\ \textquotesingle{}\textbackslash{}alpha\textquotesingle{}\ —\ це\ втрата\ гепле\ через\ поверхню\ Боха;\ \textquotesingle{}\textbackslash{}delta\textquotesingle{}\ –\ дифузія\ вздовж\ стрижня.}

\texttt{Наша\ мета\ —\ ввести\ нову\ температурну\ \textquotesingle{}w\textbackslash{}left(x,\textbackslash{}\ t\textbackslash{}right)\textquotesingle{}\ замість\ \textquotesingle{}u\textbackslash{}left(x,\textbackslash{}\ t\textbackslash{}right)\textquotesingle{}\ ,\ щоб\ рівняння\ для\ \textquotesingle{}w\textbackslash{}left(x,\textbackslash{}\ t\textbackslash{}right)\textquotesingle{}\ стало\ простішим\ за\ початкове\ рівняння}

\begin{verbatim}
u_t = \alpha^2 u_{xx} - \beta u
\end{verbatim}

.

Це досить загальний метод розв'язання диференціальних рівнянь у
частинних похідних, але вибір відповідного перетворення базується на
інтуїтивних уявленнях про поведінку розв'язків початкового рівняння.
Наприклад, у задачі (8.1) температура \texttt{u(x,\ t)} у кожній точці
\texttt{x\_0} змінюється внаслідок двох факторів:

\begin{itemize}
\tightlist
\item
  \begin{enumerate}
  \def\labelenumi{\arabic{enumi})}
  \tightlist
  \item
    дифузія тепла вздовж стрижня (через компонентний
    \texttt{\textbackslash{}alpha\^{}2\ u\_\{xx\}} ),
  \end{enumerate}
\item
  \begin{enumerate}
  \def\labelenumi{\arabic{enumi})}
  \setcounter{enumi}{1}
  \tightlist
  \item
    теплопередача через бічну поверхню стрижня (через член RI).
  \end{enumerate}
\end{itemize}

Найважливішим елементом нашого міркування є наступне: якщо не було
дифузії вздовж стрижня ( \texttt{\textbackslash{}alpha\ =\ 0} ),

тоді температура в кожній точці стрижня експоненційно знижується до нуля
згідно із законом

\begin{verbatim}
u(x_0 t) = u(x_0, 0) e^{-\beta t}
\end{verbatim}

.

Виходячи з цього спостереження, запитаємо, чи можна розв'язок проблеми
(8.1) представити як утворення двох факторів

\begin{verbatim}
(8.2) u(x, t) = e^{-\beta t}w(x, t)
\end{verbatim}

або:

\begin{verbatim}
\begin{pmatrix}
\end{verbatim}

Температура елемента з \textbackslash{}

Неізольовані боком

\begin{verbatim}
Поверхня \end{pmatrix}
\end{verbatim}

де \texttt{w(x,\ t)} описує розподіл температури, що виникає лише через
дифузію. Давайте подивимось, що станеться, якщо підставити цей вираз для
температури у задачу (8.1). Після виконання відповідних розрахунків
(читач може зробити їх сам), ми отримаємо

Нова задача майже точно така ж, як і оригінальна, але більше не має
додатку 6u. Тому замість розв'язання задачі (8.1) ми розв'язуємо
перетворену задачу (8.3), а потім множимо отриманий розв'язок
\texttt{w(x,\ t)} на \texttt{e\^{}\{-\textbackslash{}beta\ t\}} . У
нашому випадку розв'язок задачі (8.3) вже було отримано раніше методом
розділення змінних, тому ми можемо одразу записати

(8.4)

\begin{verbatim}
w(x, t) = \sum_{n=1}^{\infty} a_n e^{-(n\pi\alpha)^{\frac{n}{2}}t} \sin(n\pi x),
a_n = 2 \int_0^{\infty} \phi(\xi) \sin(n\pi \xi) d\xi,
\end{verbatim}

і, отже, розв'язок початкової задачі (8.1) має вигляд

\begin{verbatim}
u(x, t) = e^{-\beta t}w(x, t).
\end{verbatim}

Нижче, у розділі «Нотатки», наведено ще один приклад проблеми,
розв'язаної цим методом.

\section{ПРИМІТКИ}\label{ux43fux440ux438ux43cux456ux442ux43aux438}

\begin{enumerate}
\def\labelenumi{\arabic{enumi}.}
\tightlist
\item
  Вирішити проблему
\end{enumerate}

\begin{verbatim}
\begin{aligned}
u_t &= u_{xx} - u, && 0 < x < 1,\ 0 < t < \infty, \\
u(0,t) &= 0,\quad u(1,t) = 0, && 0 < t < \infty, \\
u(x,0) &= \sin(\pi x) + 0.5\sin(3\pi x), && 0 \leqslant x \leqslant 1.
\end{aligned}
\end{verbatim}

За методом, описаним вище, ми

\begin{itemize}
\tightlist
\item
  \begin{enumerate}
  \def\labelenumi{\alph{enumi})}
  \tightlist
  \item
    тимчасово відкинути конвективний член u у початковому рівнянні.
  \end{enumerate}
\item
  \begin{enumerate}
  \def\labelenumi{\alph{enumi})}
  \setcounter{enumi}{1}
  \tightlist
  \item
    записати розв'язок змішаної задачі без конвективного члена
  \end{enumerate}
\end{itemize}

\begin{verbatim}
u(x, t) = e^{-\pi^2 t} \sin(\pi x) + 0.5e^{-(3\pi)^2 t} \sin(3\pi x),
\end{verbatim}

\begin{verbatim}
c) помножимо це розв'язання на конвективний фактор e^{-\beta t}=e^{-t} отримаємо остаточне розв'язання початкової задачі
\end{verbatim}

\begin{verbatim}
u(x, t) = e^{-t} \left[ e^{-\pi^2 t} \sin(\pi x) + 0.5e^{-(3\pi)^2 t} \sin(3\pi x) \right].
\end{verbatim}

\begin{enumerate}
\def\labelenumi{\arabic{enumi}.}
\setcounter{enumi}{1}
\tightlist
\item
  Рівняння конвективної дифузії
\end{enumerate}

\begin{verbatim}
u_t = \alpha^2 u_{xx} - v u_x
\end{verbatim}

(коефіцієнт і є константою) також можна перетворити у вигляд

\begin{verbatim}
w_t = \alpha^2 w_{xx}.
\end{verbatim}

Відповідне перетворення має вигляд

\begin{verbatim}
u(x, t) = e^{v(x-vt/2)/2\alpha^2}w(x, t).
\end{verbatim}

\begin{verbatim}
Це перетворення підкреслює експоненціальний фактор, який повністю зумовлений рухом середовища. Зверніть увагу, що цей множник описує хвилю, що поширюється зліва направо зі швидкістю \upsilon/2 . Читач матиме можливість скористатися цією трансформацією.
\end{verbatim}

\section{ЗАВДАННЯ}\label{ux437ux430ux432ux434ux430ux43dux43dux44f-4}

\begin{enumerate}
\def\labelenumi{\arabic{enumi}.}
\tightlist
\item
  Розв'язати задачу конвективної дифузії
\end{enumerate}

\begin{verbatim}
\begin{aligned}
(\text{УЧП})\quad & u_t = u_{xx} - u_x, && 0 < x < 1,\ 0 < t < \infty, \\
(\text{ГУ})\quad & u(0,t) = 0,\quad u(1,t) = 0, && 0 < t < \infty, \\
(\text{НУ})\quad & u(x,0) = e^{x/2}, && 0 \leqslant x \leqslant 1.
\end{aligned}
\end{verbatim}

методом перетворення до простішого рівняння. На що схожий розв'язок?
Можна дати таку інтерпретацію цієї задачі: \texttt{u(x,t)} ---
концентрація деякої речовини в рухомому середовищі, яке рухається зліва
направо зі швидкістю v = 1; початковий розподіл концентрації задається
функцією \texttt{e\^{}\{x/2\}}, а на межах концентрація підтримується
нульовою. Чи відповідає ваш розв'язок такій інтерпретації?

\begin{enumerate}
\def\labelenumi{\arabic{enumi}.}
\setcounter{enumi}{1}
\tightlist
\item
  Розв'язати задачу
\end{enumerate}

\begin{verbatim}
\begin{aligned}
(\text{УЧП})\quad & u_t = u_{xx} - u + x, && 0 < x < 1,\ 0 < t < \infty, \\
(\text{ГУ})\quad & u(0,t) = 0,\quad u(1,t) = 1, && 0 < t < \infty, \\
(\text{НУ})\quad & u(x,0) = 0, && 0 \leqslant x \leqslant 1.
\end{aligned}
\end{verbatim}

\begin{itemize}
\item
  а) перетворивши неоднорідні граничні умови на однорідні;
\item
  б) перетворивши вихідне рівняння на нове рівняння без члена -u;
\item
  в) розв'язавши задачу, що виникає після цих перетворень.
\end{itemize}

\begin{enumerate}
\def\labelenumi{\arabic{enumi}.}
\setcounter{enumi}{2}
\tightlist
\item
  Розв'язати задачу
\end{enumerate}

\begin{verbatim}
\begin{aligned}
(\text{УЧП})\quad & u_t = u_{xx} - u, && 0 < x < 1,\ 0 < t < \infty, \\
(\text{ГУ})\quad & u(0,t) = 0,\quad u(1,t) = 0, && 0 < t < \infty, \\
(\text{НУ})\quad & u(x,0) = \sin(\pi x), && 0 \leqslant x \leqslant 1.
\end{aligned}
\end{verbatim}

застосовуючи метод розділення змінних безпосередньо до вихідного
рівняння, не виконуючи жодних попередніх перетворень. Чи збігається
отриманий розв'язок з розв'язком, знайденим раніше після заміни

\begin{verbatim}
u(x,t) = e^{-t}w(x,t)
\end{verbatim}

?

\section{РОЗВ'ЯЗАННЯ НЕОДНОРІДНИХ УЧП МЕТОДОМ РОЗКЛАДУ ЗА ВЛАСНИМИ
ФУНКЦІЯМИ}\label{ux440ux43eux437ux432ux44fux437ux430ux43dux43dux44f-ux43dux435ux43eux434ux43dux43eux440ux456ux434ux43dux438ux445-ux443ux447ux43f-ux43cux435ux442ux43eux434ux43eux43c-ux440ux43eux437ux43aux43bux430ux434ux443-ux437ux430-ux432ux43bux430ux441ux43dux438ux43cux438-ux444ux443ux43dux43aux446ux456ux44fux43cux438}

ЦІЛЬ ЛЕКЦІЇ: Навчитися розв'язувати змішані задачі вигляду

\begin{verbatim}
\begin{aligned}
(\text{УЧП})\quad & u_t = \alpha^2 u_{xx} + f(x,t), && 0 < x < 1,\ 0 < t < \infty, \\
(\text{ГУ})\quad & \alpha_1 u_x(0,t) + \beta_1 u(0,t) = 0, && 0 < t < \infty, \\
& \alpha_2 u_x(1,t) + \beta_2 u(1,t) = 0, && 0 < t < \infty, \\
(\text{НУ})\quad & u(x,0) = \varphi(x), && 0 \leqslant x \leqslant 1.
\end{aligned}
\end{verbatim}

Розв'язок неоднорідного УЧП можна шукати у вигляді ряду

\begin{verbatim}
u(x,t) = \sum_{n=1}^{\infty} T_n(t) X_n(x).
\end{verbatim}

де \texttt{X\_n}(x) --- власні функції відповідної однорідної задачі

\begin{verbatim}
\begin{aligned}
(\text{УЧП})\quad & u_t = \alpha^2 u_{xx}, && 0 < x < 1,\ 0 < t < \infty, \\
(\text{ГУ})\quad & \alpha_1 u_x(0,t) + \beta_1 u(0,t) = 0, && 0 < t < \infty, \\
& \alpha_2 u_x(1,t) + \beta_2 u(1,t) = 0, && 0 < t < \infty, \\
(\text{НУ})\quad & u(x,0) = \varphi(x), && 0 \leqslant x \leqslant 1.
\end{aligned}
\end{verbatim}

А функції \texttt{T\_n}(t) визначаються як розв'язки деяких звичайних
диференціальних рівнянь.

На лекції 6 ми ознайомилися з тим, як перетворювати неоднорідні граничні
умови на однорідні. На жаль, рівняння залишається гетерогенним, і ми
стоїмо перед завданням

\begin{verbatim}
(9.1) \quad u_{t} = \alpha^{2} u_{xx} + f(x, t), \quad 0 < x < 1, \quad 0 < t < \infty,
(9.1) \quad (\Gamma Y) \quad
\begin{cases} \alpha_{1} u_{x}(0, t) + \beta_{1} u(0, t) = 0, \\ \alpha_{2} u_{x}(1, t) + \beta_{2} u(1, t) = 0, \\ u(x, 0) = \varphi(x), \quad 0 \le x \le 1. \end{cases}
\end{verbatim}

Мета нашої лекції --- опанувати метод розв'язання таких задач. Цей метод
аналогічний методу варіації довільних констант.

теорія звичайних диференціальних рівнянь називається методом розкладу за
власними функціями.

Ідея методу надзвичайно проста. Відомо, що розв'язок задачі (9.1) при
\texttt{f(x,\ t)} = 0 (так звана відповідна однорідна задача)
записується як ряд

\begin{verbatim}
u\left(x,\ t\right)=\sum_{n=1}^{\infty}a_{n}e^{-\left(\lambda_{n}\alpha\right)^{2}}/X_{n}\left(x\right),
\end{verbatim}

де \texttt{X\_n}(x) --- власні функції задачі Штурма-Ліувіля,

(9.2)

\begin{verbatim}
X'' + \lambda^2 X = 0, \\ \alpha_1 X'(0) + \beta_1 X(0) = 0, \\ \alpha_2 X'(1) + \beta_2 X(1) = 0.
\end{verbatim}

Запитаймо себе, чи можна також шукати розв'язок гетерогенної проблеми
(9.1) у вигляді серії власних функцій, але більш загальної форми

\begin{verbatim}
u(x, t) = \sum_{n=1}^{\infty} T_n(t) X_n(x).
\end{verbatim}

Існують певні фізичні підстави для такого припущення. Насправді терміни
\texttt{f(x,t)} описують джерело тепла, розташоване всередині стрижня.
Це означає, що залежність рішення від часу не буде виражена показниками

\begin{verbatim}
e^{-(\lambda_n\alpha)^2 t}
\end{verbatim}

,

Це стосувалося однорідного завдання, але з деякими іншими функціями.

Щоб показати цей метод у дії і не ускладнювати подання деталями, ми
застосуємо його до розв'язання однієї дуже простої задачі.

\section{Построзв'язання методом демаркації за власними
функціями,}\label{ux43fux43eux441ux442ux440ux43eux437ux432ux44fux437ux430ux43dux43dux44f-ux43cux435ux442ux43eux434ux43eux43c-ux434ux435ux43cux430ux440ux43aux430ux446ux456ux457-ux437ux430-ux432ux43bux430ux441ux43dux438ux43cux438-ux444ux443ux43dux43aux446ux456ux44fux43cux438}

Розглянемо гетерогенну проблему

(9.3) (PY)

\begin{verbatim}
\begin{aligned} u_t &= \alpha^2 u_{xx} + f(x, t), & 0 < x < 1, & 0 < t < \infty, \\ u(0, t) &= 0, & 0 < t < \infty, \\ u(1, t) &= 0, & 0 < t < \infty, \\ u(x, 0) &= \varphi(x), & 0 \leqslant x \leqslant 1. \end{aligned}
\end{verbatim}

Процес розв'язання задачі подається у вигляді послідовності наступних
кроків.

КРОК 1. Основна ідея методу полягає в розкладанні густини джерела
\texttt{f(x,\ t)} на ряд власними функціями

\begin{verbatim}
f(x, t) = f_1(t) X_1(x) + f_2(t) X_2(x) + \dots + f_n(t) X_n(x) + \dots
\end{verbatim}

\begin{verbatim}
а також визначити реакції системи u_n\left(x,\,t\right) на вплив кожного компонента f_n\left(t\right)X_n\left(x\right) . Якщо знайти всі відповіді, рішення проблеми виглядатиме так:
\end{verbatim}

\begin{verbatim}
u(x, t) = \sum_{n=1}^{\infty} u_n(x, t).
\end{verbatim}

Головна складність цього методу --- розклад густини джерела на
компоненти \texttt{f\_n}(t),\texttt{X\_n}(x) . Виявляється, що
\texttt{X\_n}(x) факторами в цій задачі є власними векторами системи
Штурма--Ліувіля, які виникають при розв'язанні відповідної однорідної
задачі методом розділення змінних, а саме

\begin{verbatim}
\begin{array}{c} u_t = \alpha^2 u_{xx} \quad \text{(отметим, что } f(x, \ t) = 0), \\ u(0, \ t) = 0, \\ u(1, \ t) = 0, \\ u(x, \ 0) = \varphi(x). \end{array}
\end{verbatim}

Завдання Штурма-Ліувіля в цьому випадку виглядає так:

\begin{verbatim}
X'' + \lambda^2 X = 0,
\end{verbatim}

\texttt{X(0)} = 0, \texttt{X(1)} = 0, а його розв'язки є функціями

\begin{verbatim}
X_n(x) = \sin(n\pi x), \quad n = 1, 2, 3 \dots
\end{verbatim}

Отже, розклад на густину джерела представлений як ряд

\begin{verbatim}
(9.5) \ f(x,t) = f_1(t) \sin(\pi x) + f_2(t) \sin(2\pi x) + \dots + f_n(t) \sin(n\pi x) + \dots
\end{verbatim}

\begin{verbatim}
\mathcal{U} нарешті, щоб знайти функції f_n(t) , помножте обидві сторони останнього відношення на \sin(m\pi x) і інтегруйте від 0 до 1 за координатою x.
\end{verbatim}

У результаті ми отримуємо

\begin{verbatim}
\int_{0}^{1} f(x, t) \sin(m\pi x) dx = \sum_{n=1}^{\infty} f_{n}(t) \int_{0}^{1} \sin(m\pi x) \sin(n\pi x) dx =
= \frac{1}{2} f_{m}(t)
\end{verbatim}

або (після заміни t на p)

(9.6)

\begin{verbatim}
f_{n}(t) = 2 \int_{0}^{1} f(x, t) \sin(n\pi x) dx.
\end{verbatim}

Останнє співвідношення встановлює співвідношення між коефіцієнтами
\texttt{f\_n(t)} і часткою джерела тепла \texttt{f(x,\ t)}.

КРОК 2. (Пошук \texttt{u\_n(x,\ t)\ =\ X\_n(x)\ T\_n(t)} реакції на
вхідний вплив \texttt{f\_n(t)\ X\_n(x)} .)

Замінимо густину джерела на його розклад у рядок

\begin{verbatim}
f(x, t) = \sum_{n=1}^{\infty} f_n(t) \sin(n\pi x)
\end{verbatim}

і намагатися знайти індивідуальні відповіді, тобто функції,
\texttt{T\_\{\textbackslash{}mathbf\{n\}\}(t)} у

\begin{verbatim}
u(x, t) = \sum_{n=1}^{\infty} T_n(t) \sin(n\pi x).
\end{verbatim}

Оскільки розв'язок задачі має вигляд

\begin{verbatim}
u(x, t) = \sum_{n=1}^{\infty} T_n(t) \sin(n\pi x),
\end{verbatim}

потім підставимо цей вираз у початкову задачу

\begin{verbatim}
u_{t} = \alpha^{2} u_{xx} + \sum_{n=1}^{\infty} f_{n}(t) \sin(n\pi x),
u(0, t) = 0,
u(1, t) = 0,
u(x, 0) = \varphi(x),
\end{verbatim}

У результаті ми отримаємо

\begin{verbatim}
\sum_{n=1}^{\infty} T'_n(t) \sin(n\pi x) = -\alpha^2 \sum_{n=1}^{\infty} (n\pi)^2 T_n(t) \sin(n\pi x) + + \sum_{n=1}^{\infty} f_n(t) \sin(\pi n\lambda),
\end{verbatim}

(9.7)

\begin{verbatim}
\sum_{n=1}^{\infty} T_n(t) \sin 0 = 0
\end{verbatim}

\begin{verbatim}
(задовольняється ідентично), \sum_{n=1}^{\infty} T_n(t) \sin (n\pi) = 0 (задовольняється ідентично), \sum_{n=1}^{\infty} T_n(0) \sin (n\pi x) = \phi(x) .
\end{verbatim}

Переписування рівняння та початкової умови

(PPP)

\begin{verbatim}
\sum_{n=1}^{\infty} [T'_n + (n\pi\alpha)^2 T_n - f_n(t)] \sin(n\pi x) = 0,
\end{verbatim}

(НУ)

\begin{verbatim}
\sum_{n=1}^{\infty} T_n(0) \sin(n\pi x) = \varphi(x),
\end{verbatim}

\begin{verbatim}
можна побачити, що функції T_n(\ell) є розв'язками задачі Коші
\end{verbatim}

(9.8)

\begin{verbatim}
T_{n} + (n\pi\alpha)^{2} T_{n} = f_{n}(t),
T_{n}(0) = 2 \int_{0}^{1} \varphi(\xi) \sin(n\pi\xi) d\xi = a_{n}.
\end{verbatim}

Усі ці TAC легко розв'язати, а їхні рішення записуються у вигляді

\begin{enumerate}
\def\labelenumi{(\arabic{enumi})}
\setcounter{enumi}{98}
\tightlist
\item
\end{enumerate}

\begin{verbatim}
T_n(t) = a_n e^{-(n\pi\alpha)^2 t} + \int_0^t e^{-(n\pi\alpha)^2 (t-\tau)} f_n(\tau) d\tau,
\end{verbatim}

З того місця, де легко отримати повне розв'язання проблеми (9.3):

(9.10)

\begin{verbatim}
u(x, t) = \sum_{n=1}^{\infty} T_n(t) \sin(n\pi x) =
= \sum_{n=1}^{\infty} \left[ a_n e^{-(n\pi\alpha)^n t} \sin(n\pi x) \right] + \sum_{n=1}^{\infty} \left[ \sin(n\pi x) \int_0^t e^{-(n\pi\alpha)^n (t-\tau)} f_n(\tau) d\tau \right]
\end{verbatim}

Перехідна частина (обумовлена стаціонарною частиною (умовлена початковою
умовою) правою частиною \texttt{f(x,\ t)}).

Розв'язок (9.10) показує, що температурна реакція складається з двох
частин: перша частина зумовлена початковими умовами, а друга ---
густиною джерела \texttt{f(x,t)}. Зверніть увагу, що термін «стаціонарна
частина» не є найкращим описом другої частини температурної реакції,
оскільки він не означає, що ця частина не залежить від часу
(стаціонарний стан може бути періодичним, якщо \texttt{f(x,\ t)} є
періодичною функцією часу t).

Проблема повністю вирішена. Перш ніж перейти до наступної теми, давайте
покажемо, як використовувати описаний вище метод для розв'язання

Конкретне завдання.

70

\section{Розв'язання конкретної задачі методом декомпозиції за власними
функціями}\label{ux440ux43eux437ux432ux44fux437ux430ux43dux43dux44f-ux43aux43eux43dux43aux440ux435ux442ux43dux43eux457-ux437ux430ux434ux430ux447ux456-ux43cux435ux442ux43eux434ux43eux43c-ux434ux435ux43aux43eux43cux43fux43eux437ux438ux446ux456ux457-ux437ux430-ux432ux43bux430ux441ux43dux438ux43cux438-ux444ux443ux43dux43aux446ux456ux44fux43cux438}

Розглянемо просту задачу

(9.11)

\begin{verbatim}
\begin{aligned} u_t &= \alpha^2 u_{xx} + \sin{(3\pi x)}, & 0 < x < 1, \\ (\mathbf{FY}) & \begin{cases} u(0, t) &= 0, \\ u(1, t) &= 0, \end{cases} & 0 < t < \infty, \\ (\mathbf{HY}) & u(x, 0) &= \sin{(\pi x)}, & 0 \leq x \leq 1. \end{aligned}
\end{verbatim}

Наша мета --- обчислити коефіцієнти \texttt{T\_n}(t) при розкладі

\begin{verbatim}
u\left(x,\,t\right)=\sum_{n=1}^{\infty}\,T_{n}\left(t\right)\sin\left(n\pi x\right)
\end{verbatim}

(Власні функції цієї задачі \texttt{X\_n}(x) збігатися з власними
функціями попереднього). Якщо підставити цю декомпозицію на початкову
задачу, отримаємо наступну задачу Коші для визначення функцій
\texttt{T\_n}(t) :

\begin{verbatim}
T_n + (n\pi\alpha)^2 T_n = \int_n (t) = 2 \int_0^1 \sin(3\pi x) \sin(n\pi x) dx =
\begin{cases} 1 &n = 3, \\ 0 & n \neq 3, \end{cases}
T_n(0) = 2 \int_0^1 \sin(n\xi) \sin(n\pi\xi) d\xi =
\begin{cases} 1 &n = 1, \\ 0 & n \neq 1. \end{cases}
\end{verbatim}

Перепишем эти уравнения последовательно для n=1, 2\ldots{}

\begin{verbatim}
\begin{aligned}
(n=1)\quad & T_1' + (\pi\alpha)^2 T_1 = 0, && T_1(0) = 1, && T_1(t) = e^{-(\pi\alpha)^2 t}, \\
(n=2)\quad & T_2' + (2\pi\alpha)^2 T_2 = 0, && T_2(0) = 0, && T_2(t) = 0, \\
(n=3)\quad & T_3' + (3\pi\alpha)^2 T_3 = 4, && T_3(0) = 0, && T_3(t) = \frac{1}{(3\pi\alpha)^2}\left[1 - e^{-(3\pi\alpha)^2 t}\right], \\
(n\geqslant 4)\quad & T_n' + (n\pi\alpha)^2 T_n = 0, && T_n(0) = 0, && T_n(t) = 0.
\end{aligned}
\end{verbatim}

Отже, розв'язок задачі можна записати у вигляді

\begin{verbatim}
u(x, t) = e^{-(\pi\alpha)^2 t} \sin(\pi x) + \frac{1}{(3\pi\alpha)^2} [1 - e^{-(3\pi\alpha)^2 t}] \sin(3\pi x).
\end{verbatim}

(9.12)

\section{ЗАУВАЖЕННЯ}\label{ux437ux430ux443ux432ux430ux436ux435ux43dux43dux44f-5}

\begin{enumerate}
\def\labelenumi{\arabic{enumi}.}
\item
  Метод розкладу власних функцій є одним із найпотужніших методів
  розв'язання неоднорідних диференціальних рівнянь у похідних похідних.
  Пізніше, коли ми ознайомимося з інтегральними перетвореннями,
  побачимо, що існують й інші методи розв'язання задач такого типу.
\item
  Власні функції \texttt{X\_n(x)} , за допомогою якої виконується
  розклад, залежать від задачі, яку потрібно розв'язати, точніше від
  рівнянь і граничних умов. Читача просять розв'язати задачу 4
  наприкінці лекції, щоб перевірити, чи дійсно він навчився знаходити
  власні функції \texttt{X\_n(x)} .
\item
  Якщо читач пам'ятає основні факти з теорії звичайних диференціальних
  рівнянь, то він знає, що розв'язки, що відповідають неоднорідним
  членам вигляду
\end{enumerate}

\begin{verbatim}
P_n(x) e^{\alpha x}
\begin{cases} \sin{(\beta x)}, \\ \cos{(\beta x)}, \end{cases}
\end{verbatim}

можно найти методом неопределенных коэффициентов. То же самое и здесь.
Задача (9.11) может быть решена таким же способом.

\section{ЗАДАЧИ}\label{ux437ux430ux434ux430ux447ux438}

\begin{enumerate}
\def\labelenumi{\arabic{enumi}.}
\item
  Подставьте разложение (46.8) в задачу (46.7) и получите
  последовательность задач
  \texttt{P\_0,\textbackslash{}\ P\_1,\textbackslash{}\ P\_2,\textbackslash{}\ \textbackslash{}dots}.
\item
  Покажите, что нелинейную задачу
\end{enumerate}

\begin{verbatim}
\begin{aligned}
\Delta u + u^2 &= 0, && 0 \leqslant r < 1,\ 0 \leqslant \theta < 2\pi, \\
u(1,\theta) &= \cos\theta, && 0 \leqslant \theta < 2\pi,
\end{aligned}
\end{verbatim}

можно свести к последовательности линейных задач
\texttt{P\_0,\textbackslash{}\ P\_1,\textbackslash{}\ P\_2}.

\begin{enumerate}
\def\labelenumi{\arabic{enumi}.}
\setcounter{enumi}{2}
\item
  Подставьте приближенное решение (46.5) в нелинейную задачу 2 и оцените
  его точность.
\item
  Решите задачу \texttt{P\_1} из примера с возмущенной границей и
  проверьте, насколько хорошо приближенное решение
\end{enumerate}

\begin{verbatim}
u(r,\theta) = u_0(r,\theta) + \frac{1}{4}u_1(r,\theta)
\end{verbatim}

удовлетворяет соотношениям

\begin{verbatim}
\Delta u = 0,
u\left(1 + \frac{1}{4}\sin\theta,\ \theta\right) = \cos\theta.
\end{verbatim}

\section{Лекція 47}\label{ux43bux435ux43aux446ux456ux44f-47}

\section{РОЗВ'ЯЗАННЯ ДИФЕРЕНЦІАЛЬНИХ РІВНЯНЬ У ПОХІДНИХ ПОХІДНИХ ЗА
ДОПОМОГОЮ МЕТОДУ КОНФОРМНИХ
ВІДОБРАЖЕНЬ}\label{ux440ux43eux437ux432ux44fux437ux430ux43dux43dux44f-ux434ux438ux444ux435ux440ux435ux43dux446ux456ux430ux43bux44cux43dux438ux445-ux440ux456ux432ux43dux44fux43dux44c-ux443-ux43fux43eux445ux456ux434ux43dux438ux445-ux43fux43eux445ux456ux434ux43dux438ux445-ux437ux430-ux434ux43eux43fux43eux43cux43eux433ux43eux44e-ux43cux435ux442ux43eux434ux443-ux43aux43eux43dux444ux43eux440ux43cux43dux438ux445-ux432ux456ux434ux43eux431ux440ux430ux436ux435ux43dux44c}

МЕТА ЛЕКЦІЇ: Показати, як деякі двовимірні граничні задачі можна
перетворити на інші, простіші задачі за допомогою конформного
відображення. Припустимо, наприклад, що ми хочемо розв'язати рівняння
Лапласа

з певною граничною умовою в області комплексної форми (площина змінних x
і y). Цю задачу можна перетворити на нову, в якій потрібно знайти
розв'язок рівняння Лапласа

.

Тут z --- це комплексна змінна (з деякої області площини z), а w ---
нова комплексна змінна, яка отримає значення відповідно до формули w =
\texttt{f(z)} у новій комплексній площині w.

Наприклад, комплексна функція, визначена в \emph{першому квадранті}
комплексної z-площини, \texttt{w=z\^{}2} відображає цю область у всю
півплощину v\textgreater0 у комплексній площині w.

Таблиця 47 2

{\def\LTcaptype{none} % do not increment counter
\begin{longtable}[]{@{}
  >{\raggedright\arraybackslash}p{(\linewidth - 4\tabcolsep) * \real{0.2342}}
  >{\raggedright\arraybackslash}p{(\linewidth - 4\tabcolsep) * \real{0.2252}}
  >{\raggedright\arraybackslash}p{(\linewidth - 4\tabcolsep) * \real{0.5405}}@{}}
\toprule\noalign{}
\begin{minipage}[b]{\linewidth}\raggedright
Площина d
\end{minipage} & \begin{minipage}[b]{\linewidth}\raggedright
\texttt{z\ \textbackslash{}longrightarrow\ z\^{}2}
\end{minipage} & \begin{minipage}[b]{\linewidth}\raggedright
Площина S
\end{minipage} \\
\midrule\noalign{}
\endhead
\bottomrule\noalign{}
\endlastfoot
0 & & 0 \\
t1 → t & & -1(2) \\
Вісь x Перший квадрант & & Позитивна дійсна піввісь верхньої
півплощини \\
\end{longtable}
}

Таблиця 47.2 показує відповідність, яка встановлюється шляхом
відображення \texttt{w=z\^{}2} між деякими точками на площинах z і w.

\begin{verbatim}
x^{2} + y^{2} = 1,
\end{verbatim}

Рис. 47.1. Площина комплексної змінної
\texttt{\textbackslash{}boldsymbol\{z\}} та некольорових формул бігунів
Р И С У Н

Проблема полягає в тому, щоб знайти конформне відображення, яке
перетворює цю область на щось простіше, де нашу проблему легко
розв'язати. У цьому випадку очевидно, що ця проста ділянка має бути
кільцем. Відповідна експозиція Р И С У

\pandocbounded{\includegraphics[keepaspectratio]{images/_page_359_Figure_10.jpeg}}

НОК 47.5. Конформне відображення на кільце.

Це вже не так очевидно, але існують довідники, де можна знайти дисплей
із вказаними властивостями. Одним із найзагальних є «Довідник конформних
відображень»

G. Kober 1), яка містить сотні конформних відображень між різними
регіонами. Якщо читач витратить трохи часу і краще ознайомиться з
конформними таблицями відображення, він знайде відповідне конформне
відображення для задачі (47.3). У нашому конкретному випадку це буде
образ виду

(47.4)

ОК 47.2. Відображення першого квадранта площини
\texttt{\{\textbackslash{}bf\ z\}} у верхню півплощину
\texttt{\{\textbackslash{}bf\ w\}} .

Інакше кажучи, якщо ми говоримо про комплексісні функції змінної w =
\texttt{f(z)}, то можна розглянути питання, як криві площини z
перетворюються на криві в площині

Це питання не ставиться в теорії функцій дійсної змінної, з якою читач
уже добре знайомий.

Перш ніж рухатися далі, давайте розглянемо, як \texttt{w=z\^{}2} складне
відображення типів у еквівалентній дійсній формі. Дійсна форма показує,
як початкові координати (x,y) у площині z перетворюються на нові
координати (u,v) у площині w \texttt{w=z\^{}2} .

І давайте окремо прирівняємо реальні та уявні частини в цьому
співвідношенні. У результаті ми отримуємо

\begin{verbatim}
v = \gamma [y(x-t)],
\end{verbatim}

(действительная форма отображения w=z\^{}2 ).

Эта форма понадобится нам позже. В этом примере она интересна тем, что
показывает, как гиперболы в плоскости z переходят в координатные линии u
= const и v = const в плоскости w.

Теперь мы перейдем к изучению одного частного вида функций комплексного
переменного, которые осуществляют конформные отображения.

\section{Определение конформного
отображения}\label{ux43eux43fux440ux435ux434ux435ux43bux435ux43dux438ux435-ux43aux43eux43dux444ux43eux440ux43cux43dux43eux433ux43e-ux43eux442ux43eux431ux440ux430ux436ux435ux43dux438ux44f}

\begin{verbatim}
Відображення w=f(z) комплексної площини z у комплексну площину w називається конформним у точці  z_{\bullet}  z-площини, якщо похідна  f'(z_{\bullet}) \neq 0  . Відображення f(z) називається конформним у регіоні D, якщо  f'(z) \neq 0  в кожній точці регіону D.
\end{verbatim}

\begin{verbatim}
Наприклад, відображення  f(z) = z^2  конформно всюди, крім точки z = 0, оскільки  z'(z) = 2z \neq 0  для всіх  z \neq 0  . З іншого боку, відображення  e^z  конформно у всій площині z, оскільки скрізь  f'(z) = e^z \neq 0  . Яке призначення конформних відображень? Відповідь полягає в тому, що при розв'язанні рівняння Лапласа  \phi_{xx} + \phi_{yy} = 0  в певній області площини змінних x і y можна розглядати цю площину як площину комплексної змінної z. Розглянемо комплексне відображення w = f(z) площини z у площину w, у якій введені координати (u, v). Відображуючи w = f(z  \phi_{xy} + \phi_{uy} = 0  ), рівняння Лапласа  \varphi_{xx} + \varphi_{yy} = 0  перетворюється у нове диференціальне рівняння в похідних похідних, яке залежить від нових координат і і v. To — рівняння, яке  \phi_{xx} + \phi_{uu} = 0  повернеться у рівняння  \varphi_{nn} + \varphi_{nn} = 0  . Таким чином, ми приходимо до ідеї знайти конформне відображення, яке переносить область з комплексною межею в область з простою межею (нагадаємо, що відображення  w=z^*  переносить межу першого квадранта площини z у дійсну вісь площини w).
\end{verbatim}

Приведем несколько примеров, в которых будет показано, как можно решать
уравнения с частными производными с помощью конформных отображений.

Уравнение Лапласа в верхней полуплоскости.

Предположим, что нам необходимо решить следующую задачу Дирихле в
верхней полувлоскости (рис. 47.3):

(47.1)

Один із способів розв'язання цієї проблеми --- застосувати перетворення
Фур'є до змінної x, але набагато краще конформно відобразити верхню
половину площини у площину (u, v). Читачеві не буде складно побачити, що
конформне відображення

\begin{verbatim}
\begin{aligned} \varphi_{uu} + \varphi_{vv} &= 0, & \text{в кольче,} \\ (\Gamma Y) & & & & & & & \\ (\Gamma Y) & & & & & & \\ (\Gamma Y) & & & & & \\ & & & & & \\ (\Gamma Y) & & & & & \\ & & & & \\ & & & & \\ & & & & \\ & & & & \\ & & & \\ & & & \\ & & & \\ & & & \\ & & & \\ & & & \\ & & \\ & & \\ & & \\ & & \\ & & \\ & & \\ & & \\ & & \\ & & \\ & & \\ & & \\ & & \\ & & \\ & & \\ & & \\ & & \\ & & \\ & & \\ & & \\ & & \\ & & \\ & & \\ & & \\ & & \\ & & \\ & & \\ & & \\ & & \\ & & \\ & & \\ & & \\ & & \\ & & \\ & & \\ & & \\ & & \\ & & \\ & & \\ & & \\ & & \\ & & \\ & & \\ & & \\ & & \\ & & \\ & & \\ & & \\ & & \\ & & \\ & & \\ & & \\ & & \\ & & \\ & & \\ & & \\ & & \\ & & \\ & & \\ & & \\ & & \\ & & \\ & & \\ & & \\ & & \\ & & \\ & & \\ & & \\ & & \\ & & \\ & & \\ & & \\ & & \\ & & \\ & & \\ & & \\ & & \\ & & \\ & & \\ & & \\ & & \\ & & \\ & & \\ & & \\ & & \\ & & \\ & & \\ & & \\ & & \\ & & \\ & & \\ & & \\ & & \\ & & \\ & & \\ & & \\ & & \\ & & \\ & & \\ & & \\ & & \\ & & \\ & & \\ & & \\ & & \\ & & \\ & & \\ & & \\ & & \\ & & \\ & & \\ & & \\ & & \\ & & \\ & & \\ & & \\ & & \\ & & \\ & & \\ & & \\ & & \\ & & \\ & & \\ & & \\ & & \\ & & \\ & & \\ & & \\ & & \\ & & \\ & & \\ & & \\ & & \\ & & \\ & & \\ & & \\ & & \\ & & \\ & & \\ & & \\ & & \\ & & \\ & & \\ & & \\ & & \\ & & \\ & & \\ & & \\ & & \\ & & \\ & & \\ & & \\ & & \\ & & \\ & & \\ & & \\ & & \\ & & \\ & & \\ & & \\ & & \\ & & \\ & & \\ & & \\ & & \\ & & \\ & & \\ & & \\ & & \\ & & \\ & & \\ & & \\ & & \\ & & \\ & & \\ & & \\ & & \\ & & \\ & & \\ & & \\ & & \\ & & \\ & & \\ & & \\ & & \\ & & \\ & & \\ & & \\ & & \\ & & \\ & & \\ & & \\ & & \\ & & \\ & & \\ & & \\ & & \\ & & \\ & & \\ & & \\ & & \\ & & \\ & & \\ & & \\ & & \\ & & \\ & & \\ & & \\ & & \\ & & \\ & & \\ & & \\ & & \\ & & \\ & & \\ & & \\ & & \\ & & \\ & & \\ & & \\ & & \\ & & \\ & & \\ & & \\ & & \\ & & \\ & & \\ & & \\ & & \\ & & \\ & & \\ & & \\ & & \\ & & \\ & & \\ & & \\ & & \\ & & \\ & & \\ & & \\ & & \\ & & \\ & & \\ & & \\ & & \\ & & \\ & & \\ & & \\ & & \\ & & \\ & & \\ & & \\ & & \\ & & \\ & & \\ & & \\ & & \\ & & \\ & & \\ & & \\ & & \\ & & \\ & & \\ & & \\ & & \\ & & \\ & & \\ & & \\ & & \\ & & \\ & & \\ & & \\ & & \\ & & \\ & & \\ & & \\ & & \\ & & \\ & & \\ & & \\ & & \\ & & \\ & & \\ & & \\ & & \\ & & \\ & & \\ & & \\ & & \\ & & \\ & & \\ & & \\ & & \\ & & \\ & & \\ & & \\ & & \\ & & \\ & & \\ & & \\ & & \\ & & \\ & & \\ & & \end{aligned}
\end{verbatim}

Вона має такі властивості (див. RNS. 47.4).

\begin{enumerate}
\def\labelenumi{\arabic{enumi}.}
\tightlist
\item
  Вона відображає верхню півплощину площини z у смуги
  \texttt{-\textbackslash{}infty\ \textless{}\ u\ \textless{}\ \textbackslash{}infty}
  , \texttt{0\ \textless{}\ v\ \textless{}\ \textbackslash{}pi} , у
  площині w.
\end{enumerate}

РИС 47.3 Задача граничного значення у верхній половині площини.

\begin{verbatim}
- 2. Ребро прямої y=0, -1 < x < 1, у площині z (де потенціал дорівнює 1), переходить у пряму  v=\pi,  -  \infty < u < \infty  , у площині w. - 3. Два промені x>1, y=0, і x<-1, y=0, у площині z проходять відповідно у позитивний і негативний
\end{verbatim}

Читач повинен бути впевнений, що може підтвердити ці твердження.

\begin{verbatim}
\varphi(u, v) = 0.57 \ln(u^2 + v^2) + 1
\end{verbatim}

Рис. 47.4. Перетворення складної задачі на легку за допомогою
конформного відображення.

Важливість цих властивостей полягає в тому, що вони перетворюють
початкову вадаху (47.1) на дуже просту задачу

(47.2)

\begin{verbatim}
\varphi(x, y) = 0.57 \ln(u^2 + v^2) + 1,
\end{verbatim}

Оказывается, в этом случае бегущие волны отражаются от границ таким
образом, что результирующие колебания становятся не бегущими, а
сохраняющими форму в одном положении, т. е.

\begin{verbatim}
w = \ln \left[ \frac{z-1}{z+1} \right].
\end{verbatim}

РИС. 20.1. Три типичных стоячих волны \texttt{X(x)} \texttt{T(t)}.

Декілька прикладів стоячих хвиль показано на рис. 20.1. Якщо профілі цих
стоячих хвиль відомі \texttt{X\_n}(x) і природа коливань кожної хвилі
\texttt{T\_n}(t) , то розв'язок проблеми коливань гітарних струн є
суперпозицією найпростіших коливань \texttt{X\_n}(x),\texttt{T\_n}(t) і

сывается в виде

\begin{verbatim}
z = re^{i\theta}
\end{verbatim}

\begin{verbatim}
Единственное, что остается сделать нам,— это подобрать коэффициенты  c_n  так, чтобы при t=0 удовлетворялись начальные условия  u\left(x,\ 0\right)=f\left(x\right)  и  u_t\left(x,\ 0\right)=g\left(x\right)  .
\end{verbatim}

Задачу о колебаниях гитарной струны мы решим методом разделения
переменных.

\section{Решение задачи о колебаниях ограниченной струны методом
разделения
переменных}\label{ux440ux435ux448ux435ux43dux438ux435-ux437ux430ux434ux430ux447ux438-ux43e-ux43aux43eux43bux435ux431ux430ux43dux438ux44fux445-ux43eux433ux440ux430ux43dux438ux447ux435ux43dux43dux43eux439-ux441ux442ux440ux443ux43dux44b-ux43cux435ux442ux43eux434ux43eux43c-ux440ux430ux437ux434ux435ux43bux435ux43dux438ux44f-ux43fux435ux440ux435ux43cux435ux43dux43dux44bux445}

Чтобы решить смешанную задачу

(УЧП)

где \texttt{a\_n} , \texttt{b\_n} , \texttt{R\_n} и
\texttt{\textbackslash{}delta\_n} ---произвольные постоянные. Заметим,
что элементарные колебания удовлетворяют граничным условиям. Читателю
должно быть ясно, что каждое элементарное колебание представляет собой
стоячую волну. В стоячей волне все точки колеблются с одной и той же
частотой, а форма волны подобна волнам, изображенным на рис. 20.3.

\pandocbounded{\includegraphics[keepaspectratio]{images/_page_150_Picture_5.jpeg}}

РИС. 20.3. Стоячие волны \texttt{u\_n(x,\ t)\ =\ X\_n(x)\ T\_n(t)} .

В нашей задаче уравнение и граничные условия линейны и однородны,
поэтому произвольная сумма элементарных стоячих волн также является
решением волнового уравнения и удовлетворяет граничным условиям. Мы
должны теперь найти такую суперпозицию стоячих волн, чтобы она
удовлетворяла также и начальным условиям, тогда получим решение нашей
задачи. Подстановка суммы

\begin{verbatim}
\begin{cases} u(0, t) = 0, \\ u(L, t) = 0, \end{cases}
\end{verbatim}

\pandocbounded{\includegraphics[keepaspectratio]{images/_page_362_Figure_1.jpeg}}

\pandocbounded{\includegraphics[keepaspectratio]{images/_page_362_Figure_2.jpeg}}

якщо потенціал на внутрішньому колі дорівнює одинці, а на зовнішньому
колі дорівнює двом (рис. 47.5). Іншими словами, нам потрібно розв'язати
проблему

(47.3)

\pandocbounded{\includegraphics[keepaspectratio]{images/_page_362_Figure_5.jpeg}}

де s=-0,146, а t=-6,85. У своїй еквівалентній дійсній формі це буде

(47.5)

\pandocbounded{\includegraphics[keepaspectratio]{images/_page_362_Figure_7.jpeg}}

в начальные условия

\begin{verbatim}
\begin{cases} 0, x < a, \\ 1, x \geqslant a, \end{cases}
\end{verbatim}

найдем сначала стоячне волны, т. е. решения вида

Подставив это выражение в уравнение (20.1) и разделив переменные,
получим два обыкновенных дифференциальных уравнения

\begin{verbatim}
H(x-a) =
\end{verbatim}

Функція Хевісайду,

\begin{verbatim}
H(a-x) = \begin{cases} 1, & x \leqslant a, \\ 0, & x > a, \end{cases}
\end{verbatim}

дзеркальною функцією Хевісайду.

ТАБЛИЦЯ A. Експоненціальне перетворення Фур'є

{\def\LTcaptype{none} % do not increment counter
\begin{longtable}[]{@{}
  >{\raggedright\arraybackslash}p{(\linewidth - 2\tabcolsep) * \real{0.3797}}
  >{\raggedright\arraybackslash}p{(\linewidth - 2\tabcolsep) * \real{0.6203}}@{}}
\toprule\noalign{}
\begin{minipage}[b]{\linewidth}\raggedright
1. \texttt{f\textquotesingle{}(x)}
\end{minipage} & \begin{minipage}[b]{\linewidth}\raggedright
iωF (ω)
\end{minipage} \\
\midrule\noalign{}
\endhead
\bottomrule\noalign{}
\endlastfoot
2. \texttt{f\textquotesingle{}\textquotesingle{}(x)} &
\texttt{-\textbackslash{}omega\^{}2\ F(\textbackslash{}omega)} \\
3. fn (x), (n-та похідна) &
\texttt{(i\textbackslash{}omega)\^{}n\ F(\textbackslash{}omega)} \\
4. \texttt{f(ax),\ a\ \textgreater{}\ 0} &
\texttt{\textbackslash{}frac\{1\}\{a\}\ F\textbackslash{}left(\textbackslash{}frac\{\textbackslash{}omega\}\{a\}\textbackslash{}right)} \\
5. \texttt{f(x-a)} &
\texttt{e\^{}\{-ia\textbackslash{}omega\}F(\textbackslash{}omega)} \\
6. \texttt{e\^{}\{-a\^{}2x\^{}2\}} &
\texttt{\textbackslash{}frac\{1\}\{a\textbackslash{}sqrt\{2\}\}e\^{}\{-\textbackslash{}omega\^{}2/4a\^{}2\}} \\
7. \texttt{e\^{}\{-a\_1\ x\}} &
\texttt{\textbackslash{}sqrt\{\textbackslash{}frac\{2\}\{\textbackslash{}pi\}\}\ \textbackslash{}frac\{a\}\{a\^{}2\ +\ \textbackslash{}omega\^{}2\}} \\
\end{longtable}
}

\begin{enumerate}
\def\labelenumi{\arabic{enumi}.}
\setcounter{enumi}{19}
\item
  гріх (ax)
\item
\end{enumerate}

\begin{verbatim}
\begin{cases} 0, & m \neq n, \\ L/2, & m = n, \end{cases}
\end{verbatim}

\begin{enumerate}
\def\labelenumi{\arabic{enumi}.}
\setcounter{enumi}{10}
\tightlist
\item
\item
\item
\item
\item
\item
\end{enumerate}

\begin{verbatim}
\begin{aligned} (\text{Y} \forall \Pi) & \qquad \phi_{\text{vx}} + \phi_{yy} = 0, \quad 0 < x < \infty, \quad 0 < y < \infty, \\ \left\{ \begin{array}{l} \phi(x, \ 0) = \begin{cases} 1, & 0 < x < 1, \\ 0, & 1 \leqslant x < \infty, \end{cases} \\ \phi(0, \ y) = \begin{cases} 1, & 0 < y < 1, \\ 0, & 1 \leqslant y < \infty. \end{cases} \end{aligned}
\end{verbatim}

\begin{enumerate}
\def\labelenumi{\arabic{enumi}.}
\setcounter{enumi}{16}
\tightlist
\item
\item
\end{enumerate}

Зверніть увагу на зображення.

\begin{enumerate}
\def\labelenumi{\arabic{enumi}.}
\setcounter{enumi}{18}
\item
\item
\item
  Розв'язати змішану задачу Діріхле-Неймана з кутовим порошком за
  допомогою розчину під кутом 45°
\item
\end{enumerate}

ПРИМІТКА. Комплексна функція
\texttt{w\ =\ \textbackslash{}ln\ z\ =\ \textbackslash{}ln\ \textbar{}z\textbar{}\ +\ i\ \textbackslash{}arg(z)}
відображає промінь \texttt{\textbackslash{}theta\ =\ c\_1} у площині z у
пряму \texttt{v\ =\ c\_1} , а коло \texttt{r\ =\ c\_2} у пряму
\texttt{u\ =\ \textbackslash{}ln\ c\_2} у площині w (див. рисунки).

\begin{enumerate}
\def\labelenumi{\arabic{enumi}.}
\setcounter{enumi}{10}
\tightlist
\item
\end{enumerate}

\section{ТАБЛИЦІ ІНТЕГРАЛЬНИХ
ПЕРЕТВОРЕНЬ}\label{ux442ux430ux431ux43bux438ux446ux456-ux456ux43dux442ux435ux433ux440ux430ux43bux44cux43dux438ux445-ux43fux435ux440ux435ux442ux432ux43eux440ux435ux43dux44c}

Таблиця A: Експоненціальне перетворення Фур'є,

Таблиця B: Перетворення синуса Фур'є.

Таблиця St Cosine Фур'є-перетворення.

Таблиця Dr скінченне перетворення Фур'є, таблиця та скінченне косинусне
перетворення Фур'є.

Таблиця D: Трансформація Лапласа.

\section{\texorpdfstring{Прийняті позначення
\texttt{\textbackslash{}delta\textbackslash{}left(x\textbackslash{}right)}
---
дельта-функція,}{Прийняті позначення \textbackslash delta\textbackslash left(x\textbackslash right) --- дельта-функція,}}\label{ux43fux440ux438ux439ux43dux44fux442ux456-ux43fux43eux437ux43dux430ux447ux435ux43dux43dux44f-deltaleftxright-ux434ux435ux43bux44cux442ux430-ux444ux443ux43dux43aux446ux456ux44f}

\begin{enumerate}
\def\labelenumi{\arabic{enumi}.}
\setcounter{enumi}{11}
\tightlist
\item
\item
\item
\end{enumerate}

\begin{verbatim}
\begin{array}{llll} (\mathrm{Y} \Pi \Pi) & u_{tt} = \alpha^2 u_{xx} - \beta u_t, & 0 < x < 1, & 0 < t < \infty, \\ (\Gamma \mathrm{Y}) & \begin{cases} u \ (0, \ t) = 0, \\ u \ (1, \ t) = 0, \end{cases} & 0 < t < \infty, \\ (\mathrm{H} \mathrm{Y}) & \begin{cases} u \ (x, \ 0) = f \ (x), \\ u_1 \ (x, \ 0) = 0. \end{cases} & 0 \leqslant x \leqslant 1, \end{array}
\end{verbatim}

\begin{enumerate}
\def\labelenumi{\arabic{enumi}.}
\setcounter{enumi}{17}
\tightlist
\item
\end{enumerate}

Воспользовавшись соотношениями ортогональности

\begin{verbatim}
\begin{array}{ll} (\mathrm{Y} \mathrm{H} \Pi) & u_{tt} = \alpha^2 u_{xx} + Kx, \quad 0 < x < 1, \ 0 < t < \infty, \\ (\mathrm{\Gamma} \mathrm{Y}) & \begin{cases} u \ (0, \ t) = 0, \\ u \ (1, \ t) = 0, \end{cases} & 0 < t < \infty, \\ (\mathrm{H} \mathrm{Y}) & \begin{cases} u \ (x, \ 0) = f \ (x), \\ u_t \ (x, \ 0) = 0. \end{cases} & 0 \leqslant x \leqslant 1, \end{array}
\end{verbatim}

\begin{enumerate}
\def\labelenumi{\arabic{enumi}.}
\setcounter{enumi}{8}
\tightlist
\item
\item
\item
\end{enumerate}

\pandocbounded{\includegraphics[keepaspectratio]{images/_page_154_Figure_5.jpeg}}

находим коэффициенты \texttt{a\_n} и \texttt{b\_n} (20.5)

\begin{verbatim}
# ТАБЛИЦЯ B. Перетворення Спнуса-Фур'є i \sqrt{\frac{\pi}{2}} [\delta(\omega+a) - \delta(\omega-a)]
\end{verbatim}

Итак, задача решена. Ее решение записывается в виде

(20.6)

\begin{verbatim}
F(\omega) = \frac{?}{\pi} \int_{0}^{\infty} f(x) \sin(\omega x) dx
\end{verbatim}

а коэффициенты \texttt{a\_n} и \texttt{b\_n} находятся по формулам
(20.5). Прежде чем закончить лекцию, сделаем несколько замечаний.

\section{ЗАМЕЧАНИЯ}\label{ux437ux430ux43cux435ux447ux430ux43dux438ux44f}

\begin{enumerate}
\def\labelenumi{\arabic{enumi}.}
\tightlist
\item
  Если начальная \emph{скорость} струны равна нулю, то решение (20.6)
  принимает вид
\end{enumerate}

\begin{verbatim}
0 < \omega < x
\end{verbatim}

І це можна інтерпретувати так. Припустимо, ми розклали початкове
зміщення струни в

\pandocbounded{\includegraphics[keepaspectratio]{images/_page_155_Figure_8.jpeg}}

ряд по синусам

f'\,'(0)

Тогда каждый член такого разложения должен совершать колебания

\begin{verbatim}
\begin{aligned} & (\text{Y} \cdot \Pi) \quad u_{tt} = -u_{xxxx}, \quad 0 < x < 1, \quad 0 < t < \infty, \\ & \begin{cases} u \cdot (0, \ t) = 0, \\ u_{xx} \cdot (0, \ t) = 0, \\ u \cdot (1, \ t) = 0, \\ u_{xx} \cdot (1, \ t) = 0, \end{cases} \quad 0 < t < \infty, \\ & \begin{cases} u \cdot (x, \ 0) = f(x), \\ u_{1}(x, \ 0) = g(x). \end{cases} \quad 0 \leqslant x \leqslant 1, \end{aligned}
\end{verbatim}

Якщо ми підсумуємо всі ці коливання, отримаємо розв'язок нашої проблеми.
Наприклад, припустимо, що

\texttt{начальное\ смещение\ струны\ задается\ функцией\ \textquotesingle{}f(x)\ =\ \textbackslash{}sin(\textbackslash{}pi\ x/L)\ +\ 0.5\ \textbackslash{}sin(3\textbackslash{}pi\ x/\ L)\ +\ 0,25\ \textbackslash{}sin(5\textbackslash{}pi\ x/L)\textquotesingle{}\ .}

Полный отклик системы на начальное условие такого вида равен сумме
откликов на каждое слагаемое, т. е.

\begin{enumerate}
\def\labelenumi{\arabic{enumi}.}
\setcounter{enumi}{3}
\tightlist
\item
\item
  Слагаемое с номером n в разложении (20.6) называется \textbf{п-й модой
  колебаний} или \textbf{п-й гармоникой}. Используя тригонометрические
  формулы, можно записать эти гармоники в виде
\end{enumerate}

\begin{verbatim}
[2/\pi\omega]^{1/2}
\end{verbatim}

где \texttt{R\_n} и \texttt{\textbackslash{}delta\_n} --- новые
произвольные постоянные (амплитуда и фазовый угол). Эта новая форма
представления n-й моды оказывается более полезной при анализе колебаний.
Частота колебаний од (рад/с) и-й моды находится по формуле

\begin{verbatim}
\frac{2}{\pi\omega} [1-\cos(\omega a)]
\end{verbatim}

де T --- натяг струни, p --- лінійна густина струни. Зверніть увагу, що
частота n-ї гармоніки у n разів більша за основну частоту (n=1). Ця
властивість частоти не характерна для всіх типів коливань. Приємний звук
струни скриму або гітари зумовлений множинністю обертонів основної
частоти. У барабані частоти гармонік високого порядку не є кратними
основної частоти коливань.

\section{ЗАДАЧИ}\label{ux437ux430ux434ux430ux447ux438-1}

\begin{enumerate}
\def\labelenumi{\arabic{enumi}.}
\tightlist
\item
  Найдите решение задачи (20.1), если начальные условия задаются
  формулами
\end{enumerate}

1

\texttt{u\_t(x,\ 0)\ =\ 0.} Постройте графики решения для различных
моментов времени. Является ли это решение периодическим по времени?
Каков его период?

\begin{enumerate}
\def\labelenumi{\arabic{enumi}.}
\setcounter{enumi}{1}
\tightlist
\item
  Найдите решение задачи о колебаниях струны (20.1) при начальных
  условиях вида
\end{enumerate}

\begin{verbatim}
u_t(x, 0) = \sin(3\pi x/L).  Как выглядит график решения в различные моменты времени?
\end{verbatim}

\begin{verbatim}
- 3. Показати це, якщо на рис. 20.2 значення  \lambda \geqslant 0  , тоді розв'язки X(x) T(t) будуть або необмеженими, або ідентичними нулю.
\end{verbatim}

\begin{itemize}
\tightlist
\item
  Знайти розв'язок проблеми коливань струн за початкових умов вигляду
\end{itemize}

\begin{verbatim}
\frac{1}{2}e^{-\alpha}\sin\omega
\end{verbatim}

\begin{verbatim}
u_1(x, 0) = (3\pi \alpha/L)\sin(3\pi x/L).  5. Гитарная струна длины L=1 оттягивается за среднюю точку на величину h (см. рис.) Начальное положение струны можно задать в виде
\end{verbatim}

\begin{verbatim}
\frac{1-e^{-a\omega}}{\omega}
\end{verbatim}

\begin{enumerate}
\def\labelenumi{\arabic{enumi}.}
\setcounter{enumi}{13}
\tightlist
\item
\item
\end{enumerate}

\begin{verbatim}
F \sinh \sqrt{\omega} = 0,
\end{verbatim}

\begin{enumerate}
\def\labelenumi{\arabic{enumi}.}
\setcounter{enumi}{16}
\tightlist
\item
\end{enumerate}

\begin{verbatim}
f(x) = \int_{0}^{\infty} F(\omega) \cos(\omega x) d\omega \qquad \qquad F(x) = \frac{2}{\pi} \int_{0}^{\infty} f(x) \cos(\omega x) dx
\end{verbatim}

\begin{enumerate}
\def\labelenumi{\arabic{enumi}.}
\setcounter{enumi}{18}
\tightlist
\item
\end{enumerate}

\begin{verbatim}
D \sin \sqrt{\omega} = 0.
\end{verbatim}

\begin{enumerate}
\def\labelenumi{\arabic{enumi}.}
\setcounter{enumi}{18}
\tightlist
\item
\item
\end{enumerate}

\section{КОЛЕБАНИЯ БАЛКИ (УРАВНЕНИЕ С ЧАСТНЫМИ ПРОИЗВОДНЫМИ ЧЕТВЕРТОГО
ПОРЯДКА)}\label{ux43aux43eux43bux435ux431ux430ux43dux438ux44f-ux431ux430ux43bux43aux438-ux443ux440ux430ux432ux43dux435ux43dux438ux435-ux441-ux447ux430ux441ux442ux43dux44bux43cux438-ux43fux440ux43eux438ux437ux432ux43eux434ux43dux44bux43cux438-ux447ux435ux442ux432ux435ux440ux442ux43eux433ux43e-ux43fux43eux440ux44fux434ux43aux430}

ЦЕЛЬ ЛЕКЦИИ: Показать, каким образом задача о колебании балки приводит к
уравнению четвертого порядка и как эта задача решается в случае свободно
опирающихся концов. Проводится сравнение колебаний балки с колебаниями
скрипичной струны.

\begin{verbatim}
F = 0
\end{verbatim}

Головна різниця між поперечними коливаннями тонкої балки та поперечними
коливаннями струни полягає в тому, що промінь чинить опір згину. Не
заглиблюючись у механіку тонких балок, припустимо, що, враховуючи опір
згинання (замість рівняння хвилі) веде до рівняння четвертого порядку

(21.1)

\section{ТАБЛИЦЯ D. Скінченне перетворення синуса
Фур'є}\label{ux442ux430ux431ux43bux438ux446ux44f-d.-ux441ux43aux456ux43dux447ux435ux43dux43dux435-ux43fux435ux440ux435ux442ux432ux43eux440ux435ux43dux43dux44f-ux441ux438ux43dux443ux441ux430-ux444ux443ux440ux454}

\begin{verbatim}
У скінченному перетворенні синуса Фур'є функція f(x), визначена в інтервалі 0 \le x \le \pi , перетворюється у послідовність S_n , n = 1, 2 ... За формулою
\end{verbatim}

\texttt{rge\ \textquotesingle{}\textbackslash{}alpha\^{}2\ =\ K/\textbackslash{}rho\textquotesingle{}\ ,}

K---модуль сдвига (чем больше K, тем жестче балка, тем выше частота
колебаний),

р --- линейная плотность балки (масса/ед. длины).

Поскольку читатель впервые встречается в нашей книге с уравнением выше
чем второго порядка, будет полезно решить типовую задачу о колебаниях
балки. Позже мы познакомимся и с другими задачами теории балок.

\section{Свободно опирающаяся
балка}\label{ux441ux432ux43eux431ux43eux434ux43dux43e-ux43eux43fux438ux440ux430ux44eux449ux430ux44fux441ux44f-ux431ux430ux43bux43aux430}

Рассмотрим малые колебания тонкой балки, концы которой свободно
опираются на две опоры. Говоря «свободно опираются», мы подразумеваем,
что концы балки не перемещаются, но на-

клоны балки в концевых точках могут изменяться (концы балки вакреплены с
помощью штифтовых устройств, рис. 21.1).

Совершенно ясно, что на концах балки должны выполняться граничные
условия

\begin{verbatim}
\omega_n = (n\pi)^2,
\end{verbatim}

\texttt{u(1,\ t)\ =\ 0,} но совсем не очевидно, что на концах балки
должны удовлетворяться еще два граничных условия:

\begin{verbatim}
Коли йдеться про таблиці, зручно розглядати змінну x як належну до інтервалу [0, l]. Будь-який інший чисельний інтервал [a, b] можна перетворити на [0, \pi] за допомогою формули
\end{verbatim}

\texttt{u\_\{xx\}(1,\ t)\ =\ 0.} Використовуючи теорію тонких балок,
можна показати, що момент згинання в балці пропорційний значенню
\texttt{u\_\{xx\}} , а момент згинання на вільно лежачому кінці має
дорівнювати P I C

\begin{verbatim}
u_n(x, t) = X_n(x) T_n(t) = [a_n \sin(n\pi)^2 t + b_n \cos(n\pi)^2 t] \sin(n\pi x).
\end{verbatim}

.21.1. Свободно оппрающаяся балка.

вен нулю. Следовательно, изображенная на рис. 21.1 колеблющаяся балка
описывается следующей смешанной задачей (для простоты величина о принята
равной единице):

(21.2)

\begin{verbatim}
У скінченному синусоїдальному перетворенні немає нічого загадкового: елементи послідовності \mathcal{S}_n збігаються з коефіцієнтами в \sin(nx) у розкладі функції f(x) у ряд Фур'є, тобто
\end{verbatim}

\begin{enumerate}
\def\labelenumi{\arabic{enumi}.}
\setcounter{enumi}{1}
\tightlist
\item
  \texttt{f(\textbackslash{}alpha\ x)}
\end{enumerate}

\begin{verbatim}
u(x, t) = \sum_{n=1}^{\infty} \left[ a_n \sin(n\pi)^2 t + b_n \cos(n\pi)^2 t \right] \sin(n\pi x)
\end{verbatim}

\begin{enumerate}
\def\labelenumi{\arabic{enumi}.}
\setcounter{enumi}{2}
\tightlist
\item
  \texttt{e\^{}\{-\textbackslash{}alpha\ x\}}
\end{enumerate}

\begin{verbatim}
f(x) = \sum_{n=1}^{\infty} S_n \sin(nx)
\end{verbatim}

Для решения этой задачи воспользуемся методом разделения переменных.
Будем искать только периодические решения, т. е. колебания вида

(21.3)

\begin{verbatim}
u(x, 0) = f(x) = \sum_{n=1}^{\infty} b_n \sin(n\pi x),
\end{verbatim}

Зверніть увагу, що вибір розв'язку вигляду (21.3) фактично означає, що
ми взяли константу розділення у методі розділення змінних як від'ємну.

\texttt{Подставим\ теперь\ (21.3)\ в\ уравнение\ колебаний\ балки\ и\ получим\ уравнение\ для\ \textquotesingle{}X\textbackslash{}left(x\textbackslash{}right)\textquotesingle{}}

\begin{verbatim}
0 \le x \le \pi
\end{verbatim}

общее решение которого записывается в виде

\begin{verbatim}
u_t(x, 0) = g(x) = \sum_{n=1}^{\infty} (n\pi)^2 a_n \sin(n\pi x).
\end{verbatim}

Для определения коэффициентов C, D, E и F подставим общее решение в
граничные условия

\begin{enumerate}
\def\labelenumi{\arabic{enumi}.}
\tightlist
\item
\end{enumerate}

\begin{verbatim}
a_n = \frac{2}{(n\pi)^2} \int_0^1 g(x) \sin(n\pi x) dx,
\end{verbatim}

\begin{enumerate}
\def\labelenumi{\arabic{enumi}.}
\setcounter{enumi}{8}
\tightlist
\item
\end{enumerate}

\begin{verbatim}
b_n = 2 \int_0^1 f(x) \sin(n\pi x) dx.
\end{verbatim}

\begin{verbatim}
\frac{n}{n^2 + a^2} [1 - (-1)^n e^{ax}] 10. \frac{\sin a (\pi - x)}{\sin a\pi} \frac{2n}{\pi (n^2 + a^2)} ТАБЛИЦЯ E. Скінченне перетворення Косніуса Фур'є
\end{verbatim}

\begin{verbatim}
У скінченному косинусному перетворенні Фур'є функція f(x), визначена на інтервалі 0 \le x \le \pi , перетворюється у послідовність C_n , n=0,\ 1,\ 2,\ \ldots , за формулою
\end{verbatim}

\begin{verbatim}
u(x, 0) = \sin(\pi x) + 0.5 \sin(3\pi x),
\end{verbatim}

\begin{verbatim}
Коли йдеться про таблиці, зручно розглядати змінну x як таку, що належить до інтервалу [0,\pi] . Будь-який інший числовий інтервал [a,b] можна перетворити на [0,\pi] за допомогою формули
\end{verbatim}

\begin{verbatim}
a_n=0
\end{verbatim}

\begin{verbatim}
У скінченному косинусному перетворенні немає нічого загадкового: елементи послідовності C_n збігаються з коефіцієнтами, коли \cos{(nx)} у косинусному розкладі функції f(x) у ряд Фур'є, тобто
\end{verbatim}

\begin{verbatim}
u(x, t) = \cos(\pi^2 t) \sin(\pi x) + 0.5 \cos(9\pi^2 t) \sin(3\pi x)
f(x) = \frac{C_0}{2} + \sum_{n=1}^{\infty} C_n \cos(nx) \qquad C_n = \frac{2}{\pi} \int_{0}^{\pi} f(x) \cos nx \, dx
u(x, t) = \cos(\pi t) \sin(\pi x) + 0.5 \cos(3\pi t) \sin(3\pi x)
\end{verbatim}

\begin{enumerate}
\def\labelenumi{\arabic{enumi}.}
\tightlist
\item
\end{enumerate}

\begin{verbatim}
u(0, t) = 0,
-n^{2}C_{n} - \frac{2}{\pi} \{f'(0) - (-1)^{n} f'(n)\}
u(x, t) = \sum_{n=1}^{\infty} X_n(x) \left[ a_n \sin(\omega_n t) + b_n \cos(\omega_n t) \right],
\begin{cases} 2, &n=0 \\ 0, &n=1, 2, \dots \end{cases}
\begin{cases} 1, & n=m \\ 0, & n\neq m \end{cases}
\end{verbatim}

x

\begin{verbatim}
\quad 0 \leqslant x \leqslant 1,
u(x, 0) = \sin(\pi x),
\end{verbatim}

\begin{enumerate}
\def\labelenumi{\arabic{enumi}.}
\setcounter{enumi}{3}
\tightlist
\item
  Пусть левый конец (x=0) балки жестко закреплен в стену, а правый конец
  (x=1) свободно опирается, как показано на рис. 21.3. Решите задачу для
  балки с такими граничными условиями и объясните, как найти собственные
  частоты колебаний в этом случае. Знать собственные частоты балок очень
  важно, чтобы избежать резонанса с внешними воздействиями.
\end{enumerate}

Лекция 22

\section{ПЕРЕХОД К БЕЗРАЗМЕРНЫМ
ПЕРЕМЕННЫМ}\label{ux43fux435ux440ux435ux445ux43eux434-ux43a-ux431ux435ux437ux440ux430ux437ux43cux435ux440ux43dux44bux43c-ux43fux435ux440ux435ux43cux435ux43dux43dux44bux43c}

ЦЕЛЬ ЛЕКЦИИ: Показать, каким образом краевые, начальные и другие виды
физических задач можно записать в безразмерном виде. В конкретных
областях знания (физике, химии, биологии или экономике) уравнения
записываются по-разному.

Переход к безразмерным переменным позволяет привести уравнения к одному
и тому же виду. Именно по этой причине при изучении уравнений с частными
производными в математике стараются отвлекаться от физического смысла
входящих в них параметров. Вот почему химик, физик или биолог должен
преобразовать свое уравнение к той форме, которая принята в нашей книге.

Основная идея введения новых (безразмерных) переменных состоит в том,
что после перехода к безразмерным переменным вадача становится чисто
математической и уже не содержит

карактерных физических констант. Именно таким способом многие различные
уравнения физики, биологии и химии, содержащие всякие нюансы, связанные
с физическими параметрами, приводятся к одной и той же простой форме
(см. рис. 22.1).

Познакомимся с этим преобразованием на простом примере.

\pandocbounded{\includegraphics[keepaspectratio]{images/_page_161_Figure_4.jpeg}}

РИС. 22.1. Некоторые задачи, приводящиеся к одной и той же безразмерной
форме.

\section{Приведение диффузионной задачи к безразмерному
виду}\label{ux43fux440ux438ux432ux435ux434ux435ux43dux438ux435-ux434ux438ux444ux444ux443ux437ux438ux43eux43dux43dux43eux439-ux437ux430ux434ux430ux447ux438-ux43a-ux431ux435ux437ux440ux430ux437ux43cux435ux440ux43dux43eux43cux443-ux432ux438ux434ux443}

\texttt{Начнем\ со\ смешанной\ задачи\ для\ отрезка\ с\ начальным\ условием\ на\ температуру\ \ \textquotesingle{}u\textbackslash{}left(x,\textbackslash{},0\textbackslash{}right)=\textbackslash{}sin\textbackslash{}left(\textbackslash{}pi\ x/L\textbackslash{}right)\textquotesingle{}\ \ и\ граничными\ температурами\ на\ концах\ \ \textquotesingle{}T\_1\textquotesingle{}\ \ и\ \ \textquotesingle{}T\_2\textquotesingle{}\ .\ Другими\ словами,\ мы\ решаем\ задачу}

(22.1)

\begin{verbatim}
\begin{array}{ccccccccccccccccccccccccccccccccccc
\end{verbatim}

\pandocbounded{\includegraphics[keepaspectratio]{images/_page_161_Figure_10.jpeg}}

\begin{verbatim}
(\text{УЧП}) \quad u_{tt} + u_{xxxx} = 0, \quad 0 < x < 1, \quad 0 < t < \infty,
u \to U
(\text{ГУ}) \quad
U(x, t) = \frac{u(x, t) - T_1}{T_2 - T_1}
\end{verbatim}

\begin{verbatim}
\begin{array}{ll} \text{(УЧП)} & U_t = \alpha^2 U_{xx}, \quad 0 < x < L, \quad 0 < t < \infty, \\ U_t = 0, \quad 0 < t < \infty, \\ U_t = 0, \quad 0 < t < \infty, \\ U_t = 0, \quad 0 < t < \infty, \\ U_t = 0, \quad 0 < t < \infty, \\ U_t = 0, \quad 0 < t < \infty, \\ U_t = 0, \quad 0 < t < \infty, \\ U_t = 0, \quad 0 < t < \infty, \\ U_t = 0, \quad 0 < t < \infty, \\ U_t = 0, \quad 0 < t < \infty, \\ U_t = 0, \quad 0 < t < \infty, \\ U_t = 0, \quad 0 < t < \infty, \\ U_t = 0, \quad 0 < t < \infty, \\ U_t = 0, \quad 0 < t < \infty, \\ U_t = 0, \quad 0 < t < \infty, \\ U_t = 0, \quad 0 < t < \infty, \\ U_t = 0, \quad 0 < t < \infty, \\ U_t = 0, \quad 0 < t < \infty, \\ U_t = 0, \quad 0 < t < \infty, \\ U_t = 0, \quad 0 < t < \infty, \\ U_t = 0, \quad 0 < t < \infty, \\ U_t = 0, \quad 0 < t < \infty, \\ U_t = 0, \quad 0 < t < \infty, \\ U_t = 0, \quad 0 < t < \infty, \\ U_t = 0, \quad 0 < t < \infty, \\ U_t = 0, \quad 0 < t < \infty, \\ U_t = 0, \quad 0 < t < \infty, \\ U_t = 0, \quad 0 < t < \infty, \\ U_t = 0, \quad 0 < t < \infty, \\ U_t = 0, \quad 0 < t < \infty, \\ U_t = 0, \quad 0 < t < \infty, \\ U_t = 0, \quad 0 < t < \infty, \\ U_t = 0, \quad 0 < t < \infty, \\ U_t = 0, \quad 0 < t < \infty, \\ U_t = 0, \quad 0 < t < \infty, \\ U_t = 0, \quad 0 < t < \infty, \\ U_t = 0, \quad 0 < t < \infty, \\ U_t = 0, \quad 0 < t < \infty, \\ U_t = 0, \quad 0 < t < \infty, \\ U_t = 0, \quad 0 < t < \infty, \\ U_t = 0, \quad 0 < t < \infty, \\ U_t = 0, \quad 0 < t < \infty, \\ U_t = 0, \quad 0 < t < \infty, \\ U_t = 0, \quad 0 < t < \infty, \\ U_t = 0, \quad 0 < t < \infty, \\ U_t = 0, \quad 0 < t < \infty, \\ U_t = 0, \quad 0 < t < \infty, \\ U_t = 0, \quad 0 < t < \infty, \\ U_t = 0, \quad 0 < t < \infty, \\ U_t = 0, \quad 0 < t < \infty, \\ U_t = 0, \quad 0 < t < \infty, \\ U_t = 0, \quad 0 < t < \infty, \\ U_t = 0, \quad 0 < t < \infty, \\ U_t = 0, \quad 0 < t < \infty, \\ U_t = 0, \quad 0 < t < \infty, \\ U_t = 0, \quad 0 < t < \infty, \\ U_t = 0, \quad 0 < t < \infty, \\ U_t = 0, \quad 0 < t < \infty, \\ U_t = 0, \quad 0 < t < \infty, \\ U_t = 0, \quad 0 < t < \infty, \\ U_t = 0, \quad 0 < t < \infty, \\ U_t = 0, \quad 0 < t < \infty, \\ U_t = 0, \quad 0 < t < \infty, \\ U_t = 0, \quad 0 < t < \infty, \\ U_t = 0, \quad 0 < t < \infty, \\ U_t = 0, \quad 0 < t < \infty, \\ U_t = 0, \quad 0 < t < \infty, \\ U_t = 0, \quad 0 < t < \infty, \\ U_t = 0, \quad 0 < t < \infty, \\ U_t = 0, \quad 0 < t < \infty, \\ U_t = 0, \quad 0 < t < \infty, \\ U_t = 0, \quad 0 < t < \infty, \\ U_t = 0, \quad 0 < t < \infty, \\ U_t = 0, \quad 0 < t < \infty, \\ U_t = 0, \quad 0 < t < \infty, \\ U_t = 0, \quad 0 < t < \infty, \\ U_t = 0, \quad 0 < t < \infty, \\ U_t = 0, \quad 0 < t < \infty, \\ U_t = 0, \quad 0 < t < \infty, \\ U_t = 0, \quad 0 < t < \infty, \\ U_t = 0, \quad 0 < t < \infty, \\ U_t = 0, \quad 0 < t < \infty, \\ U_t = 0, \quad 0 < t < \infty, \\ U_t
\end{verbatim}

\begin{verbatim}
\begin{cases} u(0, t) &= 0, \\ u_x(0, t) &= 0, \\ u_{xx}(1, t) &= 0, \\ u_{xxx}(1, t) &= 0, \end{cases}
u(x, t) = T_1 + (T_2 - T_1) U(x, t).
\xi = x/L
U_x = U_{\xi} \xi_x = \frac{1}{L} U_{\xi},
\begin{cases} u(x, 0) &= f(x), \\ u_t(x, 0) &= g(x). \end{cases}
U_{xx} = \frac{1}{L^2} U_{\xi\xi}
\begin{array}{lll} (\text{YM}\Pi) & U_t = (\alpha/L)^2 \, U_{\xi\xi}, \quad 0 < \xi < 1, \quad 0 < t < \infty, \\ (22.3) & (\Gamma \text{Y}) & \begin{cases} U(0,\ t) = 0, \\ U(1,\ t) = 1, \\ U(\xi,\ 0) = \frac{\sin\left(\pi\xi\right) - T_1}{T_2 - T_1}, \quad 0 \leqslant \xi \leqslant 1. \end{cases} \end{array}
\end{verbatim}

6 M 601

УКАЗАНИЕ. Хотя собственные функции \texttt{X\_n(x)} в этой задаче уже не
обычные синусы, согласно общей теории Штурма --- Ли вилля они будут
ортогональны на отрезке {[}0, 1{]}.

\begin{enumerate}
\def\labelenumi{\arabic{enumi}.}
\setcounter{enumi}{1}
\tightlist
\item
  Найдите решение для свободно опирающейся на оба конца балки при
  начальных условиях
\end{enumerate}

\begin{verbatim}
cu_{\mathbf{t}} = [\alpha/L]^2 u_{\xi\xi}.
\end{verbatim}

\texttt{u\_t(x,\ 0)\ =\ \textbackslash{}sin(\textbackslash{}pi\ x).}
\texttt{0\ \textbackslash{}le\ x\ \textbackslash{}le\ 1,} 3. Решите
задачу 2 с начальными условиями

\begin{verbatim}
\tau = [\alpha/L]^2 t.
\quad 0 \le x \le 1,
\end{verbatim}

\begin{verbatim}
\begin{array}{cccccccccccccccccccccccccccccccccccc
\end{verbatim}

\begin{verbatim}
u(x, t) = T_1 + (T_2 - T_1) U(x/L, \alpha^2 t/L^2).
\end{verbatim}

(см. рис. 22.2).

\pandocbounded{\includegraphics[keepaspectratio]{images/_page_161_Figure_10.jpeg}}

РИС. 22.2. Задача теплопроводности в плоскости переменных (х, т).

Наша цель --- дать новую эквивалентную (22.1) постановку задачи таким
образом, чтобы

в новой задаче не было никаких физических параметров (вроде α);

\begin{enumerate}
\def\labelenumi{\arabic{enumi})}
\setcounter{enumi}{1}
\tightlist
\item
  начальные и граничные условия стали проще.
\end{enumerate}

\begin{verbatim}
Для осуществления этой цели, мы введем три новых безразмерных переменных U,   \xi   и   \tau  , чтобы заменить прежние u, x и t по схеме
\end{verbatim}

\begin{verbatim}
\begin{array}{llll} (\text{УЧП}) & u_{tt} = \alpha^2 u_{xx}, & 0 < x < L, & 0 < t < \infty, \\ u(0, t) = 0, & 0 < t < \infty, \\ u(L, t) = 0, & 0 < t < \infty, \\ u(X, 0) = \sin(\pi x/L) + (0.5)\sin(3\pi x/L), \\ u_t(x, 0) = 0. & \end{array}
\end{verbatim}

\begin{verbatim}
(безразмерная температура),   x \to \xi   (безразмерная длина),   t \to \tau   (безразмерное время).
\end{verbatim}

Мы проведем все три преобразования одновременно.

\begin{verbatim}
# Преобразование зависимой переменной  u \rightarrow U
\end{verbatim}

Определим функцию \texttt{U(x,\ t)} по формуле

\begin{verbatim}
U(x, t) = \frac{u(x, t) - T_1}{T_2 - T_1}
\xi = x/L
\end{verbatim}

.

\texttt{Ясно,\ что\ новая\ температура\ \ \textquotesingle{}U\textbackslash{}left(x,\textbackslash{},t\textbackslash{}right)\textquotesingle{}\ \ безразмерна,\ поскольку\ мы\ делим\ градусы\ Цельсия\ на\ градусы\ Цельсия.\ Новые\ граничные\ условия\ для\ \ \textquotesingle{}U\textbackslash{}left(x,\textbackslash{},t\textbackslash{}right)\textquotesingle{}\ \ при\ x=0\ и\ x=L\ примут\ вид\ \ \textquotesingle{}U\textbackslash{}left(0,\textbackslash{},t\textbackslash{}right)=0\textquotesingle{}\ \ и\ \ \textquotesingle{}U\textbackslash{}left(L,\textbackslash{},t\textbackslash{}right)=1\textquotesingle{}\ .\ Давайте\ поставим\ задачу\ для\ новой\ функции\ \ \textquotesingle{}U\textbackslash{}left(x,\textbackslash{},t\textbackslash{}right)\textquotesingle{}\ .\ После\ несложных\ преобразований\ задача\ (22.1)\ переходит\ в\ новую\ задачу}

задачу

\begin{verbatim}
(\text{YYII}) \qquad u_{\tau\tau} = u_{\xi\xi}, \quad 0 < \xi < 1, \quad 0 < \tau < \infty, 1
\end{verbatim}

\texttt{Если\ угодно,\ то\ здесь\ можно\ остановиться,\ решить\ задачу\ для\ \ \textquotesingle{}U\textbackslash{}left(x,\textbackslash{},\ t\textbackslash{}right)\textquotesingle{}\ \ и\ затем\ найти\ \ \textquotesingle{}u\textbackslash{}left(x,\textbackslash{},t\textbackslash{}right)\textquotesingle{}\ \ по\ формуле}

\begin{verbatim}
(\text{\GammaY}) \quad
u(x, t) = T_1 + (T_2 - T_1) U(x, t).
\end{verbatim}

Однако мы пойдем дальше и займемся преобразованием независимых
переменных x и t.

\begin{verbatim}
# Преобразование пространственной координаты  x \rightarrow \xi
\end{verbatim}

\begin{verbatim}
Представляется совершенно очевидным, как следует выбрать безразмерную переменную   \xi  . Поскольку   0 \leqslant x \leqslant L  , полагаем
\end{verbatim}

\begin{verbatim}
\xi = x/L
\end{verbatim}

.

После вычисления производных

\begin{verbatim}
U_x = U_{\xi} \xi_x = \frac{1}{L} U_{\xi},
U_{xx} = \frac{1}{L^2} U_{\xi\xi}
\end{verbatim}

\begin{verbatim}
\begin{array}{ll} (\text{YM}\Pi) & U_t = (\alpha/L)^2 \, U_{\xi\xi}, \quad 0 < \xi < 1, \quad 0 < t < \infty, \\ (22.3) & (\Gamma \text{Y}) & \begin{cases} U(0,\ t) = 0, \\ U(1,\ t) = 1, \\ U(\xi,\ 0) = \frac{\sin\left(\pi\xi\right) - T_1}{T_2 - T_1}, \quad 0 \leqslant \xi \leqslant 1. \end{cases}
\end{verbatim}

\begin{verbatim}
cu_{\mathbf{t}} = [\alpha/L]^2 u_{\xi\xi}.
\tau = [\alpha/L]^2 t.
\end{verbatim}

\begin{verbatim}
\begin{array}{cccccccccccccccccccccccccccccccccccc
\end{verbatim}

\begin{verbatim}
\quad 0 \leqslant \xi \leqslant 1,
u(x, t) = T_1 + (T_2 - T_1) U(x/L, \alpha^2 t/L^2).
u(\xi, \tau) = \cos(\pi \tau) \sin(\pi \xi) + 0.5 \cos(3\pi \tau) \sin(3\pi \xi).
\begin{array}{llll} (\text{УЧП}) & u_{tt} = \alpha^2 u_{xx}, & 0 < x < L, & 0 < t < \infty, \\ u(0, t) = 0, & 0 < t < \infty, \\ u(L, t) = 0, & 0 < t < \infty, \\ u(X, 0) = \sin(\pi x/L) + (0.5)\sin(3\pi x/L), \\ u_t(x, 0) = 0. & \end{array}
u(x, t) = \cos(\pi \alpha t/L)\sin(\pi x/L) + 0.5\cos(3\pi \alpha t/L)\sin(3\pi x/L).
\xi = x/L
\xi = x/L
(\text{YYII}) \qquad u_{\tau\tau} = u_{\xi\xi}, \quad 0 < \xi < 1, \quad 0 < \tau < \infty, 1
\begin{cases} u(0, \tau) = 0, \\ u(1, \tau) = 0, \\ u(\xi, 0) = \sin(\pi\xi) + 0.5\sin(3\pi\xi), \\ u_{\xi}(\xi, 0) = 0, \end{cases}
(\text{HY}) \quad
\begin{cases} u(\xi, 0) = \sin(\pi\xi) + 0.5\sin(3\pi\xi), \\ u_{\xi}(\xi, 0) = 0, \end{cases}
\quad 0 \leqslant \xi \leqslant 1,
u(\xi, \tau) = \cos(\pi \tau) \sin(\pi \xi) + 0.5 \cos(3\pi \tau) \sin(3\pi \xi).
\end{verbatim}

\begin{verbatim}
\right. \\ \end{array}
\end{verbatim}

\begin{verbatim}
u(x, t) = \cos(\pi \alpha t/L)\sin(\pi x/L) + 0.5\cos(3\pi \alpha t/L)\sin(3\pi x/L).
U(x, t) = \frac{u(x, t) - T_1}{T_2 - T_1}
\xi = x/L
u_t + vu_x = 0?
\end{verbatim}

\begin{verbatim}
\begin{array}{lll} (\text{Y} \, \Pi) & u_t = \alpha^2 u_{xx}, & 0 < x < L, \\ (\Gamma \, \text{Y}) & \left\{ \begin{array}{ll} u \, (0, \, t) = T_1, \\ u \, (L, \, t) = 0, \\ u \, (x, \, 0) = T_2, & 0 \leqslant x \leqslant L. \end{array}
\end{verbatim}

\begin{verbatim}
Au_{xx} + Bu_{xy} + Cu_{yy} + Du_x + Eu_y + Fu = G
Au_{xx} + Bu_{xy} + Cu_{yy} + Du_x + Eu_y + Fu = G,
A=1
A=1
u_{\xi\eta} = \Phi(\xi, \eta, u, u_{\xi}, u_{\eta}).
Au_{xx} + Bu_{xy} + Cu_{yy} + Du_x + Eu_y + Fu = G
\xi = \xi(x, y),
\eta = \eta(x, y)
\overline{A}u_{\xi\xi} + \overline{B}u_{\xi\eta} + \overline{C}u_{\eta\eta} + \overline{D}u_{\xi} + \overline{E}u_{\eta} + \overline{F}u = \overline{G},
\overline{A} = A\xi_{x}^{2} + B\xi_{x}\xi_{y} + C\xi_{y}^{2},
\overline{B} = 2A\xi_{x}\eta_{y} + B(\xi_{x}\eta_{y} + \xi_{y}\eta_{x}) + 2C\xi_{y}\eta_{y},
\overline{C} = A\eta_{x}^{2} + B\eta_{x}\eta_{y} + C\eta_{y}^{2},
\overline{D} = A\xi_{xx} + B\xi_{xy} + C\xi_{yy} + D\xi_{x} + E\xi_{y},
\overline{E} = A\eta_{xx} + B\eta_{xy} + C\eta_{yy} + D\eta_{x} + E\eta_{y},
\overline{F} = F,
\overline{G} = G.
\end{verbatim}

\begin{verbatim}
так, чтобы в уравнении (23.3) осталась только одна частная производная второго порядка  u_{\xi\eta}  (не следует пытаться искать такую замену переменных, с помощью которой гиперболическое уравне-
\end{verbatim}

1) Необходимо, чтобы это преобразование координат было локально обратимо
и дважды дифференцируемо. --- Прим. ред.

ние можно было бы привести и параболической или заливитической
канонической форме, это невозможно сделать).

Вычислим сначала частные производные

\begin{verbatim}
\bar{A} = A\xi_x^2 + B\xi_x\xi_y + C\xi_y^2 = 0,
\end{verbatim}

Подставляя эти соотношения в исходное уравнение (23.3), после простых,
но громоздких вычислений получаем

(23.5)

Разрешив их относительно
\texttt{{[}\textbackslash{}xi\_x/\textbackslash{}xi\_y{]}} и
\texttt{{[}\textbackslash{}eta\_x/\textbackslash{}eta\_y{]}} , получаем

(23.7)

\begin{verbatim}
= \frac{-B + \sqrt{B^2 - 4AC}}{2A}
[\eta_x/\eta_y] = \frac{-B - \sqrt{B^2 - 4AC}}{2A}
\begin{split} &A\left[\xi_{x}/\xi_{y}\right]^{2}+B\left[\xi_{x}/\xi_{y}\right]+C=0,\\ &A\left[\eta_{x}/\eta_{y}\right]^{2}+B\left[\eta_{x}/\eta_{y}\right]^{2}+C=0. \end{split}
\end{verbatim}

\pandocbounded{\includegraphics[keepaspectratio]{images/_page_170_Figure_3.jpeg}}

\begin{verbatim}
u_{xx} - 4u_{yy} + u_x = 0.
\begin{vmatrix} \frac{dy}{dx} = -[\xi_x/\xi_y] = \frac{B - \sqrt{B^2 - 4AC}}{2A} = -2, \\ \frac{dy}{dx} = -[\eta_x/\eta_y] = \frac{B^{\frac{1}{2}} + \sqrt{B^2 - 4AC}}{2A} = 2. \end{vmatrix}
\begin{bmatrix} \xi_x/\xi_y \end{bmatrix}
y = -2x + c_1,
y = 2x + c_2.
\end{verbatim}

(характеристические

\begin{verbatim}
\xi = y + 2x = c_1,
\end{verbatim}

уравнения).

\begin{verbatim}
Kaждое из квадратных уравнений для  [\xi_x/\xi_y]  и  [\eta_x/\eta_y]  имеет два кория, но мы оставляем только по одному, но так чтобы они были различны.
\end{verbatim}

\begin{verbatim}
\eta = y - 2x = c_2.
\end{verbatim}

\begin{verbatim}
РИС. 23.1. Характеристики  \xi(x, y) = c  и  \eta(x, y) = c  .
\end{verbatim}

\begin{verbatim}
Задача була зведена до знаходження двох функцій  \xi(x, y)  і  \eta(x, y)  , щоб відношення  [\xi_x/\xi_y]  і  [\eta_x/\eta_y]  задовольняли рівняння (23.7). Дуже легко знайти такі функції 1) якщо звернути увагу на рис. 23.1. Щоб зрозуміти, як це можна зробити, розглянемо просте рівняння Е
\end{verbatim}

\pandocbounded{\includegraphics[keepaspectratio]{images/_page_171_Figure_3.jpeg}}

го характеристики определяются из уравнений

\begin{verbatim}
\overline{A}u_{\xi\xi} + \overline{B}u_{\xi\eta} + \overline{C}u_{\eta\eta} + \overline{D}u_{\xi} + \overline{E}u_{\eta} + \overline{F}u = \overline{G}
\end{verbatim}

Решая эти уравнения относительно у, получаем

\begin{verbatim}
y^2u_{xx}-x^2u_{yy}=0
\frac{dy}{dx} = \frac{B + \sqrt{B^2 - 4AC}}{2A} = \frac{x}{y}
\frac{dy}{dx} = \frac{B - \sqrt{B^2 - 4AC}}{2A} = -\frac{x}{y}.
\end{verbatim}

\pandocbounded{\includegraphics[keepaspectratio]{images/_page_172_Figure_2.jpeg}}

\begin{verbatim}
y^2 - x^2 = \text{const},
\end{verbatim}

1) См., например, Тихонов А. Н., Самарский А. А. Уравнения
математической физики. --- М.: Наука, 1972 и др. издания. --- Прим. ред.

\begin{verbatim}
Очевидно, введенные таким образом функции  \xi  и  \eta  удовлетворяют уравнениям характеристик. Эти новые координаты изображены на рис. 23.2. На этом мы заканчиваем обсуждение вопроса о том, как найти новые координаты. Осталось последнее — определить новый вид исходного уравнения.
\end{verbatim}

\begin{verbatim}
y^2 + x^2 = \text{const}.
\end{verbatim}

РИС. 23.2. Характеристические координаты для уравнения
\texttt{u\_\{xx\}} - 4u\_\{uu\} + \texttt{u\_x\ =\ 0}.

\begin{verbatim}
Это совсем просто: чтобы найти каноническую форму уравнения с частными производными, подставим новые координаты  \xi(x, y)  и  \eta(x,) y)  в уравнение
\end{verbatim}

\begin{verbatim}
\xi = y^2 - x^2
\end{verbatim}

,

\begin{verbatim}
где  \overline{A}  ,  \overline{B}  ,  \overline{C}  ,  \overline{D}  ,  \overline{E}  ,  \overline{F}  и  \overline{G}  определены формулами (23.6).
\end{verbatim}

Прежде чем закончить лекцию, давайте на конкретном примере посмотрим,
как «работает» общий метод приведения к каноническому виду.

Приведение гиперболического уравнения
y\textsuperscript{2u\_\{xx\}-x}2u\_\{yy\}=0 к канонической форме

Рассмотрим уравнение

\begin{verbatim}
\vec{B} = 2A\xi_x\eta_x + B(\xi_x\eta_y + \xi_y\eta_x) + 2C\xi_y\eta_y = -16x^2y^2
,  x>  ,  0 y>0  ,
\end{verbatim}

которое является гиперболическим в первом квадранте. Найдем новые
координаты таким образом, чтобы привести исходное уравнение к
канонической форме.

ШАГ 1. Решаем уравнения характеристик

\begin{verbatim}
\bar{D} = A\xi_{xx} + B\xi_{xy} + C\xi_{yy} + D\xi_x + E\xi_y = -2(x^2 + y^2)
\end{verbatim}

(напомним, что этот шаг

\begin{verbatim}
\widetilde{E} = A\eta_{xx} + B\eta_{xy} + C\eta_{yy} + D\eta_x + E\eta_y = 2(y^2 - x^2),
\end{verbatim}

\begin{verbatim}
обеспечивает  \overline{A} = \overline{C} = 0  ),
\end{verbatim}

\begin{verbatim}
\tilde{F} = F = 0
\end{verbatim}

РИС, 23.3. Новые карактеристические координаты.

Проинтегрировав эти обыкновенные дифференциальные уравнения методом
разделения переменных, получаем два неявных соотнощения (при желании
можно было бы явно выразить у через х)

\begin{verbatim}
\overline{G} = G = 0
\overline{A}u_{\xi\xi} + \overline{B}u_{\xi\eta} + \overline{C}u_{\eta\eta} + \overline{D}u + \overline{E}u_{\eta} + \overline{F}u = \overline{G}
u_{\xi\eta} = \frac{-(x^2 + y^2) u_{\xi} + (y^2 - x^2) u_{\eta}}{8x^2y^3}.
u_{\xi\eta} = \frac{\eta u_{\xi} - \xi u_{\eta}}{2(\xi^2 - \eta^2)}.
\alpha = \alpha (\xi, \eta) = \xi + \eta,
u_{\alpha\alpha} - u_{\beta\beta} = \frac{-\beta u_{\alpha} - \alpha u_{\beta}}{2\alpha\beta}.
\alpha = \xi + \eta = (y^2 - x^2) + (y^2 + x^2) = 2y^3,
\beta = \xi - \eta = (y^2 - x^2) - (y^2 + x^2) = -2x^2.
u_{\xi\xi} - u_{\eta\eta} = \Psi(\xi, \, \eta, \, u_{\xi}, \, u, u_{\eta}).
3u_{xx} + 7u_{xy} + 2u_{yy} = 0
u_{\xi\eta} = \Phi(\xi, \eta, u, u_1, u_2).
u_{\alpha\alpha} - u_{\beta\beta} = \Psi(\alpha, \beta, u, u_{\alpha}, u_{\beta}).
u_{xx} + 4u_{xy} = 0.
u_{tt} = e^{2} (u_{xx} + u_{yy} + u_{zz}),
\begin{array}{l} u_{\xi} = u_{\alpha}\alpha_{\xi} + u_{\beta}\beta_{\xi} = u_{\alpha} + u_{\beta}. \\ u_{\eta} = u_{\alpha}\alpha_{\eta} + u_{\beta}\beta_{\eta} = u_{\alpha} - u_{\beta}, \\ u_{\xi\eta} = u_{\alpha\alpha}\alpha_{\eta} + u_{\alpha\beta}\beta_{\eta} + u_{\beta\alpha}\alpha_{\eta} + u_{\beta\beta}\beta_{\eta} = u_{\alpha\alpha} - u_{\beta\beta}, \end{array}
\end{verbatim}

(НУ)

\begin{itemize}
\tightlist
\item
  Читач може запитати, чому необхідно класифікувати та зводити їх до
  канонічної форми диференціальних рівнянь у похідних похідних.
\item
  \begin{enumerate}
  \def\labelenumi{\alph{enumi})}
  \tightlist
  \item
    Поділ рівнянь на класи гіперболічних, параболічних та еліптичних
    відповідає поділу фізичних процесів на три основні класи: хвильовий,
    дифузійний і стаціонарний. Математичні особливості розв'язків цих
    трьох типів рівнянь є абсолютно різними.
  \end{enumerate}
\item
  \begin{enumerate}
  \def\labelenumi{\alph{enumi})}
  \setcounter{enumi}{1}
  \tightlist
  \item
    У більшості праць, присвячених розв'язанню гіперболічних задач,
    припускається, що початкове рівняння записується у канонічній формі,
    тобто у вигляді
  \end{enumerate}
\end{itemize}

Якщо у нас є рівняння і ми хочемо отримати його розв'язок, ми повинні
перетворити його у канонічну форму і використати відомі результати.

в) Для численного решения канонического гиперболического уравнения
создано много различных программ для ЭВМ.

\begin{verbatim}
Функція  \Phi  (  \xi  ,  \xi \eta  , u,  u_{\xi}  ,  u_{\eta}  ) вводиться у VM  \Theta  як підпрограма, тому рівняння спочатку має бути канонічним. Як тільки розв'язок знайдено в нових координатах, його завжди можна перетворити назад на початкові координати.
\end{verbatim}

\section{ЗАДАЧИ}\label{ux437ux430ux434ux430ux447ux438-2}

\begin{itemize}
\tightlist
\item
  \begin{enumerate}
  \def\labelenumi{\arabic{enumi}.}
  \tightlist
  \item
    Выясните, являются ли приведенные ниже уравнения гыперболическими,
    параболическими или эллиптическими:
  \end{enumerate}
\end{itemize}

\begin{verbatim}
- : **a)**  u_{xx} u_{xy} = 0  ,
  - 6)  u_{tt} = u_{xx} + u_x + hu  ,
  - **B)**  u_{xx} + 3u_{yy} = \sin x  ,
  - r)  u_{xx} + u_{yy} = f(x, y)  ,
\end{verbatim}

\begin{verbatim}
-  \mathbf{D} ) \ u_{rr} + \frac{1}{r} u_r + \frac{1}{r^2} u_{\theta \theta} = f(r, \ \theta).  - 2. Получите соотношения (23.4), (23.5) и (23.6).
\end{verbatim}

\begin{itemize}
\tightlist
\item
  \begin{enumerate}
  \def\labelenumi{\arabic{enumi}.}
  \setcounter{enumi}{2}
  \tightlist
  \item
    Убедитесь в том, что уравнение
  \end{enumerate}
\end{itemize}

гиперболического типа при всех x и y и найдите
\texttt{x\ a\ p\ a\ \textbackslash{}kappa\ m\ e\ p\ u\ c} - m u + e c
\kappa u + e c \kappa u + e c \kappa u + e c \kappa u + e c \kappa u + e
c \kappa u + e c \kappa u + e c \kappa u + e c \kappa u + e c \kappa u +
e c \kappa u + e c \kappa u + e c \kappa u + e c \kappa u + e c \kappa u
+ e c \kappa u + e c \kappa u + e c \kappa u + e c \kappa u + e e c
\kappa u + e c \kappa u + e c \kappa u + e c \kappa u + e c \kappa u + e
c \kappa u + e c \kappa u + e c \kappa u + e c \kappa u + e c \kappa u +
e c \kappa u + e c \kappa u + e c \kappa u + e c \kappa u + e c \kappa u
+ e c \kappa u + e c \kappa u + e c \kappa u + e c \kappa u + e c \kap u
+ e c \kap

pa u + e c \kappa u + e c \kappa u + e c \kappa u + e c \kappa u + e c
\kappa u + e c \kappa u + e c \kappa u + e c \kappa u + e c \kappa u + e
c \kappa u + e c \kappa u + e c \kappa u + e c \kappa u + e c \kappa u +
e c \kappa u + e c \kappa u + e c \kappa u + e c \kappa u + e c \kappa u
+ e c \kappa u + e c \kappa u + e c \kappa u + e e c \kappa u + e c
\kappa u + e c \kappa u + e c \kappa u + e c \kappa u + e c \kappa u + e
c \kappa u + e c \kappa u + e c \kappa u + e c \kappa u + e c \kappa u +
e c \kappa u + e c \kappa u + e c \kappa u + e c \kappa u + e c \kappa u
+ e c \kappa u + e c \kappa u + e c \kappa u + e c \kappa u + e c
\kappa u + e c \kappa u + e c \kappa u + e c \kappa u + e c \kappa u + e
c \kappa u + e c \kappa u + e c \kappa u + e c \kappa u + e c \kappa u +
e c \kappa u + e c \kappa u + e c \kappa u + e c \kappa u + e c \kappa u
+ e c \kappa u + e c \kappa u + e c \kappa u + e c \kappa u + e c
\kappa u + e u + e c \kappa u + e c \kappa u + e c \kappa u + e c
\kappa u + e c \kappa u + e c \kappa u + e c \kappa u + e c \kappa u + e
c \kappa u + e c \kappa u + e c \kappa u + e c \kappa u + e c \kappa u +
e c \kappa u + e c \kappa u + e c \kappa u + e c \kappa u + e c \kappa u
+ e c \kappa u + e c ~kappa u + e c \kappa u + e c \kappa u + e c
\kappa u + e c \kappa u + e c \kappa u + e c \kappa u + e c \kappa u + e
c \kappa u + e c \kappa u + e c \kappa u + e c \kappa u + e c \kappa u +
e c \kappa u + e c \kappa u + e c \kappa u + e c \kappa u + e c \kappa u
+ e c \kappa u + e e c \kappa u + e c \kappa u + e c \kappa u + e c
\kappa u + e c \kappa u + e c \kappa u + e c \kappa u + e c \kappa u + e
c \kappa u + e c \kappa u + e c \kappa u + e c \kappa u + e c \kappa u +
e c \kappa u + e c \kappa u + e c \kappa u + e c \kappa u + e c \kappa u
+ e c \kappa u + e c \kappa u + e c \kappa u + e c \kappa u + e c
\kappa u + e c \kappa u + e c \kappa u + e c \kappa u + e c \kappa u + e
c \kappa u + e c \kappa u + e c \kappa u + e c \kappa u + e c \kappa u +
e c \kappa u + e c \kappa u + e c \kappa u + e c \kappa u + e c \kappa u
+ e c \kappa u + e c \kappa u + e c \kappa u + e u + e c \kappa u + e c
\kappa u + e c \kappa u + e c \kappa u + e c \kappa u + e c \kappa u + e
c \kappa u + e c \kappa u + e c \kappa u + e c \kappa u + e c \kappa u +
e c \kappa u + e c \kappa u + e c \kappa u + e c \kappa u + e c \kappa u
+ e c \kappa u + e c \kappa u + e c \kappa u + e c ~kappa u + e c
\kappa u + e c \kappa u + e c \kappa u + e c \kappa u + e c \kappa u + e
c \kappa u + e c \kappa u + e c \kappa u + e c \kappa u + e c \kappa u +
e c \kappa u + e c \kappa u + e c \kappa u + e c 4. Приведите уравнение
из предыдущей задачи к каноническому виду

\begin{verbatim}
u_{tt} = c^2 (u_{xx} + u_{yy}), \quad
\end{verbatim}

\begin{enumerate}
\def\labelenumi{\arabic{enumi}.}
\setcounter{enumi}{4}
\item
  Продолжите вадачу 4 и приведите уравнение ко второй канонической форме
\item
  Найдите характеристики уравнения
\end{enumerate}

Преобразуйте уравнение к новым координатам, решите и, вернувшись к
координатам x и y, получите решение исходной вадачи.

\section{ВОЛНОВОЕ УРАВНЕНИЕ В СВОБОДНОМ ПРОСТРАНСТВЕ (ДВУМЕРНЫЕ И
ТРЕХМЕРНЫЕ
ЗАДАЧИ)}\label{ux432ux43eux43bux43dux43eux432ux43eux435-ux443ux440ux430ux432ux43dux435ux43dux438ux435-ux432-ux441ux432ux43eux431ux43eux434ux43dux43eux43c-ux43fux440ux43eux441ux442ux440ux430ux43dux441ux442ux432ux435-ux434ux432ux443ux43cux435ux440ux43dux44bux435-ux438-ux442ux440ux435ux445ux43cux435ux440ux43dux44bux435-ux437ux430ux434ux430ux447ux438}

ЦЕЛЬ ЛЕКЦИИ: Решить задачу с начальными условиями

(УЧП)

\begin{verbatim}
\begin{cases} -\infty < x < \infty, \\ -\infty < y < \infty, \\ -\infty < z < \infty, \end{cases}
\begin{cases} -\infty < x < \infty, \\ -\infty < y < \infty, \\ -\infty < z < \infty, \end{cases}
\begin{cases} u(x, y, z, 0) = \varphi(x, y, z), \\ u_{t}(x, y, z, 0) = \psi(x, y, z) \end{cases}
P_{2}
\begin{aligned} u_{tt} &= c^2 \Delta u, \\ (HY) & \begin{cases} u(x, y, z, 0) = 0, \\ u_t(x, y, z, 0) = \psi(x, y, z), \end{cases} \end{aligned}
P_{1}
\end{verbatim}

(НУ)

\begin{verbatim}
u(x,y) = g(x,y) =
\end{verbatim}

використовується метод спуску, і доведено, що двовимірні розв'язки не
задовольняють принцип Гюйгенса. І нарешті, використовуючи метод спуску,
показано, що розв'язок одномірного рівняння виражається формулою
д'Аламбера, з якою ми познайомилися раніше.

У попередніх лекціях ми розглядали задачу з початковими умовами для
нескінченної вібруючої струни і показали, що її розв'язок задається
формулою д'Аламбера. Читачеві має бути зрозуміло, що в тривимірному
просторі одномірне хвильове рівняння описує плоскі хвилі. Наприклад,
звукові або електромагнітні хвилі на достатньо великих відстанях від
джерел можна вважати плоскими хвилями і, отже, описуватися одномірним
рівнянням. Існує така термінологія:

\begin{itemize}
\tightlist
\item
  \begin{enumerate}
  \def\labelenumi{\arabic{enumi})}
  \tightlist
  \item
    одномерные волны называются плоскими волнами,
  \end{enumerate}
\item
  \begin{enumerate}
  \def\labelenumi{\arabic{enumi})}
  \setcounter{enumi}{1}
  \tightlist
  \item
    двумерные волны называются цилиндрическими волнами.
  \end{enumerate}
\end{itemize}

\begin{enumerate}
\def\labelenumi{\arabic{enumi})}
\setcounter{enumi}{2}
\tightlist
\item
  трехмерные волны называются сферическими волнами.
\end{enumerate}

Другими словами, одномерное волновое уравнение может описывать либо
плоские волны в пространствах большего числа измерений, либо колебания
струны. Задача лекции --- обобщить формулу Даламбера на случай двух и
трех измерений.

\section{Волны в трехмерном
пространстве}\label{ux432ux43eux43bux43dux44b-ux432-ux442ux440ux435ux445ux43cux435ux440ux43dux43eux43c-ux43fux440ux43eux441ux442ux440ux430ux43dux441ux442ux432ux435}

Рассмотрим сферические волны в трехмерном пространстве при ваданных
начальных условиях, т. е. решим задачу с начальными условиями

(24.1)

где \texttt{\textbackslash{}overline\{\textbackslash{}psi\}} ---
\emph{среднее} значение начального распределения
\texttt{\textbackslash{}psi} \emph{по сфере} радиуса ct с центром в
точке (x, y, z), т. е.

(24.1)

\pandocbounded{\includegraphics[keepaspectratio]{images/_page_177_Picture_3.jpeg}}

\pandocbounded{\includegraphics[keepaspectratio]{images/_page_177_Picture_7.jpeg}}

\texttt{где\ экаком\ \textquotesingle{}\textbackslash{}Delta\textquotesingle{}\ обозначен\ дифференциальный\ оператор}

\begin{verbatim}
u_h(r, \theta) = \sum_{n=0}^{\infty} r^n [a_n \cos(n\theta) + b_n \sin(n\theta)].
u = \frac{\partial}{\partial t} [t\overline{\varphi}].
+ ct \cos \varphi  )  (ct)^2 \sin \varphi d\varphi d\theta  .
\end{verbatim}

\begin{verbatim}
Аргументи функції  \psi  проходити поверхнею сфери, якщо значення  \theta  і  \phi  змінюються відповідно в межах  [0, 2\pi]  і  [0, \pi]  (рисунок 24.1). Цю формулу можна інтерпретувати так:
\end{verbatim}

См Владимиров В. С. Уравнения математической физики. --- М.: Наука, 1971
и др. издания. --- Прим ред.

дующим образом: каждая точка пространства излучает расходящуюся (со
скоростью c) сферическую волну, через t секунд точка с координатами (x,
y, z) окажется под воздействием начального Р И С

\begin{verbatim}
u_p(r, \theta) = Ar^4 + Br^4 \cos(2\theta).
\end{verbatim}

. 24 1. Решение является средним значением начатьного распределения на
сфере.

возмущения, сосредоточенного на сфере раднуса ct с центром в заданной
точке (см. рис. 24.2).

Формула (24.3) дозволяє обчислити розв'язок задачі для більшості
початкових умов на комп'ютері. Можливо, читачу буде цікаво знайти
розв'язок для деяких простих функцій f.

\begin{verbatim}
u(x, t) = \frac{1}{2c} \int_{x-ct}^{x+ct} \varphi(s) ds.
\end{verbatim}

\begin{verbatim}
РИС. 24.2. Начальное возмущение  \psi  распространяется из каждой точки во все стороны.
\end{verbatim}

Чтобы получить полное решение, рассмотрим вторую половину задачи

(24.4)

\begin{verbatim}
u_p(r, \theta) = -\frac{r^4}{32} - \frac{r^4}{24} \cos(2\theta).
u_t(x, t) = \frac{1}{2} [\varphi(x + ct) + \varphi(x - ct)].
\sum_{n=0}^{\infty} [a_n \cos(n\theta) + b_n \sin(n\theta)] - \frac{1}{32} - \frac{1}{24} \cos(2\theta) = 0.
u_1(r, \theta) = \frac{1}{32} + \frac{1}{24}r^2\cos(2\theta) - \frac{r^4}{24}\cos(2\theta) - \frac{r^4}{32} = \frac{(r^4 - 1)}{32} - \frac{(r^4 - r^2)}{24}\cos(2\theta).
u = u_0 + u_1 = r \cos \theta - \frac{(r^4 - 1)}{32} - \frac{(r^4 - r^2)}{24} \cos (2\theta).
\Delta u = 0, \quad 0 < r < 1 + \frac{1}{4} \sin \theta,
(\text{УЧП}) \quad u_{tt} = c^2 \Delta u, \quad (x, y, z) \in \mathbb{R}^3,
u \left( 1 + \frac{1}{4} \sin \theta, \ \theta \right) = \cos \theta, \quad 0 \le \theta \le 2\pi.
(\text{HV}) \quad
u\left(1+\frac{1}{4}\sin\theta,\ \theta\right)=\cos\theta
u(1+\epsilon\sin\theta, \theta) = \cos\theta, \quad 0 \le \epsilon \le 1/4.
\begin{cases} u = \varphi, \\ u_t = \psi \end{cases}
\end{verbatim}

\pandocbounded{\includegraphics[keepaspectratio]{images/_page_350_Figure_2.jpeg}}

\begin{verbatim}
u(1+\varepsilon\sin\theta,\ \theta)=u(1,\ \theta)+u_r(1,\ \theta)(\varepsilon\sin\theta)+u_{rr}(1,\ \theta)\frac{(\varepsilon\sin\theta)^2}{21}+\ldots
\Delta u = 0, \quad 0 < r < 1 + \varepsilon \sin \theta,
u(1, \theta) + u_r(1, \theta) (\varepsilon \sin \theta) + u_{rr}(1, \theta) \frac{(\varepsilon \sin \theta)^2}{2!} + \dots = \cos \theta.
(46.8) u = u_0 + \varepsilon u_1 + \varepsilon^2 u_2 + \dots.
P_{0}
\begin{cases} \Delta u_{0} = 0, & 0 < r < 1 \text{ (внутри круга),} \\ u_{0}(1, \theta) = \cos \theta, & u_{0}(r, \theta) = r \cos \theta, \end{cases}
\begin{cases} -\infty < x < \infty, \\ -\infty < y < \infty, \end{cases}
P_{1}
\begin{cases} \Delta u_{1} = 0, & 0 < r < 1 \text{ (внутри круга),} \\ u_{1}(1, \theta) = -\sin \theta \frac{\partial u_{0}(1, \theta)}{\partial r} = -\sin \theta \cos \theta \end{cases}
\end{verbatim}

(НУ)

Тепер сформульуємо головну мету наших лекцій: знайти функцію, яка дає
мінімум (або максимум) функціоналу (44.1). Стратегія пошуку буде такою
ж, як при пошуку мінімуму функції в диференціальному численні. Там ми
знайшли критичні точки функції за умовою f'(x) = 0. У варіаційному
численні все буде трохи складніше, оскільки аргумент більше не є
числовою змінною, а функцією. Однак загальний підхід залишається
незмінним: ми обчислимо так звану функціональну похідну функції
\texttt{y(x)} і прирівняємо її до нуля. Це нове рівняння буде
аналогічним

Розглянемо розв'язок тривимірної задачі за початкових умов, які залежать
лише від двох змінних X і Y.

Сделав это, получим, что трехмерная формула

для розв'язку та опише циліндричні хвилі і, отже, дасть розв'язок
двовимірної задачі. Цей метод називається методом спуску. Після всіх
розрахунків ми отримуємо 1) (зовсім не тривіальне)

Перед нами стоїть розв'язок двовимірного хвильового рівняння, і хоча,
найімовірніше, його практичне застосування передбачає використання
комп'ютера, цікаво аналізувати його з точки зору принципу Гюйгенса. У
цьому розв'язці обидва інтеграли початкових збурень розглядаються вздовж
внутрішньої частини (тут слово «внутрішнє» відіграє вирішальну роль)
кола з радіусом cte e центр у точці (x, y). Іншими словами, якщо ми
проаналізуємо це рішення так само, як у тривимірному випадку, то
побачимо, що початкове збурення поширюється з чітко визначеним переднім
краєм і розмитим заднім краєм. Отже, принцип Гюйгенса не виконується у
просторі двох вимірів.

\begin{verbatim}
Нарешті, припустимо, що початкові збурення  \phi  і  \phi  залежать лише від однієї змінної. У цьому випадку виникають плоскі хвилі, і після застосування методу спуску отримуємо
\end{verbatim}

\&lt;sup\textgreater2) См., например, Тихонов А. Н., Самарский А. А.
Уравнения математической физики. --- М.: Наука, 1972 и др. издания. ---
Прим. ред.

хорошо известную формулу Даламбера

\begin{verbatim}
u_u + \frac{1}{4} u_t
\end{verbatim}

Слід зазначити, що в цьому випадку використання методу спуску вимагає
нетривіальних розрахунків. За формулою д'Аламбера початкове зміщення f
призводить до утворення гострого заднього краю, а початкова швидкість f
його розмиває. Іншими словами, одномірний випадок є незвичайним тим, що
принцип Гюйгенса виконується для початкових зсув і не застосовується для
початкових швидкостей. Можна стверджувати, що загалом принцип Гюйгенса
не виконується в одномірному випадку.

\section{ЗАМЕЧАНИЯ}\label{ux437ux430ux43cux435ux447ux430ux43dux438ux44f-1}

Метод спуску ми не викладали детально, оскільки неможливо зупинитися на
питанні, як можна зменшити розмірність інтегралів, включених у формули.
Усі ці розрахунки можна знайти в навчальній літературі, перелік якої
наведено наприкінці книги. Загальна ідея полягає в тому, що при
розв'язанні задач у низьковимірному просторі можна використати відоме
розв'язання подібної задачі у просторі вищої вимірності, а потім
спростити його, припускаючи, що початкові та граничні умови не залежать
від частини змінних. Читачеві має бути зрозуміло, що метод спуску можна
застосувати не лише до розглянутої проблеми.

\section{ЗАДАЧИ}\label{ux437ux430ux434ux430ux447ux438-3}

\begin{enumerate}
\def\labelenumi{\arabic{enumi}.}
\tightlist
\item
  Покажите, что в одномерном случае решение задачи
\end{enumerate}

\begin{verbatim}
u_t = (1+x) u_{xx}, \\ u(x, 0) = q(x), \quad -\infty < x < \infty,
u_1 = (1 + \varepsilon_1) u_{xx}
P_{0}
\begin{aligned} & (\forall \mathsf{Y}\Pi) \quad u_{tt} = c^2 u_{xx}, \\ & (\mathsf{H} \mathsf{Y}) \quad \begin{cases} u\left(x, \ 0\right) = 0, \\ u_t\left(x, \ 0\right) = \phi\left(x\right). \end{cases} \end{aligned}
\end{verbatim}

\texttt{,\ \textquotesingle{}(x,\ y,\ z)\ \textbackslash{}in\ R\^{}3\textquotesingle{}\ ,\ \textquotesingle{}u\ =\ 0\textquotesingle{}\ u\_t\ =}

\begin{enumerate}
\def\labelenumi{\arabic{enumi}.}
\setcounter{enumi}{3}
\tightlist
\item
  Решите аналогичную (см. задачу 3) задачу для двумерного уравнения
\end{enumerate}

(УЧП)

Зверніть увагу, що в перших двох задачах коефіцієнти є сталими. Читач
повинен розуміти, що параметр b має бути достатньо малим, а необережний
ряд може відрізнятися

\section{ЗАДАЧИ}\label{ux437ux430ux434ux430ux447ux438-4}

\begin{itemize}
\tightlist
\item
  \begin{enumerate}
  \def\labelenumi{\arabic{enumi}.}
  \tightlist
  \item
    Подставьте разложение (46.8) в задачу (46.7) и получите
    последовательность задач
    \texttt{P\_0,\textbackslash{}\ P\_1,\textbackslash{}\ P\_2,\textbackslash{}\ \textbackslash{}dots}
    . 2. Покажите, что нелинейную задачу
  \end{enumerate}
\end{itemize}

(НУ)

\begin{verbatim}
u(r, \theta) = u_0(r, \theta) + \frac{1}{4}u_1(r, \theta)
\Delta u = 0,
q_{xx} + \varphi_{uu} = 0
q_{uu} + q_{vv} = 0
\end{verbatim}

\pandocbounded{\includegraphics[keepaspectratio]{images/_page_353_Figure_9.jpeg}}

\begin{verbatim}
\begin{cases} u(x, 0) = x, \\ u_t(x, 0) = 0. \end{cases}
w = f(z)
\end{verbatim}

\pandocbounded{\includegraphics[keepaspectratio]{images/_page_355_Figure_5.jpeg}}

\begin{enumerate}
\def\labelenumi{\arabic{enumi}.}
\setcounter{enumi}{2}
\tightlist
\item
  Проиллюстрируйте и прокомментируйте решение трехмерной задачи
\end{enumerate}

(УЧП)

\pandocbounded{\includegraphics[keepaspectratio]{images/_page_355_Figure_7.jpeg}}

\begin{verbatim}
u + iv = (x + iy)^2 = x^2 - y^2 + 2ixy
\end{verbatim}

\begin{verbatim}
\left\{
\end{verbatim}

(НУ)

Чи знаєте ви, як виглядає графік цієї функції з сегментом {[}0, 1{]}?
Подумай про це. Зверніть увагу, що синусоїдальне перетворення функції
\texttt{f(x)} визначається лише для додатних цілих чисел n.~Іншими
словами, скінченний синус і косинус перетворюють функції у чисельні
послідовності.

\pandocbounded{\includegraphics[keepaspectratio]{images/_page_185_Figure_3.jpeg}}

РИС. 25.1. Графики функции \texttt{f(x)} = 1 и ее преобразования.

\section{Свойства
преобразований}\label{ux441ux432ux43eux439ux441ux442ux432ux430-ux43fux440ux435ux43eux431ux440ux430ux437ux43eux432ux430ux43dux438ux439}

Прежде чем приступить к решению задач, мы должны получить некоторые
свойства этих преобразований.

Якщо you(x, t) --- функція
\texttt{\textbackslash{}partial\ \textbackslash{}theta\ yx} змінних, і
ми виконуємо перетворення змінної x, то

\begin{enumerate}
\def\labelenumi{\arabic{enumi}.}
\setcounter{enumi}{4}
\tightlist
\item
  Решите аналогичную (см. задачу 3) задачу для одномерных плоских волн
\end{enumerate}

(УЧП)

(Зверніть увагу, що перетворення було виконано на змінній x, і отримана
послідовність залежить лише від часу t.)

А как быть с производными? Ниже приводятся несколько полезных формул

\begin{itemize}
\tightlist
\item
  \begin{enumerate}
  \def\labelenumi{\arabic{enumi}.}
  \setcounter{enumi}{5}
  \tightlist
  \item
    Какова физическая природа того факта, что принцип Гюйгенса не
    справедлив в двумерном случае?
  \end{enumerate}
\item
  \begin{enumerate}
  \def\labelenumi{\arabic{enumi}.}
  \setcounter{enumi}{6}
  \tightlist
  \item
    Воспользуйтесь формулой Лейбница
  \end{enumerate}
\end{itemize}

и продифференцируйте интеграл

\begin{verbatim}
w = \ln\left\{\frac{z-1}{z+1}\right\}
\end{verbatim}

по времени t.

\section{КОНЕЧНЫЕ ПРЕОБРАЗОВАНИЯ ФУРЬЕ (СИНУС- И
КОСИНУС-ПРЕОБРАЗОВАНИЯ)}\label{ux43aux43eux43dux435ux447ux43dux44bux435-ux43fux440ux435ux43eux431ux440ux430ux437ux43eux432ux430ux43dux438ux44f-ux444ux443ux440ux44cux435-ux441ux438ux43dux443ux441--ux438-ux43aux43eux441ux438ux43dux443ux441-ux43fux440ux435ux43eux431ux440ux430ux437ux43eux432ux430ux43dux438ux44f}

ЦЕЛЬ ЛЕКЦИИ: Ввести два новых интегральных преобразования (конечные
спиус- и косинус-преобразования):

(конечное синус-преобразование),

(конечное косинус-преобразование),

(обратное синус-преобразование),

\begin{verbatim}
(\text{УЧП}) \quad u_{tt} = u_{xx} + \sin(\pi x), \quad 0 < x < 1, \quad 0 < t < \infty,
\end{verbatim}

(обратное косинус-преобразование),

\begin{verbatim}
и показать, как с помощью этих преобразований решать краевые задачи (в частности, неоднородные). Ранее мы познакомились с регулярными преобразованиями Фурье и Лапласа и увидели, как применение этих преобразований к уравнению с частными производными сводит задачу к решению обыкновениого дифференциального уравнения. Обычное преобразование Фурье подразумевает, что преобразуемая переменная изменяется от  -\infty  до  +\infty  , и, следовательно, оно применяется для решения задач в свободном пространстве (без границ). В этой лекции мы покажем, как решать краевые задачи (с границами), преобразуя ограниченные переменные (с таким преобразованием мы встречаемся впервые).
\end{verbatim}

Давайте на время забудем о наших задачах, а займемся самими
преобразованнями. Дадим определение преобразований, приведем формулы
обращения и основные свойства. Для чего нужны преобразования и как ими
пользоваться, мы узнаем несколько позже. Впрочем, если говорить кратко,
то суть интегральных преобразований в разложении входящих в задачу
функций по

частотам, решении семейства задач для каждой частоты и после-

дующем суммировании всех полученных результатов.

Начнем с функции \texttt{f(x)}, определенной на интервале {[}0, L{]}.
Конечные синус- и косинус-преобразования этой функции определяются
формулами

\begin{verbatim}
Конечное синус-преобразование  S[f] = S_n = \frac{2}{L} \int_{\sigma}^{L} f(x) \sin(n\pi x/L) dx   n = 1, 2, \dots   C[f] = C_n = \frac{2}{L} \int_{\sigma}^{L} f(x) \cos(n\pi x/L) dx   n = 0, 1, \dots  Конечное косинус-преобразование
\end{verbatim}

Читатель, наверное, заметил, что эти преобразования не отличаются от
формул для определения коэффициентов Фурье при разложении их в ряд по
синусам или косинусам. Формулы обращения для этих преобразований---это
обычные ряды Фурье по синусам и косинусам.

Обратное синус-преобразование

\begin{verbatim}
(\text{ГУ}) \quad
\end{verbatim}

Обратное косинус-преобразование

Отметим, что в обратном косинус-преобразовании суммирование начинается с
n=0, а в синус-преобразовании --- с n=1.

\section{Примеры
синус-преобразования}\label{ux43fux440ux438ux43cux435ux440ux44b-ux441ux438ux43dux443ux441-ux43fux440ux435ux43eux431ux440ux430ux437ux43eux432ux430ux43dux438ux44f}

Пусть

тогда

\begin{verbatim}
= \ln\left[\frac{z-1}{z+1}\right] + i \arg\left[\frac{z-1}{z+1}\right]
\end{verbatim}

График функции \texttt{f(x)} и ее преобразования изображены на рис.
25.1. Обратное преобразование записывается в виде

\begin{verbatim}
= \arg\left[\frac{x^2+y^2-1+i2y}{(x+1)^2+y^2}\right] =
\end{verbatim}

(действительная форма отображения \texttt{w=z\^{}2} ).

Эта форма понадобится нам позже. В этом примере она интересна тем, что
показывает, как гиперболы в плоскости z переходят в координатные линии u
= const и v = const в плоскости w.

Теперь мы перейдем к изучению одного частного вида функций комплексного
переменного, которые осуществляют конформные отображения.

\section{Определение конформного
отображения}\label{ux43eux43fux440ux435ux434ux435ux43bux435ux43dux438ux435-ux43aux43eux43dux444ux43eux440ux43cux43dux43eux433ux43e-ux43eux442ux43eux431ux440ux430ux436ux435ux43dux438ux44f-1}

Відображення w=\texttt{f(z)} комплексної площини z у комплексну площину
w називається конформним у точці \texttt{z\_\{\textbackslash{}bullet\}}
z-площини, якщо похідна
\texttt{f\textquotesingle{}(z\_\{\textbackslash{}bullet\})\ \textbackslash{}neq\ 0}
. Відображення \texttt{f(z)} називається конформним у регіоні D, якщо
\texttt{f\textquotesingle{}(z)\ \textbackslash{}neq\ 0} в кожній точці
регіону D.

Наприклад, відображення \texttt{f(z)\ =\ z\^{}2} конформно всюди, крім z
= 0, оскільки
\texttt{z\textquotesingle{}(z)\ =\ 2z\ \textbackslash{}neq\ 0} для всіх
\texttt{z\ \textbackslash{}neq\ 0} . З іншого боку, відображення
\texttt{e\^{}z} конформно по всій z-площині, оскільки скрізь
\texttt{f\textquotesingle{}(z)\ =\ e\^{}z\ \textbackslash{}neq\ 0} . Яке
призначення конформних відображень? Відповідь полягає в тому, що при
розв'язанні рівняння Лапласа
\texttt{\textbackslash{}phi\_\{xx\}\ +\ \textbackslash{}phi\_\{yy\}\ =\ 0}
в деякій області площини змінних x і y можна розглядати цю площину як
площину комплексної змінної z. Розглянемо комплексне відображення w =
\texttt{f(z)} площини z у площину w, в якій введені координати (u, v).
Відображуючи w = \texttt{f(z} \phi\emph{\{xy\} + \phi}\{uy\} = 0
\texttt{)}, рівняння Лапласа
\texttt{\textbackslash{}varphi\_\{xx\}\ +\ \textbackslash{}varphi\_\{yy\}\ =\ 0}
перетворюється на нове диференціальне рівняння в похідних похідних, яке
залежить від нових координат і v. To --- це рівняння
\texttt{\textbackslash{}phi\_\{xx\}\ +\ \textbackslash{}phi\_\{uu\}\ =\ 0}
яке повернеться у рівняння
\texttt{\textbackslash{}varphi\_\{nn\}\ +\ \textbackslash{}varphi\_\{nn\}\ =\ 0}
. Таким чином, ми придумаємо ідею знайти конформне відображення, яке
перетворює область з комплексною межею на область з простою межею
(нагадаємо, що відображення \texttt{w=z\^{}*} переносить межу першого
квадранта площини z у дійсну вісь площини w).

Приведем несколько примеров, в которых будет показано, как можно решать
уравнения с частными производными с помощью конформных отображений.

Уравнение Лапласа в верхней полуплоскости.

Предположим, что нам необходимо решить следующую задачу Дирихле в
верхней полувлоскости (рис. 47.3):

(47.1)

\begin{verbatim}
\begin{cases} u(x, 0) = 1, \\ u_t(x, 0) = 0, \end{cases}
x^{2} + y^{2} = 1,
\begin{aligned} (Y & \Pi) & \varphi_{xx} + \varphi_{yy} = 0, & -\infty < x < \infty, & 0 < y < \infty, \\ (47.1) & \varphi(x, 0) = \begin{cases} 0, & |x| > 1, \\ 1, & |x| \le 1. \end{cases} \end{aligned}
\end{verbatim}

Решение задачи выполним за три следующих шага:

ШАГ І (Определение подходящего преобразования.)

Поскольку переменная x меняется от 0 до 1, мы воспользуемся конечным
преобразованием, а именно синус-преобразованием; позже станет ясно
почему. Мы могли бы решить эту задачу с помощью преобразования Лапласа
по переменной t (по трудоемкости этот метод не отличается от конечного
синус-преобразования).

ШАГ 2 (Выполнение преобразования).

Для удобства будем использовать обозначение
\texttt{S\_n(t)\ =\ S{[}u{]}} . Применим синус-преобразование к
исходному уравнению

\begin{verbatim}
w = \ln\left\{\frac{z-1}{z+1}\right\}
\end{verbatim}

и воспользуемся тождеством для синус-преобразования, получим

\begin{verbatim}
\begin{split} S\left[u_{t}\right] &= \frac{dS\left[u\right]}{dt}, \quad S\left[u_{tt}\right] = \frac{d^{2}S\left[u\right]}{dt^{2}}, \\ S\left[u_{xx}\right] &= -\left[n\pi/L\right]^{2}S\left[u\right] + \frac{2n\pi}{L^{2}}\left[u\left(0,\ t\right) + (-1)^{n+2}u\left(L,\ t\right)\right], \\ C\left[u_{xx}\right] &= -\left[n\pi/L\right]^{2}C\left[u\right] - \frac{2}{L}\left[u_{x}\left(0,\ t\right) + (-1)^{n+1}u_{x}\left(L,\ t\right)\right]. \end{split}
\begin{aligned} & (\text{УЧП}) & \phi_{xx} + \phi_{yy} = 0, & \text{внутри } D, \\ & (\text{ГУ}) & \begin{cases} \varphi(x, y) = 1 & \text{на } x^2 + y^2 = 1, \\ \varphi(x, y) = 2 & \text{на } (x - 1)^2 + y^2 = 9. \end{cases} \end{aligned}
\end{verbatim}

M A

\pandocbounded{\includegraphics[keepaspectratio]{images/_page_357_Figure_11.jpeg}}

THBLK1213X M A

\pandocbounded{\includegraphics[keepaspectratio]{images/_page_358_Figure_2.jpeg}}

THBLK1214X

\begin{verbatim}
\begin{cases} (\text{УЧП}) & \varphi_{uu} + \varphi_{vv} = 0, \quad -\infty < u < \infty, \quad 0 < v < \pi, \\ (\text{47.2}) & \begin{cases} \varphi(u, 0) = 0, \\ \varphi(u, \pi) = 1. \end{cases} \end{cases}
\end{verbatim}

где

\begin{verbatim}
w = u + iv = \ln\left[\frac{z-1}{z+1}\right] =
u = \gamma [(x-s)(x-t)-\gamma^2],
\begin{cases} u(0, t) = 0, \\ u(1, t) = 0, \end{cases}
\gamma = 2t/[(x-t)^2 + y^2].
= \ln\left[\frac{z-1}{z+1}\right] + i \arg\left[\frac{z-1}{z+1}\right]
\end{verbatim}

Если теперь преобразовать начальные условия краевой задачи, то получим
начальные условия для обыкновенного дифференциального уравнения

\begin{verbatim}
v = \arg\left[\frac{z-1}{z+1}\right] = \arg\left[\frac{x+iy-1}{x+1+iy}\right] =
= \arg\left[\frac{x^2+y^2-1+i2y}{(x+1)^2+y^2}\right] =
= \arg\left[\frac{2y}{x^2+y^2-1}\right]
\varphi(r) = a \ln r + b, \quad r = u^2 + v^2.
\begin{cases} u(x, 0) = 1, \\ u_t(x, 0) = 0, \end{cases}
\varphi(x, y) = 0.57 \ln(u^2 + v^2) + 1,
\varphi(x, y) = \frac{1}{\pi} \arctan\left[\frac{2y}{x^2 + y^2 - 1}\right].
\end{verbatim}

Итак, решим теперь новое семейство задач Кощи:

(ОДУ)

\begin{verbatim}
x^{2} + y^{2} = 1,
w = \ln \left[ \frac{z-1}{z+1} \right].
(x-1)^{2} + y^{2} = 9,
\begin{cases} 1, & n = 1, \\ 0, & n = 2, 3, \dots, \end{cases}
\begin{aligned} & (\text{УЧП}) & \phi_{xx} + \phi_{yy} = 0, & \text{внутри } D, \\ & (\text{ГУ}) & \begin{cases} \varphi(x, y) = 1 & \text{на } x^2 + y^2 = 1, \\ \varphi(x, y) = 2 & \text{на } (x - 1)^2 + y^2 = 9. \end{cases} \end{aligned}
\begin{split} \frac{d^2S_n\left(t\right)}{dt^2} &= -\left(\pi n\right)^2S_n\left(t\right) + 2n\pi\left[u\left(0,\ t\right) + \left(-1\right)^{n+1}u\left(1,\ t\right)\right] + D_n\left(t\right) = \\ &= -\left(n\pi\right)^2S_n\left(t\right) + D_n\left(t\right), \end{split}
\right.
\end{verbatim}

M A

\pandocbounded{\includegraphics[keepaspectratio]{images/_page_359_Figure_10.jpeg}}

THBLK1228X

\begin{verbatim}
w = 2t \left[ \frac{z-s}{z-t} \right],
u = \gamma [(x-s)(x-t)-\gamma^2],
v = \gamma [y(x-t)],
\end{verbatim}

\pandocbounded{\includegraphics[keepaspectratio]{images/_page_362_Figure_2.jpeg}}

\begin{verbatim}
\gamma = 2t/[(x-t)^2 + y^2].
\end{verbatim}

\pandocbounded{\includegraphics[keepaspectratio]{images/_page_362_Figure_5.jpeg}}

\begin{verbatim}
\begin{aligned} \varphi_{uu} + \varphi_{vv} &= 0, & \text{в кольче,} \\ (\Gamma Y) & & & & & & & \\ (\Gamma Y) & & & & & & \\ (\Gamma Y) & & & & & \\ & & & & & \\ (\Gamma Y) & & & & & \\ & & & & \\ & & & & \\ & & & & \\ & & & & \\ & & & \\ & & & \\ & & & \\ & & & \\ & & & \\ & & & \\ & & \\ & & \\ & & \\ & & \\ & & \\ & & \\ & & \\ & & \\ & & \\ & & \\ & & \\ & & \\ & & \\ & & \\ & & \\ & & \\ & & \\ & & \\ & & \\ & & \\ & & \\ & & \\ & & \\ & & \\ & & \\ & & \\ & & \\ & & \\ & & \\ & & \\ & & \\ & & \\ & & \\ & & \\ & & \\ & & \\ & & \\ & & \\ & & \\ & & \\ & & \\ & & \\ & & \\ & & \\ & & \\ & & \\ & & \\ & & \\ & & \\ & & \\ & & \\ & & \\ & & \\ & & \\ & & \\ & & \\ & & \\ & & \\ & & \\ & & \\ & & \\ & & \\ & & \\ & & \\ & & \\ & & \\ & & \\ & & \\ & & \\ & & \\ & & \\ & & \\ & & \\ & & \\ & & \\ & & \\ & & \\ & & \\ & & \\ & & \\ & & \\ & & \\ & & \\ & & \\ & & \\ & & \\ & & \\ & & \\ & & \\ & & \\ & & \\ & & \\ & & \\ & & \\ & & \\ & & \\ & & \\ & & \\ & & \\ & & \\ & & \\ & & \\ & & \\ & & \\ & & \\ & & \\ & & \\ & & \\ & & \\ & & \\ & & \\ & & \\ & & \\ & & \\ & & \\ & & \\ & & \\ & & \\ & & \\ & & \\ & & \\ & & \\ & & \\ & & \\ & & \\ & & \\ & & \\ & & \\ & & \\ & & \\ & & \\ & & \\ & & \\ & & \\ & & \\ & & \\ & & \\ & & \\ & & \\ & & \\ & & \\ & & \\ & & \\ & & \\ & & \\ & & \\ & & \\ & & \\ & & \\ & & \\ & & \\ & & \\ & & \\ & & \\ & & \\ & & \\ & & \\ & & \\ & & \\ & & \\ & & \\ & & \\ & & \\ & & \\ & & \\ & & \\ & & \\ & & \\ & & \\ & & \\ & & \\ & & \\ & & \\ & & \\ & & \\ & & \\ & & \\ & & \\ & & \\ & & \\ & & \\ & & \\ & & \\ & & \\ & & \\ & & \\ & & \\ & & \\ & & \\ & & \\ & & \\ & & \\ & & \\ & & \\ & & \\ & & \\ & & \\ & & \\ & & \\ & & \\ & & \\ & & \\ & & \\ & & \\ & & \\ & & \\ & & \\ & & \\ & & \\ & & \\ & & \\ & & \\ & & \\ & & \\ & & \\ & & \\ & & \\ & & \\ & & \\ & & \\ & & \\ & & \\ & & \\ & & \\ & & \\ & & \\ & & \\ & & \\ & & \\ & & \\ & & \\ & & \\ & & \\ & & \\ & & \\ & & \\ & & \\ & & \\ & & \\ & & \\ & & \\ & & \\ & & \\ & & \\ & & \\ & & \\ & & \\ & & \\ & & \\ & & \\ & & \\ & & \\ & & \\ & & \\ & & \\ & & \\ & & \\ & & \\ & & \\ & & \\ & & \\ & & \\ & & \\ & & \\ & & \\ & & \\ & & \\ & & \\ & & \\ & & \\ & & \\ & & \\ & & \\ & & \\ & & \\ & & \\ & & \\ & & \\ & & \\ & & \\ & & \\ & & \\ & & \\ & & \\ & & \\ & & \\ & & \\ & & \\ & & \\ & & \\ & & \\ & & \\ & & \\ & & \\ & & \\ & & \\ & & \\ & & \\ & & \\ & & \\ & & \\ & & \\ & &
\end{verbatim}

\pandocbounded{\includegraphics[keepaspectratio]{images/_page_362_Figure_7.jpeg}}

\begin{verbatim}
\begin{cases} 4/n\pi, & n = 1, 3, \dots, \\ 0, & n = 2, 4, \dots, \end{cases}
\varphi(r) = a \ln r + b, \quad r = u^2 + v^2.
\varphi(u, v) = 0.57 \ln(u^2 + v^2) + 1
H(x-a) =
\varphi(x, y) = 0.57 \ln(u^2 + v^2) + 1,
\begin{cases} 0, x < a, \\ 1, x \geqslant a, \end{cases}
w = \ln \left[ \frac{z-1}{z+1} \right].
\end{verbatim}

В результате получаем

\begin{verbatim}
z = re^{i\theta}
H(x-a) =
\begin{cases} 1, & n = 1, \\ 0, & n = 2, 3, \dots, \end{cases}
\begin{cases} 0, & x < a, \\ 1, & x \geqslant a, \end{cases}
\end{verbatim}

\begin{verbatim}
\begin{aligned} (\text{Y} \forall \Pi) & \qquad \phi_{\text{vx}} + \phi_{yy} = 0, \quad 0 < x < \infty, \quad 0 < y < \infty, \\ \left\{ \begin{array}{l} \phi(x, \ 0) = \begin{cases} 1, & 0 < x < 1, \\ 0, & 1 \leqslant x < \infty, \end{cases} \\ \phi(0, \ y) = \begin{cases} 1, & 0 < y < 1, \\ 0, & 1 \leqslant y < \infty. \end{cases} \end{aligned}
\end{verbatim}

\begin{verbatim}
f(x) = \mathcal{F}^{-1}[F] = \frac{1}{\sqrt{2\pi}} \int_{-\infty}^{+\infty} F(\omega) e^{-i\omega x} d\omega \qquad F(\omega) = \mathcal{F}[f] = \frac{1}{\sqrt{2\pi}} \int_{-\infty}^{+\infty} f(x) e^{-i\omega x} dx
\begin{cases} 1, |x| < a \\ 0, |x| > a \end{cases}
\begin{cases} S_n(0) = \begin{cases} 4/n\pi, & n = 1, 3, \dots, \\ 0, & n = 2, 4, \dots, \end{cases} \end{cases}
\end{verbatim}

\section{ЗАМЕЧ
АНИЯ}\label{ux437ux430ux43cux435ux447-ux430ux43dux438ux44f}

\begin{enumerate}
\def\labelenumi{\arabic{enumi}.}
\tightlist
\item
  Для решения задачи методом конечного синус- или косинуспреобразования
  краевые условия при x=0 и x=L должны иметь вид
\end{enumerate}

\begin{verbatim}
\begin{cases} 1, |x| < 1 \\ 0, |x| > 1 \end{cases}
\begin{cases} \frac{dS_n(0)}{dt} = 0. \end{cases}
\end{verbatim}

\texttt{f(x)} * \texttt{g(x)}

\begin{verbatim}
(1+x^2)^{-1}
\end{verbatim}

,

где

\begin{verbatim}
\frac{a}{x^2+a^2}
\frac{2ax}{(x^2+a^2)^2}
\end{verbatim}

\pandocbounded{\includegraphics[keepaspectratio]{images/_page_189_Picture_10.jpeg}}

\begin{enumerate}
\def\labelenumi{\arabic{enumi}.}
\setcounter{enumi}{17}
\tightlist
\item
\end{enumerate}

\begin{verbatim}
\begin{cases} 1-|x|, |x| < 1 \\ 0, |x| > 1 \end{cases}
\end{verbatim}

\begin{enumerate}
\def\labelenumi{\arabic{enumi}.}
\setcounter{enumi}{8}
\tightlist
\item
\end{enumerate}

\begin{verbatim}
\begin{cases} u(0, t) = f(t) \\ u(L, t) = g(t) \end{cases}
\frac{2}{\pi} \frac{\sin a\omega}{\omega}
\frac{2}{\pi} \frac{\sin a\omega}{(\omega^2+a^2)^2}
\end{verbatim}

\texttt{f(x)} * \texttt{g(x)}

\begin{enumerate}
\def\labelenumi{\arabic{enumi}.}
\setcounter{enumi}{1}
\tightlist
\item
  Решите в общем виде задачу
\end{enumerate}

\begin{verbatim}
\frac{\pi}{2} e^{-a+\omega}
\frac{\pi}{2} e^{-a+\omega}
\frac{2ax}{(x^2+a^2)^2}
\frac{\cos(ax), |x| > \pi/2a}{0, |x| > \pi/2a}
\end{verbatim}

\section{ПРИНЦИП СУПЕРПОЗИЦИИ--- ОСНОВА ТЕОРИИ ЛИНЕЙНЫХ
СИСТЕМ}\label{ux43fux440ux438ux43dux446ux438ux43f-ux441ux443ux43fux435ux440ux43fux43eux437ux438ux446ux438ux438-ux43eux441ux43dux43eux432ux430-ux442ux435ux43eux440ux438ux438-ux43bux438ux43dux435ux439ux43dux44bux445-ux441ux438ux441ux442ux435ux43c}

МЕТА ЛЕКЦІЇ: Ознайомити читача з принципом суперпозиції та показати, як
його можна використати для розкладання початкової задачі на підзадачі,
потім розв'язання всіх підзадач одночасно і, підсумувавши отримані
результати, знаходження рішення початкової задачі. Ми також хочемо
показати, що два основні методи розв'язання лінійних рівнянь --- метод
розділення змінних і метод інтегрального перетворення --- використовують
принцип суперпозиції.

Припустимо, що інженеру потрібно знайти відгук лінійної системи на
вхідний ефект f.~Загальний підхід до розв'язання задачі такий:

\begin{itemize}
\tightlist
\item
  \begin{enumerate}
  \def\labelenumi{\arabic{enumi})}
  \tightlist
  \item
    разложить f на элементарные воздействия
    \texttt{f\ =\ \textbackslash{}sum\ f\_k} ;
  \end{enumerate}
\item
  \begin{enumerate}
  \def\labelenumi{\arabic{enumi})}
  \setcounter{enumi}{1}
  \tightlist
  \item
    найти отклик системы \texttt{u\_k} на воздействие \texttt{f\_k} ;
  \end{enumerate}
\end{itemize}

\begin{enumerate}
\def\labelenumi{\arabic{enumi})}
\setcounter{enumi}{2}
\tightlist
\item
  сложить отклики
  \texttt{u\ =\ \textbackslash{}sum\_\{i=1\}\^{}\{n\}\ u\_\{i\}} .
\end{enumerate}

Если система линейная, то полученная сумма и является откликом системы
на входное воздействие f.~В этом состоит принцип суперпоэиции (см. рнс.
26.1).

РИС, 26.1. Основная идея принципа суперпоэнции; L---дифференциальный
оператор; f---вход; u---выход.

Эту основную идею можно использовать при решении сложных смешанных
задач, разбивая их на простые подзадачи, ре-

шая каждую из них и затем суммируя полученные результаты (конечно,
дифференциальное уравнение и граничные условия должны быть линейны).

\section{Разложение смешанной задачи на две более
простые}\label{ux440ux430ux437ux43bux43eux436ux435ux43dux438ux435-ux441ux43cux435ux448ux430ux43dux43dux43eux439-ux437ux430ux434ux430ux447ux438-ux43dux430-ux434ux432ux435-ux431ux43eux43bux435ux435-ux43fux440ux43eux441ux442ux44bux435}

Предположим, нам необходимо решить линейную задачу (обозначим ее P)

(УЧП)

(Р) (ГУ)

\begin{verbatim}
\cos(ax)
\begin{cases} u(0, t) = 0, \\ u(1, t) = 0, \\ u(x, 0) = \sin(2\pi x), \quad 0 \le x \le 1. \end{cases}
\frac{\pi}{2} [\delta(\omega+a) + \delta(\omega-a)]
F(\omega) = \frac{?}{\pi} \int_{0}^{\infty} f(x) \sin(\omega x) dx
\begin{cases} u(0, t) = 0, \\ u(1, t) = 0, \\ u(1, t) = 0, \end{cases}
\end{verbatim}

f'\,'(0)

\begin{verbatim}
\frac{1}{\alpha} F\left(\frac{\omega}{\alpha}\right)
\frac{2\omega}{\pi (\alpha^2 + \omega^2)}
[2/\pi\omega]^{1/2}
\frac{2}{\pi\omega} [1-\cos(\omega a)]
e^{-1}a + \omega
\frac{1}{2}e^{-\alpha}\sin\omega
\begin{cases} u_t = u_{xx} + f(x, t), \\ u(0, t) = 0, \\ u(1, t) = 0, \\ u(x, 0) = \varphi(x) \end{cases}
\cos(ax)
f(x) = \int_{0}^{\infty} F(\omega) \cos(\omega x) d\omega \qquad \qquad F(x) = \frac{2}{\pi} \int_{0}^{\infty} f(x) \cos(\omega x) dx
\begin{cases} u_t = u_{xx} + f(x, t), \\ u(0, t) = 0, \\ u(1, t) = 0, \\ u(x, 0) = 0 \end{cases}
S_n = \frac{2}{\pi} \int_0^{\pi} f(x) \sin(nx) dx.
y = \frac{\pi}{b - a} (x - a).
\begin{cases} u_t = u_{xx}, \\ u(0, t) = 0, \\ u(1, t) = 0, \\ u(x, 0) = \varphi(x). \end{cases}
f(x) = \sum_{n=1}^{\infty} S_n \sin(nx).
f(x) = \sum_{n=1}^{\infty} S_n \sin(nx)
\end{verbatim}

f'\,'(0)

\begin{verbatim}
S_n = \frac{2}{\pi} \int_{0}^{\pi} f(x) \sin(nx) dx
0 \le x \le \pi
n = 1, 2, ...
\end{verbatim}

f'\,'(x)

\begin{verbatim}
\frac{\pi}{2} e^{ax}
C_n = \frac{2}{\pi} \int_0^{\pi} f(x) \cos(nx) dx.
y = \frac{\pi}{b-a}(x-a).
f(x) = \frac{C_0}{2} + \sum_{n=1}^{\infty} C_n \cos(nx) \qquad C_n = \frac{2}{\pi} \int_{0}^{\pi} f(x) \cos nx \, dx
\end{verbatim}

f'\,'(x)

\begin{verbatim}
\begin{array}{ll} u\left(x, \ 0\right) = \sin{(\pi x)}, \\ u_{t}\left(x, \ 0\right) = 0, \end{array}
\end{verbatim}

\begin{enumerate}
\def\labelenumi{\alph{enumi})}
\tightlist
\item
\end{enumerate}

\begin{verbatim}
\begin{split} A_n(t) &= 2\int_0^1 u_t \sin(n\pi x) \, dx = \frac{d}{dt} \bigg[ 2\int_0^1 u(x, t) \sin(n\pi x) \, dx \bigg] = \frac{dU_n(t)}{dt} \,, \\ B_n(t) &= 2\int_0^1 u_{xx} \sin(n\pi x) \, dx = -(n\pi)^2 \, U_n(t) \,+ \\ &\quad + 2n\pi \, \big[ u(0, t) + (-1)^{n+1} \, u(1, t) \big], \end{split}
\end{verbatim}

\begin{verbatim}
5. \cos(mx), m=1, 2, \dots
\end{verbatim}

\texttt{\textbackslash{}mathbf\{\textbackslash{}Pi\}} для яких рівнянь
функція \texttt{u\_1\ +\ u\_2} також буде розв'язком? Які висновки ви
можете зробити зі своїх відповідей?

\begin{enumerate}
\def\labelenumi{\arabic{enumi}.}
\setcounter{enumi}{3}
\item
  Найдите четыре смещанные задачи, сумма решений которых дает решение
  следующей задачи:
\item
  Решите задачу Коши
\end{enumerate}

(OAY)

\begin{verbatim}
\begin{cases} u(0, t) = 0, \\ u(1, t) = 0 \end{cases}
\end{verbatim}

(HY) \texttt{U\_n}(0) = 0. Можете ли вы проверить полученное решение?

\begin{enumerate}
\def\labelenumi{\arabic{enumi}.}
\setcounter{enumi}{5}
\tightlist
\item
  Предположим, что функции \texttt{u\_1} и \texttt{u\_2} удовлетворяют
  линейным однородным граничным условиям
\end{enumerate}

\begin{verbatim}
\begin{cases}  u_x(1, t) + h_2 u(1, t) = 0.  Будет ли функция   u_i + u_i   удовлетворять этим условиям? , .
\end{verbatim}

Лекция 27

\section{УРАВНЕНИЯ ПЕРВОГО ПОРЯДКА (МЕТОД
ХАРАКТЕРИСТИК)}\label{ux443ux440ux430ux432ux43dux435ux43dux438ux44f-ux43fux435ux440ux432ux43eux433ux43e-ux43fux43eux440ux44fux434ux43aux430-ux43cux435ux442ux43eux434-ux445ux430ux440ux430ux43aux442ux435ux440ux438ux441ux442ux438ux43a}

ЦЕЛЬ ЛЕКЦИИ: Ввести понятие уравнения с частными производными первого
порядка (до сих пор мы рассматривали уравнения второго порядка) и
познакомиться с важным методом решения задач с начальными
условиями---методом характеристик. Задача, которой мы теперь займемся,
имеет вид

(УЧП)

Сразу же отметим, что впервые мы будем решать задачу с переменными
коэффициентами. Оказывается, если от координат (x, t) перейти к новым,
характеристическим координатам (s, \tau) , то наше УЧП превратится в
обыкновенное дифференциальное уравнение. Затем можно решить обыкновенное
дифференциальное уравнение, т. е. найти u(s, \tau) , а на последнем шаге
выразить s и t через x и t и получить \texttt{u(x,\ t)}.

Напомним читателю, что, когда мы решали уравнение дифф

константа \alpha\^{}2 играла роль коэффициента диффузии, а v---была
скоростью среды. Очевидно, что если \alpha=0 (т. е. диффузии нет), то
решение будет смещаться как целое вдоль оси x со скоростью v (у нас
осталась только конвекция). Другими словами, если начальное условие
\texttt{u(x,0)}=\varphi(x) , то соответствующее решение уравнения

будет иметь вид \texttt{u(x,\ t)} = \varphi(x-vt) .

Это дает нам основания думать, что решение уравнения первого порядка

переносится вдоль потока со скоростью

Конечно, если a и b---константы, то решение представляет собой бегущую с
постоянной скоростью волну. Если же \texttt{a(x,\ t)} и
\texttt{b(x,\ t)}

РИС. 27.1. Начальное значение в точке x оказывает влияние на решение
только вдоль характеристики в плоскости переменных (x,t). Параметр s
изменяется вдоль характеристик от нуля до бесконечности, начиная с линии
начального условия; \tau   постоянна вдоль каждой характеристики;
\tau   изменяется вдоль линии начального условия.

зависят от x и t, то скорость потока изменяется как вдоль потока, так и
во времени (читатель сможет увидеть, что начальная кривая сильно
искажается). Здесь так много аналогий с конвекцией!

Вернемся к нашей основной задаче:

(УЧП)

(НУ) \texttt{u(x,\ 0)} = \varphi(x), -\infty \textless{} x \textless{}
\infty.

Решение этого линейного уравнения первого порядка основывается на
следующем физическом факте: начальное условие в некоторой точке x
переносится в tx-плоскости вдоль линии, которая

называется характеристи-кой (см. рис. 27.1)

Этим наше уразнение отличается от других уравнений (таких, как уравнение
теплопроводности \texttt{u\_t\ =\ u\_\{xx\}\ )},для которых начальное
значение в точке x оказывает влияние на решение во всех точках
пространства и во все моменты времени. Если читатель помнит, то
начальное сме- Р И С

\texttt{.\ 27.2.\ Начальное\ возмущение\ \ \ \textquotesingle{}u\textbackslash{}left(x,0\textbackslash{}right)\textquotesingle{}\ \ \ в\ точке\ порождает\ две\ волны.\ Начальное\ возмущение\ струны\ согласно\ волновому\ уравнению\ распространяется\ вдоль\ двух\ характеристик.}

щение скрипичной струны в точке x оказывает влияние на решение вдоль
\partial syx линий в tx-плоскости (соответствующих двум бегущим волнам)
(см. рис. 27.2).

Естественно, напрашивается идея ввести две новые координаты s и
\tau   (вместо x и t) так, чтобы

\begin{itemize}
\tightlist
\item
  s изменялась вдоль характеристических кривых,
\item
  т изменялась вдоль начальной кривой (желательно вдоль линии t=0).
\end{itemize}

Рассмотрим сначала новую координату s. Если s выбрать так, что она
удовлетворяет приведенному выше условию, то уравнение

превратится в обыкновенное дифференциальное уравнение К о н

\begin{verbatim}
0 \le x \le \pi
\end{verbatim}

ечно, остается открытым вопрос, как найти эти характеристики. Ответ
прост: мы выберем характеристики \{{[}\texttt{x(s)}, \texttt{t(s)}{]}: 0
\textless{} s \textless{} \infty\} так, чтобы они удовлетворяли системе

f'\,'(x)

\begin{verbatim}
-\infty < y < +\infty -\infty < z < +\infty \begin{array}{lll}
 \begin{array}{ll} (\mathrm{Y} + \Pi) & u_{tt} = u_{xx} + \sin{(3\pi x)}, & 0 < x < 1, & 0 < t < \infty, \\ (\mathrm{\Gamma} \mathrm{Y}) & \left\{ \begin{array}{ll} u\left(0, \ t\right) = 0, \\ u\left(1, \ t\right) = 0, \\ \end{array}
\end{verbatim}

\begin{verbatim}
x(s) = s + c_1, \quad t(s) = s + c_2.
\end{verbatim}

Для каждой из подзадач выберите метод решения по своему BKYCY,

\begin{enumerate}
\def\labelenumi{\arabic{enumi}.}
\setcounter{enumi}{2}
\tightlist
\item
  Предположим, что \texttt{u\_1} и \texttt{u\_2} --- решения следующих
  уравнений:
\end{enumerate}

;

\begin{verbatim}
b)  u_t = u_{xx} + e^t  ;
\end{verbatim}

\begin{enumerate}
\def\labelenumi{\Alph{enumi})}
\setcounter{enumi}{1}
\tightlist
\item
\end{enumerate}

;

\begin{verbatim}
F)  u_t = u_{xx} + u^2  .
\frac{du}{ds} + 2u = 0, \quad 0 < s < \infty,
\begin{aligned} u_t &= u_{xx} + f(x, t), & 0 < x < 1, & 0 < t < \infty, \\ & \begin{cases} u(0, t) = g_1(t), \\ u(1, t) = g_2(t), \end{cases} & 0 < t < \infty, \\ & u(x, 0) = \varphi(x), & 0 \le x \le 1. \end{aligned}
\end{verbatim}

x

\textbar{} \_ \textbar, -\infty \textless{} x \textless{} \infty  , 0
\textless{} t \textless{} \infty  ,

\begin{verbatim}
s=t
(HV)   u(x, 0) = \varphi(x), \quad -\infty < x < \infty.
u(x, t) = \sin(x-t)e^{-xt}
\Delta u = u_{xx} + u_{yy},   -\infty < x < \infty  ,   0 < t < \infty
(\forall \Pi) \quad a(x, t) u_x + b(x, t) u_t + c(x, t) u = 0, \quad -\infty < x < \infty, \\ 0 < t < \infty,
\Delta u = u_{rr} + \frac{1}{r}u_r + \frac{1}{r^2}u_{\theta\theta}
u(x, 0) = f(x), \quad -\infty < x < \infty.
\Delta u = u_{xx} + u_{yy} + u_{zz}
\frac{dx}{ds} = a(x, t), \qquad \frac{dt}{ds} = b(x, t).
\Delta u = u_{rr} + \frac{1}{r}u_r + \frac{1}{r^2}u_{\theta\theta} + u_{zz}.
x = x (s, \tau),
\end{verbatim}

M A

\pandocbounded{\includegraphics[keepaspectratio]{images/_page_196_Figure_10.jpeg}}

THBLK1368X

\begin{verbatim}
\Delta u = u_{rr} + \frac{2}{r}u_r + \frac{1}{r^2}u_{\theta\theta} + \frac{\cot\theta}{r^2}
(\text{УЧП})
\end{verbatim}

M A

\pandocbounded{\includegraphics[keepaspectratio]{images/_page_197_Figure_5.jpeg}}

THBLK1370X

\begin{verbatim}
\Delta u = u_{rr} + \frac{2}{r}u_r + \frac{1}{r^2}u_{\theta\theta} + \frac{\cot\theta}{r^2}u_0 + \frac{1}{r^2\sin\theta}u_{\phi\phi}
\end{verbatim}

M A

\pandocbounded{\includegraphics[keepaspectratio]{images/_page_370_Figure_1.jpeg}}

THBLK1372X M A

\pandocbounded{\includegraphics[keepaspectratio]{images/_page_370_Picture_2.jpeg}}

THBLK1373X M A

\pandocbounded{\includegraphics[keepaspectratio]{images/_page_370_Figure_3.jpeg}}

THBLK1374X

\begin{verbatim}
\begin{cases} \frac{dx}{ds} = a(x, t), \\ \frac{dt}{ds} = b(x, t). \end{cases}
u(0) = F(\tau).
u(s, \tau) = F(\tau) e^{-s} \ell^2
s=t
u(x, t) = F(xe^{-t}) e^{-t^2/2}
\end{verbatim}

Уравнение одномерной диффузии Уравнение конвективной диффузии

Уравнение теплопроводности с поглощением Уравнение теплопроводности с
источником

\section{Диагонализация
матрицы}\label{ux434ux438ux430ux433ux43eux43dux430ux43bux438ux437ux430ux446ux438ux44f-ux43cux430ux442ux440ux438ux446ux44b}

Пусть A --- квадратная матрина порядка n и пусть все ее собственные
значения \texttt{\textbackslash{}lambda\_1} ,
\texttt{\textbackslash{}lambda\_2} \ldots{}
\texttt{\textbackslash{}lambda\_n} различны. Составим матрицу P так,
чтобы у нее \emph{столбцами} были координаты собственных векторов (в k-м
столбце стоит собственный вектор \texttt{X\_k} , соответствующий
собственному значению \texttt{\textbackslash{}lambda\_k} ), т. е.

\begin{verbatim}
\Lambda =
\begin{bmatrix} X_1 & X_{\tilde{2}} & \dots & X_n \\ X_{1} & X_{2} & \dots & X_n \end{bmatrix}
\begin{bmatrix} 4 & 0 \\ 0 & -4 \end{bmatrix}
\end{verbatim}

.

\begin{verbatim}
\begin{bmatrix} x_{11} & x_{1\tilde{2}} & \dots & x_{1n} \\ x_{21} & x_{22} & \dots & x_{2n} \\ \vdots & \vdots & \ddots & \vdots \\ x_{n1} & x_{n\tilde{3}} & \dots & x_{nn} \end{bmatrix}
\frac{\partial u_1}{\partial t} + 8 \frac{\partial u_2}{\partial x} = 0
\end{verbatim}

Тогда матрица Л, определяемая соотношением

где \texttt{P\^{}\{-1\}} --- обратная к матрице P, будет диагональной
матрицей вида

\begin{verbatim}
\begin{bmatrix} \frac{\partial u_1}{\partial t} \\ \frac{\partial u_2}{\partial t} \end{bmatrix}
\end{verbatim}

\begin{itemize}
\tightlist
\item
\end{itemize}

\begin{verbatim}
\begin{bmatrix} \lambda_1 & & & \\ & \lambda_2 & 0 & \\ & & \ddots & \\ & & & \lambda_n \end{bmatrix}
\begin{bmatrix} 0 & 8 \\ 2 & 0 \end{bmatrix}
\end{verbatim}

Например, у матрицы

\begin{verbatim}
\begin{bmatrix} \frac{\partial u_1}{\partial x} \\ \frac{\partial u_2}{\partial x} \end{bmatrix}
=
\begin{bmatrix} 0 \\ 0 \end{bmatrix}
in  X_2 =
\end{verbatim}

.

Значит, произведение

\begin{verbatim}
u_t + Au_s = 0
A =
\begin{bmatrix} 0 & 8 \\ 2 & 0 \end{bmatrix}
, \quad u_1 =
\end{verbatim}

равно диагональной матрице

\begin{verbatim}
\begin{bmatrix} \frac{\partial u_1}{\partial t} \\ \frac{\partial u_2}{\partial t} \end{bmatrix}
, \quad u_x =
\end{verbatim}

(Читатель может проверить это самостоятельно.)

Теперь можно решить простую систему двух уравнений с частными
производными (с соответствующими начальными условиями).

\section{\texorpdfstring{Решение линейной системы
\texttt{u\_t\ +\ Au\_x\ =\ 0}
.}{Решение линейной системы u\_t + Au\_x = 0 .}}\label{ux440ux435ux448ux435ux43dux438ux435-ux43bux438ux43dux435ux439ux43dux43eux439-ux441ux438ux441ux442ux435ux43cux44b-u_t-au_x-0-.}

Рассмотрим задачу Коши для системы, содержащей два уравнения и два
начальных условия:

(УЧП 1)

,

\begin{verbatim}
(29.3) (УЧП 2)  \frac{\partial u_2}{\partial t} + 2 \frac{\partial u_1}{\partial x} = 0  ,  -\infty < x < \infty  ,  0 < t < \infty  ,
\end{verbatim}

\begin{verbatim}
(HУ 1)  u_1(x, 0) = f(x)  ,
(HУ 2)  u_2(x, 0) = g(x)  ,  -\infty < x < \infty  .
\end{verbatim}

Эта задача может соответствовать определению давления \texttt{u\_1}(x,
t) и плотности \texttt{u\_2}(x, t) как функции пространственной
координаты x и времени t по известным распределениям этих величин в
начальный момент.

Запишем систему уравнений в матричной форме

\begin{verbatim}
\begin{bmatrix} \frac{\partial u_1}{\partial x} \\ \frac{\partial u_2}{\partial x} \end{bmatrix}
\end{verbatim}

\begin{itemize}
\tightlist
\item
\end{itemize}

\begin{verbatim}
, \quad 0 =
\begin{bmatrix} \frac{\partial u_1}{\partial x} \\ \frac{\partial u_2}{\partial x} \end{bmatrix}
\begin{bmatrix} 0 \\ 0 \end{bmatrix}
\end{verbatim}

.

\begin{verbatim}
v =
\begin{bmatrix} v_1 \\ v_2 \end{bmatrix}
\begin{bmatrix} 0 & 8 \\ 2 & 0 \end{bmatrix}
u = Pv
\begin{bmatrix} \frac{\partial u_1}{\partial t} \\ \frac{\partial u_2}{\partial t} \end{bmatrix}
u = Pv
\frac{\partial u}{\partial t} = P \frac{\partial v}{\partial t},
\begin{bmatrix} \frac{\partial u_1}{\partial x} \\ \frac{\partial u_2}{\partial x} \end{bmatrix}
\frac{\partial u}{\partial x} = P \frac{\partial v}{\partial x}.
u_t + Au_x = 0
\end{verbatim}

.

\begin{verbatim}
Pv_t + APv_x = 0.
v_t + P^{-1}APv_x = 0,
(29.6) v_t + \Lambda v_x = 0.
\end{verbatim}

с помощью преобразования

\begin{verbatim}
\Lambda =
\end{verbatim}

.

де P --- матриця, стовпці якої містять власні вектори матриці A (матриця
P, яку вже знають). Виявляється, що після такого перетворення буде
отримана дуже проста система для визначення v (два рівняння відносно
нового невідомого \texttt{v\_1} і \texttt{v\_2} виявляються
незалежними). Це означає, що \texttt{v\_1} і \texttt{v\_2} легко знайти.
Після \texttt{v\_1} і \texttt{v\_2} знайдено, згідно з формулою

находятся искомые функции и, и и2.

Сначала, однако, выясним, как выглядит система для определения v. Чтобы
разобраться в этом, читатель должен обратиться к матричному исчислению.
Продифференцировав обе части соотношения u = Pv, получаем

(29 5)

\begin{verbatim}
\begin{bmatrix} 4 & 0 \\ 0 & -4 \end{bmatrix}
\end{verbatim}

.

\begin{verbatim}
\frac{\frac{\partial v_1}{\partial t} + 4 \frac{\partial v_1}{\partial x} = 0, \\ \frac{\partial v_2}{\partial t} - 4 \frac{\partial v_2}{\partial x} = 0.
\end{verbatim}

\begin{verbatim}
v_1(x, t) = \varphi(x-4t),
u = Pv =
\begin{bmatrix} 2 & -2 \\ 1 & 1 \end{bmatrix}
\begin{bmatrix} v_1 \\ v_2 \end{bmatrix}
=
=
\end{verbatim}

.

\begin{verbatim}
\begin{bmatrix} 2 & -2 \\ 1 & 1 \end{bmatrix}
\begin{bmatrix} \varphi(x - 4t) \\ \psi(x + 4t) \end{bmatrix}
=
=
\begin{bmatrix} v_1 \\ v_2 \end{bmatrix}
\begin{bmatrix} 2\varphi(x - 4t) - 2\psi(x + 4t) \\ \varphi(x - 4t) + \psi(x + 4t) \end{bmatrix}
\end{verbatim}

.

\begin{verbatim}
u_1(x, t) = 2\varphi(x-4t) - 2\psi(x+4t), u_2(x, t) = \varphi(x-4t) + \psi(x+4t).
\phi(\xi) = \sin \xi, \ \psi(\xi) = \xi^2
u_1(x, t) = 2 \sin(x-4t) - 2(x+4t)^2,
u_1(x, 0) = f(x),
2\varphi(x) - 2\psi(x) = f(x),
\end{verbatim}

.

\begin{verbatim}
\varphi(x) + \psi(x) = g(x),
\varphi(x) = \frac{1}{4} \left[ f(x) + 2g(x) \right],
\psi(x) = \frac{1}{4} \left[ 2g(x) - f(x) \right].
u_{1}(x, t) = 2\varphi(x-4t) - 2\psi(x+4t) =
= \frac{1}{2} \left[ f(x-4t) + 2g(x-4t) \right] - \frac{1}{2} \left[ 2g(x+4t) - f(x+4t) \right],
u_{2}(x, t) = \varphi(x-4t) + \psi(x+4t) =
= \frac{1}{4} \left[ f(x-4t) + 2g(x-4t) \right] + \frac{1}{4} \left[ 2g(x+4t) - f(x+4t) \right].
u_t + A(x, t) u_x + B(x, t) u = 0
Au_t + Bu_x + Cu = 0,
A =
\end{verbatim}

откуда функци і ф и ф легко находятся:

\begin{verbatim}
\begin{bmatrix} 1 & 1 \\ 4 & 1 \end{bmatrix}
\end{verbatim}

.

\begin{verbatim}
\begin{cases} \frac{\partial u_1}{\partial t} + \frac{\partial u_1}{\partial x} + \frac{\partial u_2}{\partial x} = 0, \\ \frac{\partial u_2}{\partial t} + 4 \frac{\partial u_1}{\partial x} + \frac{\partial u_2}{\partial x} = 0. \end{cases}
\end{verbatim}

(29.9)

\begin{verbatim}
u_{tt} = c^2 \left( u_{rr} + \frac{1}{r} u_r + \frac{1}{r^2} u_{\theta\theta} \right)
u(r, \theta, t) = P(r) \Theta(\theta) T(t),
T'' + \lambda^2 c^2 T = 0
\begin{array}{llll} (\text{УЧП}) & u_{tt} = c^2 \left( u_{rr} + \frac{1}{r} u_r + \frac{1}{r^2} u_{\theta\theta} \right), & 0 < r < 1, \ 0 < t < \infty, \\ (\text{ГУ}) & u = 0 & \text{при } r = 1, & 0 < t < \infty, \\ (\text{НУ}) & \begin{cases} u = f(r, \ \theta) \\ u_t = g(r, \ \theta) \end{cases} & \text{при } t = 0. \end{array}
\end{verbatim}

.

\begin{verbatim}
u(r, \theta, t) = U(r, \theta) T(t)
\Delta U + \lambda^2 U = 0
\begin{cases} \frac{\partial u_1}{\partial t} + \frac{\partial u_1}{\partial x} + \frac{\partial u_2}{\partial x} = 0, \\ \frac{\partial u_2}{\partial t} + 4 \frac{\partial u_1}{\partial x} + \frac{\partial u_2}{\partial x} = 0. \end{cases}
T'' + \lambda^2 c^2 T = 0
\Delta U = U_{rr} + \frac{1}{r} U_r + \frac{1}{r^2} U_{\theta\theta}.
u(1, \theta, t) = U(1, \theta) T(t) = 0, \quad 0 < t < \infty,
U(1, \theta) = 0.
\Delta U + \lambda^2 U = 0,
U(1, \theta) = 0.
\begin{cases} \Delta U + \mu u = 0 & \text{в круге,} \\ U = 0 & \text{на границе круга} \end{cases}
\begin{array}{ll} (\mathsf{Y}\mathsf{H}\mathsf{\Pi}) & \Delta U + \lambda^2 U = 0, \\ (\mathsf{\Gamma}\mathsf{Y}) & U(1, \theta) = 0, \end{array}
U(r, \theta) = R(r) \Theta(\theta)
r^2R'' + rR' + (\lambda^2r^2 - n^2)R = 0
r^2R'' + rR' + (\lambda^2r^2 - n^2) R = 0
r^2R'' + rR' + (\lambda^2r^2 - n^2)R = 0
R(r) = AJ_n(\lambda r) + BY_n(\lambda r).
R(r) = AJ_n(\lambda r).
\lambda = k_{nm}
U_{nm}(r, \theta) = J_n(k_{nm}r) [A \sin(n\theta) + B \cos(n\theta)].
f_{nm} = k_{nm}c/2\pi.
T_{nm}(t) = A \sin(k_{nm}ct) + B \cos(k_{nm}ct)
T'' + k_{nm}^2 c^2 T = 0.
sU(x) = -V \frac{dU}{dx}, \quad 0 < x < \infty,
U(0) = P/s.
\begin{array}{ll} (\mathsf{Y}\mathsf{H}\mathsf{\Pi}) & \Delta U + \lambda^2 U = 0, \\ (\mathsf{\Gamma}\mathsf{Y}) & U(1, \theta) = 0, \end{array}
H(\xi) =
\begin{cases} 1, & 0 \le \theta < \pi, \\ 0, & \pi \le \theta < 2\pi. \end{cases}
\begin{cases} 0, & \xi < 0, \\ 1, & \xi \geqslant 0. \end{cases}
\cos(3\varphi) = a_0 P_0(\cos\varphi) + a_1 P_1(\cos\varphi) + \dots
\Delta u = 0, \quad 0 < r < 1, \quad 0 \le \varphi \le \pi, \quad 0 \le \theta < 2\pi.
u(1, \varphi) =
\begin{cases} 1, & 0 \le \varphi \le \pi/2, \\ -1, & \pi/2 < \varphi \le \pi. \end{cases}
\Delta u = 0
u_{rr} + \frac{1}{r} u_r + \frac{1}{r^2} u_{\theta\theta} = f(r, \theta), \quad 0 \le r < 1, \quad 0 \le \theta < 2\pi,
\Delta u = -q
u(1, \theta) = 0, \quad 0 \le \theta \le 2\pi.
= -\int_{0}^{2\pi} u_{r}(r) r d\theta = -2\pi r u_{r}(r).
-2\pi r u_r(r) = q.
u(r) = \frac{-q}{2\pi} \ln r = \frac{q}{2\pi} \ln \frac{1}{r}.
u\left(r\right)=-\frac{q}{2\pi}\ln\frac{1}{r}.
(Y \cup \Pi)
u\left(r,\;\theta\right) = \int_{0}^{2\pi} \int_{0}^{1} G\left(r,\;\theta,\;\rho,\;\varphi\right) f\left(\rho,\;\varphi\right) \rho \; d\rho \; d\varphi.
\frac{1}{2\pi} \ln \left( \frac{1}{r} \right)
\frac{1}{2\pi} \ln \frac{1}{R}
u(r, \theta) = \frac{1}{2\pi} \ln 1/R - \frac{1}{2\pi} \ln 1/\overline{R}
-\frac{1}{2\pi}\ln\rho
G(r, \theta, \rho, \phi) = \frac{1}{2\pi} \ln 1/R - \frac{1}{2\pi} \ln 1/\overline{R} + \frac{1}{2\pi} \ln \rho,
R = V \frac{r^2 - 2\rho r \cos(\theta - \phi) + \rho^2}{R},
\overline{R} = V \frac{r^2 - 2\frac{r}{\rho}\cos(\theta - \phi) + 1/\rho^2}{r^2 - 2\frac{r}{\rho}\cos(\theta - \phi) + 1/\rho^2}
u(r, \theta) = \int_{0}^{2\pi} \int_{0}^{1} G(r, \theta, \rho, \varphi) f(\rho, \varphi) \rho d\rho d\varphi
u(r, \theta) = \frac{1}{2\pi} \int_{0}^{2\pi} \int_{0}^{1} \ln (\rho \overline{R}/R) f(\rho, \varphi) \rho d\rho d\varphi.
\begin{array}{llll} (\text{УЧ\Pi}) & \Delta u = 0, & 0 < r < 1, & 0 \leqslant \theta < 2\pi, \\ (\text{ГУ}) & u(1, \theta) = g(\theta), & 0 \leqslant \theta \leqslant 2\pi, \end{array}
u(r, \theta) = \int_{0}^{2\pi} \frac{\partial G}{\partial r} (r, \theta, 1, \varphi) g(\varphi) d\varphi.
u(r, \theta) = \frac{1}{2\pi} \int_{0}^{2\pi} \left[ \frac{1 - r^2}{1 - 2r\cos(\theta - \varphi) + r^2} \right] g(\varphi) d\varphi.
\begin{array}{llll} (\text{УЧП}) & \Delta u = f(r, \theta), & 0 < r < 1, & 0 \le \theta < 2\pi, \\ (\text{ГУ}) & u(1, \theta) = g(\theta), & 0 < \theta \le 2\pi, \end{array}
\Delta u = 1
u_p = (r, \theta) = Ar^2
f(x+h) = f(x) + f'(x)h + \frac{f''(x)}{2!}h^2 + \dots
f(x+h) \cong f(x) + f'(x) h
(37.1) f'(x) \cong \frac{f(x+h)-f(x)}{h}.
f'(x) \cong \frac{f(x) - f(x - h)}{h}
f'(x) \simeq \frac{1}{2h} [f(x+h) - f(x-h)].
f''(x) \cong \frac{1}{h^2} [f(x+h) - 2f(x) + f(x-h)].
u(x+h, y) = u(x, y) + u_x(x, y)h + u_{xx}(x, y)\frac{h^2}{2!} + \dots,
u(x-h, y) = u(x, y) - u_x(x, y)h + u_{xx}(x, y)\frac{h^2}{2!} - \dots,
(Y \cup \Pi)
\begin{split} u_x(x,\ y) & \cong \frac{u(x+h,\ y)-u(x,\ y)}{h}, \\ u_{xx}(x,\ y) & \cong \frac{1}{h^2} \big[ u(x+h,\ y)-2u(x,\ y)+u(x-h,y) \big], \\ u_u(x,\ y) & \cong \frac{u(x,\ y+k)-u(x,\ y)}{k}, \\ u_{uy}(x,\ y) & \cong \frac{1}{h^2} \big[ u(x,\ y+k)-2u(x,\ y)+u(x,\ y-k) \big]. \end{split}
u_{xx} + u_{yy} = 0
\begin{split} u\left(x,\,y\right) &= u_{ij}\,,\\ u\left(x,\,y+k\right) &= u_{i+1,j}\,,\\ u\left(x,\,y-k\right) &= u_{i-1,j}\,,\\ u\left(x-h,\,y\right) &= u_{i,j-1}\,,\\ u\left(x+h,\,y\right) &= u_{i,j+1}\,,\\ u_{x}\left(x,\,y\right) &= \frac{1}{2h}\left(u_{i,j+1}-u_{i,j-1}\right)\,,\\ u_{y}\left(x,\,y\right) &= \frac{1}{2k}\left(u_{i+1,j}-u_{i-1,j}\right)\,,\\ u_{xx}\left(x,\,y\right) &= \frac{1}{h^{2}}\left(u_{i,j+1}-2u_{i,j}+u_{i,j-1}\right)\,,\\ u_{yy}\left(x,\,y\right) &= \frac{1}{h^{2}}\left(u_{i+1,j}-2u_{i,j}+u_{i-1,j}\right)\,. \end{split}
u_{xx} + u_{yy} = 0
\Delta u = \frac{1}{h^2} (u_{i,j+1} - 2u_{i,j} + u_{i,j-1}) + \frac{1}{k^2} (u_{i+1,j} - 2u_{i,j} + u_{i-1,j}) = 0.
u_{i,j} = \frac{1}{4} (u_{i+1,j} + u_{i-1,j} + u_{i,j+1} + u_{i,j-1}).
-4u_{22} + 0 + \sin(\pi/3) + u_{23} + u_{32} = 0,
-4u_{23} + u_{22} + \sin(2\pi/3) + 0 + u_{33} = 0,
-4u_{32} + 0 + u_{22} + u_{33} + 0 = 0,
\end{verbatim}

формулы, выражающие расстояние между двумя точками в полярных
координатах. Чтобы найти решение исходной задачи, осталось провести
наложение всех импульсных откликов системы. Итак, мы переходим к
последнему шагу.

ШАГ 3. Суперпозиция импульсных откликов Этот шаг совсем прост. Из
соотношения

\begin{verbatim}
-4u_{33} + u_{32} + u_{23} + 0 + 0 = 0.
\end{verbatim}

получаем

(36.4)

\begin{verbatim}
Формула (36.4) виражає розв'язок задачі Діріхле для рівняння Пуассона всередині кола через функцію Гріна цієї задачі. Якщо щільність заряду  f(r, \theta)  відома, то інтеграл можна обчислити, наприклад, чисельно.
\end{verbatim}

\section{ЗАМЕЧАНИЯ}\label{ux437ux430ux43cux435ux447ux430ux43dux438ux44f-2}

\begin{enumerate}
\def\labelenumi{\arabic{enumi}.}
\tightlist
\item
  Задачу
\end{enumerate}

також можна розв'язати за допомогою методу \emph{функції Гріна}. У цьому
випадку рішення фіксується у формі

\begin{verbatim}
\begin{bmatrix} -4 & 1 & 1 \\ 1 & -4 & 0 & 1 \\ 1 & 0 & -4 & 1 \\ 0 & 1 & 1 & -4 \end{bmatrix}
\end{verbatim}

Останнє відношення можна надати більш обчислювально корисну форму, явно
виставляючи вираз у
\texttt{\textbackslash{}partial\ G/\textbackslash{}partial\ r} .
Внаслідок цього ми отримуємо інтегральну формулу Пуассона, знайдену
раніше в лекції 33.

(36.5)

\begin{enumerate}
\def\labelenumi{\arabic{enumi}.}
\setcounter{enumi}{1}
\tightlist
\item
  Решение общей задачи Дирихле
\end{enumerate}

\begin{verbatim}
\begin{bmatrix} u_{22} \\ u_{23} \\ u_{32} \\ u_{33} \end{bmatrix}
\end{verbatim}

находится как сумма решений (36.4) и (36.5).

\begin{enumerate}
\def\labelenumi{\arabic{enumi}.}
\setcounter{enumi}{2}
\item
  Метод функцій Гріна дозволяє отримати розв'язки багатьох задач у
  областях різних форм. Однак для кожної області і (точніше, для кожного
  оператора зліва від граничної умови) для кожного рівняння потрібно
  знайти власну функцію Гріна, і це не завжди легко зробити.
\item
  Чтобы реально найти решение по формуле (36.4), в большинстве случаев
  необходимо вычислять интегралы на ЭВМ.
\end{enumerate}

\section{ЗАДАЧИ}\label{ux437ux430ux434ux430ux447ux438-5}

Найти потенциал точечного источника в трехмерном пространстве.

\begin{verbatim}
2. Найти функцию Грина  G(x, y, \xi, \eta)  для уравнения Лапласа в верхней полуплоскости y > 0. Другими словами, найти потенциал в точке (x, y) в верхней полуплоскости, если заряд расположен в точке  (\xi, \eta)  , а потенциал на границе y = 0 равен нулю (см. рис.).
\end{verbatim}

УКАЗАНИЕ. Если отрицательный заряд поместить в точку
\texttt{\textbackslash{}overline\{Q\}\ =\ (\textbackslash{}xi,\ -\textbackslash{}eta)}
, очевидно, що загальний потенціал на лінії y = 0 буде нульовим. Отже,
функція Гріна збігається з потенціалом цих двох зарядів.

\texttt{-\ 3.\ Використовуючи\ результати\ задачі\ 2,\ знайдемо\ розв\textquotesingle{}язок\ рівняння\ Пуассона\ \textquotesingle{}\textbackslash{}Delta\ u\ =\ -k\textquotesingle{}\ у\ верхній\ півплощині\ з\ умовою\ нульової\ межі.}
- 4. Як би ви побудували функцію Зеленого для першого квадранта
\texttt{x\textgreater{}0,\textbackslash{}\ y\textgreater{}0} ? - 5.
Інший підхід до розв'язання рівняння Пуассона такий. Припустимо, ви
хочете розв'язати проблему

(УЧП)

\begin{verbatim}
=
\end{verbatim}

\texttt{,\ \textquotesingle{}0\ \textless{}\ r\ \textless{}\ 1\textquotesingle{}\ ,\ \textquotesingle{}0\ \textbackslash{}le\ \textbackslash{}theta\ \textless{}\ 2\textbackslash{}pi\textquotesingle{}\ ,\ \textquotesingle{}(\textbackslash{}Gamma\ Y)\textquotesingle{}\ \textquotesingle{}u(1,\ \textbackslash{}theta)\ =\ \textbackslash{}sin\ \textbackslash{}theta\textquotesingle{}\ ,\ \textquotesingle{}0\ \textbackslash{}le\ \textbackslash{}theta\ \textbackslash{}le\ 2\textbackslash{}pi\textquotesingle{}\ .}

Сначала попытайтесь найти частное решение уравнения с помощью
подстановки

.

\begin{verbatim}
Затем подставьте это выражение в уравнение Пуассона и определите константу  A_{\cdot}  Після того, як  u_p  знайдено часткове рішення, необхідно шукати повне рішення у вигляді  u=w+u_p  . Яку граничну умову має задовольняти функція w? Знайдіть  w(r, \theta)  . Яке рішення початкової проблеми  u(r, \theta)  ? Подивись. Ретельно проаналізуйте відповідь. Дайте інтерпретацію кожного члена.
\end{verbatim}

\section{Часть 5. ЧИСЛЕННЫЕ И ПРИБЛИЖЕННЫЕ
МЕТОДЫ}\label{ux447ux430ux441ux442ux44c-5.-ux447ux438ux441ux43bux435ux43dux43dux44bux435-ux438-ux43fux440ux438ux431ux43bux438ux436ux435ux43dux43dux44bux435-ux43cux435ux442ux43eux434ux44b}

Лекция 37

\section{ЧИСЛЕННЫЕ РЕШЕНИЯ (ЭЛЛИПТИЧЕСКИЕ
ЗАДАЧИ)}\label{ux447ux438ux441ux43bux435ux43dux43dux44bux435-ux440ux435ux448ux435ux43dux438ux44f-ux44dux43bux43bux438ux43fux442ux438ux447ux435ux441ux43aux438ux435-ux437ux430ux434ux430ux447ux438}

МЕТА ЛЕКЦІЇ: Показати, як диференціальні рівняння в похідних похідних
можна звести до систем алгебраїчних рівнянь, якщо похідні з частковими
похідними замінити на наближення з скінченними різницями. Розв'язання
цієї системи алгебраїчних рівнянь ітеративним методом дозволяє отримати
наближене розв'язок

Дается информация о пакете прикладных программ ELLPACK, предназначенном
для решения эллиптических

задач общего вида на ЭВМ.

К настоящему моменту мы познакомились с несколькими методами решения
линейных уравнений с частными производными. Однако большинство
уравнений, с которыми мы имели дело, были очень простыми. Так же просты
были и формы областей, в которых решались задачи. Тем не менее многие
задачи нельзя упростить до такой степени, чтобы они свелись к заданному
набору шаблонных задач. Такие задачи следует решать приближенно
численными методами. Прогресс в области высокопроизводительных ЭВМ
привел к созданию новых численных методов. Сегодня уже решены такие
нелинейные задачи динамики жидкости, теории упругости и теории
потенциала, о решении которых десять лет назад никто даже не
задумывался.

Под общим названием «численные методы» объединены несколько различных
подходов к решению задач. Подробное описание всех этих подходов можно
найти в книге {[}4{]} из рекомендованного списка литературы. В этой и
двух последующих лекциях будет показано, как решать эллиптические,
гиперболические и параболические уравнения методом конечных разностей.

Сначала познакомимся с понятием конечных разностей, а затем покажем, как
с их помощью решить задачу Дирихле в квадрате.

\section{Конечно-разностные
аппроксимации}\label{ux43aux43eux43dux435ux447ux43dux43e-ux440ux430ux437ux43dux43eux441ux442ux43dux44bux435-ux430ux43fux43fux440ux43eux43aux441ux438ux43cux430ux446ux438ux438}

Вспомним ряд Тейлора для функции \texttt{f(x)}:

\begin{verbatim}
\begin{bmatrix} -\sin{(\pi/3)} \\ -\sin{(2\pi/3)} \\ 0 \\ 0 \end{bmatrix}
\end{verbatim}

Если оборвать этот ряд на втором члене, то получим

.

Откуда

\begin{verbatim}
=
\end{verbatim}

Выражение, стоящее в правой части, называется правой разностной
производной. Она аппроксимирует первую производную f'(x) в точке x.

В разложении Тейлора функции \texttt{f(\textbackslash{}mathbf\{r\})}
можно заменить h на -h и получить левую разностную производную

(37.2)

.

\begin{verbatim}
Buyutan  f(x - h) \cong f(x) - f'(x) h  \nif  f(x + h) \cong f( x) + f'(x) h  ,
\end{verbatim}

получаем центральную разностную производную

(37.3)

\begin{verbatim}
\begin{bmatrix} -0.86 \\ -0.86 \\ 0 \\ 0 \end{bmatrix}
\end{verbatim}

Якщо в ряду Тейлора залишиться ще один член, то центральну різницю
похідну для наближення f'\,'(x) можна отримати точно таким самим
способом

(37.4)

.

Тепер ми можемо розширити поняття наближення скінченних різниць до
\emph{часткових похідних}. Якщо виходити з розкладу Тейлора функції двох
змінних

\begin{verbatim}
u_{xx} + u_{yy} = f(x, y)
xu_{xx} + u_{yy} + 2u = \sin(x - y)
\sin xu_{xx} + u_{xy} + 3u = 0
\end{verbatim}

\pandocbounded{\includegraphics[keepaspectratio]{images/_page_285_Figure_8.jpeg}}

\begin{verbatim}
\begin{split} u_x(x,\ y) & \cong \frac{u(x+h,\ y)-u(x,\ y)}{h}, \\ u_{xx}(x,\ y) & \cong \frac{1}{h^2} \big[ u(x+h,\ y)-2u(x,\ y)+u(x-h,y) \big], \\ u_u(x,\ y) & \cong \frac{u(x,\ y+k)-u(x,\ y)}{k}, \\ u_{uy}(x,\ y) & \cong \frac{1}{h^2} \big[ u(x,\ y+k)-2u(x,\ y)+u(x,\ y-k) \big]. \end{split}
\end{verbatim}

\texttt{(УЧ\textbackslash{}Pi)}

\begin{verbatim}
u_{xx} + u_{yy} + 2u = 0
u_{xx} + u_{yy} = 0
\begin{cases} u = 0 & \text{квадрата,} \\ \frac{\partial u}{\partial x}(1, y) = 1, & 0 \leqslant y \leqslant 1, \end{cases}
\begin{split} u\left(x,\,y\right) &= u_{ij}\,,\\ u\left(x,\,y+k\right) &= u_{i+1,j}\,,\\ u\left(x,\,y-k\right) &= u_{i-1,j}\,,\\ u\left(x-h,\,y\right) &= u_{i,j-1}\,,\\ u\left(x+h,\,y\right) &= u_{i,j+1}\,,\\ u_{x}\left(x,\,y\right) &= \frac{1}{2h}\left(u_{i,j+1}-u_{i,j-1}\right)\,,\\ u_{y}\left(x,\,y\right) &= \frac{1}{2k}\left(u_{i+1,j}-u_{i-1,j}\right)\,,\\ u_{xx}\left(x,\,y\right) &= \frac{1}{h^{2}}\left(u_{i,j+1}-2u_{i,j}+u_{i,j-1}\right)\,,\\ u_{yy}\left(x,\,y\right) &= \frac{1}{h^{2}}\left(u_{i+1,j}-2u_{i,j}+u_{i-1,j}\right)\,. \end{split}
(\text{УЧП})
u_t = u_{xx}
\end{verbatim}

\pandocbounded{\includegraphics[keepaspectratio]{images/_page_288_Figure_3.jpeg}}

\begin{verbatim}
u_t = u_{xx}, \quad 0 < x < 1, \quad 0 < t < \infty,
\begin{cases} u(0, t) = 1, \\ u_x(1, t) = -[u(1, t) - g(t)], \\ u(x, 0) = 0, \quad 0 \le x \le 1. \end{cases}
x_j = jh
u_{t} = \frac{1}{k} [u(x, t+k) - u(x, t)] = \frac{1}{k} [u_{t+1,j} - u_{t,j}],
u_{xx} = \frac{1}{h^{2}} [u(x+h, t) - 2u(x, t) + u(x-h, t)] = \frac{1}{h^{2}} [u_{t,j+1} - 2_{t,j} + u_{t,j-1}].
\end{verbatim}

Це потрібно вирішувати у зв'язку з \texttt{u\_\{22\}} ,
\texttt{u\_\{21\}} , \texttt{u\_\{32\}} та \texttt{u\_\{33\}} .
Розв'язок такої системи можна знайти за допомогою \emph{ітеративного
методу}. Метод Лібмана також є одним із ітеративних методів.

\begin{itemize}
\tightlist
\item
  \begin{enumerate}
  \def\labelenumi{\arabic{enumi}.}
  \setcounter{enumi}{1}
  \tightlist
  \item
    Якщо зменшити кроки сітки на h і k (кількість вузлів сітки
    збільшиться), то отримаємо систему рівнянь типу (37,7), але більших
    розмірів.
  \end{enumerate}
\item
  \begin{enumerate}
  \def\labelenumi{\arabic{enumi}.}
  \setcounter{enumi}{2}
  \tightlist
  \item
    Систему рівнянь (37.7) можна переписати у матричній формі
  \end{enumerate}
\end{itemize}

\begin{verbatim}
u_{t+1,j} = u_{t,j} + \frac{k}{h^2} \left[ u_{t,j+1} - 2u_{t,j} + u_{t,j-1} \right].
a_x(1, t) = -[u(1, t) - g(t)].
\begin{bmatrix} u_{22} \\ u_{23} \\ u_{32} \\ u_{33} \end{bmatrix}
\frac{1}{h}[u_{t,n} - u_{t,n-1}] = -[u_{t,n} - g_t].
u_{t,n} = \frac{u_{t,n-1} - hg_t}{1 + h}.
u_{2,j} = u_{1,j} + \frac{k}{h^2} [u_{1,j+1} - 2u_{1,j} + u_{1,j-1}], \quad j = 2, 3, \ldots, n-1.
u_{2,n} = \frac{u_{2,n-1} + hg_2}{1+h}.
u_{tt} = u_{xx}, \quad 0 < x < 1, \quad 0 < t < \infty,
\end{verbatim}

С ростом числа уравнений (например, до 1000) матрица коэффициентов
становится разреженной, т. е. содержит много нулей. Системы с
разреженными матрицами решаются специальными методами. Обычно
используются итерационные методы, такие, как метод Якоби, метод
Гаусса---Зейделя или метод последовательной верхней релаксации.

\begin{enumerate}
\def\labelenumi{\arabic{enumi}.}
\setcounter{enumi}{3}
\item
  При розв'язанні задачі Неймана ті \emph{похідні}, що входять до
  граничної умови, також слід замінити різницею апроксимаціями.
\item
  Подобным образом можно решать задачи
\end{enumerate}

\begin{enumerate}
\def\labelenumi{\alph{enumi})}
\tightlist
\item
\end{enumerate}

\begin{enumerate}
\def\labelenumi{\arabic{enumi})}
\setcounter{enumi}{5}
\tightlist
\item
\end{enumerate}

\begin{verbatim}
\begin{cases} u(0, t) = g_1(t), \\ u(1, t) = g_2(t), \end{cases}
\end{verbatim}

\begin{enumerate}
\def\labelenumi{\Alph{enumi})}
\setcounter{enumi}{1}
\tightlist
\item
\end{enumerate}

(неоднородные уравнения с переменными коэффициентами), (уравнения с
переменными коэффициентами).

\begin{enumerate}
\def\labelenumi{\arabic{enumi}.}
\setcounter{enumi}{5}
\tightlist
\item
  Якщо область, в якій розв'язується задача, має неправильну форму, її
  можна покрити сіткою і наближити розв'язок у точках, найближчих до
  межі, інтерполюючи граничні умови. Після цього задачу вирішують
  звичайним способом (див. рисунок 37.2).
\end{enumerate}

\begin{verbatim}
\quad 0 < t < \infty,
\end{verbatim}

□ новые ГУ находятся интерполяцией ГУ на кривой

РИС, 37 2.

\begin{itemize}
\item
  \begin{enumerate}
  \def\labelenumi{\arabic{enumi}.}
  \setcounter{enumi}{6}
  \tightlist
  \item
    В некоторых журналах публикуются тексты программ для решения
    уравнений с частными производными на ЭВМ. Ниже приведен краткий
    перечень этих журналов:
  \end{enumerate}

  \begin{itemize}
  \tightlist
  \item
    \begin{enumerate}
    \def\labelenumi{\alph{enumi})}
    \tightlist
    \item
      ACM Transactions on Mathematical Software,
    \end{enumerate}
  \item
    \begin{enumerate}
    \def\labelenumi{\arabic{enumi})}
    \setcounter{enumi}{5}
    \tightlist
    \item
      Computer Journal,
    \end{enumerate}
  \item
    \begin{enumerate}
    \def\labelenumi{\Alph{enumi})}
    \setcounter{enumi}{1}
    \tightlist
    \item
      Numerische Mathematik,
    \end{enumerate}
  \item
    \begin{enumerate}
    \def\labelenumi{\alph{enumi})}
    \setcounter{enumi}{17}
    \tightlist
    \item
      BIT
    \end{enumerate}
  \end{itemize}
\end{itemize}

Кроме того, недавно был разработан пакет прикладных программ ELLPACK,
предназначенный для решения достаточно широкого класса краевых
эллиптических задач. Этот пакет позволит решать разнообразные двух- и
трехмерные задачи в

различных системах координат для произвольных границ и общих граничных
условий. Пользователю представляется большой ассортимент различных
методов решения 1).

\section{ЗАДАЧИ}\label{ux437ux430ux434ux430ux447ux438-6}

\begin{itemize}
\tightlist
\item
  1 Получите аппроксимацию (39.4) для второй производной f'\,'(x)
  \texttt{f\textquotesingle{}\textquotesingle{}(x)\ =\ \textbackslash{}\{f(x+h)\ 2f(x)\ +\ f(x-h)\textbackslash{}\}/h\^{}2}
  .
\item
  \begin{enumerate}
  \def\labelenumi{\arabic{enumi}.}
  \setcounter{enumi}{1}
  \tightlist
  \item
    Проделайте две итерации для задачи Дирихле (37.5), используя
    итерационный метод Либманна.
  \end{enumerate}
\item
  \begin{enumerate}
  \def\labelenumi{\arabic{enumi}.}
  \setcounter{enumi}{2}
  \tightlist
  \item
    К какой алгебранческой системе сведется задача Дирихле для уравнений
    Пуассона в квадрате
  \end{enumerate}
\end{itemize}

\begin{verbatim}
u_{xx} + u_{yy} = f(x,y), \quad 0 < x < 1, \quad 0 < y < 1,   (\Gamma Y)   u(x,y) = g(x,y)  на границе, если ее решать ме-
\end{verbatim}

тодом конечных разностей? 4. Решите задачу 3, если

(УЧП)

, \texttt{0\ \textless{}\ x\ \textless{}\ 1} ,
\texttt{0\ \textless{}\ y\ \textless{}\ 1} , (ГУ)
\texttt{u(x,\ y)\ =\ g(x,\ y)} на границе.

\begin{enumerate}
\def\labelenumi{\arabic{enumi}.}
\setcounter{enumi}{4}
\tightlist
\item
  Как бы вы решали задачу Неймана внутри квадрата
\end{enumerate}

(УЧП)

\begin{verbatim}
\begin{cases} u(x, 0) = \varphi(x), \\ u_t(x, 0) = \psi(x), \end{cases}
\end{verbatim}

, \texttt{0\ \textless{}\ x\ \textless{}\ 1} ,
\texttt{0\ \textless{}\ y\ \textless{}\ 1} , на верхней, нижней и левой
стороне квадрата,

методом конечных разностей?

\begin{enumerate}
\def\labelenumi{\arabic{enumi}.}
\setcounter{enumi}{5}
\tightlist
\item
  Постройте блок-схему алеоритма решения задачи Дирихле в квадрате
\end{enumerate}

\begin{verbatim}
\quad 0 \leqslant x \leqslant 1,
\end{verbatim}

\begin{verbatim}
u_{xx} + u_{yy} = f(x, y), \quad 0 < x < 1, \quad 0 < y < 1,   (\text{ГУ})   u(x, y) = g(x, y)  на границе,
\end{verbatim}

\begin{itemize}
\item
  з довільною кількістю вузлів сітки. Якщо ви знаєте мову програмування,
  напишіть програму для виконання цих обчислень.
\item
  \begin{enumerate}
  \def\labelenumi{\arabic{enumi})}
  \tightlist
  \item
    Російською мовою публікуються статті з чисельних методів розв'язання
    диференціальних рівнянь у численних журналах і збірках. Наприклад:
    «Журнал обчислювальної математики та математичної фізики»,
    видавництво «Наука», «Журнал інженерної фізики», видавництво
    Академії наук БРСР, збірки серії «Обчислювальні методи і
    програмування», Видавництво Московського державного університету.
    Багато публікацій, переважно статей, присвячені наборам прикладних
    проблем. Як приклад, можна вказати: Горбунов-Посадов М. М., Карпов
    В. Я, Корягін та ін. Пакет SAFRA. Програмне забезпечення для
    обчислювальних експериментів. У: Пакети прикладних програм.
    Обчислювальний експеримент. Москва, Наука, 1983. Примітка. Переклад.
  \end{enumerate}
\end{itemize}

\section{ЯВНЫЕ РАЗНОСТНЫЕ
СХЕМЫ}\label{ux44fux432ux43dux44bux435-ux440ux430ux437ux43dux43eux441ux442ux43dux44bux435-ux441ux445ux435ux43cux44b}

МЕТА ЛЕКЦІЇ: Ввести концепцію явної різницької схеми та показати, як її
можна використовувати для розв'язання гіперболічних і параболічних
задач. Основна ідея полягає в тому, що після заміни рівняння типу

\begin{verbatim}
u_{tt} \simeq \frac{1}{h^2} [u(x, t+k) - 2u(x, t) + u(x, t-k)],
\end{verbatim}

его конечно-разностной аппроксимацией получаются формулы, явно
выражающие значения решения для одного момента времени через значения
решения в предыдущий момент времени. Таким образом, смешанную задачу для
параболического или гиперболического уравнения можно решать,
последовательно вычисляя решения для всех последующих моментов времени.

Явна схема не позбавлена недоліків. Якщо необхідно підвищити точність
наближення похідних, то при зміні сегмента збільшується не лише обсяг
обчислень, а й похибки округлення.

Попередня лекція була присвячена розв'язанню еліптичних задач на межі
лінії (стаціонарних задач). У таких задачах необхідно знайти розв'язок
диференціального рівняння в похідних у заданій області простору, якщо
розв'язок або його похідна задано на межі області визначення Приблизне
розв'язання еліптичних задач зводиться до розв'язку системи алгебраїчних
рівнянь для значень функції у внутрішніх вузлах сітки. Іншими словами,
значення розв'язку у всіх внутрішніх вузлах визначаються одночасно.

\begin{verbatim}
У цій лекції ми розглянемо диференціальні методи розв'язання задач, залежних від часу. Основна ідея полягає в тому, що якщо розв'язок відомий у початковий момент часу, то можливо знайти розв'язок за допомогою схеми підрахунку тика, коли  t=\Delta t,\ 2\Delta t,\ 3\Delta t,\ldots  Замінюючи часткові похідні часу та просторову змінну на скінченні різниці, можна отримати явні вирази для  u_{tj}  за допомогою значень функції u у попередні часи. Такий процес називається явною схемою підрахунку з тикуванням.
\end{verbatim}

Чтобы показать этот метод в действии, рассмотрим типичную задачу
теплопроводности.

\section{Явная схема для уравнения
теплопроводности}\label{ux44fux432ux43dux430ux44f-ux441ux445ux435ux43cux430-ux434ux43bux44f-ux443ux440ux430ux432ux43dux435ux43dux438ux44f-ux442ux435ux43fux43bux43eux43fux440ux43eux432ux43eux434ux43dux43eux441ux442ux438}

Рассмотрим задачу теплопроводности в стержне, начальная температура
которого равна нулю. Пусть температура левого конца фиксирована, а на
правом конце происходит теплообмен с окружающей средой, так что тепловой
полок пропорционален Р И

\begin{verbatim}
u_{xx} \simeq \frac{1}{h^2} [u(x+h, t) - 2u(x, t) + u(x-h, t)].
\end{verbatim}

С. 38.1. Сетка для уравнения теплопроводности.

\begin{verbatim}
разности температур конца стержия и среды. Пусть температура среды определяется функцией  g\left(t\right)  . Другими словами, мы решаем задачу
\end{verbatim}

задачу (УЧП)

\begin{verbatim}
u_t(x, 0) \approx \frac{1}{k} [u(x, k) - u(x, 0)] = \frac{1}{k} [u(x, k) - \varphi(x)].
\end{verbatim}

(38.1) (ГУ)

\begin{verbatim}
u(x, t+k) = 2u(x, t) - u(x, t-k) + + (k/h)^2 [u(x+h, t) - 2u(x, t) + u(x-h, t)].
\end{verbatim}

Для решения этой задачи методом конечных разностей построим
прямоугольную сетку, узлы которой определяются формулами (см. рис. 38.1)

\begin{verbatim}
\frac{1}{k}\left[u\left(x,\,k\right)-\varphi\left(x\right)\right]=\psi\left(x\right),
,  j = 0, 1, 2, ..., n  ,  y_i = ik  ,  i = 0, 1, 2, ..., m  .
\end{verbatim}

Отметим, что значения \texttt{u\_\{ij\}} на левой и нижней сторонах
сетки на рис. 38.1 известны из граничных и начальных условий и наша
вадача состоит в отыскании остальных значений \texttt{u\_\{ij\}} . Для
решения этой задачи заменим частные производные в уравнении
теплопроводности их конечно-разностными аппроксимациями

\begin{verbatim}
u_t = u_{xx}
\begin{cases} u(0, t) = 0, \\ u(1, t) = 0, \\ u(x, 0) = \sin(\pi x), \quad 0 \le x \le 1. \end{cases}
\end{verbatim}

\pandocbounded{\includegraphics[keepaspectratio]{images/_page_293_Figure_6.jpeg}}

\begin{verbatim}
\begin{array}{ccccccccccccccccccccccccccccccccccc
\end{verbatim}

\begin{verbatim}
\frac{\Delta t}{(\Delta x)^2} \leqslant 0.5.
\end{verbatim}

\begin{verbatim}
\begin{array}{ccccccccccccccccccccccccccccccccccc
\end{verbatim}

\begin{verbatim}
\begin{split} u_t(x,\ t) &= \frac{1}{k} \big[ u(x,\ t+k) - u(x,\ t) \big], \\ u_{xx}(x,\ t) &= \frac{\lambda}{h^2} \big[ u(x+h,\ t+k) - 2u(x,\ t+k) + u(x-h,\ t+k) \big] + \\ &\quad + \frac{(1-\lambda)}{h^2} \big[ u(x+h,\ t) - 2u(x,\ t) + u(x-h,\ t) \big], \end{split}
\frac{1}{k} (u_{i+1, j} - u_{i, j}) = \frac{\lambda}{h^2} (u_{i+1, j+1} - 2u_{i+1, j} + u_{i+i, j-1}) + \frac{(1-\lambda)}{h^2} (u_{i, j+1} - 2u_{i, j} + u_{i, j-1}),
\end{verbatim}

(39.3)

\begin{verbatim}
(\Gamma Y) \begin{cases} u_{i, 1} = 0, & i = 1, 2, ..., m, \\ u_{i, n} = 0, & i = 1, 2, ..., m, \\ u_{i, j} = 1, & j = 2, ..., n-1. \end{cases}
-\lambda r u_{i+1, j+1} + (1+2r\lambda) u_{i+1, j} - \lambda r u_{i+1, j-1} = = r(1-\lambda) u_{i, j+1} + [1-2r(1-\lambda)] u_{i, j} + r(1-\lambda) u_{i, j-1},
\end{verbatim}

\pandocbounded{\includegraphics[keepaspectratio]{images/_page_296_Figure_5.jpeg}}

\pandocbounded{\includegraphics[keepaspectratio]{images/_page_296_Figure_7.jpeg}}

\begin{verbatim}
\begin{array}{l} -u_{21}+3u_{22}-u_{23}=u_{11}-u_{12}+u_{13}=1,\\ -u_{22}+3u_{23}-u_{24}=u_{12}-u_{13}+u_{14}=1,\\ -u_{23}+3u_{24}-u_{25}=u_{13}-u_{14}+u_{15}=1,\\ -u_{24}+3u_{25}-u_{26}=u_{14}-u_{15}+u_{16}=1. \end{array}
\begin{cases} u(0, t) = g_1(t), \\ u(1, t) = g_2(t), \end{cases}
\end{verbatim}

(HY)

\begin{verbatim}
\begin{bmatrix} 3 & -1 & 0 & 0 \\ -1 & 3 & -1 & 0 \\ 0 & -1 & 3 & -1 \\ 0 & 0 & -1 & 3 \end{bmatrix}
\quad 0 \leqslant x \leqslant 1,
\begin{bmatrix} u_{2\bar{z}} \\ u_{\bar{z}a} \\ u_{24} \\ u_{\bar{z}x} \end{bmatrix}
\end{verbatim}

а начальное условие--- по схеме

\begin{verbatim}
=
\end{verbatim}

В результате для вычисления величины \texttt{u(x,\ t+k)} получаем
следующую явную схему:

(38.5)

З (38.5) видно, що для обчислення розв'язку на заданому часовому шарі
необхідно знати розв'язок на двох попередніх шарах. Отже, щоб почати
рахувати, потрібно використовувати P И С

\begin{verbatim}
\begin{bmatrix} 1 \\ 1 \\ 1 \\ 1 \end{bmatrix}
\end{verbatim}

. 38.4. Использование массивов при реализации явной схемы.

начальным условнем на скорость

.

\begin{verbatim}
з якої ми виводимо  u(x, k) = \varphi(x) + k\psi(x)  , тобто цінність розв'язку на  t = \Delta t  . Розв'язок у наступні часи можна знайти за явною формулою (38.5).
\end{verbatim}

\section{ЗАДАЧИ}\label{ux437ux430ux434ux430ux447ux438-7}

\begin{enumerate}
\def\labelenumi{\arabic{enumi}.}
\tightlist
\item
  Постройте явную конечно-разностную схему для задачи
\end{enumerate}

(УЧП)

\texttt{,\ \textquotesingle{}0\ \textless{}\ x\ \textless{}\ 1\textquotesingle{}\ ,\ \textquotesingle{}0\ \textless{}\ t\ \textless{}\ \textbackslash{}infty\textquotesingle{}\ ,}
(ГУ)

\begin{verbatim}
\begin{bmatrix} b_1 & c_1 & 0 & 0 & \dots & & & 0 \\ a_1 & b_2 & c_2 & 0 & \dots & & & 0 \\ 0 & a_{\bar{z}} & b_3 & c_3 & \dots & & & 0 \\ & & & & & & & & \\ \vdots & & & & & & & & \\ 0 & 0 & 0 & \dots & & & & a_{n-1} & b_n \end{bmatrix}
\end{verbatim}

\texttt{(HУ)\ \textquotesingle{}u(x,\ 0)\ =\ \textbackslash{}sin(\textbackslash{}pi\ x),\ \textbackslash{}quad\ 0\ \textbackslash{}le\ x\ \textbackslash{}le\ 1.\textquotesingle{}\ Полагая\ \textquotesingle{}h\ =\ \textbackslash{}Delta\ x\ =\ 0.1\textquotesingle{}\ ,\ найдите\ решение\ при\ \textquotesingle{}t\_1\ =\ 0.005\textquotesingle{}\ ,\ \textquotesingle{}t\_2\ =\ 0.010\textquotesingle{}\ ,\ \textquotesingle{}t\_3\ =\ 0.015\textquotesingle{}\ .\ Постройте\ график\ полученного\ решения\ на\ сетке\ x\ =\ 0;\ 0.1;\ 0.2;\ 0.3;\ .\ .;\ 0.9:\ 1\ при\ t\ =\ 0.015.}

РИС. 38.5. Зависимость полной ошибки от шага сетки.

\begin{itemize}
\tightlist
\item
  \begin{enumerate}
  \def\labelenumi{\arabic{enumi}.}
  \setcounter{enumi}{1}
  \tightlist
  \item
    Решите задачу 1 аналитически (методом разделения переменных).
    Найдите значение решения в узлах сетки x=0; 0,1; 0,2; 0,3; \ldots;
    0,9; 1 при t=0,015. Сравните эти результаты с результатами
    численного решения задачи 1. (Для табулирования аналитического
    решения можно воспользоваться калькулятором или написать маленькую
    программу для ЭВМ.)
  \end{enumerate}
\item
  \begin{enumerate}
  \def\labelenumi{\arabic{enumi}.}
  \setcounter{enumi}{2}
  \tightlist
  \item
    В соответствии с замечанием 3 постройте блок-схему решения
    гиперболической задачи.
  \end{enumerate}
\item
  \begin{enumerate}
  \def\labelenumi{\arabic{enumi}.}
  \setcounter{enumi}{3}
  \tightlist
  \item
    Решите задачу 1 с новым граничным условием при x=1 \texttt{u\_x}(1,
    t) = -{[}\texttt{u(1,\ t)} 1{]}. \# НЕЯВНЫЕ РАЗНОСТНЫЕ СХЕМЫ (СХЕМА
    КРАНКА---НИКОЛЬСОНА)
  \end{enumerate}
\end{itemize}

ЦЕЛЬ ЛЕКЦИИ: Показать, как можно решать задачи, содержащие функции,
зависящие от времени, с помощью неявной разностной схемы. В этой схеме,
как и рачьше, частные производные заменяются конечно-разностными
аппроксимациями, но \texttt{u\_\{l+1,\textbackslash{}\ J\}} уже не
выражаются в явном виде через значения на предыдущих слоях. Теперь для
определения величин \texttt{u\_\{l+1,\textbackslash{}\ J\}} необходимо
решить систему уравнений. Другими словами, для определения решения на
каждом временном наге необходимо решать систему уравнений.

Перевага неявних схем над явними полягає в тому, що в неявних схемах
танг сітки можна зробити достатньо великим без страху, що помилки
округлення «знищать» розв'язок.

Для решения параболической задачи мы воспользуемся широко известной
\emph{схемой Кранка} --- \emph{Никольсона}.

Как уже говорилось в предыдущей лекции, для надежной работы явной схемы
необходимо, чтобы временной шаг сетки был мал. В частности, если простая
задача для уравнения теплопроводности

(39.1)

\begin{verbatim}
\begin{bmatrix} x_1 \\ x_{\bar{z}} \\ x_{\bar{s}} \\ \vdots \\ x_n \end{bmatrix}
\end{verbatim}

\begin{verbatim}
решается по явной схеме, то шаги сетки  \Delta t  и  \Delta x  должны удовлетворять условию
\end{verbatim}

\texttt{Інакше\ метод\ буде\ чисельно\ нестабільним\ (а\ помилки\ округлення\ зростуть\ нескінченно).\ (Питання\ чисельної\ стабільності\ розглядаються\ у\ {[}53{]}.)\ Іншими\ словами,\ якщо\ відстань\ сітки\ на\ осі\ x\ вибрана\ \textquotesingle{}\textbackslash{}Delta\ x=0,1\textquotesingle{}\ ,\ то\ часовий\ крок\ \textquotesingle{}\textbackslash{}Delta\ t\textquotesingle{}\ не\ може\ бути\ більшим\ за\ \textquotesingle{}0.5\textbackslash{}Delta\ x\^{}2=0.005\textquotesingle{}\ .\ Це\ означає,\ що\ ви\ можете\ пройти\ від\ t=0\ до\ t=1,\ зробивши\ 200\ кроків.}

Существуют, однако, такие схемы (неявные схемы), которые позволяют вести
вычисления с гораздо большим шагом, при этом,

правда, растет объем вычислений, совершаемых на одном шаге. По таким
схемам можно вести вычисления с достаточно большим шагом, но на каждом
шаге приходится решать систему алгебраимеских уравнений. Для иллюстрации
этого метода решим следующую задачу теплопроводности.

\section{Неявная схема для уравнения
теплопроводности}\label{ux43dux435ux44fux432ux43dux430ux44f-ux441ux445ux435ux43cux430-ux434ux43bux44f-ux443ux440ux430ux432ux43dux435ux43dux438ux44f-ux442ux435ux43fux43bux43eux43fux440ux43eux432ux43eux434ux43dux43eux441ux442ux438}

Рассмотрим задачу

(39 2)

\begin{verbatim}
=
\end{verbatim}

Воспользуемся следующими конечно-разностными аппроксимациями для частных
производных \texttt{u\_t} и \texttt{u\_\{xx\}} :

\begin{verbatim}
где  \lambda  выбирается из отрезка [0,1]. Отметим, что  u_{xx}  аппроксимируется взвешенным средним центральных разностных производных в момент t и t+k. При  \lambda=0.5  получается обычное среднее этих двух центральных производных, а при  \lambda=0.75  —одна из разностных производных берется с весом 0.75, а вторая—с весом 0.25. При  \lambda=0  получается обычная явная схема, о которой шла речь в предыдущей лекции.
\end{verbatim}

После замены частных производных \texttt{u\_t} и \texttt{u\_\{xx\}} в
задаче (39.2) мы получаем разностную задачу (см. рис. 39.1).

Разностное уравнение

\begin{verbatim}
\begin{bmatrix} d_t \\ d_{\bar{z}} \\ d_{\bar{s}} \\ \vdots \\ \vdots \\ d_n \end{bmatrix}
\end{verbatim}

(39.3)

,

298

Перенесем все неизвестные значения u с верхнего временного слоя (с
индексом i+1) в левую часть уравнения (39.3) и получим (39.4)

де вводиться позначення \texttt{r\ =\ k/h\^{}2} . Зверніть увагу, що
якщо i фіксоване, а j змінюється від 2 до n-1, співвідношення (39.4)
визначають систему n-2 рівнянь з n-2 невідомим
\texttt{u\_\{t+1,2\},\ u\_\{t+1,3\},}

\begin{verbatim}
\begin{bmatrix} 1 & c_1^* & 0 & 0 & \dots & 0 \\ 0 & 1 & c_2^* & 0 & \dots & 0 \\ 0 & 0 & 1 & c_2^* & \dots & 0 \\ \vdots & \vdots & & & \vdots & \vdots \\ 0 & 0 & \dots & & 1 \end{bmatrix}
\end{verbatim}

РИС. 39.1. Сетка для неявной разностной схемы.

\begin{verbatim}
РИС. 39.2. Шаблон для неявной схемы.  u_{i+1, 4}, \ldots, u_{i+1, n-1}  , которые являются решением задачи во внут-
\end{verbatim}

\begin{verbatim}
ренних узлах сетки на временном слое  t=(i+1)\Delta t  .
\end{verbatim}

Рис. 39.2 дает наглядное представление о ссруктуре каждого уравнения
системы (39.4). Перейдем теперу к решению системы (39.4).

\section{Алгоритм решения задачи
(39.2)}\label{ux430ux43bux433ux43eux440ux438ux442ux43c-ux440ux435ux448ux435ux43dux438ux44f-ux437ux430ux434ux430ux447ux438-39.2}

\begin{verbatim}
ШАГ 1. Выбираем некоторое значение  \lambda  (  0 \le \lambda \le 1  ). Если  \lambda = 0  , то уравнения (39.4), переходят в явные формулы из лекции 38.
\end{verbatim}

\begin{verbatim}
ШАГ 2. Пусть, например,  h=\Delta x=0,2  и  k=\Delta t=0,08  (при этом  r=k/h^2=2  ). В данном случае сетка содержит 6 узлов вдоль осн x (4 внутренних узла) (см. рис. 39.1). Возьчем весовой параметр  \lambda=0,5  (получающаяся при этом схема называется схемой Кранка—Никольсона). В соответствии с вычислительной схемой (ее принято называть шаблоном), изображенной на рис. 39.2, двигаясь слева направо (j=2,3,4,5) по первым двум слоям (i=1), получаем следующие четыре уравнения:
\end{verbatim}

\begin{verbatim}
\begin{bmatrix} x_1 \\ x_2 \\ x_3 \\ \vdots \\ x_n \end{bmatrix}
\end{verbatim}

Перепишем их в матричной форме

(39.5)

\begin{verbatim}
=
\begin{bmatrix} u_{2\bar{z}} \\ u_{\bar{z}a} \\ u_{24} \\ u_{\bar{z}x} \end{bmatrix}
\begin{bmatrix} d_1^* \\ d_2^* \\ d_3^* \\ \vdots \\ d_n^* \end{bmatrix}
\end{verbatim}

.

Матрица этой системы называется трехдиагональной. Для того чтобы решить
трехдиагональную систему

\begin{verbatim}
c_1^* = c_1/b_1
x_n = d_n^*, \quad x_j = d_j^* - c_j^* x_{j+1}, \quad j = n-1, \ n-2, \ \dots, \ 2, \ 1.
\begin{bmatrix} x_1 \\ x_{\bar{z}} \\ x_{\bar{s}} \\ \vdots \\ x_n \end{bmatrix}
\end{verbatim}

\begin{verbatim}
\begin{array}{ll} (\mathrm{Y}\mathrm{H}\Pi) & u_t = u_{xx}, \quad 0 < x < 1, \quad 0 < t < \infty, \\ (\Gamma\mathrm{Y}) & \left\{ \begin{array}{ll} u_{\,0}(0, \ t) = 1, \\ u_{\,x}(1, \ t) + u_{\,}(1, \ t) = g_{\,}(t), \\ u_{\,x}(1, \ t) = 0, \quad 0 \leqslant \kappa \leqslant 1. \end{array}
\end{verbatim}

\begin{verbatim}
\begin{bmatrix} d_t \\ d_{\bar{z}} \\ d_{\bar{s}} \\ \vdots \\ \vdots \\ d_n \end{bmatrix}
\end{verbatim}

\begin{verbatim}
\right. & 0 < t < \infty, \end{array}
\end{verbatim}

преобразуем ее в эквивалентную систему вида

\begin{verbatim}
u_t = u_{xx} + u, \quad 0 < x < 1,
\begin{bmatrix} x_1 \\ x_2 \\ x_3 \\ \vdots \\ x_n \end{bmatrix}
u(x, t) = e^{-\pi^2 t} \sin(\pi x).
\begin{bmatrix} d_1^* \\ d_2^* \\ d_3^* \\ \vdots \\ d_n^* \end{bmatrix}
\end{verbatim}

где

Розв'язок якого показано на рис. 40.1. Що краще --- мати аналітичне
рішення цієї проблеми?

(40.2)

\texttt{,\ \textquotesingle{}c\_\{j+1\}\^{}*\ =\ \textbackslash{}frac\{c\_\{j+1\}\}\{b\_\{j+1\}\ -\ a\_j\ c\_j\^{}*\}\textquotesingle{}\ ,\ \textquotesingle{}j\ =\ 1,\ 2,\ \textbackslash{}ldots,\ n-2\textquotesingle{}\ ,\ и\ \textquotesingle{}d\_1\^{}*\ =\ d\_1/b\_1\textquotesingle{}\ ,\ \textquotesingle{}d\_\{j+1\}\^{}*\ =\ \textbackslash{}frac\{d\_\{j+1\}\ -\ a\_j\ d\_j\^{}*\}\{b\_\{j+1\}\ -\ a\_j\ c\_j\^{}*\}\textquotesingle{}\ ,\ \textquotesingle{}j\ =\ 1,\ 2,\ \textbackslash{}ldots,\ n-1\textquotesingle{}\ .}

В этом преобразовании нет ничего таниственного, поскольку вторая система
полностью эквивалентна первой. Матрица новой системы устроена так, что
эта система очень легко решается. Решая уравнения последовательно снизу
вверх, получаем

\begin{verbatim}
или численное, представленное при  \alpha = 1  в табл. 40.1?
\end{verbatim}

Вопрос хороший, а ответ на него зависит от того, что мы хотим делать с
этим решением дальше. Однако каждому из типов решений присущи очевидные
преимущества.

\section{Преимущества аналитического
решения}\label{ux43fux440ux435ux438ux43cux443ux449ux435ux441ux442ux432ux430-ux430ux43dux430ux43bux438ux442ux438ux447ux435ux441ux43aux43eux433ux43e-ux440ux435ux448ux435ux43dux438ux44f}

\begin{enumerate}
\def\labelenumi{\arabic{enumi}.}
\tightlist
\item
  Очевидно, що розв'язок (40.2) є більш інформативним, ніж таблиця
  чисел. Якщо ми хочемо обчислити розв'язок у певній точці (x, t), ми
  можемо зробити це так точно, як хочемо, просто збільшуючи кількість
  членів ряду, які розглядаються. Водночас легко отримати оцінку зверху
  щодо значення допущеної помилки.
\end{enumerate}

\pandocbounded{\includegraphics[keepaspectratio]{images/_page_301_Figure_4.jpeg}}

РИС. 40.1. Решение уравнения теплопроводности (40.1) в различные моменты
времени.

\begin{itemize}
\tightlist
\item
  \begin{enumerate}
  \def\labelenumi{\arabic{enumi}.}
  \setcounter{enumi}{1}
  \tightlist
  \item
    Аналітичне розв'язання завжди дозволяє обчислити значення розв'язку
    в одній точці (x, t) без необхідності обчислювати значення
    розв'язків в інших точках, як це відбувається при розв'язанні задач
    із явними або неявними різницею схем.
  \end{enumerate}
\item
  \begin{enumerate}
  \def\labelenumi{\arabic{enumi}.}
  \setcounter{enumi}{2}
  \tightlist
  \item
    Аналітичне рішення дозволяє визначити розв'язок у будь-якій точці, а
    не лише на вузлах сітки.
  \end{enumerate}
\item
  \begin{enumerate}
  \def\labelenumi{\arabic{enumi}.}
  \setcounter{enumi}{3}
  \tightlist
  \item
    Найважливішою перевагою аналітичного розв'язку для нас є можливість
    простежити вплив фізичних параметрів, початкових і граничних умов на
    природу розв'язку.
  \end{enumerate}
\end{itemize}

Чисельні методи не виявляють цих регулярностей, оскільки дозволяють
знаходити розв'язок лише за заданих параметрів --- початкових і
граничних умовах. Іноді важливо знати взаємозв'язок між параметрами
моделі та розв'язком, особливо коли йдеться про оцінку фізичних
параметрів за типом розв'язку. Наприклад, припустимо, що розв'язок і
визначається експериментально, і ми знаємо аналітичний розв'язок

u = функция от параметров.

Тогда можно поставить вопрос об определении параметров нак функции
входных данных по схеме

Параметры = функция от u = функция исходных данных.

Такая задача называется параметрической идентификацией. Для
параметрической идентификации необходимо уметь решать уравнение с
частными производными. Несколько позже в этой лекции мы приведем пример
параметрической идентификации в биологии, но сначала посмотрим, чем же
хороши численные решения.

\section{Преимущества численных
решений}\label{ux43fux440ux435ux438ux43cux443ux449ux435ux441ux442ux432ux430-ux447ux438ux441ux43bux435ux43dux43dux44bux445-ux440ux435ux448ux435ux43dux438ux439}

Головна перевага чисельних розв'язків полягає в тому, що їх можна
отримати навіть якщо аналітичне рішення неможливо. Майже всі нелінійні
диференціальні рівняння в похідних похідних потрібно розв'язувати
чисельними методами, і більшість реальних моделей фізики, хімії,
біології тощо є нелінійними. Як правило, лінійні моделі наближаються до
нелінійних лише якщо в них відкинути нелінійні члени. Наприклад, ось
кілька нелінійних рівнянь:

\begin{enumerate}
\def\labelenumi{\arabic{enumi})}
\tightlist
\item
  нелинейное волновое уравнение \texttt{u\_\{tt\}\ =\ u\_\{xx\}\ +}f(u)
\end{enumerate}

\begin{verbatim}
;

2) уравнение реакции с нелинейной диффузией  u_t = u_{xx} +f(u)
\end{verbatim}

;

\begin{enumerate}
\def\labelenumi{\arabic{enumi})}
\setcounter{enumi}{2}
\tightlist
\item
  система Ходжкина---Хаксли
\end{enumerate}

\texttt{Застосовуючи\ цей\ метод\ до\ системи\ чотирьох\ рівнянь\ (39,5),\ отримуємо\ розв\textquotesingle{}язок\ \textquotesingle{}u\_\{22\}=0,60\textquotesingle{}\ ;\ \textquotesingle{}u\_\{23\}=0,80\textquotesingle{}\ ;\ \textquotesingle{}u\_\{24\}=0,80\textquotesingle{}\ ;\ \textquotesingle{}u\_\{23\}=0,60\textquotesingle{}\ .\ Ми\ отримали\ наближені\ значення\ розв\textquotesingle{}язку\ у\ внутрішніх\ точках\ сітки\ на\ \textquotesingle{}t=\textbackslash{}Delta\ t\textquotesingle{}\ .\ Тепер\ ми\ можемо\ зробити\ наступний\ крок\ у\ часі,\ але\ для\ цього\ нам\ доведеться\ розв\textquotesingle{}язати\ нову\ систему\ рівнянь.}

В неявной схеме объем вычислений на каждом шаге больше, чем в явной, но
хорошую точность можно получить даже при гораздо большем шаге.

\section{ЗАДАЧИ}\label{ux437ux430ux434ux430ux447ux438-8}

\begin{itemize}
\tightlist
\item
  \begin{enumerate}
  \def\labelenumi{\arabic{enumi}.}
  \tightlist
  \item
    Из дифференциального уравнения (39.3) получите уравнение (39.4).
  \end{enumerate}
\item
  \begin{enumerate}
  \def\labelenumi{\arabic{enumi}.}
  \setcounter{enumi}{1}
  \tightlist
  \item
    Постройте неявную конечно-разностную схему для задачи
  \end{enumerate}
\end{itemize}

\begin{verbatim}
u(x, 0) = Ae^{-x^2/a}
\begin{aligned} (\text{УЧ}\Pi) & u_t = Du_{xx} - Vu_x, & -\infty < x < \infty, \\ (\text{HY}) & u(x, 0) = Ae^{-x^2, a}, & -\infty < x < \infty. \end{aligned}
u(x, t) = \frac{A}{2\sqrt{\pi Dt}} \int_{-\infty}^{+\infty} e^{-\xi^2/a} e^{-(x-\xi)^2/4Dt} d\xi =
= \frac{A\sqrt{a}}{\sqrt{a+4Dt}} e^{-[x^2/(a+4Dt)]},
\end{verbatim}

\begin{verbatim}
- 4. Как выглядит шаблон для уравнения (39.4), когда  \lambda = 1  ?
\end{verbatim}

\begin{itemize}
\tightlist
\item
  \begin{enumerate}
  \def\labelenumi{\arabic{enumi}.}
  \setcounter{enumi}{4}
  \tightlist
  \item
    Постройте блок-схему розв'язки рівняння теплопровідності (39.2).
    Якщо есть возможность, то напишите программу для ЭВМ Було бы очень
    полезно решить эту задачу численно при различных значениях λ с
    начальным условием и (x, 0) = sin (πx). Сравните полученные при
    различных λ численные результаты с аналитическим решением, яке в
    данном случае имеет вид
  \end{enumerate}
\end{itemize}

,

\texttt{\textbackslash{}overline\{C\}\ =\ 0} (по той же причине, что и
\texttt{\textbackslash{}overline\{A\}\ =\ 0} ),

6 Воспользовавшись формулами, приведенными в тексте, решите систему
алгебраических уравнений (39.5).

Лекція 40

\section{СРАВНЕНИЕ АНАЛИТИЧЕСКИХ РЕШЕНИЙ С
ЧИСЛЕННЫМИ}\label{ux441ux440ux430ux432ux43dux435ux43dux438ux435-ux430ux43dux430ux43bux438ux442ux438ux447ux435ux441ux43aux438ux445-ux440ux435ux448ux435ux43dux438ux439-ux441-ux447ux438ux441ux43bux435ux43dux43dux44bux43cux438}

ЦЕЛЬ ЛЕКЦИИ: Сравнить достоинства и недостатки аналитических и численных
методов розв'язки уравнений з частинними похідними. Приведена важная
математическая задача об идентификации физических величин
(параметрическая идентификация). В связи с цієї задачей рассматривается
один важный пример из биологии.

Вероятно, пришло время обсудить достоинства и недостатки аналитических и
численных решений уравнений з частинними похідними. Сначала выясним, що
мы подразумеваем под этими двумя типами решений.

\section{Аналитические
розв'язки}\label{ux430ux43dux430ux43bux438ux442ux438ux447ux435ux441ux43aux438ux435-ux440ux43eux437ux432ux44fux437ux43aux438}

Под аналитическими решениями мы подразумеваем такие розв'язки, в которых
неизвестная функція и выражена через независимые змінні и параметры
системи в виде формул, бесконечных рядов и интегралов.

\section{Численные
розв'язки}\label{ux447ux438ux441ux43bux435ux43dux43dux44bux435-ux440ux43eux437ux432ux44fux437ux43aux438}

Під чисельним ми зрозуміємо чисельні розв'язки, отримані чисельно після
приблизної заміни початкового рівняння іншим, простішим рівнянням.
Наприклад, у методі скінченної різниці похідні виражаються через
скінченні різниці, а розв'язок диференціального рівняння в похідних у
похідних апроксимується розв'язком різницького рівняння. Результатом
такої процедури зазвичай є таблиця значень розв'язків, а для деяких
значень пояснювальних змінних.

Таблина 40.1

\textbar{} × \textbar{} \textbar{} \textbar{} \textbar{} \textbar{}
\textbar{} \textbar{} \textbar{} \textbar{} \textbar{} \textbar{}

\textbar---------------------------\textbar-------------\textbar------------------\textbar-----------\textbar----------\textbar----------\textbar-----\textbar-----\textbar-----\textbar-----------\textbar------------------\textbar------------------\textbar{}

\textbar{} 0 \textbar{} 0,1 \textbar{} 0,2 \textbar{} 0,3 \textbar{} 0,4
\textbar{} 0,5 \textbar{} 0,6 \textbar{} 0,7 \textbar{} 0,8 \textbar{}
0,9 \textbar{} 1 \textbar{}

00,010,020,03 \textbar{} 100 \textbar{} 10,20,15 \textbar{} 10,34
\textbar{} l: \textbar{} : \textbar{} : \textbar{} ! \textbar{} :
\textbar{} 10,34 \textbar{} 10,20,15 \textbar{} 1000 \textbar{}

; \textbar{} \textbar{} : \textbar{} : \textbar{} : \textbar{} :
\textbar{} : \textbar{} : \textbar{} : \textbar{} : \textbar{} :
\textbar{} : \textbar{}

Теперь, когда мы знаем, що подразумевается под решением каждого из цих
типов, давайте сравним их между собой.

\section{\texorpdfstring{\textbf{Сравнительный анализ численного} и
аналитического
решений}{Сравнительный анализ численного и аналитического решений}}\label{ux441ux440ux430ux432ux43dux438ux442ux435ux43bux44cux43dux44bux439-ux430ux43dux430ux43bux438ux437-ux447ux438ux441ux43bux435ux43dux43dux43eux433ux43e-ux438-ux430ux43dux430ux43bux438ux442ux438ux447ux435ux441ux43aux43eux433ux43e-ux440ux435ux448ux435ux43dux438ux439}

Рассмотрим простейшую смещанную параболическую задачу

(40.1)

\begin{verbatim}
\bar{D} = A\xi_{xx} + B\xi_{xy} + C\xi_{yy} + D\xi_x + E\xi_y = -2(x^2 + y^2)
\end{verbatim}

Розв'язок якого показано на рис. 40.1. Що краще --- мати аналітичне
рішення цієї проблеми?

(40.2)

\begin{verbatim}
\widetilde{E} = A\eta_{xx} + B\eta_{xy} + C\eta_{yy} + D\eta_x + E\eta_y = 2(y^2 - x^2),
\end{verbatim}

\begin{verbatim}
или численное, представленное при  \alpha = 1  в табл. 40.1?
\end{verbatim}

Вопрос хороший, а ответ на него зависит от того, що мы хотим делать с
этим решением дальше. Однако каждому из типов решений присущи очевидные
преимущества.

\section{Преимущества аналитического
розв'язки}\label{ux43fux440ux435ux438ux43cux443ux449ux435ux441ux442ux432ux430-ux430ux43dux430ux43bux438ux442ux438ux447ux435ux441ux43aux43eux433ux43e-ux440ux43eux437ux432ux44fux437ux43aux438}

\begin{enumerate}
\def\labelenumi{\arabic{enumi}.}
\tightlist
\item
  Очевидно, що розв'язок (40.2) є більш інформативним, ніж таблиця
  чисел. Якщо ми хочемо обчислити розв'язок у певній точці (x, t), ми
  можемо зробити це так точно, як хочемо, просто збільшуючи кількість
  членів ряду, які розглядаються. Водночас легко отримати оцінку зверху
  щодо значення допущеної помилки.
\end{enumerate}

\pandocbounded{\includegraphics[keepaspectratio]{images/_page_301_Figure_4.jpeg}}

Рис. 40.1. Розв'язок рівняння теплопровідності (40.1) в различные
моменты времени.

\begin{itemize}
\tightlist
\item
  \begin{enumerate}
  \def\labelenumi{\arabic{enumi}.}
  \setcounter{enumi}{1}
  \tightlist
  \item
    Аналітичне розв'язання завжди дозволяє обчислити значення розв'язку
    в одній точці (x, t) без необхідності обчислювати значення
    розв'язків в інших точках, як це відбувається при розв'язанні задач
    із явними або неявними різницею схем.
  \end{enumerate}
\item
  \begin{enumerate}
  \def\labelenumi{\arabic{enumi}.}
  \setcounter{enumi}{2}
  \tightlist
  \item
    Аналітичне рішення дозволяє визначити розв'язок у будь-якій точці, а
    не лише на вузлах сітки.
  \end{enumerate}
\item
  \begin{enumerate}
  \def\labelenumi{\arabic{enumi}.}
  \setcounter{enumi}{3}
  \tightlist
  \item
    Найважливішою перевагою аналітичного розв'язку для нас є можливість
    простежити вплив фізичних параметрів, початкових і граничних умов на
    природу розв'язку.
  \end{enumerate}
\end{itemize}

Чисельні методи не виявляють цих регулярностей, оскільки дозволяють
знаходити розв'язок лише за заданих параметрів --- початкових і
граничних умовах. Іноді важливо знати взаємозв'язок між параметрами
моделі та розв'язком, особливо коли йдеться про оцінку фізичних
параметрів за типом розв'язку. Наприклад, припустимо, що розв'язок і
визначається експериментально, і ми знаємо аналітичний розв'язок

u = функція от параметров.

Тогда можна поставить вопрос об определении параметров нак функції
входных данных по схеме

Параметры = функція от u = функція исходных данных.

Такая задача називається параметрической идентификацией. Для
параметрической идентификации необходимо уметь решать уравнение з
частинними похідними. Несколько позже в цієї лекции мы приведем пример
параметрической идентификации в биологии, но сначала посмотрим, чем же
хороши численные розв'язки.

\section{Преимущества численных
решений}\label{ux43fux440ux435ux438ux43cux443ux449ux435ux441ux442ux432ux430-ux447ux438ux441ux43bux435ux43dux43dux44bux445-ux440ux435ux448ux435ux43dux438ux439-1}

Головна перевага чисельних розв'язків полягає в тому, що їх можна
отримати навіть якщо аналітичне рішення неможливо. Майже всі нелінійні
диференціальні рівняння в похідних похідних потрібно розв'язувати
чисельними методами, і більшість реальних моделей фізики, хімії,
біології тощо є нелінійними. Як правило, лінійні моделі наближаються до
нелінійних лише якщо в них відкинути нелінійні члени. Наприклад, ось
кілька нелінійних рівнянь:

\begin{enumerate}
\def\labelenumi{\arabic{enumi})}
\tightlist
\item
  нелинейное хвильове рівняння \texttt{u\_\{tt\}\ =\ u\_\{xx\}\ +}f(u)
\end{enumerate}

\begin{verbatim}
;

2) уравнение реакции с нелинейной диффузией  u_t = u_{xx} +f(u)
\end{verbatim}

;

\begin{enumerate}
\def\labelenumi{\arabic{enumi})}
\setcounter{enumi}{2}
\tightlist
\item
  система Ходжкина---Хаксли
\end{enumerate}

\begin{verbatim}
u(x, t) = \frac{A\sqrt{a}}{\sqrt{a+4Dt}}e^{-[(x-Vt)^{2}/(a+4Dt)]}.
\end{verbatim}

Аналитические розв'язки для цих уравнений не известны ни при каких
нелинейных функциях f и \texttt{g\^{}\{(1)\}} . Поэтому общий подход к
решению нелинейных (а в ряде случаев и линейных) задач обычно базируется
на численных решениях.

Тепер давайте розглянемо приклад того, як можна використати аналітичне
рішення для визначення фізичних параметрів. Детальніший опис цих завдань
можна знайти у рекомендованій літературі.

\section{Параметрическая идентификация (в
биологии)}\label{ux43fux430ux440ux430ux43cux435ux442ux440ux438ux447ux435ux441ux43aux430ux44f-ux438ux434ux435ux43dux442ux438ux444ux438ux43aux430ux446ux438ux44f-ux432-ux431ux438ux43eux43bux43eux433ux438ux438}

Припустимо, що біолог намагається визначити швидкість, з якою (K\^{}+)
іони калію дифундують у екзоплазмовому розчині. Якщо ви знаєте цей
коефіцієнт дифузії, можна багато сказати про те, як нервові імпульси
передаються вздовж аксонів. Проблема в тому, що це співвідношення майже
неможливо знайти прямим вимірюванням. Однак можливо знайти математичний

\begin{verbatim}
<sup>1)</sup> Це не так. В качестве примера можна привести уравнение  u_{ii} = u_{xx} - \sin u_x - \Pi pum  , ped.
\end{verbatim}

\begin{verbatim}
связь между концентрацией поташа  u\left(x,\,t\right)  и коэффициентом диффузии  D,\,\,a\,\,  затем по измерениям величины  u\left(x,\,t\right)  определить величину D.
\end{verbatim}

Біологи Ходжкін і Кейз виявили, що після розміщення гігантського аксону
у спеціальному соляному розчині концентрація радіохктивного калію (**2K)
вздовж аксону приблизно описується кривою в початковий момент

,

.

Отже, параметри A і A слід обирати так, щоб крива відповідала
експериментальним даним (наприклад, у сенсі методу найменших квадратів).
Також було встановлено, що P И С

\pandocbounded{\includegraphics[keepaspectratio]{images/_page_303_Figure_6.jpeg}}

\begin{verbatim}
. 40.2. Начальная концентрация ионов (  ^{42}K  ): a—теоретическая кривая;  \delta  —экспериментальные точки.
\end{verbatim}

с течением времени ноны поташа под действием диффузии и конвекции
растекаются вдоль аксона. Ходжкин и Кейс предположили, що концентрация и
является решением следующей задачи с начальными данными для уравнения
диффузии:

(40.3)

\begin{verbatim}
\overline{G} = G = 0
\end{verbatim}

Цю проблему вирішили у лекції 15, переключившись на рухому систему
координат. Якщо читач пам'ятає, загальна схема методу виглядає так.
Спочатку відкидаємо конвективний член і отримуємо розв'язок суто
дифузійної задачі (при V = 0):

,

подставим их в уравнение

\begin{verbatim}
= \frac{A\sqrt{a}}{\sqrt{a+4Dt}} e^{-[x^2/(a+4Dt)]},
\overline{A}u_{\xi\xi} + \overline{B}u_{\xi\eta} + \overline{C}u_{\eta\eta} + \overline{D}u + \overline{E}u_{\eta} + \overline{F}u = \overline{G}
u(x, t) = \frac{A\sqrt{a}}{\sqrt{a+4Dt}}e^{-[(x-Vt)^{2}/(a+4Dt)]}.
\end{verbatim}

и получим

\begin{verbatim}
SS = \sum_{i=1}^{n} [y_i - (a + bx_i)]^2.
u_{\xi\eta} = \frac{-(x^2 + y^2) u_{\xi} + (y^2 - x^2) u_{\eta}}{8x^2y^3}.
u_t = \frac{1}{r} \frac{\partial}{\partial r} \left( Dr u_r - s \omega^2 r^2 u \right), \quad 0 < r < 1,
\end{verbatim}

ШАГ 2. Выражая, наконец, x и y через \texttt{\textbackslash{}xi} и
\texttt{\textbackslash{}eta} , окончательно получаем

(23.8)

\begin{verbatim}
Au_{xx} + Bu_{xy} + Cu_{yy} + Du_x + Eu_y + Fu = G
u_{\xi\eta} = \frac{\eta u_{\xi} - \xi u_{\eta}}{2(\xi^2 - \eta^2)}.
\begin{cases} 0, & t < x/V, \\ P, & t \ge x/V. \end{cases}
\end{verbatim}

Задача решена. Теперь мы знаем, что нужно сделать, чтобы найти новые
координаты и привести уравнение к канонической форме.

\section{ЗАМЕЧАНИЯ}\label{ux437ux430ux43cux435ux447ux430ux43dux438ux44f-3}

\begin{enumerate}
\def\labelenumi{\arabic{enumi}.}
\tightlist
\item
  На самом деле для гиперболических уравнений существует две
  канонические формы. Вторую можно получить из первов заменой переменных
  вида
\end{enumerate}

Конечно, все це очень просто! Не сложнее, чем сбрасывать що-нибудь на
ленту транспортера, а затем наблюдать за его движением. Однако все
становится очень интересным, якщо примесь диффундирует в среде. Чтобы
разобраться в том, що про-исходит с движущейся волной при наличии
диффузии, решим

следующую задачу

(15.2)

\begin{verbatim}
\alpha = \alpha (\xi, \eta) = \xi + \eta,
\end{verbatim}

\begin{verbatim}
З формальною заміною змінних  \xi = \xi(x, y)  та  \eta = \eta(x, y)  (41.1) зводиться до тієї ж форми
\end{verbatim}

(41.2)

\texttt{\textbackslash{}beta\ =\ \textbackslash{}beta\ (\textbackslash{}xi,\ \textbackslash{}eta)\ =\ \textbackslash{}xi\ -\ \textbackslash{}eta.}
Воспользуемся этой заменой для преобразования уравнения (23.8). Так как

где, как обычно, \texttt{H(x)} --- функція Хевисайда Начальное
распределение концентрации изображено на рис 15 4.

\pandocbounded{\includegraphics[keepaspectratio]{images/_page_114_Figure_4.jpeg}}

\begin{verbatim}
Рис. 15.3 Конвективная волна a—передиції фронт примеси,  \delta  —на одну единицу вперед волнового фронта
\end{verbatim}

\begin{verbatim}
Обратим внимание на то, що в задаче (15.2) мы отодвинули границу в —  \infty  (теперь решается задача Коши), так що она не оказывает теперь влияние на оценку розв'язок. Чтобы решить задачу (15.2), можна воспользоваться преобразованием Лапласа Р И С
\end{verbatim}

\pandocbounded{\includegraphics[keepaspectratio]{images/_page_114_Figure_7.jpeg}}

. 15 4 початкові умови в задаче конвективной диффузии.

\begin{verbatim}
ло переменной t или преобразованием Фурье по переменной x. Однако в цьому случае гораздо интереснее поступить совсем по-другому Введем новую систему координат, яка движется вдоль старой со скоростью V. Другими словами, вместо системи координат, привязанной к берегу реки, мы рассмотрим систему координат, яка движется со скоростью фронта примеси (конечно, при налични диффузии фронт буде «смазан»). С точки зрения математики це означает, що мы заменяем пространственную координату x на новую  \xi = x - Vt  .
\end{verbatim}

Теперь ясно, що

\begin{verbatim}
когда  \xi = 0  , мы находимся па фронте распространяющейся примеси.
\end{verbatim}

\begin{verbatim}
когда  \xi = 1  , мы находимся на одну единицу длины впереди фронта.
\end{verbatim}

\begin{verbatim}
когда  \xi = -1  , мы находимся на одну единицу длины позади фронта.
\end{verbatim}

Наша задача --- преобразовать исходную задачу с НУ

\begin{verbatim}
\overline{D} = A\xi_{xx} + B\xi_{xy} + C\xi_{yy} + D\xi_{x} + E\xi_{y},
[\xi_x/\xi_y] = -B/2A
\end{verbatim}

TO

(23 9)

\begin{verbatim}
\xi(x, y) = c
u_{\alpha\alpha} - u_{\beta\beta} = \frac{-\beta u_{\alpha} - \alpha u_{\beta}}{2\alpha\beta}.
\frac{dy}{dx} = -\left[\xi_x/\xi_y\right] = B_t 2A.
\end{verbatim}

При желании можно выразить \texttt{\textbackslash{}alpha} и
\texttt{\textbackslash{}beta} через исходные координаты x и y. В
результате получаем

\begin{verbatim}
\overline{B} = 2A\xi_x\eta_x + B(\xi_x\eta_y + \xi_y\eta_x) + 2C\xi_y\eta_y,
\alpha = \xi + \eta = (y^2 - x^2) + (y^2 + x^2) = 2y^3,
\end{verbatim}

1) Завдяки рівності

І з B\^{}2 - 4AC = 0 році , то

\begin{verbatim}
\beta = \xi - \eta = (y^2 - x^2) - (y^2 + x^2) = -2x^2.
\end{verbatim}

\begin{verbatim}
Оскільки  \xi_x/\xi_y = -B \ 2A = -2 \sqrt{A\tilde{C}}/2A = -\sqrt{C/A}  останній вираз для B буде записаний як
\end{verbatim}

\begin{verbatim}
u_{\xi\xi} - u_{\eta\eta} = \Psi(\xi, \, \eta, \, u_{\xi}, \, u, u_{\eta}).
\end{verbatim}

\begin{verbatim}
І оскільки  \overline{A}  дорівнює нулю, то  \overline{B}  , звісно, також дорівнює нулю.
\end{verbatim}

\begin{verbatim}
Оскільки при нашому виборі  \xi  обидва коефіцієнти  \overline{A}  і  \overline{B}  повертаються до нуля, ми можемо взяти будь-яку функцію як  \eta  (якщо вона не пропорційна координаті  \xi  ). У прикладі можна вибрати  \eta = y  .
\end{verbatim}

\begin{verbatim}
Залишається знайти канонічну форму рівняння в нових координатах. Для цього потрібно просто підставити формули  \xi  і  \eta  (41.3) і знайти всі шанси  \overline{A}  ,  \overline{B}  ,  \overline{C}  ,  \overline{D}  ,  \overline{E}  ,  \overline{F}  і  \overline{G}  . На цьому вивчення параболічного випадку завершується, але перш ніж перейти до іншої теми, розглянемо простий приклад.
\end{verbatim}

\section{\texorpdfstring{Перетворення параболічного рівняння
\texttt{u\_\{xx\}} + 2u\_\{xy\} + \texttt{u\_\{yy\}\ =\ 0}у канонічну
форму}{Перетворення параболічного рівняння u\_\{xx\} + 2u\_\{xy\} + u\_\{yy\} = 0у канонічну форму}}\label{ux43fux435ux440ux435ux442ux432ux43eux440ux435ux43dux43dux44f-ux43fux430ux440ux430ux431ux43eux43bux456ux447ux43dux43eux433ux43e-ux440ux456ux432ux43dux44fux43dux43dux44f-u_xx-2u_xy-u_yy-0ux443-ux43aux430ux43dux43eux43dux456ux447ux43dux443-ux444ux43eux440ux43cux443}

Розглянемо рівняння

\begin{verbatim}
3u_{xx} + 7u_{xy} + 2u_{yy} = 0
\end{verbatim}

\begin{verbatim}
Його коефіцієнти мають такі значення: A=1, B=2, C=1, D=E=F=G=0. Отже,  B^2-4AC=0  для всіх значень x і y. Щоб знайти нові координати  \xi  і  \eta  та канонічну форму рівняння, ми рухаємося наступним чином.
\end{verbatim}

КРОК 1. Запишемо \emph{характеристичне рівняння} (тепер лише одне, а не
два)

\begin{verbatim}
u_{\xi\eta} = \Phi(\xi, \eta, u, u_1, u_2).
\end{verbatim}

Розв'язуючи це рівняння відносно y, отримуємо

\begin{verbatim}
u_{\alpha\alpha} - u_{\beta\beta} = \Psi(\alpha, \beta, u, u_{\alpha}, u_{\beta}).
\end{verbatim}

,

\begin{verbatim}
Де — характеристика  \xi = y - x  визначена. Використання цієї характеристики забезпечує досягнення рівності  \overline{A} = 0  . Координату  \eta  можна вибрати довільно (якщо вона не залежить від  \xi  ). Ми оберемо його наступним чином:
\end{verbatim}

\begin{verbatim}
u_{xx} + 4u_{xy} = 0.
\end{verbatim}

.

Нові координати

\begin{verbatim}
u_{tt} = e^{2} (u_{xx} + u_{yy} + u_{zz}),
\end{verbatim}

изображены на рис. 41 1.

\pandocbounded{\includegraphics[keepaspectratio]{images/_page_309_Figure_3.jpeg}}

\begin{verbatim}
Рис. 41.1. Новая система координат  \xi = y - x  ,  \eta = y  .
\end{verbatim}

\begin{verbatim}
ШАГ 2. Осталось найти каноническую форму уравнения. Подставляя  \xi  и  \eta  в формулы для коефіцієнтів  \overline{A}  ,  \overline{B}  ,  \overline{C}  ,  \overline{D}  ,  \overline{E}  ,  \overline{F}  и  \overline{G}  , получаем  \overline{A} = 0  (так и должно быть; мы находили  \xi  из условия равенства нулю цього коэффициента),  \overline{B}=0  (мы уже раньше показали, що этог коефіцієнт також обращается в нуль),
\end{verbatim}

(41.5)

\begin{verbatim}
\overline{D} = A \xi_{xx} + B \xi_{xy} + C \eta_{yy} + D \xi_x + E \xi_y = 0,
\end{verbatim}

Отже, нове рівняння

(НУ)

\begin{verbatim}
\xi = x - Vt
u(x, t) =
\end{verbatim}

є загальним розв'язком рівняння

в трехмерном пространстве и показать, что полученное решение
удовлетворяет \emph{принципу Гюйгенса}. Для решения соответствующей
двумерной задачи

(УЧП)

Перед нами розв'язок задачи конвективной диффузии (15.2). Его очень
легко интерпретировать, якщо представить, що мы движемся относительно
графика, изображенного на рис. 15.6.

\pandocbounded{\includegraphics[keepaspectratio]{images/_page_117_Figure_3.jpeg}}

Рис. 15.6. Диффузня из области высокой концентрации в область низкой
концентрации. Чем больше коефіцієнт дифузії, тем быстрее установится
стационарное значение.

Другими словами, в зависимости от относительной величины D (коэффициента
диффузии) и V (скорости потока) розв'язок движется слева направо со
скоростью V и в то же самое время передний фронт расплывается со
скоростью, определяемой величиной D.

\pandocbounded{\includegraphics[keepaspectratio]{images/_page_117_Figure_6.jpeg}}

Рис. 15.7. Розв'язок задачи конвективнои диффузии. Вещество одновременно
движется и диффундирует.

(На рис. 15.7 показано, как расплывается передний фронт концентрации.)

\section{ЗАМЕЧАНИЯ}\label{ux437ux430ux43cux435ux447ux430ux43dux438ux44f-4}

Преобразование координат является важным методом розв'язки уравнений з
частинними похідними. Выбрав подходящую систему координат, можна
существенно упростить уравнение.

\section{ЗАДАЧИ}\label{ux437ux430ux434ux430ux447ux438-9}

\begin{enumerate}
\def\labelenumi{\arabic{enumi}.}
\tightlist
\item
  Решите задачу Коши:
\end{enumerate}

\begin{verbatim}
u_{tt} = c^2 (u_{xx} + u_{yy}), \quad
\end{verbatim}

Але тепер у B\^{}2-4AC\textless0 . Переходячи до нових пояснювальних
змінних, ми хочемо перетворити її у нову форму

\begin{verbatim}
\begin{cases} -\infty < x < \infty, \\ -\infty < y < \infty, \end{cases}
\end{verbatim}

\begin{enumerate}
\def\labelenumi{\arabic{enumi}.}
\setcounter{enumi}{1}
\tightlist
\item
  Найдите розв'язок задачи Коши для уравнения конвективной диффузии:
\end{enumerate}

\begin{verbatim}
\begin{cases} -\infty < x < \infty, \\ -\infty < y < \infty, \\ -\infty < z < \infty, \end{cases}
\end{verbatim}

\begin{verbatim}
Це можливо, але лише шляхом введення складних координат. Щоб знайти ці комплексні координати  \xi  і  \eta  , розв'яжемо характеристичні рівняння, так само як у гіперболічному випадку:
\end{verbatim}

\begin{verbatim}
\begin{cases} u(x, y, z, 0) = \varphi(x, y, z), \\ u_{t}(x, y, z, 0) = \psi(x, y, z). \end{cases}
\end{verbatim}

Воспользуйтесь преобразованием из лекции 8.

\begin{enumerate}
\def\labelenumi{\arabic{enumi}.}
\setcounter{enumi}{2}
\tightlist
\item
  Найдите розв'язок задачи переноса:
\end{enumerate}

\begin{verbatim}
P_{2}
\frac{dy}{dx} = \sqrt{-4x^2} = 2ix
\begin{aligned} u_{tt} &= c^2 \Delta u, \\ (HY) & \begin{cases} u(x, y, z, 0) = 0, \\ u_t(x, y, z, 0) = \psi(x, y, z), \end{cases} \end{aligned}
\\ -\infty < x < \infty, \quad 0 < t < \infty,
\end{verbatim}

Цю задачу можна розв'язати методом \texttt{\textbackslash{}Phi} Юр'єра
\texttt{\^{}\{1\}} ), і виявляється, що розв'язок можна представити у
вигляді

\begin{verbatim}
u_t = Du_{xx} - Vu_x, \quad -\infty < x < \infty, \quad 0 < t < \infty,
P_{1}
u(x, 0) = e^{-x^2}, \quad -\infty < x < \infty.
\end{verbatim}

где \texttt{\textbackslash{}overline\{\textbackslash{}psi\}} ---
\emph{среднее} значение начального распределения
\texttt{\textbackslash{}psi} \emph{по сфере} радиуса ct с центром в
точке (x, y, z), т. е.

\begin{verbatim}
y^2 u_{xx} + x^2 u_{yy} = 0,
\end{verbatim}

\pandocbounded{\includegraphics[keepaspectratio]{images/_page_177_Picture_3.jpeg}}

\begin{verbatim}
\frac{dy}{dx} = \frac{B - \sqrt{B^2 - 4AC}}{2A} = -\frac{\sqrt{-4x^2y^2}}{2y^2} = -i\frac{x}{y},
\end{verbatim}

\pandocbounded{\includegraphics[keepaspectratio]{images/_page_177_Picture_7.jpeg}}

\begin{verbatim}
u(\xi, \tau) = \frac{1}{2\sqrt{D\pi\tau}} \int_{-\infty}^{+\infty} e^{-\beta^2} e^{-(\xi-\beta)^2/4D\tau} d\beta.
\begin{aligned} (\text{YAII}) & u_{tt} = c^2 \Delta u, \\ (\text{HY}) & \begin{cases} u = 0, \\ u_t = \varphi, \end{cases} \end{aligned}
u_{tt} = \alpha^2 u_{xx}
u_h(r, \theta) = \sum_{n=0}^{\infty} r^n [a_n \cos(n\theta) + b_n \sin(n\theta)].
u_{\xi\eta} = \psi(\xi, \eta, u, u_{\xi}, u_{\eta}).
\Delta u = -\frac{r^2}{2} - \frac{r^3}{2} \cos{(2\theta)}
\xi(x, y) = y^2, \quad \eta(x, y) = x^2.
u = \frac{\partial}{\partial t} [t\overline{\varphi}].
T \sin \theta_2 - T \sin \theta_1 =
u_p(r, \theta) = Ar^4 + Br^4 \cos(2\theta).
\Delta x \rho u_{tt} = T \left[ u_x \left( x + \Delta x, \ t \right) - u_x \left( x, \ t \right) \right] + \Delta x F \left( x, \ t \right) - \frac{1}{\Delta x \rho u_t} \left( x, \ t \right) - \Delta x \gamma \dot{u} \left( x, \ t \right)
u(x, t) = \frac{1}{2c} \int_{x-ct}^{x+ct} \varphi(s) ds.
u_{tt} = \alpha^{3} u_{xx} - \beta u_{t} - \gamma u + F(x, t).
u_p(r, \theta) = -\frac{r^4}{32} - \frac{r^4}{24} \cos(2\theta).
(16.2) u_{tt} = \alpha^2 u_{xx},
u_t(x, t) = \frac{1}{2} [\varphi(x + ct) + \varphi(x - ct)].
u_{tt} = ku_{xx}
\sum_{n=0}^{\infty} [a_n \cos(n\theta) + b_n \sin(n\theta)] - \frac{1}{32} - \frac{1}{24} \cos(2\theta) = 0.
u_{tt} = \frac{\partial}{\partial x} [\alpha^2(x) u_x],
u_1(r, \theta) = \frac{1}{32} + \frac{1}{24}r^2\cos(2\theta) - \frac{r^4}{24}\cos(2\theta) - \frac{r^4}{32} = \frac{(r^4 - 1)}{32} - \frac{(r^4 - r^2)}{24}\cos(2\theta).
u(x, 0) = f(x)
u = u_0 + u_1 = r \cos \theta - \frac{(r^4 - 1)}{32} - \frac{(r^4 - r^2)}{24} \cos (2\theta).
i_x + Cv_t + Gv = 0, \\ v_x + Li_t + Ri = 0,
i_{xx} + Gv_x - CLi_{tt} - CRi_t = 0.
v_{x} = -Li_{t} - Ri,
P_{2}
i_{t,s} = CLi_{tt} + (CR + GL)i_t + GRi.
\Delta u = 0, \quad 0 < r < 1 + \frac{1}{4} \sin \theta,
v_{xx} = CLv_{tt} + (CR + GL)v_t + GRv.
\end{verbatim}

записывается в виде

\begin{verbatim}
v_{tt} = \alpha^2 v_{xx}, \\ i_{tt} = \alpha^2 i_{xx}, \qquad \alpha^2 = 1/CL,
u \left( 1 + \frac{1}{4} \sin \theta, \ \theta \right) = \cos \theta, \quad 0 \le \theta \le 2\pi.
u_{tt} = u_{xx} - u_t
\end{verbatim}

где \texttt{\textbackslash{}overline\{\textbackslash{}phi\}} и
\texttt{\textbackslash{}overline\{\textbackslash{}psi\}} --- средние
значения функций \texttt{\textbackslash{}phi} и
\texttt{\textbackslash{}psi} на сфере раднусом ct с центром в точке (x,
y, z).

Цей розв'язок зазвичай називають формулою Пуассона для хвильового
рівняння у тривимірному просторі. Це природне узагальнення формули
д'Аламбера на тривимірний випадок. Найважливішим у формулі Пуассона є
те, що обидва інтеграли в
\texttt{\textbackslash{}bar\{\textbackslash{}phi\}} і
\texttt{\textbackslash{}bar\{\textbackslash{}psi\}} взяті з поверхні
сфери. Цей факт дозволяє нам надати таку важливу інтерпретацію рішення.

У момент \texttt{t=t\_1} розв'язок u у (x,y,z) залежить лише від
початкових розподілів сфери з радіусом \texttt{ct\_1} центрованим на
(x,y,z) (рисунок 24.3). Тепер припустимо, що початкові розподіли
\texttt{\textbackslash{}phi} і \texttt{\textbackslash{}psi} дорівнюють
нулю всюди, крім малої сфери (див. рисунок 24.3). З часом радіус P И С

\pandocbounded{\includegraphics[keepaspectratio]{images/_page_179_Picture_3.jpeg}}

. 24.3. Схема поясияет, как начальное возмущение влияет на точку (к, ч,
г).

Сфера, що оточує точку (x, y, z), зростає зі швидкістю c, так що через
\texttt{t\_\{\textbackslash{}hat\{z\}\}} секунди ця сфера починає
перетинати область початкового збурення, і тому розв'язок
\texttt{u(x,\ y,\ z,\ t)} стає ненульовим. У
\texttt{t\_\{\textbackslash{}hat\{z\}\}\ \textless{}\ t\ \textless{}\ t\_\{\textbackslash{}hat\{a\}\}}
розв'язок U у (x, y, z) залишається ненульовим, оскільки сфера продовжує
перетинати область початкового збурення, але на
\texttt{t=t\_\{\textbackslash{}hat\{a\}\}} розв'язок знову стає
нульовим. Іншими словами, розв'язок, що поширюється з початкової області
збурення, має чітко визначену задню кромку. Цей загальний принцип
відомий як принцип Гюйгенса для тривимірного простору. Виявляється, що
передній край хвилі завжди чітко окреслений, але задній край чітко
окреслений лише в просторах розмірності 3, 5, 7\ldots. Ми вже знаємо з
формули д'Аламбера, що початкові умови

\pandocbounded{\includegraphics[keepaspectratio]{images/_page_349_Figure_5.jpeg}}

Що ви можете сказати про точність?

\begin{enumerate}
\def\labelenumi{\arabic{enumi}.}
\setcounter{enumi}{3}
\tightlist
\item
  Як отримати послідовність випадкових точок, що лежать всередині
  трикутника T?
\end{enumerate}

\pandocbounded{
\includegraphics[keepaspectratio]{images/_page_319_Figure_13.jpeg}}

\pandocbounded{\includegraphics[keepaspectratio]{images/_page_319_Figure_14.jpeg}}

\begin{enumerate}
\def\labelenumi{\arabic{enumi}.}
\setcounter{enumi}{4}
\tightlist
\item
  Як отримати послідовність випадкових чисел, розподіл яких показаний на
  рисунку?
\end{enumerate}

\pandocbounded{\includegraphics[keepaspectratio]{images/_page_319_Figure_16.jpeg}}

Іншими словами, як отримати послідовність чисел \{0, 1, 2\}, якщо
ймовірність виникнення 0 і 2 дорівнює 0,25, а ймовірність появи 1 ---
0,5?

Лекція 43

\section{РОЗВ'ЯЗАННЯ ДИФЕРЕНЦІАЛЬНИХ РІВНЯНЬ У ПОХІДНИХ ПОХІДНИХ
МОНТЕ-КАРЛО}\label{ux440ux43eux437ux432ux44fux437ux430ux43dux43dux44f-ux434ux438ux444ux435ux440ux435ux43dux446ux456ux430ux43bux44cux43dux438ux445-ux440ux456ux432ux43dux44fux43dux44c-ux443-ux43fux43eux445ux456ux434ux43dux438ux445-ux43fux43eux445ux456ux434ux43dux438ux445-ux43cux43eux43dux442ux435-ux43aux430ux440ux43bux43e}

МЕТА ЛЕКЦІЇ: Показати, як побудувати таку гру, результатом якої є
приблизне розв'язання диференціального рівняння. Одна з цих ігор (модель
випадкової ходьби) веде до скінченної різниці апроксимації задачі
Діріхле в квадраті. Цю модель можна узагальнити для отримання розв'язків
інших задач.

У попередній лекції було сказано, що можливо винайти такі ігри на удачу,
результатом яких є розв'язки (можливо, приблизні) задач для рівнянь з
частковими коефіцієнтами проникнення. У цій лекції ми покажемо, як
побудувати таку гру для наближеного розв'язання задачі Діріхле (PPP)
\texttt{u\_\{xx\}} +
\texttt{u\_\{yy\}\ =\ 0\ ,\ 0\ \textless{}\ x\ \textless{}\ 1,\ 0\ \textless{}\ y\ \textless{}\ 1},

(GI)

\begin{verbatim}
u\left(1+\frac{1}{4}\sin\theta,\ \theta\right)=\cos\theta
R(A) = g_1 P_A(p_1) + g_2 P_A(p_2) + \dots + g_{12} P_A(p_{12}).
\end{verbatim}

приводят, вообще говоря, к нерезкому заднему фронту (поскольку формула
Даламбера подразумевает интегрирование функции
\texttt{\textbackslash{}psi(x)} в пределах от x-ct до x+ct).

Покажем теперь, что принцип Гюйгенса неприменим к цилиндрическим волнам.
Такого рода решения описывают волны на воде от точечного источника,
когда задний фронт волны не резкий, а постепенно затухает до нуля.

\section{Двумерное волновое
уравнение}\label{ux434ux432ux443ux43cux435ux440ux43dux43eux435-ux432ux43eux43bux43dux43eux432ux43eux435-ux443ux440ux430ux432ux43dux435ux43dux438ux435}

Чтобы решить двумерную волновую задачу

(УЧП)

\begin{verbatim}
R(A) = 1 \cdot (0.04) + 1 \cdot (0.15) + 1 \cdot (0.03) + 0 \cdot (0.06) + \dot + 0 \cdot (0.04).
u(1+\epsilon\sin\theta, \theta) = \cos\theta, \quad 0 \le \epsilon \le 1/4.
R(A) = \frac{1}{4} [R(B) + R(C) + R(D) + R(E)]
f(x+h) = f(x) + f'(x)h + f''(x)\frac{h^2}{2!} + \dots,
u_{i,j} = \frac{1}{4} \left[ u_{i-1,j} + u_{i+1,j} + u_{i,j-1} + u_{i,j+1} \right],
\end{verbatim}

(PPP)

\begin{verbatim}
u_{xx}=u_{l,\ j+1}-2u_{l,\ j}+u_{l,\ j-1},
\end{verbatim}

\pandocbounded{\includegraphics[keepaspectratio]{images/_page_350_Figure_2.jpeg}}

\begin{verbatim}
u_{i, j} = \frac{u_{i, j+1} + u_{i, j-1} + \sin x_j (u_{i+1, j} + u_{i-1, j})}{2 (1 + \sin x_j)}.
u(1+\varepsilon\sin\theta,\ \theta)=u(1,\ \theta)+u_r(1,\ \theta)(\varepsilon\sin\theta)+u_{rr}(1,\ \theta)\frac{(\varepsilon\sin\theta)^2}{21}+\ldots
\end{verbatim}

(i, j+1)

\begin{verbatim}
\Delta u = 0, \quad 0 < r < 1 + \varepsilon \sin \theta,
u(A) = g_1 P_A(p_1) + g_2 P_A(p_2) + \dots + g_N P_A(p_N),
u(1, \theta) + u_r(1, \theta) (\varepsilon \sin \theta) + u_{rr}(1, \theta) \frac{(\varepsilon \sin \theta)^2}{2!} + \dots = \cos \theta.
u_{xx} + u_{yy} = 0
(46.8) u = u_0 + \varepsilon u_1 + \varepsilon^2 u_2 + \dots.
u_{xx} + x^2 u_{yy} = 0
P_{0}
(\text{UBP}) \qquad
\begin{cases} \Delta u_{0} = 0, & 0 < r < 1 \text{ (внутри круга),} \\ u_{0}(1, \theta) = \cos \theta, & u_{0}(r, \theta) = r \cos \theta, \end{cases}
\end{verbatim}

\begin{enumerate}
\def\labelenumi{\arabic{enumi}.}
\setcounter{enumi}{4}
\tightlist
\item
  Разработайте схему метода Монте-Карло для розв'язки смешанной
  параболической задачи
\end{enumerate}

\begin{verbatim}
\begin{cases} u(x, \ 0) = \varphi(x), \\ u_t(x, \ 0) = 0 \end{cases}
\end{verbatim}

\begin{verbatim}
\right. \end{array}
\end{verbatim}

\begin{verbatim}
P_{1}
J[y] = \int_{b}^{a} F(x, y, y') dx,
\begin{cases} \Delta u_{1} = 0, & 0 < r < 1 \text{ (внутри круга),} \\ u_{1}(1, \theta) = -\sin \theta \frac{\partial u_{0}(1, \theta)}{\partial r} = -\sin \theta \cos \theta \end{cases}
J[y] = \int_{0}^{1} [y^{2}(x) + y'^{2}(x)] dx.
\frac{df(x)}{dx} = 0
T = \int_{0}^{T} dt = \int_{a}^{L} \frac{dt}{ds} ds = \int_{0}^{L} \frac{ds}{v} = \frac{1}{\sqrt{2mg}} \int_{0}^{L} \frac{ds}{\sqrt{y}} = \frac{1}{\sqrt{2mg}} \int_{a}^{b} \sqrt{\frac{1+y'^{2}}{y}} dy
u = u_0 + \frac{1}{4}u_1 + \left(\frac{1}{4}\right)^2 u_1 + \cdots
J[y] = \int_{a}^{b} F(x, y, y') dx.
u_u + \frac{1}{4} u_t
\frac{df(x)}{dx} = 0
\end{verbatim}

\begin{enumerate}
\def\labelenumi{\arabic{enumi}.}
\setcounter{enumi}{1}
\tightlist
\item
  Результаты задачи 1 используйте для решения следующей задачи:
\end{enumerate}

(УЧП)

\begin{verbatim}
J[y] = \int_{a}^{y} F(x, y, y') dx
u_t = (1+x) u_{xx}, \\ u(x, 0) = q(x), \quad -\infty < x < \infty,
J[y] = \int_{a}^{b} F(x, y, y') dx
\end{verbatim}

(НУ)

\begin{verbatim}
y(a) = A
u_1 = (1 + \varepsilon_1) u_{xx}
J[\overline{y}] \leqslant J[\overline{y} + \varepsilon \eta]
u = u_1 + \varepsilon u_1 + \varepsilon^2 u_2 + \cdots
\varphi(\varepsilon) = J[\bar{y} + \varepsilon \eta]
\frac{d\varphi\left(\varepsilon\right)}{d\varepsilon} = \frac{d}{d\varepsilon} J\left[\overline{y} + \varepsilon\eta\right]|_{\varepsilon=0} =
\end{verbatim}

(Читач має робити це самостійно.) Інтегруючи на частини, ми отримуємо

\begin{verbatim}
P_{0}
\end{verbatim}

\begin{verbatim}
Оскільки цей інтеграл є нульовим у будь-якій функції  \eta(x)  , що задовольняє граничні умови  \eta(a) = \eta(b) = 0  , ми Р и с
\end{verbatim}

\pandocbounded{\includegraphics[keepaspectratio]{images/_page_330_Figure_6.jpeg}}

\begin{verbatim}
. 44 3. Граф функції  J[y+] \in \eta  ] в околиці  \varepsilon = 0  .
\end{verbatim}

Ми приходимо до висновку, що друга частина підоб'єктивного виразу має
дорівнювати нулю, тобто

(44 2)

\texttt{,\ \textquotesingle{}(x,\ y,\ z)\ \textbackslash{}in\ R\^{}3\textquotesingle{}\ ,\ \textquotesingle{}u\ =\ 0\textquotesingle{}\ u\_t\ =}

(Рівняння Ейлера--Лагранжа).

\begin{verbatim}
Рівняння (44 2) називається рівнянням Ейлера-Лагранжа, і хоча загалом воно здається комплексним, коли в нього підставляється конкретна функція F(x, y, y'), воно стає звичайним диференціальним рівнянням другого порядку відносно невідомої функції  \overline{y}(x)  . Отже, щоб визначити мінімізуючу функцію  \overline{y}  , необхідно розв'язати рівняння Ейлера-Лагранжа.
\end{verbatim}

Отже, ми показали, що

\begin{verbatim}
якщо функція y(x) мінімізує функціональний  J[y] = \int_a^b F(x, y, y') dx  (у класі гладких функцій з граничними умовами y(a) = A; y(b) = B), то вона повинна задовольняти рівняння
\end{verbatim}

\begin{verbatim}
g(x) = e^{-x^2}.
\end{verbatim}

(з граничною умовою).

\begin{verbatim}
(Щоб спростити позначення, ми прибрали лінію над  \bar{y}  .) Щоб проілюструвати теорію, наведемо приклад.
\end{verbatim}

\begin{verbatim}
Пошук мінімальної функціональності  J[y] = \int_{0}^{1} [y^2 + y'^2] dx  .
\end{verbatim}

Спробуємо знайти функцію \texttt{y(x)}, яка проходить через точки (0,0)
і (1,1) і мінімізує функціональність

\begin{enumerate}
\def\labelenumi{\arabic{enumi}.}
\setcounter{enumi}{3}
\tightlist
\item
  Решите аналогичную (см. задачу 3) задачу для двумерного уравнения
\end{enumerate}

(УЧП)

\begin{verbatim}
З контексту зрозуміло, що шукана функція має бути диференційовною, оскільки інтегративний вираз залежить від y'. Щоб визначити  \bar{y}(x)  , запишемо рівняння Ейлера–Лагранжа (разом із граничною умовою y(0)=0 та y(1)=1).
\end{verbatim}

Зверніть увагу, що в перших двох задачах коефіцієнти є сталими. Читач
повинен розуміти, що параметр b має бути достатньо малим, а необережний
ряд може відрізнятися

\section{ЗАДАЧИ}\label{ux437ux430ux434ux430ux447ux438-10}

\begin{itemize}
\tightlist
\item
  \begin{enumerate}
  \def\labelenumi{\arabic{enumi}.}
  \tightlist
  \item
    Подставьте разложение (46.8) в задачу (46.7) и получите
    последовательность задач
    \texttt{P\_0,\textbackslash{}\ P\_1,\textbackslash{}\ P\_2,\textbackslash{}\ \textbackslash{}dots}
    . 2. Покажите, что нелинейную задачу
  \end{enumerate}
\end{itemize}

\texttt{y(0)} = 0, \texttt{y(1)} = 1. Але F(x, y, y') = y\^{}2 + y'\^{}2
і, отже,

(НУ)

(Продифференцировали F по y и y'. )

Уравнение Эйлера --- Лагранжа принимает вид

\begin{verbatim}
u(r, \theta) = u_0(r, \theta) + \frac{1}{4}u_1(r, \theta)
\end{verbatim}

откуда

\begin{verbatim}
\Delta u = 0,
\end{verbatim}

Решая це простое дифференциальное уравнение с граничными условиями
\texttt{y(0)}=0 и \texttt{y(1)}=1, получаем

\begin{verbatim}
q_{xx} + \varphi_{uu} = 0
\end{verbatim}

.

График цієї функції изображен на рис. 44.4. Ясно, що любая другая
гладкая функція, удовлетворяющая тем же граничным условиям, даст большее
значение функционала \texttt{J(y)}.

\pandocbounded{\includegraphics[keepaspectratio]{images/_page_332_Figure_5.jpeg}}

\begin{verbatim}
Рис. 44 4. Все допустимые гладкие кривые удовлетворяют граничным условиям y(0) = 0 и y(1) = 1: a—допустимая кривая;  \delta  —минимизирующая функція  y(x) = 0.42e^x - 0.42e^{-x}  .
\end{verbatim}

\section{ЗАМЕЧАНИЯ}\label{ux437ux430ux43cux435ux447ux430ux43dux438ux44f-5}

\begin{enumerate}
\def\labelenumi{\arabic{enumi}.}
\item
  Уравнение Эйлера---Лагранжа аналогично равенству нулю производной в
  дифференциальном исчислении. Якщо читатель помнит, то на самом деле
  равенство нулю производной не является достаточным условием
  экстремума. Наприклад, производная функції \texttt{f(x)} = x\^{}3 при
  x = 0 равна нулю, но ця точка не является нп точкой минимума, ни
  точкой максимума. Так же и с уравнением Эйлера---Лагранжа: оно
  является лише необходимым условием экстремума, но не достаточным.
  Значит, экстремум может достигаться на других функциях. Очень часто,
  однако, розв'язок уравнения Эйлера---Лагранжа по существу задачи дает
  локальный или даже глобальный минимум.
\item
  Якщо у (х) минимизирует функционал
\end{enumerate}

\begin{verbatim}
q_{uu} + q_{vv} = 0
\end{verbatim}

то площадь поверхности, образованной вращением кривой \texttt{y(x)}
вокруг оси x, буде минимальной (рис. 44.5). Решением уравнения Эйлера
--- Лагранжа в цьому случае буде

цепная линия г д

\pandocbounded{\includegraphics[keepaspectratio]{images/_page_353_Figure_9.jpeg}}

\begin{verbatim}
е константы  \alpha  и  \beta  определяются из условий y(a) = A и y(b) = B (т. е. цепная линия должна проходить через заданные точки). Решить уравнение Эйлера—Лагранжа в данном случае достаточно трудно, поэтому мы рекомендуем читателю лише проверить, що цепная линия ему удовлетворяет.
\end{verbatim}

\begin{enumerate}
\def\labelenumi{\arabic{enumi}.}
\setcounter{enumi}{2}
\tightlist
\item
  В цієї лекции мы нашли, що функція
\end{enumerate}

\begin{verbatim}
w = f(z)
\end{verbatim}

доставляет минимум функционалу П р

\pandocbounded{\includegraphics[keepaspectratio]{images/_page_355_Figure_5.jpeg}}

\begin{verbatim}
и цьому оказывается, що  J[\overline{y}] = 0.46  . Якщо подставить в функционал любую другую гладкую функцию, график которой проходит через точки (0, 0) и (1, 1), то значение функционала J[y] станет больше.
\end{verbatim}

\pandocbounded{\includegraphics[keepaspectratio]{images/_page_333_Picture_8.jpeg}}

Рис. 44 5. Минимальная поверхность вращения.

\begin{enumerate}
\def\labelenumi{\arabic{enumi}.}
\setcounter{enumi}{3}
\tightlist
\item
  Основные законы физики чаще всего формулируются на языке вариационных
  принципов, а не дифференциальных уравнений. В качестве примеров можна
  привести принцип Ферма (свет при распространении из одной точки в
  другую выбирает путь, которому соответствует наименьшее время
  распространения) или принцип Гамильтона (в консервативном поле частица
  движется так, що интеграл действия
\end{enumerate}

\pandocbounded{\includegraphics[keepaspectratio]{images/_page_355_Figure_7.jpeg}}

(кинетическая энергия--- потенциальная энергия) dt буде минимальным).
Таким образом, физические явления развиваются лише так, що эти
функционалы принимают минимальные значения.

Основные идеи вариационного исчисления можна распространить на
функционалы, зависящие от функції нескольких независимых змінних, вида

\begin{verbatim}
u + iv = (x + iy)^2 = x^2 - y^2 + 2ixy
\end{verbatim}

Соответствующее уравнение Эйлера --- Лагранжа для такого функционала
буде уже уравнением з частинними похідними

\begin{verbatim}
\left\{
\end{verbatim}

Такого рода функционалы мы рассмотрим в следующей лекции. Однако общая
философия метода буде иной. Некоторым новым способом мы найдем функцию
\texttt{u(x,\ y)}, яка минимизирует функционал, а поскольку она является
решением уравнения Эйлера --- Лагранжа, то тем самым мы найдем розв'язок
уравнения з частинними похідними. Другими словами, мы будем искать
розв'язок уравнения з частинними похідними методом минимизации
функционала. В цієї лекции все було наоборот: задачу минимизации
функционала, зависящего от функції одной переменной, мы сводили к
решению уравнения Эйлера --- Лагранжа. Методы розв'язки дифференциальных
уравнений путем минимизации соответствующих функционалов принято
называть прямыми методами вариационного исчисления. Методы минимизации
функционалов путем розв'язки соответствующих уравнений Эйлера ---
Лагранжа називаються непрямыми методами вариационного исчисления.
Следующая лекция посвящена хорошо известному прямому методу Ритца.

\section{ЗАДАЧИ}\label{ux437ux430ux434ux430ux447ux438-11}

\begin{enumerate}
\def\labelenumi{\arabic{enumi}.}
\tightlist
\item
  Среди всех кривых, удовлетворяющих условиям \texttt{y(0)} = 0, и
  \texttt{y(1)} = 1, найдите ту, яка минимизирует функционал
\end{enumerate}

Чи знаєте ви, як виглядає графік цієї функції з сегментом {[}0, 1{]}?
Подумай про це. Зверніть увагу, що синусоїдальне перетворення функції
\texttt{f(x)} визначається лише для додатних цілих чисел n.~Іншими
словами, скінченний синус і косинус перетворюють функції у чисельні
послідовності.

\pandocbounded{\includegraphics[keepaspectratio]{images/_page_185_Figure_3.jpeg}}

РИС. 25.1. Графики функции \texttt{f(x)} = 1 и ее преобразования.

\section{Свойства
преобразований}\label{ux441ux432ux43eux439ux441ux442ux432ux430-ux43fux440ux435ux43eux431ux440ux430ux437ux43eux432ux430ux43dux438ux439-1}

Прежде чем приступить к решению задач, мы должны получить некоторые
свойства этих преобразований.

Якщо you(x, t) --- функція
\texttt{\textbackslash{}partial\ \textbackslash{}theta\ yx} змінних, і
ми виконуємо перетворення змінної x, то

\begin{verbatim}
Как вы интерпретируете получелный результат? Чему равна величина  J[\overline{y}]  ? Каков смысл величины  J[\overline{y}]  ?
\end{verbatim}

\begin{verbatim}
2. Кинетическая энергия колеблющейся материальной точки определяется выражением  KE = \frac{1}{2} m\dot{y}^2  , где y = dy/dt, а потенциальная равна  PE = \frac{1}{2} ky^2  .
\end{verbatim}

\pandocbounded{\includegraphics[keepaspectratio]{images/_page_335_Picture_2.jpeg}}

Принцип Гамильтона утверждает, що материальная точка буде двигаться так,
що интеграл

(Зверніть увагу, що перетворення було виконано на змінній x, і отримана
послідовність залежить лише від часу t.)

А как быть с производными? Ниже приводятся несколько полезных формул

буде минимален. Якщо цей закоп справедлив, то получите дифференциальное
уравнение движения материальной точки.

\begin{verbatim}
3. Покажите, що в классе гладких функцій, удовлетворяющих граничным условиям y(0)=0 и  y(\pi/2)=1  , минимум функционала
\end{verbatim}

достигается на функції
\texttt{\textbackslash{}overline\{y\}(x)\ =\ \textbackslash{}sin\ x} .
Вычислите \texttt{J(\textbackslash{}sin\ x)} . 4. Исходя из функционала

получите уравнение Эйлера --- Лагранжа

\# ВАРИАЦИОННЫЕ МЕТОДЫ РОЗВ'ЯЗКИ УРАВНЕНИЙ З ЧАСТИННИМИ ПОХІДНИМИ

ЦЕЛЬ ЛЕКЦИИ: Показать, как можна решать дифференциальное уравнение,
рассматривая его как уравнение Эйлера---Лагранжа для некоторого
функционала и находя некоторым новым методом функцию, минимизирующую цей
функционал. Тогда минимизирующая функція буде решением уравнения з
частинними похідними. Задача, конечно, состоит в том, чтобы найти такой
функционал, для которого исходное уравнение является уравнением
Эйлера---Лагранжа. Хорошо известен результат (теорема о минимуме
энергии), який гласит, що нахождение розв'язки и некоторой эллиптической
краевой задачи

, в области D , u = 0 , на границе D ,

эквивалентно нахождению такой функції u (тоже обращающейся в нуль на
границе области D), яка минимизирует функционал потенинальной энергии

\texttt{То\ есть,\ \textquotesingle{}\textbackslash{}Delta\ u\ =\ f\textquotesingle{}\ —уравнение\ Эйлера—Лагранжа\ для\ J{[}u{]}.\ Функція,\ приближенно\ минимизирующая\ функционал\ J{[}u{]},\ ищется\ методом\ Ритца;\ тем\ самым\ мы\ получаем\ розв\textquotesingle{}язок\ (приближениое)\ уравнения\ з\ частинними\ похідними.\ В\ цієї\ лекции\ мы\ познакомимся\ с\ методом\ Ритца\ и\ посмотрим,\ как\ с\ помощью\ цього\ метода\ можна\ и\ и\ и\ и\ и\ и\ и\ уравнения.}

Очень хорошо решать граничные задачи (как, наприклад, задачу о
закрепленной мембране) путем отыскания гладкой поверхности, яка
минимизирует потенциальную энергию мембраны. То есть, якщо мы
рассматриваем исходное уравнение как уравнение Эйлера --- Лагранжа для
искоторого функционала J{[}u{]}, то розв'язок дифференциального
уравнения можна получить путем минимизации функционала (поскольку
минимизирующая функція функционала является решением соответствующего
уравнения

Эйлера --- Лагранжа). В лекции 44 мы рассматривали такие функционалы, у
которых уравнение Эйлера --- Лагранжа було обыкновенным дифференциальным
уравнением. В цієї лекции рассматриваются функционалы, для которых
уравнение Эйлера --- Лагранжа относится к классу уравнений з частинними
похідними. Наприклад для функционала

\begin{verbatim}
(\text{УЧП}) \quad u_{tt} = u_{xx} + \sin(\pi x), \quad 0 < x < 1, \quad 0 < t < \infty,
\end{verbatim}

уравнение Эйлера---Лагранжа (подобно случаю обыкновенного
дифференциального уравнения в последней лекции) имеет вид

\begin{verbatim}
(\text{ГУ}) \quad
\end{verbatim}

и, отже, для розв'язки задачи Дирихле в единичном квадрате

(УЧП)

, 0 \textless{} x \textless{} 1 , 0 \textless{} y \textless{} 1 , (ГУ) u
= g , на границе квадрата,

можна, наоборот, найти такую функцию \texttt{u(x,\ y)}, яка минимизирует
J{[}u{]} и равна g на границе. Вероятно, не стоит удивляться тому, що
функционал

представляет собой потенциальную эпергию мембраны и фактически мы
находим поверхность с минимальной потенциальной энерешей. Вопрос,
конечно, состоит в следующем: якщо дано дифференциальное уравнение, как
мы находим функционал, выражающий потенциальную энергию розв'язки? Ответ
на цей вопрос дает хорошо известная теорема (теорема о минимуме
энергии), яка утверждает:

Розв'язок и задачи Дирихле

\begin{verbatim}
^{\Lambda}u = f  в области  ^{D}   (ГУ)   ^{U}=0  на границе  ^{D}  является той же функциен, яка минимизирует (среди функцій, удовлетворяющих граничному условию u=0) энергегический функционал
\end{verbatim}

\begin{verbatim}
\begin{cases} u_{tx}, \quad 0 = \begin{cases} 1, \quad -1 < x < 1, \\ 0 \quad \text{в остальных точках,} \end{cases}
\end{verbatim}

\begin{verbatim}
и, отже, для розв'язки задачи (45.1) достаточно найти функцию  \overline{u}  (среди функцій, обращающихся в нуль на граннце), минимизирующую  J\{u\}  . Этим замечанием завершается первая часть лекции. В оставшейся части мы покажем, как находить минимизирующую функцию  \overline{u}  для J[u], используя метод Ритца.
\end{verbatim}

\section{Метод Ритца минимизации
функционалов}\label{ux43cux435ux442ux43eux434-ux440ux438ux442ux446ux430-ux43cux438ux43dux438ux43cux438ux437ux430ux446ux438ux438-ux444ux443ux43dux43aux446ux438ux43eux43dux430ux43bux43eux432}

Цей метод --- один из многих вариационных методов, описанных в
литературе по вариационному исчислению. Идея, предложенная математиком
В. Ритцем, очень проста. Метод состоит из следующих шагов.

ШАГ 1. Выбираем n и функцию u, минимизирующую функционал

ищем в виде

\begin{verbatim}
где функції  \phi_1, \phi_2, \phi_3, \ldots, \phi_n  принадлежат к классу достаточно хороших функцій и все обращаются в нуль на границе, так що из них можна построить разумное приближение для розв'язки задачи. Эти функції принято называть пробными функциями. Типичный набор пробных функцій для задачи Дирихле в единичном квадрате приведен ниже:
\end{verbatim}

\texttt{—\ обращаются\ в\ нуль\ на\ границе,\ \textquotesingle{}\textbackslash{}varphi\_2(x,\ y)\ =\ x\textbackslash{}varphi\_1(x,\ y),\textquotesingle{}\ \textquotesingle{}\textbackslash{}varphi\_3(x,\ y)\ =\ y\textbackslash{}varphi\_1(x,\ y),\textquotesingle{}\ \textquotesingle{}\textbackslash{}varphi\_4(x,\ y)\ =\ x\^{}2\textbackslash{}varphi\_1(x,\ y),\textquotesingle{}\ \textquotesingle{}\textbackslash{}varphi\_5(x,\ y)\ =\ xy\textbackslash{}varphi\_1(x,\ y),\textquotesingle{}\ \textquotesingle{}\textbackslash{}varphi\_6(x,\ y)\ =\ y\^{}2\textbackslash{}varphi\_1(x,\ y),\textquotesingle{}\ \textquotesingle{}\textbackslash{}vdots\textquotesingle{}\ \textquotesingle{}\textbackslash{}vdots\textquotesingle{}\ Иными\ словами,\ первые\ четыре\ приближения\ имеют\ вид}

Решение задачи выполним за три следующих шага:

ШАГ І (Определение подходящего преобразования.)

Поскольку переменная x меняется от 0 до 1, мы воспользуемся конечным
преобразованием, а именно синус-преобразованием; позже станет ясно
почему. Мы могли бы решить эту задачу с помощью преобразования Лапласа
по переменной t (по трудоемкости этот метод не отличается от конечного
синус-преобразования).

ШАГ 2 (Выполнение преобразования).

Для удобства будем использовать обозначение \texttt{S\_n}(t) = S{[}u{]}
. Применим синус-преобразование к исходному уравнению

\begin{verbatim}
\begin{aligned} & u_{ti} = c^{2}u_{xx}, & -\infty < x < +\infty, & 0 < t < \infty, \\ & (\text{НУ}) & \begin{cases} u_{t}(x, 0) = 0, \\ u_{t}(x, 0) = \begin{cases} 1, & -1 < x < 1, \\ 0, & \text{в остальных случаях.} \end{cases} \end{aligned}
\end{verbatim}

и воспользуемся тождеством для синус-преобразования, получим

в координатах (x, t).

\pandocbounded{\includegraphics[keepaspectratio]{images/_page_134_Figure_7.jpeg}}

Рис. 18.4. Розв'язок задачи (18.4) в xt-плоскости.

Задача (18.4) описывает поведение струны, которой сообщается рачальная
единичная скорость на отрезке -1 \textless{} x \textless{} 1. Смещение .
.

описывается формулой Даламбера

\begin{verbatim}
является теперь функцией коефіцієнтів  a_1a_2\dots a_n  . Отже, для того чтобы найти минимум функционала J, приравняем нулю частные производные
\end{verbatim}

\begin{verbatim}
Aa = b
b_i = -\int_0^1 \int_0^1 f(x, y) \varphi_i(x, y) dx dy,
\end{verbatim}

где

Це розв'язок в различные моменты времени изображено на рис. 18.5.

\pandocbounded{\includegraphics[keepaspectratio]{images/_page_135_Figure_5.jpeg}}

Рис. 18.5. Розв'язок задачи (18.4) в различные моменты времени.

\begin{verbatim}
Этим завершается наша интерпретация формулы Даламбера в плоскости змінних  \boldsymbol{x}  и t. В оставшейся части лекции мы займемся решением смешанной задачи для полубесконечной струны:
\end{verbatim}

(18.6)

Единственное отличие в том, що аппроксимирующие функції
\texttt{\textbackslash{}phi\_1,\ \textbackslash{}phi\_2,\ \textbackslash{}phi\_3,\ \textbackslash{}ldots,\ \textbackslash{}phi\_n}
повинні удовлетворять граничным условиям:

является иравнение Лапласа, но це доказательство проводится аналогично
тому, как це делалось в лекции 44 для функционала, зависящего от функції
одной переменной.

\begin{enumerate}
\def\labelenumi{\arabic{enumi}.}
\setcounter{enumi}{2}
\tightlist
\item
  В книгах по вариационному исчислению показано, как строить
  энергетический функционал не лише для уравнения Пуассона с граничным
  условием Дирихле u=0, но и для многих других дифференциальных
  уравнений и многих других типов граничных условий. Отже, большое число
  краевых задач можна решить, минимизируя соответствующие энергетические
  функционалы.
\end{enumerate}

\begin{verbatim}
4. Чем большее значение п мы выбираем, тем меньшее значение  J[\bar{u}_n]  получается. Значит, с увеличением n растет точность л приближенного розв'язки дифференциального уравнения. Один из методов определения и, основан на последовательном вычислении функционала  J[u_n]  для возрастающих значений n. Обчислення прекращаются, якщо с ростом п функционал перестает практически уменьшаться.
\end{verbatim}

\begin{enumerate}
\def\labelenumi{\arabic{enumi}.}
\setcounter{enumi}{4}
\tightlist
\item
  Якщо и велико, то проведение вычислений по методу Ритца требует
  привлечения ЭВМ. На рис. 45.2 (с. 345) приведена блок-
\end{enumerate}

схема розв'язки краевой задачи методом Ритца.

Пользователь повинен позаботиться о подпрограмме, яка должна вычислять
значения пробных функцій. Такая подпрограмма может иметь вид:

SUBROUTINE BC (X.Y.PHI)

\begin{itemize}
\tightlist
\item
  SUBROUTINE PROVIDED BY THE USER TO C - EVALUATE THE - FUNCTIONS
  PHI(1), PHI(2), \ldots, PHI(N) DIMENSION PHI(20)
\end{itemize}

PHI(1) = X*Y*(1-X)*(1-Y) PHI(2) = X\^{}*PHI(1) (используемые функції)

\begin{verbatim}
PHI(3) = Y\*PHI(1)
\end{verbatim}

PHI(N) = (що получится)RETURN END

\section{\texorpdfstring{\textbf{ЗА}ДАЧИ}{ЗАДАЧИ}}\label{ux437ux430ux434ux430ux447ux438-12}

Найдите энергетический функционал для задачи

(УЧП)

Если теперь преобразовать начальные условия краевой задачи, то получим
начальные условия для обыкновенного дифференциального уравнения

, 0 \textless{} x \textless{} 1 , 0 \textless{} y \textless{} 1 , (ГУ) u
= 0 на границе.

\begin{enumerate}
\def\labelenumi{\arabic{enumi}.}
\setcounter{enumi}{1}
\tightlist
\item
  Как минимизировать функционал
\end{enumerate}

методом Ритца?

УКАЗАНИЕ. Якщо ввести новую функцию \texttt{z(x)}:

то можна заметить, що она удовлетворяет граничным условиям \texttt{z(0)}
= 0 и \texttt{z(1)} = 0.1

\begin{itemize}
\tightlist
\item
  \begin{enumerate}
  \def\labelenumi{\arabic{enumi}.}
  \setcounter{enumi}{2}
  \tightlist
  \item
    Напишите программу для проведения вычислений по блоксхеме,
    изображенной на рис. 45 2.
  \end{enumerate}
\item
  \begin{enumerate}
  \def\labelenumi{\arabic{enumi}.}
  \setcounter{enumi}{3}
  \tightlist
  \item
    Покажите, що уравнение
  \end{enumerate}
\end{itemize}

является уравнением Эйлера --- Лагранжа для функционала

\begin{enumerate}
\def\labelenumi{\arabic{enumi}.}
\setcounter{enumi}{4}
\tightlist
\item
  Задачу Дирихле
\end{enumerate}

Итак, решим теперь новое семейство задач Кощи:

(ОДУ)

\begin{verbatim}
где A=0.06 и B=0.04. Как найти потенциальную энергию цього розв'язки? Мы рекомендуем читателю вспомнить синуспреобразование и самостоятельно получить розв'язок  u\left(x,y\right)  .
\end{verbatim}

*) При такой замене возникает особенность в т. x=1. Лучше ввести функцию
\texttt{z(x)} по формуле \texttt{z(x)}=\texttt{y(x)}=x. Прим. ред.

\pandocbounded{\includegraphics[keepaspectratio]{images/_page_343_Figure_2.jpeg}}

\begin{verbatim}
Далсе. Приближенное розв'язок  u_n(x,y) = a_1 \varphi_1 + a_2 \varphi_2 + \dots + a_n \varphi_n  .
\end{verbatim}

В це время можна численно найти эначение функционала
J{[}\texttt{u\_n}{]} . Якщо необходимо найти значение розв'язки
\texttt{u\_n}(x,y) в некоторой точке, можна написать небольшую программу
эля быполнения отой операции. При необходимости можна повторить
обчислення по цієї программе с другими значениями п

РИС 452. Блок-схема метода Ритца.

\section{РОЗВ'ЯЗОК УРАВНЕНИЙ З ЧАСТИННИМИ ПОХІДНИМИ МЕТОДАМИ ТЕОРИИ
ВОЗМУЩЕНИЙ}\label{ux440ux43eux437ux432ux44fux437ux43eux43a-ux443ux440ux430ux432ux43dux435ux43dux438ux439-ux437-ux447ux430ux441ux442ux438ux43dux43dux438ux43cux438-ux43fux43eux445ux456ux434ux43dux438ux43cux438-ux43cux435ux442ux43eux434ux430ux43cux438-ux442ux435ux43eux440ux438ux438-ux432ux43eux437ux43cux443ux449ux435ux43dux438ux439}

ЦЕЛЬ ЛЕКЦИИ: Показать, как различные сложные задачи (нелинейные
уравнения с змінними коэффициентами, области неправильной формы и т. д.)
можна решить методом возмущения более простых задач. Иными словами, мы
хотим показать, как следует изменить розв'язок простой задачи, чтобы оно
давало приближенное розв'язок трудной задачи.

Наприклад, мы покажем, как нелинейную задачу

(УЧП)

, \texttt{0\ \textless{}\ r\ \textless{}\ 1} ,
\texttt{(ГУ)\ \textquotesingle{}u(1,\ \textbackslash{}theta)\ =\ \textbackslash{}cos\ \textbackslash{}theta\textquotesingle{}\ ,\ \textquotesingle{}0\ \textbackslash{}le\ \textbackslash{}theta\ \textbackslash{}le\ 2\textbackslash{}pi\textquotesingle{}\ .}

Очень часто задачу, розв'язок которой не известно, можна путем
непрерывного изменения каких-то параметров свести к близкой, но легко
решаемой задаче. Тогда и розв'язок исходной задачи буде близко решению
модифицирующей задачи.

Наприклад, рівняння Лапласа

в помощью семейства уравнений

\texttt{,\ \textquotesingle{}0\ \textbackslash{}le\ \textbackslash{}varepsilon\ \textbackslash{}le\ 1\textquotesingle{}\ ,}

\begin{verbatim}
,  0 < r < 1  ,  0 \le \theta < 2\pi  ,  u(1, \theta) = g(\theta)  ,  0 \le \theta < 2\pi  ,
\end{verbatim}

\begin{verbatim}
при одних граничных условиях (наприклад,  g(\theta) = \cos \theta  ) имеет два различных вещественных розв'язки, а при других (наприклад,  g(\theta) = A \cos \theta  с A > 20,65) ця задача не имеет ни одного вещественного розв'язки. — Прим. ред.
\end{verbatim}

4) Краевая задача

\begin{verbatim}
можна непрерывно модифицировать в нелинейное уравнение  \Delta u = u^2 = 0  1).
\end{verbatim}

На рис. 46.1 изображена схема розв'язки нелинейной задачи путем сведения
ее к решению уравнения Лапласа2)

\begin{verbatim}
\right.
\end{verbatim}

Рис. 46.1. Схема метода возмущении.

Чтобы найти розв'язок возмущенного уравнения Лапласа u

\pandocbounded{\includegraphics[keepaspectratio]{images/_page_362_Figure_1.jpeg}}

\begin{verbatim}
(0, t) = f(t),
u_{x}(0, t) = 0.
x + ct = \text{const},
\end{verbatim}

\pandocbounded{\includegraphics[keepaspectratio]{images/_page_362_Figure_2.jpeg}}

\begin{verbatim}
0 < r < 1,   0 \le \theta < 2\pi,   u(1, 0) = 0,   0 \le \theta < 2\pi,  нмеет два розв'язки (одно — нулсьое, другое — нетривиальное), в то время как задача u
\end{verbatim}

\pandocbounded{\includegraphics[keepaspectratio]{images/_page_362_Figure_5.jpeg}}

\begin{verbatim}
(1, \theta) = 0, \quad 0 \le \theta < 2\pi,
\end{verbatim}

\pandocbounded{\includegraphics[keepaspectratio]{images/_page_362_Figure_7.jpeg}}

\begin{verbatim}
\begin{array}{ccccccccccccccccccccccccccccccccccc
\end{verbatim}

\begin{verbatim}
\begin{array}{ccc} (\text{УЧ}\Pi) & \Delta u = 0, & 0 < r < 1, \\ (\text{ГУ}) & u(1, \theta) = \cos \theta, & 0 \leqslant \theta \leqslant 2\pi, \end{array}
u_{tt} = c^2 u_{xx}, \quad 0 < x < L, \quad 0 < t < \infty,
H(x-a) =
u(0, t) = g_1(t),
\begin{cases} 0, x < a, \\ 1, x \geqslant a, \end{cases}
u_x(0, t) = g_1(t),
\end{verbatim}

В результате получаем

\begin{verbatim}
u_x(0, t) - \gamma_1 u(0, t) = g_1(t),
H(x-a) =
u(0, t) = g_1(t),
\begin{cases} 0, & x < a, \\ 1, & x \geqslant a, \end{cases}
u_x(0, t) = g_1(t),
\begin{cases} 0, & n = 2, 4, \dots, \\ (4/n\pi)\cos(n\pi t), & n = 3, 5, 7, \dots \end{cases}
u_x(0, t) - \gamma_1 u(0, t) = g_1(t),
f(x) = \mathcal{F}^{-1}[F] = \frac{1}{\sqrt{2\pi}} \int_{-\infty}^{+\infty} F(\omega) e^{-i\omega x} d\omega \qquad F(\omega) = \mathcal{F}[f] = \frac{1}{\sqrt{2\pi}} \int_{-\infty}^{+\infty} f(x) e^{-i\omega x} dx
\begin{cases} \Delta u_{1} = -r^{2} \cos^{2} \theta = -\frac{r^{2}}{2} \left[1 + \cos(2\theta)\right] = -\frac{r^{2}}{2} - \frac{r^{2}}{2} \cos(2\theta), \\ u_{1}(1, \theta) = 0. \end{cases}
\end{verbatim}

Следовательно, решение нашей задачи можно представить в виде

\begin{verbatim}
\Delta u = -\frac{r^2}{2} - \frac{r^3}{2} \cos{(2\theta)}
\begin{cases} 1, |x| < a \\ 0, |x| > a \end{cases}
u_p(r, \theta) = Ar^4 + Br^4 \cos(2\theta).
\end{verbatim}

\section{ЗАМЕЧ
АНИЯ}\label{ux437ux430ux43cux435ux447-ux430ux43dux438ux44f-1}

\begin{enumerate}
\def\labelenumi{\arabic{enumi}.}
\tightlist
\item
  Для решения задачи методом конечного синус- или косинуспреобразования
  краевые условия при x=0 и x=L должны иметь вид
\end{enumerate}

\begin{verbatim}
u_p(r, \theta) = -\frac{r^4}{32} - \frac{r^4}{24} \cos(2\theta).
\begin{cases} 1, |x| < 1 \\ 0, |x| > 1 \end{cases}
\sum_{n=0}^{\infty} [a_n \cos(n\theta) + b_n \sin(n\theta)] - \frac{1}{32} - \frac{1}{24} \cos(2\theta) = 0.
\end{verbatim}

Отметим, що якщо величина \texttt{u(0,\ t)} положительна, то величина
\texttt{u\_x(0,\ t)} також положительна, якщо \texttt{u(L,\ t)}
положительна, то \texttt{u\_x(L,\ t)}

\pandocbounded{\includegraphics[keepaspectratio]{images/_page_144_Picture_4.jpeg}}

Рис. 19.6. Возникновение неоднородных граничных условий упругого
закрепления: h---коефіцієнт упругости пружины.

отрицательна. Однородные граничні умови (19.2) можна записать в виде

(19.3)

\begin{verbatim}
P_{2}
\delta(x-a)
\Delta u = 0, \quad 0 < r < 1 + \frac{1}{4} \sin \theta,
\end{verbatim}

\begin{verbatim}
\begin{array}{ll} (\mathrm{Y} \mathrm{U} \Pi) & u_t = u_{xx}, \quad 0 < x < 1, \quad 0 < t < \infty, \\ (\mathrm{\Gamma} \mathrm{Y} & \left\{ \begin{array}{ll} u_x \left( 0, \ t \right) = 0 \\ u_x \left( 1, \ t \right) = 0, \end{array}
\end{verbatim}

досягається в кінцевій точці сімейства граничних умов виду (див. рисунок
46 3)

\begin{verbatim}
(1+x^2)^{-1}
Tu_x(0, t) = -\beta u_t(0, t).
xe^{-a+x^4}, a > 0
Tu_{x}(0, t) = \varphi[u(0, t)],
\end{verbatim}

\texttt{H(x+a)} - \texttt{H(x-a)}

\begin{verbatim}
Tu_{x}(0, t) = -hu^{3}(0, t)
\frac{a}{x^2+a^2}
mu_{tt}(L, t) = -ku_x(L, t) + mg.
\end{verbatim}

\begin{verbatim}
\right. & 0 < t < \infty, \\ (\text{НУ}) & u(x, 0) = 0, \quad 0 \leqslant x \leqslant 1. \end{array}
\end{verbatim}

\begin{verbatim}
P_{0}
\frac{2ax}{(x^2+a^2)^2}
\end{verbatim}

Отже, розв'язуючи кожну з задач Дірікла всередині кола для функцій
\texttt{u\_0} , \texttt{u\_1} , \texttt{u\_2} , отримуємо розв'язок

(46.9)

\begin{verbatim}
\begin{cases} \cos(ax), |x| > \pi/2a \\ 0, |x| > \pi/2a \end{cases}
\end{verbatim}

Задачі Діріхле в деформованій області (зверніть увагу, що тут

\begin{verbatim}
Ми вже взяли  \epsilon = 1/4  :)
\end{verbatim}

Як вправу, читача просять знайти пертурбацію першого порядку
\texttt{u\_1} і перевірити, наскільки добре є приблизне розв'язання я к

\pandocbounded{\includegraphics[keepaspectratio]{images/_page_189_Picture_10.jpeg}}

що розглянути рівняння з параметром

(46.11)

\begin{enumerate}
\def\labelenumi{\arabic{enumi}.}
\setcounter{enumi}{17}
\tightlist
\item
\end{enumerate}

і побачити його розв'язок у вигляді

\begin{verbatim}
\begin{cases} 1-|x|, |x| < 1 \\ 0, |x| > 1 \end{cases}
\end{verbatim}

Підставляючи цю декомпозицію на (46.11), отримуємо таку послідовність
завдань

\begin{verbatim}
u_{x}(0, t) = \frac{h}{T} [u(0, t) - \theta_{1}(t)]
u_{tt} = u_{xx}, \quad 0 < x < 1, \quad 0 < t < \infty,
\cos(ax)
\Delta u = 0,
\frac{2}{\pi} \frac{\sin a\omega}{\omega}
q_{uu} + q_{vv} = 0
\end{verbatim}

Перед нами неоднородное уравнение теплопроводности, так что метод
разделения переменных неприменим. Можпо, конечно, воспользоваться
конечным синус-преобразованием по переменной x или преобразованием
Лапласа по переменной t, но вместо этого мы рассмотрим две подзадачи

\begin{verbatim}
w = f(z)
\frac{2}{\pi} \frac{\sin a\omega}{(\omega^2+a^2)^2}
\end{verbatim}

\begin{verbatim}
\left\{
\end{verbatim}

\begin{verbatim}
\frac{\pi}{2} e^{-a+\omega}
\Rightarrow.
\end{verbatim}

\begin{verbatim}
\begin{aligned} (Y & \Pi) & \varphi_{xx} + \varphi_{yy} = 0, & -\infty < x < \infty, & 0 < y < \infty, \\ (47.1) & \varphi(x, 0) = \begin{cases} 0, & |x| > 1, \\ 1, & |x| \le 1. \end{cases}
\end{verbatim}

\begin{verbatim}
w = \ln\left\{\frac{z-1}{z+1}\right\}
\frac{\pi}{2} e^{-a+\omega}
\begin{cases} (\text{УЧП}) & \varphi_{uu} + \varphi_{vv} = 0, \quad -\infty < u < \infty, \quad 0 < v < \pi, \\ (\text{47.2}) & \begin{cases} \varphi(u, 0) = 0, \\ \varphi(u, \pi) = 1. \end{cases} \end{cases}
\frac{\pi}{2} e^{-a+\omega}
u(x, t) = \sum_{n=1}^{\infty} c_n X_n(x) T_n(t),
\frac{\pi}{2} e^{-a+\omega}
u_{tt} = \alpha^2 u_{xx}, \quad 0 < x < L, \quad 0 < t < \infty,
\frac{\pi}{2} e^{-a+\omega}
(\Gamma Y)
\end{verbatim}

Каждая из этих задач легко решается и в то же время ясно, что
\texttt{\textbackslash{}mathit\{suma\}} решений задач \texttt{P\_1} и
\texttt{P\_2} является решением исходной вадачи, т. е.

Щоб візуально зрозуміти, як працює ця функція, ми радимо читачеві
побудувати її за кількома лініями y=c.

У другому прикладі ми перетворюємо область між двома неконцентричними
колами у кільце.

Задача Діріхле в області між двома некопоцентричними колами

Припустимо, ви хочете знайти потенціал між двома колами

\begin{verbatim}
\frac{2ax}{(x^2+a^2)^2}
\quad 0 < t < \infty,
\end{verbatim}

Решение задачи \texttt{P\_\{i\}} Розв'язання задачі \texttt{P\_i}
Розв'язання задачі \texttt{P\_i} Можна показати, що в загальному випадку
розв'язок задачі

(HY)

\begin{verbatim}
\frac{\cos(ax), |x| > \pi/2a}{0, |x| > \pi/2a}
u = \gamma [(x-s)(x-t)-\gamma^2],
\frac{1-|x|, |x| < 1}{0, |x| > 1}
\quad 0 \le x \le L,
u(x, t) = X(x) T(t).
T'' - \alpha^2 \lambda T = 0,
\cos(ax)
X'' - \lambda X = 0,
\frac{\pi}{2} [\delta(\omega+a) + \delta(\omega-a)]
u(x, t) = [C \sin(\beta x) + D \cos(\beta x)][A \sin(\alpha \beta t) + B \cos(\alpha \beta t)]
\beta_n = \frac{n\pi}{L}, \quad n = 0, 1, 2, \ldots
u_n(x, t) = X_n(x) T_n(t) = \sin(n\pi x/L) [a_n \sin(n\pi at/L) + b_n \cos(n\pi at/L)],
f(x) = \int_{0}^{\infty} F(\omega) \sin(\omega x) d\omega
u_n(x, t) = R_n \sin(n\pi x/L) \cos(n\pi\alpha (t - \delta_n)/L],
F(\omega) = \frac{?}{\pi} \int_{0}^{\infty} f(x) \sin(\omega x) dx
u(x, t) = \sum_{n=1}^{\infty} \sin(n\pi x/L) \left[ a_n \sin(n\pi \alpha t/L) + b_n \cos(n\pi \alpha t/L) \right]
0 < x < \infty
u(x, \theta) = f(x),
\sum_{n=1}^{\infty} b_n \sin(n\pi x/L) = f(x),
\sum_{n=1}^{\infty} a_n (n\pi \alpha/L) \sin(n\pi x/L) = g(x).
0 < \omega < x
\int_{0}^{L} \sin(m\pi x/L) \sin(n\pi x/L) dx =
-\omega^2 F(\omega) + \frac{2}{\pi} \omega f(0)
\begin{cases} 1, |x| < 1 \\ 0, |x| > 1 \end{cases}
\frac{1}{\alpha} F\left(\frac{\omega}{\alpha}\right)
\end{verbatim}

\texttt{f(x)} * \texttt{g(x)}

проследим, как оно решается методом конечного синус-преобразования. Вот
что мы должны сделать: разложить еходное воздействие \texttt{f(x,\ t)}
на компоненты, найти отклик \texttt{U\_n} на каждую компоненту, а затем
сложить эти отклики. Математически это может быть и не совсем очевидно,
но давайте внимательно проследим за ходом решения. Разложим входящие в
уравнение с частными производными члены в ряды Фурье по синусам:

\begin{verbatim}
(1+x^2)^{-1}
\frac{2\omega}{\pi (\alpha^2 + \omega^2)}
xe^{-a+x^4}, a > 0
x^{-1/2}
\end{verbatim}

\texttt{H(x+a)} - \texttt{H(x-a)}

\begin{verbatim}
[2/\pi\omega]^{1/2}
\frac{a}{x^2+a^2}
\end{verbatim}

where

\begin{verbatim}
\frac{2ax}{(x^2+a^2)^2}
\frac{2}{\pi\omega} [1-\cos(\omega a)]
\begin{cases} \cos(ax), |x| > \pi/2a \\ 0, |x| > \pi/2a \end{cases}
\end{verbatim}

1

\begin{verbatim}
\begin{cases} 1-|x|, |x| < 1 \\ 0, |x| > 1 \end{cases}
e^{-1}a + \omega
\cos(ax)
\end{verbatim}

Зверніть увагу, по-перше, що коефіцієнти \texttt{A\_n} , \texttt{B\_n}
\texttt{F\_n} дійсно є функціями часу t, оскільки ми застосували
синусоїдальне перетворення до функції двох змінних x, t, u, а по-друге,
ми розклали вхідний ефект \texttt{f(x,t)} на прості компоненти
\texttt{F\_n(t)} . Тепер знайдіть відповіді \texttt{U\_n(t)} для кожного
компонента \texttt{F\_n(t)} , потім додайте всі відповіді
\texttt{U\_n(t)} і отримайте розв'язок \texttt{u(x,t)}.

Щоб знайти відповіді \texttt{U\_n(t)} , необхідно робити невеликі
розрахунки з коефіцієнтами \texttt{A\_n(t)} і \texttt{B\_n(t)} , щоб у
вирази інтегрування була включена лише функція u, а не її похідні
\texttt{u\_i} і \texttt{u\_\{xx\}} . Інтегруючи на частини, ми отримуємо

\begin{verbatim}
\frac{2}{\pi} \frac{\sin a\omega}{\omega}
\frac{1}{2}e^{-\alpha}\sin\omega
\frac{2}{\pi} \frac{\sin a\omega}{(\omega^2+a^2)^2}
\frac{1-e^{-a\omega}}{\omega}
\frac{\pi}{2} e^{-a+\omega}
\frac{\pi}{2} e^{-a+\omega}
\frac{\pi}{2} e^{-a+\omega}
\frac{2}{\pi}F''(\omega)
\frac{\pi}{2} e^{-a+\omega}
f(x) = \int_{0}^{\infty} F(\omega) \cos(\omega x) d\omega \qquad \qquad F(x) = \frac{2}{\pi} \int_{0}^{\infty} f(x) \cos(\omega x) dx
\frac{\pi}{2} e^{-a+\omega}
0 < x < \infty \qquad \qquad 0 < \omega < \infty
\frac{\cos(ax), |x| > \pi/2a}{0, |x| > \pi/2a}
\frac{1-|x|, |x| < 1}{0, |x| > 1}
\cos(ax)
\end{verbatim}

j'\,'(x)

\begin{verbatim}
u_{tt} = -\alpha^2 u_{xxxx},
S_n = \frac{2}{\pi} \int_0^{\pi} f(x) \sin(nx) dx.
u(0, t) = 0,
\end{verbatim}

в последнее соотношение, получаем

\begin{verbatim}
u_{xx}(0, t) = 0,
y = \frac{\pi}{b - a} (x - a).
-\omega^2 F(\omega) + \frac{2}{\pi} \omega f(0)
\end{verbatim}

и, следовательно, (26.1) принимает вид

\begin{verbatim}
\frac{2\omega}{\pi (\alpha^2 + \omega^2)}
f(x) = \sum_{n=1}^{\infty} S_n \sin(nx).
u(x, t) = X(x) [A \sin(\omega t) + B \cos(\omega t)].
\end{verbatim}

Поскольку это разложение является тождеством по x, его коэффициенты
должны быть равны нулю, \texttt{\textbackslash{}tau} . е.

\begin{verbatim}
X^{(1V)} - \omega^2 X = 0,
f(x) = \sum_{n=1}^{\infty} S_n \sin(nx)
X(x) = C\cos\sqrt{\omega}x + D\sin\sqrt{\omega}x + E \cosh\sqrt{\omega}x + F \sinh\sqrt{\omega}x.
\end{verbatim}

.

Таким образом, мы получили связь между входом и выходом системы, между
коэффициентами \texttt{F\_n} и \texttt{U\_n} . Прежде чем решать
уравнения для \texttt{U\_n(t)} , обратимся к начальному условию

\begin{verbatim}
8.  \frac{x}{x^4+4} 
\end{verbatim}

\begin{verbatim}
S_n = \frac{2}{\pi} \int_{0}^{\pi} f(x) \sin(nx) dx
\end{verbatim}

\begin{verbatim}
10.  -x^2/(x) 
\end{verbatim}

Если разложить \texttt{u(x,\ 0)} в ряд Фурье по синусам и приравнять его
нулю:

\section{ТАБЛИЦЯ C. Перетворення
Фур'є}\label{ux442ux430ux431ux43bux438ux446ux44f-c.-ux43fux435ux440ux435ux442ux432ux43eux440ux435ux43dux43dux44f-ux444ux443ux440ux454}

\texttt{то\ определяются\ начальные\ условия\ для\ функции\ \textquotesingle{}U\_n\textbackslash{}left(t\textbackslash{}right)\textquotesingle{}\ ,\ т.\ е.}

Решая их, находим

\begin{verbatim}
n = 1, 2, ...
\end{verbatim}

\begin{verbatim}
,  \sin \sqrt{\omega} = 0 \Rightarrow \omega = (n\pi)^2  ,  n = 1, 2, \dots  Другими словами, *собственные частоты* свободно опирающейся балки равны
\end{verbatim}

, \texttt{n\ =\ 1,\ 2,\ \textbackslash{}dots} Мы получили разложение
нашей смешанной задачи на ряд простых задач об отклике системы на
элементарные входные воздействия:

(OДУ)

а фундаментальные розв'язки и, (удовлетворяют уравнению и граничным
условиям (21.2)) записываются в виде

(HУ) \texttt{U\_n(0)\ =\ 0,\ n\ =\ 1,\ 2,\ ...} Каждое из этих уравнений
легко решается методом интегрирующего множителя (или с помощью
преобразования Лапласа). В любом случае получаем

А теперь, поскольку уравнение и граничні умови линейны и однородны,
можна утверждать, що сумма

(21.4)

\begin{verbatim}
\frac{\pi}{2} e^{ax}
\end{verbatim}

також удовлетворяет уравнению и граничным условиям. Единственное, що нам
осталось---це выбрать коефіцієнти \texttt{a\_n} и \texttt{b\_n} так,
чтобы удовлетворить начальным условиям. Подставляя (21.4) в початкові
умови, получаем

(21.5)

Таким образом, мы получили отклики \texttt{U\_n(t)} на простейшие
входные воздействия \texttt{F\_n(t)} . Для получения решения исходной
задачи осталось совершить последний шаг---сложить все отклики

\begin{verbatim}
n = 1, 2, ...
C_n = \frac{2}{\pi} \int_0^{\pi} f(x) \cos(nx) dx.
\end{verbatim}

f'\,'(x)

\texttt{Читатель,\ должно\ быть,\ заметил,\ что\ каждый\ отклик\ \textquotesingle{}U\_n(t)\textquotesingle{}\ умножается\ на\ весовой\ множитель\ \textquotesingle{}\textbackslash{}sin(n\textbackslash{}pi\ x)\textquotesingle{}\ .\ Конечно,\ в\ разложении\ функции\ \textquotesingle{}f(x,\ t)\textquotesingle{}\ на\ компоненты\ \textquotesingle{}F\_n(t)\textquotesingle{}\ присутствуют\ те\ же\ весовые\ множители.}

\section{ЗАМЕЧАНИЕ}\label{ux437ux430ux43cux435ux447ux430ux43dux438ux435}

Використання скінченного синусоїдального перетворення дозволяє записати
розв'язок у вигляді ряду, але для більшості інших інтегральних
перетворень розв'язки представлені як інтеграли (неперервний розклад).
Читача можна попросити дати інтерпретацію проблеми

Таким образом, розв'язок задачи имеет вид (21.4), а коефіцієнти
\texttt{a\_n} и \texttt{b\_n} вычисляются по формулам (21.6)

Чтобы помочь читателю разобраться в цієї задаче, приведем простой
пример.

\section{Модель колеблющейся
балки}\label{ux43cux43eux434ux435ux43bux44c-ux43aux43eux43bux435ux431ux43bux44eux449ux435ux439ux441ux44f-ux431ux430ux43bux43aux438}

Рассмотрим свободно опирающуюся балку, изображенную на рис. 21.2. Пусть
початкові умови таковы:

\begin{verbatim}
y = \frac{\pi}{b-a}(x-a).
\end{verbatim}

\texttt{u\_1}(x, 0) = 0 Мы могли бы найти розв'язок, подставив функції
\texttt{f(x)} и \texttt{g(x)} в формулы (21.6), однако легко видеть, що

\begin{verbatim}
f(x) = \frac{C_0}{2} + \sum_{n=1}^{\infty} C_n \cos(nx).
\end{verbatim}

\begin{verbatim}
для всех  n=1,\,2,\,\ldots,   b_1=1,\,b_2=0,\,b_3=0,5   b_n=0,\,n=4,\,5,\,\ldots  Отже, розв'язок имеет вид
\end{verbatim}

\begin{verbatim}
f(x) = \frac{C_0}{2} + \sum_{n=1}^{\infty} C_n \cos(nx) \qquad C_n = \frac{2}{\pi} \int_{0}^{\pi} f(x) \cos nx \, dx
\end{verbatim}

.

Интересно сравнить це розв'язок с решением для струны при тех же
начальных условиях. Якщо заглянуть в лекцию 20, там можна найти
розв'язок соответствующей задачи для струны:

\begin{verbatim}
0 \le x \le \pi \qquad n = 0, 1, 2, \dots
\end{verbatim}

.

\pandocbounded{\includegraphics[keepaspectratio]{images/_page_158_Figure_1.jpeg}}

РИС 21.2. Простые колебания свободно опирающейся балки.

Другими словами, балка колеблется на более высоких частотах, чем струна.
Отметим, що в том и другом случаях высшие частоты кратны основным
частотам колебаний.

\section{ЗАМЕЧАНИЯ}\label{ux437ux430ux43cux435ux447ux430ux43dux438ux44f-6}

\begin{itemize}
\item
  \begin{enumerate}
  \def\labelenumi{\arabic{enumi}.}
  \tightlist
  \item
    Вообще говоря, балку можна закрепить тремя способами:
  \end{enumerate}
\item
  а) оставить свободной (незакрепленной),
\item
  б) свободно опереть.
\item
  в) жестко закрепить.
\end{itemize}

Схематически эти способы закрепления изображены на рис. 21.3.

\begin{enumerate}
\def\labelenumi{\arabic{enumi}.}
\setcounter{enumi}{1}
\tightlist
\item
  Еще одна важная задача --- колебания балки с заделанным концом (см.
  рис. 213). Благодаря нестандартным граничным условиям
\end{enumerate}

\texttt{u\_x(0,\ t)\ =\ 0,} \texttt{u\_\{xx\}(1,\ t)\ =\ 0,}
\texttt{u\_\{xx\}(1,\ t)\ =\ 0,} розв'язок уже не буде выражаться в виде
ряда из произведений синусов и косинусов. Оно буде иметь более сложную
структуру

где власні функції \texttt{X\_n}(x) являются линейными комбинациями
синусов, косинусов, гиперболических синусов и гиперболических косинусов.
Розв'язок цієї задачи можна найта в руководствах по теории упругости.

\pandocbounded{\includegraphics[keepaspectratio]{images/_page_159_Figure_2.jpeg}}

Рис. 21.3. Типичные задачи для балки: a---защемленная балка; b---балка,
жестко заделанная с обоих концов; a---левый конец жестко заделан, правый
свободно опирается.

\section{ЗАДАЧИ}\label{ux437ux430ux434ux430ux447ux438-13}

\begin{enumerate}
\def\labelenumi{\arabic{enumi}.}
\tightlist
\item
  Решите задачу для балки с закрепленным концом.
\end{enumerate}

f'\,'(x)

\begin{enumerate}
\def\labelenumi{\arabic{enumi}.}
\setcounter{enumi}{5}
\tightlist
\item
\end{enumerate}

\begin{verbatim}
\begin{cases}
\end{verbatim}

\end{document}
